% Options for packages loaded elsewhere
% Options for packages loaded elsewhere
\PassOptionsToPackage{unicode}{hyperref}
\PassOptionsToPackage{hyphens}{url}
\PassOptionsToPackage{dvipsnames,svgnames,x11names}{xcolor}
%
\documentclass[
  letterpaper,
  DIV=11,
  numbers=noendperiod]{scrreprt}
\usepackage{xcolor}
\usepackage{amsmath,amssymb}
\setcounter{secnumdepth}{5}
\usepackage{iftex}
\ifPDFTeX
  \usepackage[T1]{fontenc}
  \usepackage[utf8]{inputenc}
  \usepackage{textcomp} % provide euro and other symbols
\else % if luatex or xetex
  \usepackage{unicode-math} % this also loads fontspec
  \defaultfontfeatures{Scale=MatchLowercase}
  \defaultfontfeatures[\rmfamily]{Ligatures=TeX,Scale=1}
\fi
\usepackage{lmodern}
\ifPDFTeX\else
  % xetex/luatex font selection
  \setmainfont[]{DejaVu Sans}
\fi
% Use upquote if available, for straight quotes in verbatim environments
\IfFileExists{upquote.sty}{\usepackage{upquote}}{}
\IfFileExists{microtype.sty}{% use microtype if available
  \usepackage[]{microtype}
  \UseMicrotypeSet[protrusion]{basicmath} % disable protrusion for tt fonts
}{}
\makeatletter
\@ifundefined{KOMAClassName}{% if non-KOMA class
  \IfFileExists{parskip.sty}{%
    \usepackage{parskip}
  }{% else
    \setlength{\parindent}{0pt}
    \setlength{\parskip}{6pt plus 2pt minus 1pt}}
}{% if KOMA class
  \KOMAoptions{parskip=half}}
\makeatother
% Make \paragraph and \subparagraph free-standing
\makeatletter
\ifx\paragraph\undefined\else
  \let\oldparagraph\paragraph
  \renewcommand{\paragraph}{
    \@ifstar
      \xxxParagraphStar
      \xxxParagraphNoStar
  }
  \newcommand{\xxxParagraphStar}[1]{\oldparagraph*{#1}\mbox{}}
  \newcommand{\xxxParagraphNoStar}[1]{\oldparagraph{#1}\mbox{}}
\fi
\ifx\subparagraph\undefined\else
  \let\oldsubparagraph\subparagraph
  \renewcommand{\subparagraph}{
    \@ifstar
      \xxxSubParagraphStar
      \xxxSubParagraphNoStar
  }
  \newcommand{\xxxSubParagraphStar}[1]{\oldsubparagraph*{#1}\mbox{}}
  \newcommand{\xxxSubParagraphNoStar}[1]{\oldsubparagraph{#1}\mbox{}}
\fi
\makeatother


\usepackage{longtable,booktabs,array}
\usepackage{calc} % for calculating minipage widths
% Correct order of tables after \paragraph or \subparagraph
\usepackage{etoolbox}
\makeatletter
\patchcmd\longtable{\par}{\if@noskipsec\mbox{}\fi\par}{}{}
\makeatother
% Allow footnotes in longtable head/foot
\IfFileExists{footnotehyper.sty}{\usepackage{footnotehyper}}{\usepackage{footnote}}
\makesavenoteenv{longtable}
\usepackage{graphicx}
\makeatletter
\newsavebox\pandoc@box
\newcommand*\pandocbounded[1]{% scales image to fit in text height/width
  \sbox\pandoc@box{#1}%
  \Gscale@div\@tempa{\textheight}{\dimexpr\ht\pandoc@box+\dp\pandoc@box\relax}%
  \Gscale@div\@tempb{\linewidth}{\wd\pandoc@box}%
  \ifdim\@tempb\p@<\@tempa\p@\let\@tempa\@tempb\fi% select the smaller of both
  \ifdim\@tempa\p@<\p@\scalebox{\@tempa}{\usebox\pandoc@box}%
  \else\usebox{\pandoc@box}%
  \fi%
}
% Set default figure placement to htbp
\def\fps@figure{htbp}
\makeatother





\setlength{\emergencystretch}{3em} % prevent overfull lines

\providecommand{\tightlist}{%
  \setlength{\itemsep}{0pt}\setlength{\parskip}{0pt}}



 


\usepackage{cancel}
\usepackage{fontspec}
\KOMAoption{captions}{tableheading}
\makeatletter
\@ifpackageloaded{tcolorbox}{}{\usepackage[skins,breakable]{tcolorbox}}
\@ifpackageloaded{fontawesome5}{}{\usepackage{fontawesome5}}
\definecolor{quarto-callout-color}{HTML}{909090}
\definecolor{quarto-callout-note-color}{HTML}{0758E5}
\definecolor{quarto-callout-important-color}{HTML}{CC1914}
\definecolor{quarto-callout-warning-color}{HTML}{EB9113}
\definecolor{quarto-callout-tip-color}{HTML}{00A047}
\definecolor{quarto-callout-caution-color}{HTML}{FC5300}
\definecolor{quarto-callout-color-frame}{HTML}{acacac}
\definecolor{quarto-callout-note-color-frame}{HTML}{4582ec}
\definecolor{quarto-callout-important-color-frame}{HTML}{d9534f}
\definecolor{quarto-callout-warning-color-frame}{HTML}{f0ad4e}
\definecolor{quarto-callout-tip-color-frame}{HTML}{02b875}
\definecolor{quarto-callout-caution-color-frame}{HTML}{fd7e14}
\makeatother
\makeatletter
\@ifpackageloaded{bookmark}{}{\usepackage{bookmark}}
\makeatother
\makeatletter
\@ifpackageloaded{caption}{}{\usepackage{caption}}
\AtBeginDocument{%
\ifdefined\contentsname
  \renewcommand*\contentsname{Table of contents}
\else
  \newcommand\contentsname{Table of contents}
\fi
\ifdefined\listfigurename
  \renewcommand*\listfigurename{List of Figures}
\else
  \newcommand\listfigurename{List of Figures}
\fi
\ifdefined\listtablename
  \renewcommand*\listtablename{List of Tables}
\else
  \newcommand\listtablename{List of Tables}
\fi
\ifdefined\figurename
  \renewcommand*\figurename{Figure}
\else
  \newcommand\figurename{Figure}
\fi
\ifdefined\tablename
  \renewcommand*\tablename{Table}
\else
  \newcommand\tablename{Table}
\fi
}
\@ifpackageloaded{float}{}{\usepackage{float}}
\floatstyle{ruled}
\@ifundefined{c@chapter}{\newfloat{codelisting}{h}{lop}}{\newfloat{codelisting}{h}{lop}[chapter]}
\floatname{codelisting}{Listing}
\newcommand*\listoflistings{\listof{codelisting}{List of Listings}}
\usepackage{amsthm}
\theoremstyle{definition}
\newtheorem{exercise}{Exercise}[chapter]
\theoremstyle{remark}
\AtBeginDocument{\renewcommand*{\proofname}{Proof}}
\newtheorem*{remark}{Remark}
\newtheorem*{solution}{Solution}
\newtheorem{refremark}{Remark}[chapter]
\newtheorem{refsolution}{Solution}[chapter]
\makeatother
\makeatletter
\makeatother
\makeatletter
\@ifpackageloaded{caption}{}{\usepackage{caption}}
\@ifpackageloaded{subcaption}{}{\usepackage{subcaption}}
\makeatother
\usepackage{bookmark}
\IfFileExists{xurl.sty}{\usepackage{xurl}}{} % add URL line breaks if available
\urlstyle{same}
\hypersetup{
  pdftitle={Math-Booklet},
  pdfauthor={Maths Department},
  colorlinks=true,
  linkcolor={blue},
  filecolor={Maroon},
  citecolor={Blue},
  urlcolor={Blue},
  pdfcreator={LaTeX via pandoc}}


\title{Math-Booklet}
\author{Maths Department}
\date{2026-02-02}
\begin{document}
\maketitle

\renewcommand*\contentsname{Table of contents}
{
\hypersetup{linkcolor=}
\setcounter{tocdepth}{2}
\tableofcontents
}

\bookmarksetup{startatroot}

\chapter*{Welcome to MYP 5
Mathematics}\label{welcome-to-myp-5-mathematics}
\addcontentsline{toc}{chapter}{Welcome to MYP 5 Mathematics}

\markboth{Welcome to MYP 5 Mathematics}{Welcome to MYP 5 Mathematics}

Welcome to your comprehensive revision hub. This isn't just a textbook;
it's a \textbf{living document} designed to evolve with you as we move
toward your final MYP 5 assessments.

\begin{tcolorbox}[enhanced jigsaw, arc=.35mm, breakable, toptitle=1mm, title={🚀 A Growing Resource}, colframe=quarto-callout-note-color-frame, rightrule=.15mm, opacityback=0, titlerule=0mm, colback=white, coltitle=black, bottomtitle=1mm, toprule=.15mm, leftrule=.75mm, bottomrule=.15mm, colbacktitle=quarto-callout-note-color!10!white, left=2mm, opacitybacktitle=0.6]

This booklet is updated \textbf{ongoingly}. New chapters and advanced
units will appear as we cover them in class, ensuring your revision
material is always synchronized with your learning.

\end{tcolorbox}

\section*{How to use this Booklet}\label{how-to-use-this-booklet}
\addcontentsline{toc}{section}{How to use this Booklet}

\markright{How to use this Booklet}

\begin{figure}

\begin{minipage}{0.50\linewidth}

\subsection*{🎯 Mastery Tracking}\label{mastery-tracking}
\addcontentsline{toc}{subsection}{🎯 Mastery Tracking}

Every unit starts with clear \textbf{Learning Targets}. Use these as
your personal checklist. If you can't tick it off, you aren't ready for
the test!\end{minipage}%
%
\begin{minipage}{0.50\linewidth}

\subsection*{🗂️ Formula Cards}\label{formula-cards}
\addcontentsline{toc}{subsection}{🗂️ Formula Cards}

No more hunting through notebooks. Quick-access digital cards provide
the identities and rules you need at a glance.\end{minipage}%
\newline
\begin{minipage}{0.50\linewidth}

\subsection*{⚡ Dynamic Practice}\label{dynamic-practice}
\addcontentsline{toc}{subsection}{⚡ Dynamic Practice}

Our \textbf{``Test Your Knowledge''} sections feature a ``Try Again''
button that generates new numbers for the same problem type. Practice
the \emph{method}, not the memory.\end{minipage}%
%
\begin{minipage}{0.50\linewidth}

\subsection*{🔍 Instant Feedback}\label{instant-feedback}
\addcontentsline{toc}{subsection}{🔍 Instant Feedback}

Every exercise includes a hidden solution. Toggle it only after you've
put pen to paper.\end{minipage}%

\end{figure}%

\section*{📈 The Journey: Tiered
Practice}\label{the-journey-tiered-practice}
\addcontentsline{toc}{section}{📈 The Journey: Tiered Practice}

\markright{📈 The Journey: Tiered Practice}

We use a color-coded system to help you scale your skills from
foundational to exam-ready:

\begin{longtable}[]{@{}
  >{\centering\arraybackslash}p{(\linewidth - 4\tabcolsep) * \real{0.0986}}
  >{\raggedright\arraybackslash}p{(\linewidth - 4\tabcolsep) * \real{0.2394}}
  >{\raggedright\arraybackslash}p{(\linewidth - 4\tabcolsep) * \real{0.6620}}@{}}
\toprule\noalign{}
\begin{minipage}[b]{\linewidth}\centering
Level
\end{minipage} & \begin{minipage}[b]{\linewidth}\raggedright
Focus
\end{minipage} & \begin{minipage}[b]{\linewidth}\raggedright
Goal
\end{minipage} \\
\midrule\noalign{}
\endhead
\bottomrule\noalign{}
\endlastfoot
🟢 & \textbf{Easy} & Building fluency and basic recall. \\
🟡 & \textbf{Medium} & Multi-step processing and application. \\
🔴 & \textbf{Challenging} & Exam-style logic and complex problem
solving. \\
\end{longtable}

\begin{tcolorbox}[enhanced jigsaw, arc=.35mm, breakable, toptitle=1mm, title=\textcolor{quarto-callout-tip-color}{\faLightbulb}\hspace{0.5em}{Pro Tip for Success}, colframe=quarto-callout-tip-color-frame, rightrule=.15mm, opacityback=0, titlerule=0mm, colback=white, coltitle=black, bottomtitle=1mm, toprule=.15mm, leftrule=.75mm, bottomrule=.15mm, colbacktitle=quarto-callout-tip-color!10!white, left=2mm, opacitybacktitle=0.6]

Don't just read the solutions. Mathematics is a ``doing'' subject. Use
the \textbf{Interactive Testing} feature until you can solve the 🔴
Challenging problems without looking at the hints!

\end{tcolorbox}

\part{Algebra}

\chapter{Unit 1: Operations}\label{unit-1-operations}

\section{🎯 Learning Targets}\label{learning-targets}

1

\textbf{Calculate with Integers} -- Add, subtract, multiply, and divide
positive and negative numbers accurately.

2

\textbf{Apply Order of Operations} -- Evaluate numerical expressions
using PEMDAS/BODMAS, including brackets and exponents.

3

\textbf{Operate with Fractions} -- Add, subtract, multiply, and divide
fractions and mixed numbers fluently.

4

\textbf{Work with Decimals} -- Perform calculations involving decimals
and convert between fractions, decimals, and percentages.

5

\textbf{Simplify Complex Calculations} -- Combine integers, fractions,
and decimals in multi-step numerical problems.

\section{🗂️ Recall: Formula Card}\label{recall-formula-card}

\subsection*{Essential Operation Rules}\label{essential-operation-rules}
\addcontentsline{toc}{subsection}{Essential Operation Rules}

\subsubsection*{Simplifying Fractions}\label{simplifying-fractions}
\addcontentsline{toc}{subsubsection}{Simplifying Fractions}

\textbf{💡 Pro Tip} A fraction is fully simplified when GCF = 1

\textbf{How to Simplify}

\begin{enumerate}
\def\labelenumi{\arabic{enumi}.}
\tightlist
\item
  Find the Greatest Common Factor (GCF) of numerator and denominator
\item
  Divide both top and bottom by the GCF
\end{enumerate}

\[\frac{a}{b} = \frac{a \div \text{GCF}}{b \div \text{GCF}}\]

\textbf{Examples}: -
\(\frac{12}{18} = \frac{12 \div 6}{18 \div 6} = \frac{2}{3}\) -
\(\frac{15}{25} = \frac{15 \div 5}{25 \div 5} = \frac{3}{5}\)

\subsubsection*{Rules of Signs (Multiplication \&
Division)}\label{rules-of-signs-multiplication-division}
\addcontentsline{toc}{subsubsection}{Rules of Signs (Multiplication \&
Division)}

\textbf{🎯 Memory Trick} Same signs = Happy (+)\\
Different signs = Sad (-)

\textbf{Same Signs → Positive} - \((+) \times (+) = (+)\) -
\((-) \times (-) = (+)\)

\textbf{Different Signs → Negative} - \((+) \times (-) = (-)\) -
\((-) \times (+) = (-)\)

\subsubsection*{Adding/Subtracting
Fractions}\label{addingsubtracting-fractions}
\addcontentsline{toc}{subsubsection}{Adding/Subtracting Fractions}

\textbf{⚠️ Common Mistake} Don't add denominators!\\
WRONG: \(\frac{1}{2} + \frac{1}{3} = \frac{2}{5}\)

\textbf{Same Denominator}: Keep denominator, add/subtract numerators
\[\frac{a}{c} + \frac{b}{c} = \frac{a + b}{c}\]

\subsubsection*{Multiplying \& Dividing
Fractions}\label{multiplying-dividing-fractions}
\addcontentsline{toc}{subsubsection}{Multiplying \& Dividing Fractions}

\textbf{📝 Quick Method} Cancel BEFORE you multiply - saves time!

\textbf{Multiply}:
\(\frac{a}{b} \times \frac{c}{d} = \frac{a \times c}{b \times d}\)

\textbf{Divide (Keep-Change-Flip)}:
\(\frac{a}{b} \div \frac{c}{d} = \frac{a}{b} \times \frac{d}{c}\)

\subsubsection*{Order of Operations
(PEMDAS/BODMAS)}\label{order-of-operations-pemdasbodmas}
\addcontentsline{toc}{subsubsection}{Order of Operations
(PEMDAS/BODMAS)}

\begin{enumerate}
\def\labelenumi{\arabic{enumi}.}
\tightlist
\item
  \textbf{P}arentheses / \textbf{B}rackets
\item
  \textbf{E}xponents / \textbf{O}rders
\item
  \textbf{M/D} Multiplication \& Division (Left to Right)
\item
  \textbf{A/S} Addition \& Subtraction (Left to Right)
\end{enumerate}

\section{📝 Practice Problems}\label{practice-problems}

\subsection{🟢 Easy}

\begin{tcolorbox}[enhanced jigsaw, arc=.35mm, breakable, toptitle=1mm, title={Guidance}, colframe=quarto-callout-important-color-frame, rightrule=.15mm, opacityback=0, titlerule=0mm, colback=white, coltitle=black, bottomtitle=1mm, toprule=.15mm, leftrule=.75mm, bottomrule=.15mm, colbacktitle=quarto-callout-important-color!10!white, left=2mm, opacitybacktitle=0.6]

\begin{enumerate}
\def\labelenumi{\arabic{enumi}.}
\tightlist
\item
  Read each question carefully before you begin answering it.
\item
  Check your answers seem right.
\item
  Always show your workings.
\end{enumerate}

\end{tcolorbox}

\begin{exercise}[Conversion
Table]\protect\hypertarget{exr-u1-conv-table}{}\label{exr-u1-conv-table}

Complete the following table:

\begin{longtable}[]{@{}lll@{}}
\toprule\noalign{}
Fraction & Decimal & Percentage \\
\midrule\noalign{}
\endhead
\bottomrule\noalign{}
\endlastfoot
\(\frac{1}{2}\) & \ldots{} & \ldots{} \\
\ldots{} & 0.75 & 75\% \\
\end{longtable}

\begin{tcolorbox}[enhanced jigsaw, arc=.35mm, breakable, toptitle=1mm, title=\textcolor{quarto-callout-tip-color}{\faLightbulb}\hspace{0.5em}{🔍 View Solution}, colframe=quarto-callout-tip-color-frame, rightrule=.15mm, opacityback=0, titlerule=0mm, colback=white, coltitle=black, bottomtitle=1mm, toprule=.15mm, leftrule=.75mm, bottomrule=.15mm, colbacktitle=quarto-callout-tip-color!10!white, left=2mm, opacitybacktitle=0.6]

\begin{itemize}
\tightlist
\item
  \textbf{Row 1:} \(\frac{1}{2} = 0.5 = \mathbf{50\%}\)
\item
  \textbf{Row 2:} \(\mathbf{\frac{3}{4}} = 0.75 = 75\%\)
\end{itemize}

\end{tcolorbox}

{Topic: Conversion Table} {Skills: Basic Conversions}

\end{exercise}

∑ × π ÷ √ ◆ ∞ ± ∫ ≠ Δ

\begin{exercise}[Basic
Conversions]\protect\hypertarget{exr-u1-basic-conv}{}\label{exr-u1-basic-conv}

~

\begin{enumerate}
\def\labelenumi{\alph{enumi})}
\tightlist
\item
  Write 20\% as a decimal.\\
\item
  Write 9\% as a fraction.
\end{enumerate}

\begin{tcolorbox}[enhanced jigsaw, arc=.35mm, breakable, toptitle=1mm, title=\textcolor{quarto-callout-tip-color}{\faLightbulb}\hspace{0.5em}{🔍 View Solution}, colframe=quarto-callout-tip-color-frame, rightrule=.15mm, opacityback=0, titlerule=0mm, colback=white, coltitle=black, bottomtitle=1mm, toprule=.15mm, leftrule=.75mm, bottomrule=.15mm, colbacktitle=quarto-callout-tip-color!10!white, left=2mm, opacitybacktitle=0.6]

\begin{enumerate}
\def\labelenumi{\alph{enumi})}
\tightlist
\item
  \(20 \div 100 = \mathbf{0.2}\)\\
\item
  \(\mathbf{\frac{9}{100}}\)
\end{enumerate}

\end{tcolorbox}

{Topic: Basic Conversions} {Skills: Percentage-to-Decimal/Fraction}

\end{exercise}

∑ × π ÷ √ ◆ ∞ ± ∫ ≠ Δ

\begin{tcolorbox}[enhanced jigsaw, arc=.35mm, breakable, toptitle=1mm, title={Level 1: Basic Operations}, colframe=quarto-callout-important-color-frame, rightrule=.15mm, opacityback=0, titlerule=0mm, colback=white, coltitle=black, bottomtitle=1mm, toprule=.15mm, leftrule=.75mm, bottomrule=.15mm, colbacktitle=quarto-callout-important-color!10!white, left=2mm, opacitybacktitle=0.6]

Calculate

\begin{longtable}[]{@{}
  >{\centering\arraybackslash}p{(\linewidth - 6\tabcolsep) * \real{0.2500}}
  >{\centering\arraybackslash}p{(\linewidth - 6\tabcolsep) * \real{0.2500}}
  >{\centering\arraybackslash}p{(\linewidth - 6\tabcolsep) * \real{0.2500}}
  >{\centering\arraybackslash}p{(\linewidth - 6\tabcolsep) * \real{0.2500}}@{}}
\toprule\noalign{}
\begin{minipage}[b]{\linewidth}\centering
a
\end{minipage} & \begin{minipage}[b]{\linewidth}\centering
b
\end{minipage} & \begin{minipage}[b]{\linewidth}\centering
c
\end{minipage} & \begin{minipage}[b]{\linewidth}\centering
d
\end{minipage} \\
\midrule\noalign{}
\endhead
\bottomrule\noalign{}
\endlastfoot
\(\frac{3}{4}\) of 36 & \(\frac{1}{3}\) of 24 & \(\frac{2}{5}\) of 35 &
\(\frac{4}{9}\) of 72 \\
\end{longtable}

\end{tcolorbox}

\begin{tcolorbox}[enhanced jigsaw, arc=.35mm, breakable, toptitle=1mm, title=\textcolor{quarto-callout-tip-color}{\faLightbulb}\hspace{0.5em}{🔍 Solutions}, colframe=quarto-callout-tip-color-frame, rightrule=.15mm, opacityback=0, titlerule=0mm, colback=white, coltitle=black, bottomtitle=1mm, toprule=.15mm, leftrule=.75mm, bottomrule=.15mm, colbacktitle=quarto-callout-tip-color!10!white, left=2mm, opacitybacktitle=0.6]

\begin{enumerate}
\def\labelenumi{\alph{enumi})}
\tightlist
\item
  \((36 \div 4) \times 3 = 27\)\\
\item
  \((24 \div 3) \times 1 = 8\)\\
\item
  \((35 \div 5) \times 2 = 14\)\\
\item
  \((72 \div 9) \times 4 = 32\)
\end{enumerate}

\end{tcolorbox}

{Topic: Fractions of Amounts} {Skills: Multiplicative Reasoning}

∑ × π ÷ √ ◆ ∞ ± ∫ ≠ Δ

\begin{tcolorbox}[enhanced jigsaw, arc=.35mm, breakable, toptitle=1mm, title={Level 2: Applied Context}, colframe=quarto-callout-important-color-frame, rightrule=.15mm, opacityback=0, titlerule=0mm, colback=white, coltitle=black, bottomtitle=1mm, toprule=.15mm, leftrule=.75mm, bottomrule=.15mm, colbacktitle=quarto-callout-important-color!10!white, left=2mm, opacitybacktitle=0.6]

\textbf{Calculate:}\\
a) Zaid earns 800 QAR a month. He spends \(\frac{1}{4}\) on rent and
\(\frac{2}{5}\) on food. How much money does he have left?\\
b) Fatima is buying shoes that normally cost 200 QAR. They are
\(\frac{1}{5}\) off. What is the sale price?\\
c) Amir has 300 marbles. He gives \(\frac{1}{2}\) to his brother and
\(\frac{1}{3}\) of the remaining marbles to his friend. How many does he
keep?

\end{tcolorbox}

\begin{tcolorbox}[enhanced jigsaw, arc=.35mm, breakable, toptitle=1mm, title=\textcolor{quarto-callout-tip-color}{\faLightbulb}\hspace{0.5em}{🔍 Solutions}, colframe=quarto-callout-tip-color-frame, rightrule=.15mm, opacityback=0, titlerule=0mm, colback=white, coltitle=black, bottomtitle=1mm, toprule=.15mm, leftrule=.75mm, bottomrule=.15mm, colbacktitle=quarto-callout-tip-color!10!white, left=2mm, opacitybacktitle=0.6]

\begin{enumerate}
\def\labelenumi{\alph{enumi})}
\item
  Rent: \((800 \div 4) \times 1 = 200\)\\
  Food: \((800 \div 5) \times 2 = 320\)\\
  Total spent: \(200 + 320 = 520\)\\
  Left: \(800 - 520 = 280\) QAR
\item
  Discount: \((200 \div 5) \times 1 = 40\)\\
  Sale Price: \(200 - 40 = 160\) QAR
\item
  To brother: \(300 \times \frac{1}{2} = 150\)\\
  Remaining: \(150\)\\
  To friend: \(150 \times \frac{1}{3} = 50\)\\
  Kept: \(150 - 50 = 100\)
\end{enumerate}

\end{tcolorbox}

{Topic: Fractions of Amounts}\\
{Skills: Word Problems}

\subsection{🟡 Medium}

\begin{exercise}[Time and
Fractions]\protect\hypertarget{exr-u1-time-frac}{}\label{exr-u1-time-frac}

A train is late. It should arrive at 5 pm but arrives at 5:15 pm.

\begin{enumerate}
\def\labelenumi{\alph{enumi})}
\tightlist
\item
  How many minutes late is the train?\\
\item
  Write your answer as a fraction of an hour.
\end{enumerate}

\begin{tcolorbox}[enhanced jigsaw, arc=.35mm, breakable, toptitle=1mm, title=\textcolor{quarto-callout-tip-color}{\faLightbulb}\hspace{0.5em}{🔍 View Solution}, colframe=quarto-callout-tip-color-frame, rightrule=.15mm, opacityback=0, titlerule=0mm, colback=white, coltitle=black, bottomtitle=1mm, toprule=.15mm, leftrule=.75mm, bottomrule=.15mm, colbacktitle=quarto-callout-tip-color!10!white, left=2mm, opacitybacktitle=0.6]

\begin{enumerate}
\def\labelenumi{\alph{enumi})}
\tightlist
\item
  \textbf{15 minutes} b) \(\frac{15}{60} = \mathbf{\frac{1}{4}}\) of an
  hour
\end{enumerate}

\end{tcolorbox}

{Topic: Time and Fractions} {Skills: Fraction Simplification}

\end{exercise}

∑ × π ÷ √ ◆ ∞ ± ∫ ≠ Δ

\begin{exercise}[Savings and
Values]\protect\hypertarget{exr-u1-savings-perc}{}\label{exr-u1-savings-perc}

For every £200 Mrs.~Joy earns, she saves £34.

\begin{enumerate}
\def\labelenumi{\alph{enumi})}
\tightlist
\item
  Work out £34 as a percentage of £200.\\
\item
  Last month she earned £1000. How much does she save?
\end{enumerate}

\begin{tcolorbox}[enhanced jigsaw, arc=.35mm, breakable, toptitle=1mm, title=\textcolor{quarto-callout-tip-color}{\faLightbulb}\hspace{0.5em}{🔍 View Solution}, colframe=quarto-callout-tip-color-frame, rightrule=.15mm, opacityback=0, titlerule=0mm, colback=white, coltitle=black, bottomtitle=1mm, toprule=.15mm, leftrule=.75mm, bottomrule=.15mm, colbacktitle=quarto-callout-tip-color!10!white, left=2mm, opacitybacktitle=0.6]

\begin{enumerate}
\def\labelenumi{\alph{enumi})}
\tightlist
\item
  \(\frac{34}{200} \times 100 = \mathbf{17\%}\)\\
\item
  \(17\%\) of \(£1000 = 0.17 \times 1000 = \mathbf{£170}\)
\end{enumerate}

\end{tcolorbox}

{Topic: Savings} {Skills: Values as Percentages}

\end{exercise}

∑ × π ÷ √ ◆ ∞ ± ∫ ≠ Δ

\subsection{🔴 Challenging}

\begin{tcolorbox}[enhanced jigsaw, arc=.35mm, breakable, toptitle=1mm, title={Multi-Step Reasoning}, colframe=quarto-callout-important-color-frame, rightrule=.15mm, opacityback=0, titlerule=0mm, colback=white, coltitle=black, bottomtitle=1mm, toprule=.15mm, leftrule=.75mm, bottomrule=.15mm, colbacktitle=quarto-callout-important-color!10!white, left=2mm, opacitybacktitle=0.6]

\begin{enumerate}
\def\labelenumi{\alph{enumi})}
\item
  Maryam works 8 hours a week and earns 50 QAR per hour. She decides to
  save \(\frac{2}{5}\) of her total weekly earnings. How many weeks will
  it take her to save 640 QAR?
\item
  There are 40 students in a class. \(\frac{3}{5}\) are girls.
  \(\frac{1}{4}\) of the girls wear glasses. How many girls in the class
  wear glasses?
\end{enumerate}

\end{tcolorbox}

\begin{tcolorbox}[enhanced jigsaw, arc=.35mm, breakable, toptitle=1mm, title=\textcolor{quarto-callout-tip-color}{\faLightbulb}\hspace{0.5em}{🔍 Solutions}, colframe=quarto-callout-tip-color-frame, rightrule=.15mm, opacityback=0, titlerule=0mm, colback=white, coltitle=black, bottomtitle=1mm, toprule=.15mm, leftrule=.75mm, bottomrule=.15mm, colbacktitle=quarto-callout-tip-color!10!white, left=2mm, opacitybacktitle=0.6]

\begin{enumerate}
\def\labelenumi{\alph{enumi})}
\item
  Weekly Pay: \(8 \times 50 = 400\)\\
  Weekly Savings: \((400 \div 5) \times 2 = 160\)\\
  Weeks needed: \(640 \div 160 = 4\) weeks
\item
  Total Girls: \((40 \div 5) \times 3 = 24\)\\
  Girls with glasses: \((24 \div 4) \times 1 = 6\)
\end{enumerate}

\end{tcolorbox}

{Topic: Fractions of Amounts}\\
{Skills: Logic \& Savings}

\section{⚡ Test Your Knowledge}\label{sec-ope-worksheet}

\chapter{Unit 2: Indices}\label{unit-2-indices}

\section{🎯 Learning Targets}\label{learning-targets-1}

1

\textbf{Apply Multiplication Rules} -- Simplify expressions by adding
powers when multiplying terms with the same base.

2

\textbf{Apply Division Rules} -- Simplify expressions by subtracting
powers when dividing terms with the same base.

3

\textbf{Master Power of a Power} -- Use the rule to simplify expressions
where a power is raised to another power.

4

\textbf{Manage Coefficients} -- Accurately multiply or divide numerical
coefficients separately from the algebraic index laws.

5

\textbf{Handle Negative Indices} -- Correctly perform addition and
subtraction of exponents when calculations involve negative numbers.

6

\textbf{Simplify Complex Fractions} -- Combine multiple index laws to
simplify multi-step algebraic fractions involving multiplication and
division.

\section{🗂️ Formula Card}\label{formula-card}

\subsection*{Algebraic Indices \& Expansion
Rules}\label{algebraic-indices-expansion-rules}
\addcontentsline{toc}{subsection}{Algebraic Indices \& Expansion Rules}

\subsubsection*{Laws of Indices (The Core
Rules)}\label{laws-of-indices-the-core-rules}
\addcontentsline{toc}{subsubsection}{Laws of Indices (The Core Rules)}

\textbf{💡 Pro Tip} These rules ONLY apply when the \textbf{base} (the
big letter) is the same.

\textbf{Multiplication \& Division} 1. \textbf{Multiplication}: Add the
powers \(\rightarrow\) \[x^a \times x^b = x^{a+b}\] 2.
\textbf{Division}: Subtract the powers \(\rightarrow\)
\[\frac{x^a}{x^b} = x^{a-b}\]

\textbf{Examples}: \[y^5 \times y^3 = y^{5+3} = y^8\]
\[10w^7 \div 2w^3 = (10 \div 2)w^{7-3} = 5w^4\]

\subsubsection*{Power of a Power \& Zero}\label{power-of-a-power-zero}
\addcontentsline{toc}{subsubsection}{Power of a Power \& Zero}

\textbf{🎯 Memory Trick} Brackets mean multiply!\\
Anything to the power of \(0\) is a hero (it becomes \(1\)).

\textbf{The Bracket Rule}: \[(x^a)^b = x^{a \times b} = x^{ab}\]

\textbf{Examples}: \[(k^3)^4 = k^{3 \times 4} = k^{12}\]
\[(5x^2)^3 = 5^3 \times x^{2 \times 3} = 125x^6\]

\textbf{The Zero Rule}: \[x^0 = 1\]

\subsubsection*{Negative \& Fractional
Indices}\label{negative-fractional-indices}
\addcontentsline{toc}{subsubsection}{Negative \& Fractional Indices}

\textbf{⚠️ Common Mistake} A negative power does NOT make the number
negative; it makes it a reciprocal (fraction)!

\textbf{Negative Powers}: \[x^{-a} = \frac{1}{x^a}\]\\
\textbf{Fractional Powers}: \[x^{\frac{1}{n}} = \sqrt[n]{x}\]

\textbf{Examples}: \[4^{-2} = \frac{1}{4^2} = \frac{1}{16}\]
\[9^{\frac{1}{2}} = \sqrt{9} = 3\]

\subsubsection*{Expanding Brackets with
Indices}\label{expanding-brackets-with-indices}
\addcontentsline{toc}{subsubsection}{Expanding Brackets with Indices}

\textbf{📝 Quick Method} Multiply the ``Big'' numbers (coefficients),
add the ``Small'' numbers (indices).

\textbf{Single Brackets}: \[a(b + c) = ab + ac\]

\textbf{Examples}:
\[3x^2(2x^3 + 4) = (3 \times 2)x^{2+3} + (3 \times 4)x^2 = 6x^5 + 12x^2\]
\[x^5(x^2 - x^{-3}) = x^{5+2} - x^{5-3} = x^7 - x^2\]

\subsubsection*{Algebraic Fractions with
Indices}\label{algebraic-fractions-with-indices}
\addcontentsline{toc}{subsubsection}{Algebraic Fractions with Indices}

\textbf{Simplifying Fractions}:

\begin{enumerate}
\def\labelenumi{\arabic{enumi}.}
\item
  Simplify the numerator (top) first using multiplication laws.
\item
  Simplify the denominator (bottom).
\item
  Use the division law (subtract) to find the final result.
\end{enumerate}

\textbf{Example}:
\[\frac{a^5 \times a^4}{a^2} = \frac{a^{5+4}}{a^2} = \frac{a^9}{a^2} = a^{9-2} = a^7\]

\textbf{Example with Coefficients}:
\[\frac{(2x^3)^2}{4x^4} = \frac{4x^6}{4x^4} = 1x^{6-4} = x^2\]

\section{📝 Practice Problems}\label{practice-problems-1}

\subsection{🟢 Easy}

\begin{exercise}[Multiplication
Rule]\protect\hypertarget{exr-u2-1}{}\label{exr-u2-1}

Simplify. Give your answer in index form.

\begin{enumerate}
\def\labelenumi{\alph{enumi})}
\tightlist
\item
  \(w^{2} \times w^{2}\)\\
\item
  \(n^{3} \times n^{3}\)\\
\item
  \(c^{7} \times c^{2}\)\\
\item
  \(w^{7} \times w^{6}\)\\
\item
  \(b^{6} \times b^{5}\)\\
\item
  \(a^{2} \times a^{2}\)
\end{enumerate}

\begin{tcolorbox}[enhanced jigsaw, arc=.35mm, breakable, toptitle=1mm, title=\textcolor{quarto-callout-tip-color}{\faLightbulb}\hspace{0.5em}{🔍 View Solution}, colframe=quarto-callout-tip-color-frame, rightrule=.15mm, opacityback=0, titlerule=0mm, colback=white, coltitle=black, bottomtitle=1mm, toprule=.15mm, leftrule=.75mm, bottomrule=.15mm, colbacktitle=quarto-callout-tip-color!10!white, left=2mm, opacitybacktitle=0.6]

\begin{enumerate}
\def\labelenumi{\alph{enumi})}
\tightlist
\item
  \(w^{4}\)\\
\item
  \(n^{6}\)\\
\item
  \(c^{9}\)\\
\item
  \(w^{13}\)\\
\item
  \(b^{11}\)\\
\item
  \(a^{4}\)\\
\end{enumerate}

\end{tcolorbox}

{Index notation} {Same base identification} {Addition of integers}

\end{exercise}

\begin{exercise}[Multiplication Rule --- Three
Terms]\protect\hypertarget{exr-u2-2}{}\label{exr-u2-2}

Simplify. Give your answer in index form.

\begin{enumerate}
\def\labelenumi{\alph{enumi})}
\tightlist
\item
  \(m^{3} \times m^{6} \times m^{6}\)\\
\item
  \(a^{6} \times a^{3} \times a^{6}\)\\
\item
  \(y^{3} \times y^{5} \times y^{6}\)\\
\item
  \(n^{2} \times n^{3} \times n^{5}\)\\
\item
  \(x^{4} \times x^{3} \times x^{3}\)\\
\item
  \(x^{2} \times x^{2} \times x^{5}\)
\end{enumerate}

\begin{tcolorbox}[enhanced jigsaw, arc=.35mm, breakable, toptitle=1mm, title=\textcolor{quarto-callout-tip-color}{\faLightbulb}\hspace{0.5em}{🔍 View Solution}, colframe=quarto-callout-tip-color-frame, rightrule=.15mm, opacityback=0, titlerule=0mm, colback=white, coltitle=black, bottomtitle=1mm, toprule=.15mm, leftrule=.75mm, bottomrule=.15mm, colbacktitle=quarto-callout-tip-color!10!white, left=2mm, opacitybacktitle=0.6]

\begin{enumerate}
\def\labelenumi{\alph{enumi})}
\tightlist
\item
  \(m^{15}\)\\
\item
  \(a^{15}\)\\
\item
  \(y^{14}\)\\
\item
  \(n^{10}\)\\
\item
  \(x^{10}\)\\
\item
  \(x^{9}\)\\
\end{enumerate}

\end{tcolorbox}

{Index notation} {Same base identification} {Addition of integers}
{Multi step simplification}

\end{exercise}

\begin{exercise}[Division
Rule]\protect\hypertarget{exr-u2-7}{}\label{exr-u2-7}

Simplify. Give your answer in index form.

\begin{enumerate}
\def\labelenumi{\alph{enumi})}
\tightlist
\item
  \(w^{8} \div w^{2}\)\\
\item
  \(p^{10} \div p^{4}\)\\
\item
  \(x^{5} \div x^{2}\)\\
\item
  \(y^{7} \div y^{3}\)\\
\item
  \(a^{9} \div a^{4}\)\\
\item
  \(c^{8} \div c^{2}\)
\end{enumerate}

\begin{tcolorbox}[enhanced jigsaw, arc=.35mm, breakable, toptitle=1mm, title=\textcolor{quarto-callout-tip-color}{\faLightbulb}\hspace{0.5em}{🔍 View Solution}, colframe=quarto-callout-tip-color-frame, rightrule=.15mm, opacityback=0, titlerule=0mm, colback=white, coltitle=black, bottomtitle=1mm, toprule=.15mm, leftrule=.75mm, bottomrule=.15mm, colbacktitle=quarto-callout-tip-color!10!white, left=2mm, opacitybacktitle=0.6]

\begin{enumerate}
\def\labelenumi{\alph{enumi})}
\tightlist
\item
  \(w^{6}\)\\
\item
  \(p^{6}\)\\
\item
  \(x^{3}\)\\
\item
  \(y^{4}\)\\
\item
  \(a^{5}\)\\
\item
  \(c^{6}\)\\
\end{enumerate}

\end{tcolorbox}

{Index notation} {Same base identification} {Subtraction of integers}

\end{exercise}

\begin{exercise}[Power of a
Power]\protect\hypertarget{exr-u2-10}{}\label{exr-u2-10}

Simplify. Give your answer in index form.

\begin{enumerate}
\def\labelenumi{\alph{enumi})}
\tightlist
\item
  \((k^{2})^{2}\)\\
\item
  \((a^{5})^{4}\)\\
\item
  \((b^{3})^{3}\)\\
\item
  \((m^{6})^{5}\)\\
\item
  \((c^{5})^{3}\)\\
\item
  \((n^{5})^{3}\)
\end{enumerate}

\begin{tcolorbox}[enhanced jigsaw, arc=.35mm, breakable, toptitle=1mm, title=\textcolor{quarto-callout-tip-color}{\faLightbulb}\hspace{0.5em}{🔍 View Solution}, colframe=quarto-callout-tip-color-frame, rightrule=.15mm, opacityback=0, titlerule=0mm, colback=white, coltitle=black, bottomtitle=1mm, toprule=.15mm, leftrule=.75mm, bottomrule=.15mm, colbacktitle=quarto-callout-tip-color!10!white, left=2mm, opacitybacktitle=0.6]

\begin{enumerate}
\def\labelenumi{\alph{enumi})}
\tightlist
\item
  \(k^{4}\)\\
\item
  \(a^{20}\)\\
\item
  \(b^{9}\)\\
\item
  \(m^{30}\)\\
\item
  \(c^{15}\)\\
\item
  \(n^{15}\)\\
\end{enumerate}

\end{tcolorbox}

{Index notation} {Multiplication of integers} {Power of power law}

\end{exercise}

\begin{exercise}[Zero
Index]\protect\hypertarget{exr-u2-21}{}\label{exr-u2-21}

Evaluate.

\begin{enumerate}
\def\labelenumi{\alph{enumi})}
\tightlist
\item
  \(1\)\\
\item
  \(1\)\\
\item
  \(1\)\\
\item
  \(1\)\\
\item
  \(1\)\\
\item
  \(1\)
\end{enumerate}

\begin{tcolorbox}[enhanced jigsaw, arc=.35mm, breakable, toptitle=1mm, title=\textcolor{quarto-callout-tip-color}{\faLightbulb}\hspace{0.5em}{🔍 View Solution}, colframe=quarto-callout-tip-color-frame, rightrule=.15mm, opacityback=0, titlerule=0mm, colback=white, coltitle=black, bottomtitle=1mm, toprule=.15mm, leftrule=.75mm, bottomrule=.15mm, colbacktitle=quarto-callout-tip-color!10!white, left=2mm, opacitybacktitle=0.6]

\begin{enumerate}
\def\labelenumi{\alph{enumi})}
\tightlist
\item
  \(1\)\\
\item
  \(1\)\\
\item
  \(1\)\\
\item
  \(1\)\\
\item
  \(1\)\\
\item
  \(1\)\\
\end{enumerate}

\end{tcolorbox}

{Zero exponent rule} {Index notation}

\end{exercise}

∑ × π ÷ √ ◆ ∞ ± ∫ ≠ Δ

\subsection{🟡 Medium}

\begin{exercise}[Multiplication with
Coefficients]\protect\hypertarget{exr-u2-3}{}\label{exr-u2-3}

Simplify. Give your answer in index form.

\begin{enumerate}
\def\labelenumi{\alph{enumi})}
\tightlist
\item
  \(4b^{6} \times 4b^{4}\)\\
\item
  \(5a^{2} \times 6a^{5}\)\\
\item
  \(6b^{7} \times 4b^{6}\)\\
\item
  \(6x^{7} \times 3x^{2}\)\\
\item
  \(3a^{2} \times 4a^{3}\)\\
\item
  \(5b^{5} \times 4b^{7}\)
\end{enumerate}

\begin{tcolorbox}[enhanced jigsaw, arc=.35mm, breakable, toptitle=1mm, title=\textcolor{quarto-callout-tip-color}{\faLightbulb}\hspace{0.5em}{🔍 View Solution}, colframe=quarto-callout-tip-color-frame, rightrule=.15mm, opacityback=0, titlerule=0mm, colback=white, coltitle=black, bottomtitle=1mm, toprule=.15mm, leftrule=.75mm, bottomrule=.15mm, colbacktitle=quarto-callout-tip-color!10!white, left=2mm, opacitybacktitle=0.6]

\begin{enumerate}
\def\labelenumi{\alph{enumi})}
\tightlist
\item
  \(16b^{10}\)\\
\item
  \(30a^{7}\)\\
\item
  \(24b^{13}\)\\
\item
  \(18x^{9}\)\\
\item
  \(12a^{5}\)\\
\item
  \(20b^{12}\)\\
\end{enumerate}

\end{tcolorbox}

{Index notation} {Coefficient multiplication} {Addition of integers}
{Separating coefficients from indices}

\end{exercise}

\begin{exercise}[Combining Same Powers --- (ab)ⁿ =
aⁿ×bⁿ]\protect\hypertarget{exr-u2-5}{}\label{exr-u2-5}

Write as a single power by expressing each number as a prime power.

\begin{enumerate}
\def\labelenumi{\alph{enumi})}
\tightlist
\item
  \(4 \times 49\)\\
\item
  \(32 \times 243\)\\
\item
  \(32 \times 243\)\\
\item
  \(8 \times 27\)\\
\item
  \(9 \times 25\)\\
\item
  \(32 \times 243\)
\end{enumerate}

\begin{tcolorbox}[enhanced jigsaw, arc=.35mm, breakable, toptitle=1mm, title=\textcolor{quarto-callout-tip-color}{\faLightbulb}\hspace{0.5em}{🔍 View Solution}, colframe=quarto-callout-tip-color-frame, rightrule=.15mm, opacityback=0, titlerule=0mm, colback=white, coltitle=black, bottomtitle=1mm, toprule=.15mm, leftrule=.75mm, bottomrule=.15mm, colbacktitle=quarto-callout-tip-color!10!white, left=2mm, opacitybacktitle=0.6]

\begin{enumerate}
\def\labelenumi{\alph{enumi})}
\tightlist
\item
  \(14^{2}\)\\
\item
  \(6^{5}\)\\
\item
  \(6^{5}\)\\
\item
  \(6^{3}\)\\
\item
  \(15^{2}\)\\
\item
  \(6^{5}\)\\
\end{enumerate}

\end{tcolorbox}

{Prime factorisation} {Perfect square recognition} {Perfect cube
recognition} {Power distribution law}

\end{exercise}

\begin{exercise}[Multiplication ---
Multi-variable]\protect\hypertarget{exr-u2-6}{}\label{exr-u2-6}

Simplify. Give your answer in index form.

\begin{enumerate}
\def\labelenumi{\alph{enumi})}
\tightlist
\item
  \(3c^{5}t^{3} \times 3c^{5}t^{2}\)\\
\item
  \(2w^{2}x^{4} \times 3w^{4}x^{3}\)\\
\item
  \(4b^{5}n^{4} \times 3b^{5}n^{2}\)\\
\item
  \(3n^{5}c^{3} \times 4n^{5}c^{2}\)\\
\item
  \(5c^{2}t^{2} \times 2c^{2}t^{2}\)\\
\item
  \(2w^{5}y^{3} \times 5w^{5}y\)
\end{enumerate}

\begin{tcolorbox}[enhanced jigsaw, arc=.35mm, breakable, toptitle=1mm, title=\textcolor{quarto-callout-tip-color}{\faLightbulb}\hspace{0.5em}{🔍 View Solution}, colframe=quarto-callout-tip-color-frame, rightrule=.15mm, opacityback=0, titlerule=0mm, colback=white, coltitle=black, bottomtitle=1mm, toprule=.15mm, leftrule=.75mm, bottomrule=.15mm, colbacktitle=quarto-callout-tip-color!10!white, left=2mm, opacitybacktitle=0.6]

\begin{enumerate}
\def\labelenumi{\alph{enumi})}
\tightlist
\item
  \(9c^{10}t^{5}\)\\
\item
  \(6w^{6}x^{7}\)\\
\item
  \(12b^{10}n^{6}\)\\
\item
  \(12n^{10}c^{5}\)\\
\item
  \(10c^{4}t^{4}\)\\
\item
  \(10w^{10}y^{4}\)\\
\end{enumerate}

\end{tcolorbox}

{Multi variable expressions} {Coefficient multiplication} {Addition of
integers} {Grouping like terms}

\end{exercise}

\begin{exercise}[Division with
Coefficients]\protect\hypertarget{exr-u2-8}{}\label{exr-u2-8}

Simplify. Give your answer in index form.

\begin{enumerate}
\def\labelenumi{\alph{enumi})}
\tightlist
\item
  \(12n^{5} \div 4n^{3}\)\\
\item
  \(20p^{7} \div 4p^{2}\)\\
\item
  \(12a^{5} \div 4a^{2}\)\\
\item
  \(4a^{8} \div 2a^{3}\)\\
\item
  \(9k^{7} \div 3k^{2}\)\\
\item
  \(12w^{10} \div 3w^{5}\)
\end{enumerate}

\begin{tcolorbox}[enhanced jigsaw, arc=.35mm, breakable, toptitle=1mm, title=\textcolor{quarto-callout-tip-color}{\faLightbulb}\hspace{0.5em}{🔍 View Solution}, colframe=quarto-callout-tip-color-frame, rightrule=.15mm, opacityback=0, titlerule=0mm, colback=white, coltitle=black, bottomtitle=1mm, toprule=.15mm, leftrule=.75mm, bottomrule=.15mm, colbacktitle=quarto-callout-tip-color!10!white, left=2mm, opacitybacktitle=0.6]

\begin{enumerate}
\def\labelenumi{\alph{enumi})}
\tightlist
\item
  \(3n^{2}\)\\
\item
  \(5p^{5}\)\\
\item
  \(3a^{3}\)\\
\item
  \(2a^{5}\)\\
\item
  \(3k^{5}\)\\
\item
  \(4w^{5}\)\\
\end{enumerate}

\end{tcolorbox}

{Coefficient division} {Subtraction of integers} {Separating
coefficients from indices}

\end{exercise}

\begin{exercise}[Division ---
Multi-variable]\protect\hypertarget{exr-u2-9}{}\label{exr-u2-9}

Simplify. Give your answer in index form.

\begin{enumerate}
\def\labelenumi{\alph{enumi})}
\tightlist
\item
  \(\dfrac{12p^{6}w^{3}}{3p^{4}w}\)\\
\item
  \(\dfrac{4w^{8}x^{3}}{2w^{4}x^{3}}\)\\
\item
  \(\dfrac{4b^{3}y^{4}}{2by^{3}}\)\\
\item
  \(\dfrac{8w^{3}a^{3}}{2w^{2}a}\)\\
\item
  \(\dfrac{6m^{6}x^{4}}{2m^{4}x}\)\\
\item
  \(\dfrac{9y^{5}m^{3}}{3ym}\)
\end{enumerate}

\begin{tcolorbox}[enhanced jigsaw, arc=.35mm, breakable, toptitle=1mm, title=\textcolor{quarto-callout-tip-color}{\faLightbulb}\hspace{0.5em}{🔍 View Solution}, colframe=quarto-callout-tip-color-frame, rightrule=.15mm, opacityback=0, titlerule=0mm, colback=white, coltitle=black, bottomtitle=1mm, toprule=.15mm, leftrule=.75mm, bottomrule=.15mm, colbacktitle=quarto-callout-tip-color!10!white, left=2mm, opacitybacktitle=0.6]

\begin{enumerate}
\def\labelenumi{\alph{enumi})}
\tightlist
\item
  \(4p^{2}w^{2}\)\\
\item
  \(2w^{4}1\)\\
\item
  \(2b^{2}y\)\\
\item
  \(4wa^{2}\)\\
\item
  \(3m^{2}x^{3}\)\\
\item
  \(3y^{4}m^{2}\)\\
\end{enumerate}

\end{tcolorbox}

{Multi variable expressions} {Coefficient division} {Subtraction of
integers} {Grouping like terms}

\end{exercise}

\begin{exercise}[Power of a Power with
Coefficient]\protect\hypertarget{exr-u2-11}{}\label{exr-u2-11}

Expand and simplify.

\begin{enumerate}
\def\labelenumi{\alph{enumi})}
\tightlist
\item
  \((5b^{2})^{2}\)\\
\item
  \((2w^{2})^{2}\)\\
\item
  \((5c^{5})^{3}\)\\
\item
  \((5m^{2})^{2}\)\\
\item
  \((2y^{5})^{3}\)\\
\item
  \((4k^{5})^{4}\)
\end{enumerate}

\begin{tcolorbox}[enhanced jigsaw, arc=.35mm, breakable, toptitle=1mm, title=\textcolor{quarto-callout-tip-color}{\faLightbulb}\hspace{0.5em}{🔍 View Solution}, colframe=quarto-callout-tip-color-frame, rightrule=.15mm, opacityback=0, titlerule=0mm, colback=white, coltitle=black, bottomtitle=1mm, toprule=.15mm, leftrule=.75mm, bottomrule=.15mm, colbacktitle=quarto-callout-tip-color!10!white, left=2mm, opacitybacktitle=0.6]

\begin{enumerate}
\def\labelenumi{\alph{enumi})}
\tightlist
\item
  \(25b^{4}\)\\
\item
  \(4w^{4}\)\\
\item
  \(125c^{15}\)\\
\item
  \(25m^{4}\)\\
\item
  \(8y^{15}\)\\
\item
  \(256k^{20}\)\\
\end{enumerate}

\end{tcolorbox}

{Power of power law} {Coefficient exponentiation} {Multiplication of
integers}

\end{exercise}

\begin{exercise}[Power of a Power ---
Multi-variable]\protect\hypertarget{exr-u2-12}{}\label{exr-u2-12}

Expand and simplify.

\begin{enumerate}
\def\labelenumi{\alph{enumi})}
\tightlist
\item
  \((2p^{2}k^{3})^{2}\)\\
\item
  \((4aw^{3})^{2}\)\\
\item
  \((3a^{2}w)^{2}\)\\
\item
  \((4cb^{2})^{3}\)\\
\item
  \((2bw)^{3}\)\\
\item
  \((4w^{3}t^{3})^{2}\)
\end{enumerate}

\begin{tcolorbox}[enhanced jigsaw, arc=.35mm, breakable, toptitle=1mm, title=\textcolor{quarto-callout-tip-color}{\faLightbulb}\hspace{0.5em}{🔍 View Solution}, colframe=quarto-callout-tip-color-frame, rightrule=.15mm, opacityback=0, titlerule=0mm, colback=white, coltitle=black, bottomtitle=1mm, toprule=.15mm, leftrule=.75mm, bottomrule=.15mm, colbacktitle=quarto-callout-tip-color!10!white, left=2mm, opacitybacktitle=0.6]

\begin{enumerate}
\def\labelenumi{\alph{enumi})}
\tightlist
\item
  \(4p^{4}k^{6}\)\\
\item
  \(16a^{2}w^{6}\)\\
\item
  \(9a^{4}w^{2}\)\\
\item
  \(64c^{3}b^{6}\)\\
\item
  \(8b^{3}w^{3}\)\\
\item
  \(16w^{6}t^{6}\)\\
\end{enumerate}

\end{tcolorbox}

{Power of power law} {Coefficient exponentiation} {Multi variable
expressions}

\end{exercise}

\begin{exercise}[Evaluate Negative
Indices]\protect\hypertarget{exr-u2-14}{}\label{exr-u2-14}

Evaluate. Write as a fraction.

\begin{enumerate}
\def\labelenumi{\alph{enumi})}
\tightlist
\item
  \(4^{-1}\)\\
\item
  \(3^{-3}\)\\
\item
  \(4^{-2}\)\\
\item
  \(5^{-3}\)\\
\item
  \(7^{-1}\)\\
\item
  \(10^{-3}\)
\end{enumerate}

\begin{tcolorbox}[enhanced jigsaw, arc=.35mm, breakable, toptitle=1mm, title=\textcolor{quarto-callout-tip-color}{\faLightbulb}\hspace{0.5em}{🔍 View Solution}, colframe=quarto-callout-tip-color-frame, rightrule=.15mm, opacityback=0, titlerule=0mm, colback=white, coltitle=black, bottomtitle=1mm, toprule=.15mm, leftrule=.75mm, bottomrule=.15mm, colbacktitle=quarto-callout-tip-color!10!white, left=2mm, opacitybacktitle=0.6]

\begin{enumerate}
\def\labelenumi{\alph{enumi})}
\tightlist
\item
  \(\dfrac{1}{4}\)\\
\item
  \(\dfrac{1}{27}\)\\
\item
  \(\dfrac{1}{16}\)\\
\item
  \(\dfrac{1}{125}\)\\
\item
  \(\dfrac{1}{7}\)\\
\item
  \(\dfrac{1}{1000}\)\\
\end{enumerate}

\end{tcolorbox}

{Negative index as reciprocal} {Integer exponentiation} {Fraction
notation}

\end{exercise}

\begin{exercise}[Negative Index --- Fraction
Base]\protect\hypertarget{exr-u2-15}{}\label{exr-u2-15}

Evaluate. Simplify your answer.

\begin{enumerate}
\def\labelenumi{\alph{enumi})}
\tightlist
\item
  \(\left(\dfrac{2}{3}\right)^{-1}\)\\
\item
  \(\left(\dfrac{2}{5}\right)^{-1}\)\\
\item
  \(\left(\dfrac{2}{3}\right)^{-3}\)\\
\item
  \(\left(\dfrac{4}{5}\right)^{-1}\)\\
\item
  \(\left(\dfrac{2}{5}\right)^{-2}\)\\
\item
  \(\left(\dfrac{4}{5}\right)^{-1}\)
\end{enumerate}

\begin{tcolorbox}[enhanced jigsaw, arc=.35mm, breakable, toptitle=1mm, title=\textcolor{quarto-callout-tip-color}{\faLightbulb}\hspace{0.5em}{🔍 View Solution}, colframe=quarto-callout-tip-color-frame, rightrule=.15mm, opacityback=0, titlerule=0mm, colback=white, coltitle=black, bottomtitle=1mm, toprule=.15mm, leftrule=.75mm, bottomrule=.15mm, colbacktitle=quarto-callout-tip-color!10!white, left=2mm, opacitybacktitle=0.6]

\begin{enumerate}
\def\labelenumi{\alph{enumi})}
\tightlist
\item
  \(\dfrac{3}{2}\)\\
\item
  \(\dfrac{5}{2}\)\\
\item
  \(\dfrac{27}{8}\)\\
\item
  \(\dfrac{5}{4}\)\\
\item
  \(\dfrac{25}{4}\)\\
\item
  \(\dfrac{5}{4}\)\\
\end{enumerate}

\end{tcolorbox}

{Negative index as reciprocal} {Fraction reciprocal} {Integer
exponentiation}

\end{exercise}

\begin{exercise}[Simplify with Negative
Indices]\protect\hypertarget{exr-u2-16}{}\label{exr-u2-16}

Simplify. Give your answer in index form.

\begin{enumerate}
\def\labelenumi{\alph{enumi})}
\tightlist
\item
  \(x^{3} \times x^{-3}\)\\
\item
  \(p^{5} \times p^{-3}\)\\
\item
  \(a^{2} \times a^{-2}\)\\
\item
  \(n^{2} \times n^{-5}\)\\
\item
  \(x^{3} \times x^{-3}\)\\
\item
  \(c^{5} \times c^{-1}\)
\end{enumerate}

\begin{tcolorbox}[enhanced jigsaw, arc=.35mm, breakable, toptitle=1mm, title=\textcolor{quarto-callout-tip-color}{\faLightbulb}\hspace{0.5em}{🔍 View Solution}, colframe=quarto-callout-tip-color-frame, rightrule=.15mm, opacityback=0, titlerule=0mm, colback=white, coltitle=black, bottomtitle=1mm, toprule=.15mm, leftrule=.75mm, bottomrule=.15mm, colbacktitle=quarto-callout-tip-color!10!white, left=2mm, opacitybacktitle=0.6]

\begin{enumerate}
\def\labelenumi{\alph{enumi})}
\tightlist
\item
  \(1\)\\
\item
  \(p^{2}\)\\
\item
  \(1\)\\
\item
  \(n^{-3}\)\\
\item
  \(1\)\\
\item
  \(c^{4}\)\\
\end{enumerate}

\end{tcolorbox}

{Addition of integers} {Negative integer addition} {Index notation}

\end{exercise}

\begin{exercise}[Division Giving Negative
Index]\protect\hypertarget{exr-u2-17}{}\label{exr-u2-17}

Simplify. Give your answer in index form.

\begin{enumerate}
\def\labelenumi{\alph{enumi})}
\tightlist
\item
  \(y^{2} \div y^{3}\)\\
\item
  \(b^{3} \div b^{4}\)\\
\item
  \(x^{3} \div x^{7}\)\\
\item
  \(c^{2} \div c^{5}\)\\
\item
  \(x^{2} \div x^{5}\)\\
\item
  \(m^{3} \div m^{4}\)
\end{enumerate}

\begin{tcolorbox}[enhanced jigsaw, arc=.35mm, breakable, toptitle=1mm, title=\textcolor{quarto-callout-tip-color}{\faLightbulb}\hspace{0.5em}{🔍 View Solution}, colframe=quarto-callout-tip-color-frame, rightrule=.15mm, opacityback=0, titlerule=0mm, colback=white, coltitle=black, bottomtitle=1mm, toprule=.15mm, leftrule=.75mm, bottomrule=.15mm, colbacktitle=quarto-callout-tip-color!10!white, left=2mm, opacitybacktitle=0.6]

\begin{enumerate}
\def\labelenumi{\alph{enumi})}
\tightlist
\item
  \(y^{-1}\)\\
\item
  \(b^{-1}\)\\
\item
  \(x^{-4}\)\\
\item
  \(c^{-3}\)\\
\item
  \(x^{-3}\)\\
\item
  \(m^{-1}\)\\
\end{enumerate}

\end{tcolorbox}

{Subtraction of integers} {Negative index as reciprocal} {Index
notation}

\end{exercise}

\begin{exercise}[Fractional Index --- Unit
Fraction]\protect\hypertarget{exr-u2-18}{}\label{exr-u2-18}

Evaluate.

\begin{enumerate}
\def\labelenumi{\alph{enumi})}
\tightlist
\item
  \(49^{\frac{1}{2}}\)\\
\item
  \(16^{\frac{1}{4}}\)\\
\item
  \(8^{\frac{1}{3}}\)\\
\item
  \(32^{\frac{1}{5}}\)\\
\item
  \(81^{\frac{1}{4}}\)\\
\item
  \(64^{\frac{1}{2}}\)
\end{enumerate}

\begin{tcolorbox}[enhanced jigsaw, arc=.35mm, breakable, toptitle=1mm, title=\textcolor{quarto-callout-tip-color}{\faLightbulb}\hspace{0.5em}{🔍 View Solution}, colframe=quarto-callout-tip-color-frame, rightrule=.15mm, opacityback=0, titlerule=0mm, colback=white, coltitle=black, bottomtitle=1mm, toprule=.15mm, leftrule=.75mm, bottomrule=.15mm, colbacktitle=quarto-callout-tip-color!10!white, left=2mm, opacitybacktitle=0.6]

\begin{enumerate}
\def\labelenumi{\alph{enumi})}
\tightlist
\item
  \(7\)\\
\item
  \(2\)\\
\item
  \(2\)\\
\item
  \(2\)\\
\item
  \(3\)\\
\item
  \(8\)\\
\end{enumerate}

\end{tcolorbox}

{Fractional index as root} {Perfect square recognition} {Perfect cube
recognition} {Nth root evaluation}

\end{exercise}

\begin{exercise}[Complex
Fraction]\protect\hypertarget{exr-u2-22}{}\label{exr-u2-22}

Simplify. Give your answer in index form.

\begin{enumerate}
\def\labelenumi{\alph{enumi})}
\tightlist
\item
  \(\dfrac{n^{3} \times n^{4}}{n}\)\\
\item
  \(\dfrac{b^{6} \times b^{5}}{b^{6}}\)\\
\item
  \(\dfrac{x^{5} \times x^{2}}{x^{4}}\)\\
\item
  \(\dfrac{t^{3} \times t^{4}}{t}\)\\
\item
  \(\dfrac{y^{2} \times y^{3}}{y^{3}}\)\\
\item
  \(\dfrac{y^{2} \times y^{4}}{y^{5}}\)
\end{enumerate}

\begin{tcolorbox}[enhanced jigsaw, arc=.35mm, breakable, toptitle=1mm, title=\textcolor{quarto-callout-tip-color}{\faLightbulb}\hspace{0.5em}{🔍 View Solution}, colframe=quarto-callout-tip-color-frame, rightrule=.15mm, opacityback=0, titlerule=0mm, colback=white, coltitle=black, bottomtitle=1mm, toprule=.15mm, leftrule=.75mm, bottomrule=.15mm, colbacktitle=quarto-callout-tip-color!10!white, left=2mm, opacitybacktitle=0.6]

\begin{enumerate}
\def\labelenumi{\alph{enumi})}
\tightlist
\item
  \(n^{6}\)\\
\item
  \(b^{5}\)\\
\item
  \(x^{3}\)\\
\item
  \(t^{6}\)\\
\item
  \(y^{2}\)\\
\item
  \(y\)\\
\end{enumerate}

\end{tcolorbox}

{Combining multiplication and division laws} {Multi step simplification}
{Index notation}

\end{exercise}

\begin{exercise}[Numeric Index Laws ---
Multiply]\protect\hypertarget{exr-u2-26}{}\label{exr-u2-26}

Simplify. Give your answer as a single power.

\begin{enumerate}
\def\labelenumi{\alph{enumi})}
\tightlist
\item
  \(3^{5} \times 3^{2}\)\\
\item
  \(3^{4} \times 3^{5}\)\\
\item
  \(5^{3} \times 5^{5}\)\\
\item
  \(2^{6} \times 2^{6}\)\\
\item
  \(2^{5} \times 2^{6}\)\\
\item
  \(5^{2} \times 5^{2}\)
\end{enumerate}

\begin{tcolorbox}[enhanced jigsaw, arc=.35mm, breakable, toptitle=1mm, title=\textcolor{quarto-callout-tip-color}{\faLightbulb}\hspace{0.5em}{🔍 View Solution}, colframe=quarto-callout-tip-color-frame, rightrule=.15mm, opacityback=0, titlerule=0mm, colback=white, coltitle=black, bottomtitle=1mm, toprule=.15mm, leftrule=.75mm, bottomrule=.15mm, colbacktitle=quarto-callout-tip-color!10!white, left=2mm, opacitybacktitle=0.6]

\begin{enumerate}
\def\labelenumi{\alph{enumi})}
\tightlist
\item
  \(3^{7}\)\\
\item
  \(3^{9}\)\\
\item
  \(5^{8}\)\\
\item
  \(2^{12}\)\\
\item
  \(2^{11}\)\\
\item
  \(5^{4}\)\\
\end{enumerate}

\end{tcolorbox}

{Numeric base recognition} {Addition of integers} {Expressing as single
power}

\end{exercise}

\begin{exercise}[Numeric Index Laws ---
Divide]\protect\hypertarget{exr-u2-27}{}\label{exr-u2-27}

Simplify. Give your answer as a single power.

\begin{enumerate}
\def\labelenumi{\alph{enumi})}
\tightlist
\item
  \(5^{7} \div 5^{5}\)\\
\item
  \(3^{4} \div 3^{3}\)\\
\item
  \(3^{8} \div 3^{5}\)\\
\item
  \(2^{8} \div 2^{5}\)\\
\item
  \(3^{8} \div 3^{5}\)\\
\item
  \(3^{5} \div 3^{4}\)
\end{enumerate}

\begin{tcolorbox}[enhanced jigsaw, arc=.35mm, breakable, toptitle=1mm, title=\textcolor{quarto-callout-tip-color}{\faLightbulb}\hspace{0.5em}{🔍 View Solution}, colframe=quarto-callout-tip-color-frame, rightrule=.15mm, opacityback=0, titlerule=0mm, colback=white, coltitle=black, bottomtitle=1mm, toprule=.15mm, leftrule=.75mm, bottomrule=.15mm, colbacktitle=quarto-callout-tip-color!10!white, left=2mm, opacitybacktitle=0.6]

\begin{enumerate}
\def\labelenumi{\alph{enumi})}
\tightlist
\item
  \(5^{2}\)\\
\item
  \(3^{1}\)\\
\item
  \(3^{3}\)\\
\item
  \(2^{3}\)\\
\item
  \(3^{3}\)\\
\item
  \(3^{1}\)\\
\end{enumerate}

\end{tcolorbox}

{Numeric base recognition} {Subtraction of integers} {Expressing as
single power}

\end{exercise}

\begin{exercise}[Write as Power of
n]\protect\hypertarget{exr-u2-29}{}\label{exr-u2-29}

Write as a single power of the given base.

\begin{enumerate}
\def\labelenumi{\alph{enumi})}
\tightlist
\item
  \(\text{Write } 1 \text{ as a power of } 5\)\\
\item
  \(\text{Write } \frac{1}{2} \text{ as a power of } 2\)\\
\item
  \(\text{Write } 16 \text{ as a power of } 2\)\\
\item
  \(\text{Write } \frac{1}{2} \text{ as a power of } 2\)\\
\item
  \(\text{Write } 125 \text{ as a power of } 5\)\\
\item
  \(\text{Write } 16 \text{ as a power of } 2\)
\end{enumerate}

\begin{tcolorbox}[enhanced jigsaw, arc=.35mm, breakable, toptitle=1mm, title=\textcolor{quarto-callout-tip-color}{\faLightbulb}\hspace{0.5em}{🔍 View Solution}, colframe=quarto-callout-tip-color-frame, rightrule=.15mm, opacityback=0, titlerule=0mm, colback=white, coltitle=black, bottomtitle=1mm, toprule=.15mm, leftrule=.75mm, bottomrule=.15mm, colbacktitle=quarto-callout-tip-color!10!white, left=2mm, opacitybacktitle=0.6]

\begin{enumerate}
\def\labelenumi{\alph{enumi})}
\tightlist
\item
  \(5^{0}\)\\
\item
  \(2^{-1}\)\\
\item
  \(2^{4}\)\\
\item
  \(2^{-1}\)\\
\item
  \(5^{3}\)\\
\item
  \(2^{4}\)\\
\end{enumerate}

\end{tcolorbox}

{Prime factorisation} {Index notation} {Negative index as reciprocal}
{Zero exponent rule}

\end{exercise}

\begin{exercise}[Find Missing
Index]\protect\hypertarget{exr-u2-30}{}\label{exr-u2-30}

Find the missing index.

\begin{enumerate}
\def\labelenumi{\alph{enumi})}
\tightlist
\item
  \(t^{4} \times t^{\square} = t^{6}\)\\
\item
  \(t^{5} \times t^{\square} = t^{12}\)\\
\item
  \(y^{3} \times y^{\square} = y^{10}\)\\
\item
  \(b^{7} \times b^{\square} = b^{13}\)\\
\item
  \(w^{2} \times w^{\square} = w^{5}\)\\
\item
  \(n^{6} \times n^{\square} = n^{13}\)
\end{enumerate}

\begin{tcolorbox}[enhanced jigsaw, arc=.35mm, breakable, toptitle=1mm, title=\textcolor{quarto-callout-tip-color}{\faLightbulb}\hspace{0.5em}{🔍 View Solution}, colframe=quarto-callout-tip-color-frame, rightrule=.15mm, opacityback=0, titlerule=0mm, colback=white, coltitle=black, bottomtitle=1mm, toprule=.15mm, leftrule=.75mm, bottomrule=.15mm, colbacktitle=quarto-callout-tip-color!10!white, left=2mm, opacitybacktitle=0.6]

\begin{enumerate}
\def\labelenumi{\alph{enumi})}
\tightlist
\item
  \(2\)\\
\item
  \(7\)\\
\item
  \(7\)\\
\item
  \(6\)\\
\item
  \(3\)\\
\item
  \(7\)\\
\end{enumerate}

\end{tcolorbox}

{Reverse index law} {Subtraction of integers} {Algebraic reasoning}

\end{exercise}

\begin{exercise}[Find Missing Index ---
Division]\protect\hypertarget{exr-u2-31}{}\label{exr-u2-31}

Find the missing index.

\begin{enumerate}
\def\labelenumi{\alph{enumi})}
\tightlist
\item
  \(\dfrac{b^{11}}{b^{\square}} = b^{6}\)\\
\item
  \(\dfrac{x^{14}}{x^{\square}} = x^{9}\)\\
\item
  \(\dfrac{w^{7}}{w^{\square}} = w^{4}\)\\
\item
  \(\dfrac{p^{7}}{p^{\square}} = p^{4}\)\\
\item
  \(\dfrac{y^{9}}{y^{\square}} = y^{5}\)\\
\item
  \(\dfrac{y^{12}}{y^{\square}} = y^{9}\)
\end{enumerate}

\begin{tcolorbox}[enhanced jigsaw, arc=.35mm, breakable, toptitle=1mm, title=\textcolor{quarto-callout-tip-color}{\faLightbulb}\hspace{0.5em}{🔍 View Solution}, colframe=quarto-callout-tip-color-frame, rightrule=.15mm, opacityback=0, titlerule=0mm, colback=white, coltitle=black, bottomtitle=1mm, toprule=.15mm, leftrule=.75mm, bottomrule=.15mm, colbacktitle=quarto-callout-tip-color!10!white, left=2mm, opacitybacktitle=0.6]

\begin{enumerate}
\def\labelenumi{\alph{enumi})}
\tightlist
\item
  \(5\)\\
\item
  \(5\)\\
\item
  \(3\)\\
\item
  \(3\)\\
\item
  \(4\)\\
\item
  \(3\)\\
\end{enumerate}

\end{tcolorbox}

{Reverse index law} {Subtraction of integers} {Algebraic reasoning}

\end{exercise}

∑ × π ÷ √ ◆ ∞ ± ∫ ≠ Δ

\subsection{🔴 Challenging}

\begin{exercise}[Fractional Index ---
Non-unit]\protect\hypertarget{exr-u2-4}{}\label{exr-u2-4}

Evaluate.

\begin{enumerate}
\def\labelenumi{\alph{enumi})}
\tightlist
\item
  \(25^{\frac{3}{2}}\)\\
\item
  \(27^{\frac{2}{3}}\)\\
\item
  \(25^{\frac{3}{2}}\)\\
\item
  \(25^{\frac{3}{2}}\)\\
\item
  \(16^{\frac{3}{2}}\)\\
\item
  \(125^{\frac{2}{3}}\)
\end{enumerate}

\begin{tcolorbox}[enhanced jigsaw, arc=.35mm, breakable, toptitle=1mm, title=\textcolor{quarto-callout-tip-color}{\faLightbulb}\hspace{0.5em}{🔍 View Solution}, colframe=quarto-callout-tip-color-frame, rightrule=.15mm, opacityback=0, titlerule=0mm, colback=white, coltitle=black, bottomtitle=1mm, toprule=.15mm, leftrule=.75mm, bottomrule=.15mm, colbacktitle=quarto-callout-tip-color!10!white, left=2mm, opacitybacktitle=0.6]

\begin{enumerate}
\def\labelenumi{\alph{enumi})}
\tightlist
\item
  \(125\)\\
\item
  \(9\)\\
\item
  \(125\)\\
\item
  \(125\)\\
\item
  \(64\)\\
\item
  \(25\)\\
\end{enumerate}

\end{tcolorbox}

{Fractional index as root and power} {Nth root evaluation} {Integer
exponentiation} {Order of operations}

\end{exercise}

\begin{exercise}[Negative Outer
Power]\protect\hypertarget{exr-u2-13}{}\label{exr-u2-13}

Simplify. Give your answer with positive exponents.

\begin{enumerate}
\def\labelenumi{\alph{enumi})}
\tightlist
\item
  \((4w^{3})^{-1}\)\\
\item
  \((3n^{2})^{-3}\)\\
\item
  \((3w^{3})^{-2}\)\\
\item
  \((2b^{3})^{-3}\)\\
\item
  \((2t^{2})^{-3}\)\\
\item
  \((4m^{3})^{-1}\)
\end{enumerate}

\begin{tcolorbox}[enhanced jigsaw, arc=.35mm, breakable, toptitle=1mm, title=\textcolor{quarto-callout-tip-color}{\faLightbulb}\hspace{0.5em}{🔍 View Solution}, colframe=quarto-callout-tip-color-frame, rightrule=.15mm, opacityback=0, titlerule=0mm, colback=white, coltitle=black, bottomtitle=1mm, toprule=.15mm, leftrule=.75mm, bottomrule=.15mm, colbacktitle=quarto-callout-tip-color!10!white, left=2mm, opacitybacktitle=0.6]

\begin{enumerate}
\def\labelenumi{\alph{enumi})}
\tightlist
\item
  \(\dfrac{1}{4w^{3}}\)\\
\item
  \(\dfrac{1}{27n^{6}}\)\\
\item
  \(\dfrac{1}{9w^{6}}\)\\
\item
  \(\dfrac{1}{8b^{9}}\)\\
\item
  \(\dfrac{1}{8t^{6}}\)\\
\item
  \(\dfrac{1}{4m^{3}}\)\\
\end{enumerate}

\end{tcolorbox}

{Power of power law} {Negative index as reciprocal} {Coefficient
exponentiation}

\end{exercise}

\begin{exercise}[Negative Fractional
Index]\protect\hypertarget{exr-u2-19}{}\label{exr-u2-19}

Evaluate.

\begin{enumerate}
\def\labelenumi{\alph{enumi})}
\tightlist
\item
  \(32^{\frac{-4}{5}}\)\\
\item
  \(25^{\frac{-1}{2}}\)\\
\item
  \(81^{\frac{-3}{4}}\)\\
\item
  \(81^{\frac{-3}{4}}\)\\
\item
  \(64^{\frac{-1}{3}}\)\\
\item
  \(81^{\frac{-3}{4}}\)
\end{enumerate}

\begin{tcolorbox}[enhanced jigsaw, arc=.35mm, breakable, toptitle=1mm, title=\textcolor{quarto-callout-tip-color}{\faLightbulb}\hspace{0.5em}{🔍 View Solution}, colframe=quarto-callout-tip-color-frame, rightrule=.15mm, opacityback=0, titlerule=0mm, colback=white, coltitle=black, bottomtitle=1mm, toprule=.15mm, leftrule=.75mm, bottomrule=.15mm, colbacktitle=quarto-callout-tip-color!10!white, left=2mm, opacitybacktitle=0.6]

\begin{enumerate}
\def\labelenumi{\alph{enumi})}
\tightlist
\item
  \(1/16\)\\
\item
  \(1/5\)\\
\item
  \(1/27\)\\
\item
  \(1/27\)\\
\item
  \(1/4\)\\
\item
  \(1/27\)\\
\end{enumerate}

\end{tcolorbox}

{Fractional index as root and power} {Negative index as reciprocal} {Nth
root evaluation} {Combining index laws}

\end{exercise}

\begin{exercise}[Fractional Index of a
Fraction]\protect\hypertarget{exr-u2-20}{}\label{exr-u2-20}

Evaluate. Simplify your answer.

\begin{enumerate}
\def\labelenumi{\alph{enumi})}
\tightlist
\item
  \(\left(\dfrac{9}{25}\right)^{\frac{3}{2}}\)\\
\item
  \(\left(\dfrac{4}{9}\right)^{\frac{1}{2}}\)\\
\item
  \(\left(\dfrac{16}{25}\right)^{\frac{1}{2}}\)\\
\item
  \(\left(\dfrac{9}{25}\right)^{\frac{3}{2}}\)\\
\item
  \(\left(\dfrac{27}{1000}\right)^{\frac{2}{3}}\)\\
\item
  \(\left(\dfrac{4}{9}\right)^{\frac{1}{2}}\)
\end{enumerate}

\begin{tcolorbox}[enhanced jigsaw, arc=.35mm, breakable, toptitle=1mm, title=\textcolor{quarto-callout-tip-color}{\faLightbulb}\hspace{0.5em}{🔍 View Solution}, colframe=quarto-callout-tip-color-frame, rightrule=.15mm, opacityback=0, titlerule=0mm, colback=white, coltitle=black, bottomtitle=1mm, toprule=.15mm, leftrule=.75mm, bottomrule=.15mm, colbacktitle=quarto-callout-tip-color!10!white, left=2mm, opacitybacktitle=0.6]

\begin{enumerate}
\def\labelenumi{\alph{enumi})}
\tightlist
\item
  \(27/125\)\\
\item
  \(2/3\)\\
\item
  \(4/5\)\\
\item
  \(27/125\)\\
\item
  \(9/100\)\\
\item
  \(2/3\)\\
\end{enumerate}

\end{tcolorbox}

{Fractional index as root and power} {Fraction exponentiation} {Nth root
of fraction}

\end{exercise}

\begin{exercise}[Complex Fraction with
Coefficients]\protect\hypertarget{exr-u2-23}{}\label{exr-u2-23}

Simplify fully.

\begin{enumerate}
\def\labelenumi{\alph{enumi})}
\tightlist
\item
  \(\dfrac{3x^{2} \times 3x^{4}}{9x^{3}}\)\\
\item
  \(\dfrac{2p^{3} \times 4p^{2}}{8p^{4}}\)\\
\item
  \(\dfrac{3n^{5} \times 4n^{4}}{12n}\)\\
\item
  \(\dfrac{4w^{4} \times 3w^{3}}{3w^{4}}\)\\
\item
  \(\dfrac{3w^{2} \times 2w^{3}}{2w^{3}}\)\\
\item
  \(\dfrac{3c^{2} \times 3c^{2}}{3c}\)
\end{enumerate}

\begin{tcolorbox}[enhanced jigsaw, arc=.35mm, breakable, toptitle=1mm, title=\textcolor{quarto-callout-tip-color}{\faLightbulb}\hspace{0.5em}{🔍 View Solution}, colframe=quarto-callout-tip-color-frame, rightrule=.15mm, opacityback=0, titlerule=0mm, colback=white, coltitle=black, bottomtitle=1mm, toprule=.15mm, leftrule=.75mm, bottomrule=.15mm, colbacktitle=quarto-callout-tip-color!10!white, left=2mm, opacitybacktitle=0.6]

\begin{enumerate}
\def\labelenumi{\alph{enumi})}
\tightlist
\item
  \(x^{3}\)\\
\item
  \(p\)\\
\item
  \(n^{8}\)\\
\item
  \(4w^{3}\)\\
\item
  \(3w^{2}\)\\
\item
  \(3c^{3}\)\\
\end{enumerate}

\end{tcolorbox}

{Coefficient division} {Combining multiplication and division laws}
{Multi step simplification} {Separating coefficients from indices}

\end{exercise}

\begin{exercise}[Complex Fraction ---
Multi-variable]\protect\hypertarget{exr-u2-24}{}\label{exr-u2-24}

Simplify fully.

\begin{enumerate}
\def\labelenumi{\alph{enumi})}
\tightlist
\item
  \(\dfrac{8y^{6}p^{4}}{4y^{4}p}\)\\
\item
  \(\dfrac{8b^{6}k^{2}}{4b^{2}k}\)\\
\item
  \(\dfrac{16b^{5}p^{4}}{4b^{2}p^{2}}\)\\
\item
  \(\dfrac{6p^{8}y^{3}}{3p^{4}y}\)\\
\item
  \(\dfrac{6w^{6}n^{5}}{3w^{4}n^{3}}\)\\
\item
  \(\dfrac{4n^{6}m^{5}}{2n^{3}m^{2}}\)
\end{enumerate}

\begin{tcolorbox}[enhanced jigsaw, arc=.35mm, breakable, toptitle=1mm, title=\textcolor{quarto-callout-tip-color}{\faLightbulb}\hspace{0.5em}{🔍 View Solution}, colframe=quarto-callout-tip-color-frame, rightrule=.15mm, opacityback=0, titlerule=0mm, colback=white, coltitle=black, bottomtitle=1mm, toprule=.15mm, leftrule=.75mm, bottomrule=.15mm, colbacktitle=quarto-callout-tip-color!10!white, left=2mm, opacitybacktitle=0.6]

\begin{enumerate}
\def\labelenumi{\alph{enumi})}
\tightlist
\item
  \(2y^{2}p^{3}\)\\
\item
  \(2b^{4}k\)\\
\item
  \(4b^{3}p^{2}\)\\
\item
  \(2p^{4}y^{2}\)\\
\item
  \(2w^{2}n^{2}\)\\
\item
  \(2n^{3}m^{3}\)\\
\end{enumerate}

\end{tcolorbox}

{Multi variable expressions} {Coefficient division} {Multi step
simplification} {Grouping like terms}

\end{exercise}

\begin{exercise}[Nested Power in
Fraction]\protect\hypertarget{exr-u2-25}{}\label{exr-u2-25}

Simplify fully.

\begin{enumerate}
\def\labelenumi{\alph{enumi})}
\tightlist
\item
  \(\dfrac{c^{3} \times (c^{3})^{2}}{c^{4}}\)\\
\item
  \(\dfrac{b^{5} \times (b^{3})^{3}}{b^{9}}\)\\
\item
  \(\dfrac{k^{5} \times (k^{2})^{2}}{k^{7}}\)\\
\item
  \(\dfrac{y^{2} \times (y^{4})^{3}}{y^{8}}\)\\
\item
  \(\dfrac{a^{4} \times (a^{3})^{3}}{a^{7}}\)\\
\item
  \(\dfrac{p^{3} \times (p^{3})^{2}}{p^{5}}\)
\end{enumerate}

\begin{tcolorbox}[enhanced jigsaw, arc=.35mm, breakable, toptitle=1mm, title=\textcolor{quarto-callout-tip-color}{\faLightbulb}\hspace{0.5em}{🔍 View Solution}, colframe=quarto-callout-tip-color-frame, rightrule=.15mm, opacityback=0, titlerule=0mm, colback=white, coltitle=black, bottomtitle=1mm, toprule=.15mm, leftrule=.75mm, bottomrule=.15mm, colbacktitle=quarto-callout-tip-color!10!white, left=2mm, opacitybacktitle=0.6]

\begin{enumerate}
\def\labelenumi{\alph{enumi})}
\tightlist
\item
  \(c^{5}\)\\
\item
  \(b^{5}\)\\
\item
  \(k^{2}\)\\
\item
  \(y^{6}\)\\
\item
  \(a^{6}\)\\
\item
  \(p^{4}\)\\
\end{enumerate}

\end{tcolorbox}

{Power of power law} {Combining multiplication and division laws} {Multi
step simplification} {Order of operations}

\end{exercise}

\begin{exercise}[Numeric Index Laws ---
Mixed]\protect\hypertarget{exr-u2-28}{}\label{exr-u2-28}

Simplify. Give your answer as a single power.

\begin{enumerate}
\def\labelenumi{\alph{enumi})}
\tightlist
\item
  \(\dfrac{3^{4} \times 3^{-1}}{3^{1}}\)\\
\item
  \(\dfrac{3^{5} \times 3^{-1}}{3^{1}}\)\\
\item
  \(\dfrac{5^{4} \times 5^{-3}}{5^{2}}\)\\
\item
  \(\dfrac{2^{4} \times 2^{-1}}{2^{2}}\)\\
\item
  \(\dfrac{2^{5} \times 2^{-2}}{2^{1}}\)\\
\item
  \(\dfrac{5^{5} \times 5^{-2}}{5^{3}}\)
\end{enumerate}

\begin{tcolorbox}[enhanced jigsaw, arc=.35mm, breakable, toptitle=1mm, title=\textcolor{quarto-callout-tip-color}{\faLightbulb}\hspace{0.5em}{🔍 View Solution}, colframe=quarto-callout-tip-color-frame, rightrule=.15mm, opacityback=0, titlerule=0mm, colback=white, coltitle=black, bottomtitle=1mm, toprule=.15mm, leftrule=.75mm, bottomrule=.15mm, colbacktitle=quarto-callout-tip-color!10!white, left=2mm, opacitybacktitle=0.6]

\begin{enumerate}
\def\labelenumi{\alph{enumi})}
\tightlist
\item
  \(3^{2}\)\\
\item
  \(3^{3}\)\\
\item
  \(5^{-1}\)\\
\item
  \(2^{1}\)\\
\item
  \(2^{2}\)\\
\item
  \(5^{0}\)\\
\end{enumerate}

\end{tcolorbox}

{Numeric base recognition} {Combining multiplication and division laws}
{Negative integer addition} {Multi step simplification}

\end{exercise}

\begin{exercise}[Find Missing
Power]\protect\hypertarget{exr-u2-32}{}\label{exr-u2-32}

Find the missing power.

\begin{enumerate}
\def\labelenumi{\alph{enumi})}
\tightlist
\item
  \((k^{3})^{\square} = k^{9}\)\\
\item
  \((c^{4})^{\square} = c^{16}\)\\
\item
  \((p^{5})^{\square} = p^{15}\)\\
\item
  \((y^{4})^{\square} = y^{12}\)\\
\item
  \((m^{2})^{\square} = m^{6}\)\\
\item
  \((k^{3})^{\square} = k^{9}\)
\end{enumerate}

\begin{tcolorbox}[enhanced jigsaw, arc=.35mm, breakable, toptitle=1mm, title=\textcolor{quarto-callout-tip-color}{\faLightbulb}\hspace{0.5em}{🔍 View Solution}, colframe=quarto-callout-tip-color-frame, rightrule=.15mm, opacityback=0, titlerule=0mm, colback=white, coltitle=black, bottomtitle=1mm, toprule=.15mm, leftrule=.75mm, bottomrule=.15mm, colbacktitle=quarto-callout-tip-color!10!white, left=2mm, opacitybacktitle=0.6]

\begin{enumerate}
\def\labelenumi{\alph{enumi})}
\tightlist
\item
  \(3\)\\
\item
  \(4\)\\
\item
  \(3\)\\
\item
  \(3\)\\
\item
  \(3\)\\
\item
  \(3\)\\
\end{enumerate}

\end{tcolorbox}

{Reverse power of power law} {Division of integers} {Algebraic
reasoning}

\end{exercise}

\begin{exercise}[Identify and Correct
Error]\protect\hypertarget{exr-u2-33}{}\label{exr-u2-33}

Identify the mistake and write the correct answer.

\begin{enumerate}
\def\labelenumi{\alph{enumi})}
\tightlist
\item
  \(27^{\frac{1}{3}} = 9\text{ (claim: } \frac{1}{3} \text{ of } 27 = 9\text{). Correct?}\)\\
\item
  \(27^{\frac{1}{3}} = 9\text{ (claim: } \frac{1}{3} \text{ of } 27 = 9\text{). Correct?}\)\\
\item
  \(27^{\frac{1}{3}} = 9\text{ (claim: } \frac{1}{3} \text{ of } 27 = 9\text{). Correct?}\)\\
\item
  \(9^{-2} = -81\text{. Is this correct? Explain.}\)\\
\item
  \(9^{-2} = -81\text{. Is this correct? Explain.}\)\\
\item
  \(y^8 \times y^3 = y^{24}\text{. Explain the mistake.}\)
\end{enumerate}

\begin{tcolorbox}[enhanced jigsaw, arc=.35mm, breakable, toptitle=1mm, title=\textcolor{quarto-callout-tip-color}{\faLightbulb}\hspace{0.5em}{🔍 View Solution}, colframe=quarto-callout-tip-color-frame, rightrule=.15mm, opacityback=0, titlerule=0mm, colback=white, coltitle=black, bottomtitle=1mm, toprule=.15mm, leftrule=.75mm, bottomrule=.15mm, colbacktitle=quarto-callout-tip-color!10!white, left=2mm, opacitybacktitle=0.6]

\begin{enumerate}
\def\labelenumi{\alph{enumi})}
\tightlist
\item
  \(\text{No. } 27^{\frac{1}{3}} = \sqrt[3]{27} = 3\)\\
\item
  \(\text{No. } 27^{\frac{1}{3}} = \sqrt[3]{27} = 3\)\\
\item
  \(\text{No. } 27^{\frac{1}{3}} = \sqrt[3]{27} = 3\)\\
\item
  \(\text{No. } 9^{-2} = \dfrac{1}{9^2} = \dfrac{1}{81}\)\\
\item
  \(\text{No. } 9^{-2} = \dfrac{1}{9^2} = \dfrac{1}{81}\)\\
\item
  \(\text{Powers should be added, not multiplied. Correct: } y^{8+3} = y^{11}\)\\
\end{enumerate}

\end{tcolorbox}

{Mathematical reasoning} {Error analysis} {Verbal mathematical
explanation} {Index notation}

\end{exercise}

∑ × π ÷ √ ◆ ∞ ± ∫ ≠ Δ

\section{⚡ Test Your Knowledge}\label{sec-indices-worksheet}

\chapter{Unit 3: Mensuration}\label{unit-3-mensuration}

\section{🎯 Learning Targets}\label{learning-targets-2}

\subsection{🔵 Volume}\label{volume}

1

\textbf{Volume --- Prisms}\\
Calculate the volume of right prisms using\\
\(V = A_{\text{base}} \times l\)

Includes rectangular, triangular, trapezoidal, and parallelogram
cross-sections.

2

\textbf{Volume --- Cylinders}\\
Calculate the volume of a full cylinder using\\
\(V = \pi r^2 h\)

Use either radius or diameter.

3

\textbf{Volume --- Portions}\\
Find volumes of partial cylinders:\\
- Half: \(\tfrac{1}{2}\pi r^2 h\)\\
- Quarter\\
- Sector: \(\tfrac{\theta}{360}\pi r^2 h\)

4

\textbf{Volume --- Reverse}\\
Rearrange volume formulas to find unknown dimensions such as length,
height, or radius.

5

\textbf{Volume --- Context}\\
Solve real-world problems involving volume, including converting m³ to
litres.

6

\textbf{Composite Volume}\\
Break complex solids into simpler shapes and add or subtract their
volumes.

\begin{center}\rule{0.5\linewidth}{0.5pt}\end{center}

\subsection{🟢 Surface Area}\label{surface-area}

7

\textbf{SA --- Prisms}\\
Calculate surface area of cuboids, triangular and trapezoidal prisms,
and open boxes.

8

\textbf{SA --- Cylinders}\\
Calculate surface area of: - Closed cylinders\\
- Open-top cylinders\\
- Tubes and half-cylinders

9

\textbf{SA --- Composite}\\
Find surface area of composite solids, including added or removed faces.

\begin{center}\rule{0.5\linewidth}{0.5pt}\end{center}

\subsection{🟣 Advanced \& Algebra}\label{advanced-algebra}

10

\textbf{Unit Conversions}\\
Convert between length, area, and volume units using correct powers of
conversion.

11

\textbf{Pyramids, Cones \& Spheres}\\
- Pyramid: \(V=\tfrac{1}{3}Ah\)\\
- Cone: \(V=\tfrac{1}{3}\pi r^2h\), \(SA=\pi r^2+\pi rl\)\\
- Sphere: \(V=\tfrac{4}{3}\pi r^3\), \(SA=4\pi r^2\)

Includes composite solids.

12

\textbf{Algebraic Volume}\\
Work with algebraic dimensions:\\
- Expand and simplify expressions\\
- Solve for \(x\) given volume

\begin{center}\rule{0.5\linewidth}{0.5pt}\end{center}

\section{🗂️ Formula Cards}\label{formula-cards-1}

\subsection*{Unit 3: Mensuration --- Formula
Reference}\label{unit-3-mensuration-formula-reference}
\addcontentsline{toc}{subsection}{Unit 3: Mensuration --- Formula
Reference}

\begin{center}\rule{0.5\linewidth}{0.5pt}\end{center}

\subsubsection*{📏 Unit Conversions}\label{unit-conversions}
\addcontentsline{toc}{subsubsection}{📏 Unit Conversions}

\textbf{💡 Key Rule} When converting units:\\
• \textbf{Length} → multiply or divide by the factor\\
• \textbf{Area} → use the factor \textbf{squared}\\
• \textbf{Volume} → use the factor \textbf{cubed}

\textbf{Length chain:}
\(\text{km} \xrightarrow{\times 1000} \text{m} \xrightarrow{\times 100} \text{cm} \xrightarrow{\times 10} \text{mm}\)

\textbf{Area chain:}
\(\text{km}^2 \xrightarrow{\times 10^6} \text{m}^2 \xrightarrow{\times 10^4} \text{cm}^2 \xrightarrow{\times 100} \text{mm}^2\)

\textbf{Volume chain:}
\(\text{km}^3 \xrightarrow{\times 10^9} \text{m}^3 \xrightarrow{\times 10^6} \text{cm}^3 \xrightarrow{\times 10^3} \text{mm}^3\)

\textbf{Capacity:} \(1\text{ m}^3 = 1000\text{ L}\) \(\qquad\)
\(1\text{ cm}^3 = 1\text{ mL}\)

\textbf{⚠️ Common Error} \(1\text{ m}^2 = 100\text{ cm}^2\) ✗ --- should
be \textbf{10 000} cm².\\
\(1\text{ m}^3 = 100\text{ cm}^3\) ✗ --- should be \textbf{1 000 000}
cm³.

\begin{center}\rule{0.5\linewidth}{0.5pt}\end{center}

\subsubsection*{📦 Volume Formulae}\label{volume-formulae}
\addcontentsline{toc}{subsubsection}{📦 Volume Formulae}

\textbf{🎯 Core Idea} For ANY prism or cylinder:\\
\[V = A_{\text{cross-section}} \times \text{length}\]

\begin{longtable}[]{@{}
  >{\raggedright\arraybackslash}p{(\linewidth - 4\tabcolsep) * \real{0.2414}}
  >{\raggedright\arraybackslash}p{(\linewidth - 4\tabcolsep) * \real{0.4483}}
  >{\raggedright\arraybackslash}p{(\linewidth - 4\tabcolsep) * \real{0.3103}}@{}}
\toprule\noalign{}
\begin{minipage}[b]{\linewidth}\raggedright
Solid
\end{minipage} & \begin{minipage}[b]{\linewidth}\raggedright
Cross-section
\end{minipage} & \begin{minipage}[b]{\linewidth}\raggedright
Formula
\end{minipage} \\
\midrule\noalign{}
\endhead
\bottomrule\noalign{}
\endlastfoot
Rectangular prism & Rectangle & \(V = l \times w \times h\) \\
Triangular prism & Triangle & \(V = \tfrac{1}{2}\,b\,h\,l\) \\
Trapezoidal prism & Trapezoid & \(V = \tfrac{1}{2}(a+b)\,h\,l\) \\
Cylinder & Circle & \(V = \pi r^2 h\) \\
Half-cylinder & Semicircle & \(V = \tfrac{1}{2}\pi r^2 h\) \\
Quarter-cylinder & Quarter circle & \(V = \tfrac{1}{4}\pi r^2 h\) \\
Cylindrical sector & Sector (\(\theta°\)) &
\(V = \dfrac{\theta}{360}\,\pi r^2 h\) \\
\end{longtable}

\textbf{Composite solids:} \(V_{\text{total}} = V_1 + V_2\) (add) or
\(V_1 - V_2\) (subtract for holes)

\textbf{⚠️ Common Errors}\\
- Using diameter instead of radius: \(r = d \div 2\) first!\\
- Missing the \(\tfrac{1}{2}\) for triangular cross-sections.\\
- Adding volumes when you should subtract (hollow solids).

\begin{center}\rule{0.5\linewidth}{0.5pt}\end{center}

\subsubsection*{🎨 Surface Area Formulae}\label{surface-area-formulae}
\addcontentsline{toc}{subsubsection}{🎨 Surface Area Formulae}

\textbf{📝 Strategy} Surface Area = \textbf{sum of areas of ALL outer
faces}.\\
Sketch the solid, label every face, calculate each, then add.

\textbf{Cuboid (all 6 faces):} \[SA = 2(lw + lh + wh)\]

\textbf{Open-top box (5 faces --- no lid):} \[SA = lw + 2lh + 2wh\]

\textbf{Triangular prism (5 faces):}
\[SA = 2 \times A_{\triangle} + (b + s_1 + s_2) \times l\]

\textbf{Trapezoidal prism (6 faces):}
\[SA = 2 \times \tfrac{1}{2}(a+b)h + (a + b + s_1 + s_2) \times l\]

\textbf{Closed cylinder (curved + 2 circles):}
\[SA = 2\pi rh + 2\pi r^2\]

\textbf{Open-top cylinder (curved + 1 circle):}
\[SA = 2\pi rh + \pi r^2\]

\textbf{Tube/pipe (curved only, both ends open):} \[SA = 2\pi rh\]

\textbf{Half-cylinder (3 regions):}
\[SA = \underbrace{\pi r h}_{\text{half-curved}} + \underbrace{2rh}_{\text{flat rect}} + \underbrace{\pi r^2}_{\text{2 semicircles}}\]

\textbf{💡 Cylinder Unrolled} The curved surface of a cylinder unrolls
into a rectangle:\\
width \(= 2\pi r\) (the circumference), height \(= h\).\\
So curved \(SA = 2\pi r \times h = 2\pi rh\).

\textbf{⚠️ Common Errors --- Surface Area}\\
- Using the full cylinder SA formula when the cylinder is open at one or
both ends.\\
- Forgetting to include \textbf{both} triangular ends of a triangular
prism.\\
- Counting an internal join face that is \textbf{not} on the outside
surface.\\
- Confusing SA (cm²) with Volume (cm³).

\begin{center}\rule{0.5\linewidth}{0.5pt}\end{center}

\subsubsection*{🔺 Pyramids, Cones and
Spheres}\label{pyramids-cones-and-spheres}
\addcontentsline{toc}{subsubsection}{🔺 Pyramids, Cones and Spheres}

\textbf{🎯 Key Pattern} Pyramids and cones are exactly \textbf{⅓} the
volume of the prism or cylinder with the same base and height.

\begin{longtable}[]{@{}
  >{\raggedright\arraybackslash}p{(\linewidth - 4\tabcolsep) * \real{0.2500}}
  >{\raggedright\arraybackslash}p{(\linewidth - 4\tabcolsep) * \real{0.2857}}
  >{\raggedright\arraybackslash}p{(\linewidth - 4\tabcolsep) * \real{0.4643}}@{}}
\toprule\noalign{}
\begin{minipage}[b]{\linewidth}\raggedright
Solid
\end{minipage} & \begin{minipage}[b]{\linewidth}\raggedright
Volume
\end{minipage} & \begin{minipage}[b]{\linewidth}\raggedright
Surface Area
\end{minipage} \\
\midrule\noalign{}
\endhead
\bottomrule\noalign{}
\endlastfoot
Pyramid (any base) & \(V = \dfrac{1}{3} A_{\text{base}} \times h\) & sum
of all faces \\
Rectangular pyramid & \(V = \dfrac{1}{3} lwh\) & \\
Square pyramid & \(V = \dfrac{1}{3} s^2 h\) & \\
Cone & \(V = \dfrac{1}{3}\pi r^2 h\) & \(SA = \pi r^2 + \pi r l\) \\
Sphere & \(V = \dfrac{4}{3}\pi r^3\) & \(SA = 4\pi r^2\) \\
Hemisphere & \(V = \dfrac{2}{3}\pi r^3\) & \(SA = 3\pi r^2\) \\
\end{longtable}

\textbf{Cone slant height:} \(l = \sqrt{r^2 + h^2}\) (Pythagoras ---
slant is the hypotenuse)\\
\textbf{Remember:} \(r = d \div 2\) always before substituting.

\textbf{⚠️ Common Errors --- Curved Solids}\\
- Forgetting the \(\tfrac{1}{3}\) in pyramid and cone formulas.\\
- Squaring \(r\) when you should cube it (sphere formula uses
\(r^3\)).\\
- Using \(\tfrac{4}{3}\) instead of \(\tfrac{2}{3}\) for hemisphere
volume.\\
- Using the perpendicular height \(h\) instead of slant \(l\) in cone
SA.

\begin{center}\rule{0.5\linewidth}{0.5pt}\end{center}

\subsubsection*{🔣 Algebraic Volume}\label{algebraic-volume}
\addcontentsline{toc}{subsubsection}{🔣 Algebraic Volume}

\textbf{📝 Method} Substitute the algebraic expressions into the volume
formula, then expand brackets and collect like terms using index laws.

\textbf{Steps:}

\begin{enumerate}
\def\labelenumi{\arabic{enumi}.}
\tightlist
\item
  Write the formula:
  e.g.~\(V = \tfrac{1}{2} \times \text{base} \times h \times l\)
\item
  Substitute each expression: e.g.~base \(= 2x\), \(h = x\),
  \(l = (x+4)\)
\item
  Multiply step by step:
  \(V = \tfrac{1}{2} \times 2x \times x \times (x+4) = x^2(x+4)\)
\item
  Expand: \(V = x^3 + 4x^2\)
\end{enumerate}

\textbf{Index law reminder:} \((ax)^2 = a^2x^2\) --- square both the
coefficient AND the variable.

\textbf{Example --- solve for x:} If \(V = 3x^3 = 375\), then
\(x^3 = 125\), so \(x = \sqrt[3]{125} = 5\).

\textbf{⚠️ Common Errors --- Algebraic Volume}\\
- \((2x)^2 = 4x^2\) ✓ not \(2x^2\) ✗\\
- Stopping before expanding all brackets.\\
- Forgetting \(V =\) on the left-hand side of the formula.\\
- Square rooting when the equation requires a \textbf{cube root}.

\begin{center}\rule{0.5\linewidth}{0.5pt}\end{center}

\section{📏 Part A --- Unit Conversions}\label{sec-conversions}

\begin{tcolorbox}[enhanced jigsaw, arc=.35mm, breakable, toptitle=1mm, title={🔑 Key Conversion Facts}, colframe=quarto-callout-note-color-frame, rightrule=.15mm, opacityback=0, titlerule=0mm, colback=white, coltitle=black, bottomtitle=1mm, toprule=.15mm, leftrule=.75mm, bottomrule=.15mm, colbacktitle=quarto-callout-note-color!10!white, left=2mm, opacitybacktitle=0.6]

\begin{longtable}[]{@{}
  >{\raggedright\arraybackslash}p{(\linewidth - 2\tabcolsep) * \real{0.5000}}
  >{\raggedright\arraybackslash}p{(\linewidth - 2\tabcolsep) * \real{0.5000}}@{}}
\toprule\noalign{}
\begin{minipage}[b]{\linewidth}\raggedright
Dimension
\end{minipage} & \begin{minipage}[b]{\linewidth}\raggedright
Conversion
\end{minipage} \\
\midrule\noalign{}
\endhead
\bottomrule\noalign{}
\endlastfoot
\textbf{Length} & \(\times10\): cm→mm · \(\times100\): m→cm ·
\(\times1000\): km→m \\
\textbf{Area} & \(\times100\): cm²→mm² · \(\times10{,}000\): m²→cm² ·
\(\times10^6\): km²→m² \\
\textbf{Volume} & \(\times1000\): cm³→mm³ · \(\times10^6\): m³→cm³ ·
\(1\text{ m}^3=1000\text{ L}\) \\
\end{longtable}

\emph{To go the other way: divide by the same factor.}

\end{tcolorbox}

\subsection{🟢 Easy}

\begin{exercise}[Converting Units of
Length]\protect\hypertarget{exr-conv-1}{}\label{exr-conv-1}

Convert the measurement.

\begin{enumerate}
\def\labelenumi{\alph{enumi})}
\item
  \(5.1\,\text{cm} = \square\,\text{mm}\)
\item
  \(3.2\,\text{cm} = \square\,\text{mm}\)
\item
  \(4373\,\text{mm} = \square\,\text{m}\)
\item
  \(6.7\,\text{km} = \square\,\text{m}\)
\item
  \(1.3\,\text{m} = \square\,\text{cm}\)
\item
  \(2342\,\text{mm} = \square\,\text{m}\)
\end{enumerate}

\begin{tcolorbox}[enhanced jigsaw, arc=.35mm, breakable, toptitle=1mm, title=\textcolor{quarto-callout-tip-color}{\faLightbulb}\hspace{0.5em}{🔍 View Solution}, colframe=quarto-callout-tip-color-frame, rightrule=.15mm, opacityback=0, titlerule=0mm, colback=white, coltitle=black, bottomtitle=1mm, toprule=.15mm, leftrule=.75mm, bottomrule=.15mm, colbacktitle=quarto-callout-tip-color!10!white, left=2mm, opacitybacktitle=0.6]

\begin{enumerate}
\def\labelenumi{\alph{enumi})}
\item
  \(51\,\text{mm}\)
\item
  \(32\,\text{mm}\)
\item
  \(4.373\,\text{m}\)
\item
  \(6700\,\text{m}\)
\item
  \(130\,\text{cm}\)
\item
  \(2.342\,\text{m}\)
\end{enumerate}

\end{tcolorbox}

{Unit conversion length} {Multiply by powers of 10}

\end{exercise}

∑ × π ÷ √ ◆ ∞ ± ∫ ≠ Δ

\subsection{🟡 Medium}

\begin{exercise}[Converting Units of
Area]\protect\hypertarget{exr-conv-2}{}\label{exr-conv-2}

Convert the measurement.

\begin{enumerate}
\def\labelenumi{\alph{enumi})}
\item
  \(3.6\,\text{cm}^{2} = \square\,\text{mm}^{2}\)
\item
  \(4029618\,\text{m}^{2} = \square\,\text{km}^{2}\)
\item
  \(5.0\,\text{cm}^{2} = \square\,\text{mm}^{2}\)
\item
  \(717\,\text{mm}^{2} = \square\,\text{cm}^{2}\)
\item
  \(4.5\,\text{m}^{2} = \square\,\text{cm}^{2}\)
\item
  \(7.4\,\text{m}^{2} = \square\,\text{cm}^{2}\)
\end{enumerate}

\begin{tcolorbox}[enhanced jigsaw, arc=.35mm, breakable, toptitle=1mm, title=\textcolor{quarto-callout-tip-color}{\faLightbulb}\hspace{0.5em}{🔍 View Solution}, colframe=quarto-callout-tip-color-frame, rightrule=.15mm, opacityback=0, titlerule=0mm, colback=white, coltitle=black, bottomtitle=1mm, toprule=.15mm, leftrule=.75mm, bottomrule=.15mm, colbacktitle=quarto-callout-tip-color!10!white, left=2mm, opacitybacktitle=0.6]

\begin{enumerate}
\def\labelenumi{\alph{enumi})}
\item
  \(360\,\text{mm}^{2}\)
\item
  \(4.029618\,\text{km}^{2}\)
\item
  \(500\,\text{mm}^{2}\)
\item
  \(7.17\,\text{cm}^{2}\)
\item
  \(45{,}000\,\text{cm}^{2}\)
\item
  \(74{,}000\,\text{cm}^{2}\)
\end{enumerate}

\end{tcolorbox}

{Unit conversion area} {Squaring conversion factors}

\end{exercise}

\begin{exercise}[Converting Units of
Volume]\protect\hypertarget{exr-conv-3}{}\label{exr-conv-3}

Convert the measurement.

\begin{enumerate}
\def\labelenumi{\alph{enumi})}
\item
  \(3.1\,\text{m}^{3} = \square\,\text{cm}^{3}\)
\item
  \(3483\,\text{L} = \square\,\text{m}^{3}\)
\item
  \(3970\,\text{L} = \square\,\text{m}^{3}\)
\item
  \(4743\,\text{mm}^{3} = \square\,\text{cm}^{3}\)
\item
  \(295\,\text{cm}^{3} = \square\,\text{mL}\)
\item
  \(6144958\,\text{cm}^{3} = \square\,\text{m}^{3}\)
\end{enumerate}

\begin{tcolorbox}[enhanced jigsaw, arc=.35mm, breakable, toptitle=1mm, title=\textcolor{quarto-callout-tip-color}{\faLightbulb}\hspace{0.5em}{🔍 View Solution}, colframe=quarto-callout-tip-color-frame, rightrule=.15mm, opacityback=0, titlerule=0mm, colback=white, coltitle=black, bottomtitle=1mm, toprule=.15mm, leftrule=.75mm, bottomrule=.15mm, colbacktitle=quarto-callout-tip-color!10!white, left=2mm, opacitybacktitle=0.6]

\begin{enumerate}
\def\labelenumi{\alph{enumi})}
\item
  \(3{,}100{,}000\,\text{cm}^{3}\)
\item
  \(3.483\,\text{m}^{3}\)
\item
  \(3.97\,\text{m}^{3}\)
\item
  \(4.743\,\text{cm}^{3}\)
\item
  \(295\,\text{mL}\)
\item
  \(6.144958\,\text{m}^{3}\)
\end{enumerate}

\end{tcolorbox}

{Unit conversion volume} {Cubing conversion factors}

\end{exercise}

\begin{exercise}[Mixed Unit
Conversions]\protect\hypertarget{exr-conv-4}{}\label{exr-conv-4}

Convert the measurement.

\begin{enumerate}
\def\labelenumi{\alph{enumi})}
\item
  \(6.4\,\text{cm}^{2} = \square\,\text{mm}^{2}\)
\item
  \(2.7\,\text{m} = \square\,\text{cm}\)
\item
  \(3.3\,\text{km}^{2} = \square\,\text{m}^{2}\)
\item
  \(7461\,\text{m} = \square\,\text{km}\)
\item
  \(4657692\,\text{cm}^{3} = \square\,\text{m}^{3}\)
\item
  \(7.0\,\text{cm}^{2} = \square\,\text{mm}^{2}\)
\end{enumerate}

\begin{tcolorbox}[enhanced jigsaw, arc=.35mm, breakable, toptitle=1mm, title=\textcolor{quarto-callout-tip-color}{\faLightbulb}\hspace{0.5em}{🔍 View Solution}, colframe=quarto-callout-tip-color-frame, rightrule=.15mm, opacityback=0, titlerule=0mm, colback=white, coltitle=black, bottomtitle=1mm, toprule=.15mm, leftrule=.75mm, bottomrule=.15mm, colbacktitle=quarto-callout-tip-color!10!white, left=2mm, opacitybacktitle=0.6]

\begin{enumerate}
\def\labelenumi{\alph{enumi})}
\item
  \(640\,\text{mm}^{2}\)
\item
  \(270\,\text{cm}\)
\item
  \(3{,}300{,}000\,\text{m}^{2}\)
\item
  \(7.461\,\text{km}\)
\item
  \(4.657692\,\text{m}^{3}\)
\item
  \(700\,\text{mm}^{2}\)
\end{enumerate}

\end{tcolorbox}

{Unit conversion length} {Unit conversion area} {Unit conversion volume}

\end{exercise}

∑ × π ÷ √ ◆ ∞ ± ∫ ≠ Δ

\begin{center}\rule{0.5\linewidth}{0.5pt}\end{center}

\section{📦 Part B --- Volume of Prisms and Cylinders}\label{sec-volume}

\subsection{🟢 Easy}

\begin{exercise}[Volume of Rectangular
Prism]\protect\hypertarget{exr-vol-1}{}\label{exr-vol-1}

Find the volume of each solid.

\begin{enumerate}
\def\labelenumi{\alph{enumi})}
\item
  \(\text{Cuboid: }l=14\text{ cm},\ w=4\text{ cm},\ h=4\text{ cm}\)
\item
  \(\text{Cuboid: }l=13\text{ cm},\ w=3\text{ cm},\ h=8\text{ cm}\)
\item
  \(\text{Cuboid: }l=8\text{ cm},\ w=4\text{ cm},\ h=2\text{ cm}\)
\item
  \(\text{Cuboid: }l=9\text{ cm},\ w=8\text{ cm},\ h=2\text{ cm}\)
\item
  \(\text{Cuboid: }l=4\text{ cm},\ w=4\text{ cm},\ h=4\text{ cm}\)
\item
  \(\text{Cuboid: }l=10\text{ cm},\ w=9\text{ cm},\ h=5\text{ cm}\)
\end{enumerate}

\begin{tcolorbox}[enhanced jigsaw, arc=.35mm, breakable, toptitle=1mm, title=\textcolor{quarto-callout-tip-color}{\faLightbulb}\hspace{0.5em}{🔍 View Solution}, colframe=quarto-callout-tip-color-frame, rightrule=.15mm, opacityback=0, titlerule=0mm, colback=white, coltitle=black, bottomtitle=1mm, toprule=.15mm, leftrule=.75mm, bottomrule=.15mm, colbacktitle=quarto-callout-tip-color!10!white, left=2mm, opacitybacktitle=0.6]

\begin{enumerate}
\def\labelenumi{\alph{enumi})}
\item
  \(224\text{ cm}^3\)
\item
  \(312\text{ cm}^3\)
\item
  \(64\text{ cm}^3\)
\item
  \(144\text{ cm}^3\)
\item
  \(64\text{ cm}^3\)
\item
  \(450\text{ cm}^3\)
\end{enumerate}

\end{tcolorbox}

{Volume formula prism} {Multiplication of integers} {Unit recognition}

\end{exercise}

\begin{exercise}[Volume from Cross-Section
Area]\protect\hypertarget{exr-vol-2}{}\label{exr-vol-2}

Find the volume. The cross-section area is given.

\begin{enumerate}
\def\labelenumi{\alph{enumi})}
\item
  \(\text{Prism: }A=21\,\text{cm}^{2},\ l=6\text{ cm}\)
\item
  \(\text{Prism: }A=28\,\text{cm}^{2},\ l=13\text{ cm}\)
\item
  \(\text{Prism: }A=29\,\text{cm}^{2},\ l=8\text{ cm}\)
\item
  \(\text{Prism: }A=35\,\text{cm}^{2},\ l=4\text{ cm}\)
\item
  \(\text{Prism: }A=48\,\text{cm}^{2},\ l=11\text{ cm}\)
\item
  \(\text{Prism: }A=39\,\text{cm}^{2},\ l=9\text{ cm}\)
\end{enumerate}

\begin{tcolorbox}[enhanced jigsaw, arc=.35mm, breakable, toptitle=1mm, title=\textcolor{quarto-callout-tip-color}{\faLightbulb}\hspace{0.5em}{🔍 View Solution}, colframe=quarto-callout-tip-color-frame, rightrule=.15mm, opacityback=0, titlerule=0mm, colback=white, coltitle=black, bottomtitle=1mm, toprule=.15mm, leftrule=.75mm, bottomrule=.15mm, colbacktitle=quarto-callout-tip-color!10!white, left=2mm, opacitybacktitle=0.6]

\begin{enumerate}
\def\labelenumi{\alph{enumi})}
\item
  \(126\,\text{cm}^{3}\)
\item
  \(364\,\text{cm}^{3}\)
\item
  \(232\,\text{cm}^{3}\)
\item
  \(140\,\text{cm}^{3}\)
\item
  \(528\,\text{cm}^{3}\)
\item
  \(351\,\text{cm}^{3}\)
\end{enumerate}

\end{tcolorbox}

{Volume formula prism} {Area identification}

\end{exercise}

\begin{exercise}[Volume of Triangular
Prism]\protect\hypertarget{exr-vol-3}{}\label{exr-vol-3}

Find the volume of each triangular prism.

\begin{enumerate}
\def\labelenumi{\alph{enumi})}
\item
  \(\text{Triangular prism: }b=4,\ h=4,\ l=14\text{ cm}\)
\item
  \(\text{Triangular prism: }b=6,\ h=3,\ l=8\text{ cm}\)
\item
  \(\text{Triangular prism: }b=4,\ h=3,\ l=8\text{ cm}\)
\item
  \(\text{Triangular prism: }b=7,\ h=8,\ l=9\text{ cm}\)
\item
  \(\text{Triangular prism: }b=7,\ h=8,\ l=9\text{ cm}\)
\item
  \(\text{Triangular prism: }b=5,\ h=7,\ l=6\text{ cm}\)
\end{enumerate}

\begin{tcolorbox}[enhanced jigsaw, arc=.35mm, breakable, toptitle=1mm, title=\textcolor{quarto-callout-tip-color}{\faLightbulb}\hspace{0.5em}{🔍 View Solution}, colframe=quarto-callout-tip-color-frame, rightrule=.15mm, opacityback=0, titlerule=0mm, colback=white, coltitle=black, bottomtitle=1mm, toprule=.15mm, leftrule=.75mm, bottomrule=.15mm, colbacktitle=quarto-callout-tip-color!10!white, left=2mm, opacitybacktitle=0.6]

\begin{enumerate}
\def\labelenumi{\alph{enumi})}
\item
  \(112.0\,\text{cm}^{3}\)
\item
  \(72.0\,\text{cm}^{3}\)
\item
  \(48.0\,\text{cm}^{3}\)
\item
  \(252.0\,\text{cm}^{3}\)
\item
  \(252.0\,\text{cm}^{3}\)
\item
  \(105.0\,\text{cm}^{3}\)
\end{enumerate}

\end{tcolorbox}

{Area triangle} {Volume formula prism} {Multi step calculation}

\end{exercise}

\begin{exercise}[Volume of Cylinder
(radius)]\protect\hypertarget{exr-vol-6}{}\label{exr-vol-6}

Find the volume. Round to 1 decimal place.

\begin{enumerate}
\def\labelenumi{\alph{enumi})}
\item
  \(\text{Cylinder: }r=3\text{ cm},\ h=4\text{ cm}\)
\item
  \(\text{Cylinder: }r=7\text{ cm},\ h=7\text{ cm}\)
\item
  \(\text{Cylinder: }r=3\text{ cm},\ h=13\text{ cm}\)
\item
  \(\text{Cylinder: }r=6\text{ cm},\ h=13\text{ cm}\)
\item
  \(\text{Cylinder: }r=2\text{ cm},\ h=15\text{ cm}\)
\item
  \(\text{Cylinder: }r=2\text{ cm},\ h=6\text{ cm}\)
\end{enumerate}

\begin{tcolorbox}[enhanced jigsaw, arc=.35mm, breakable, toptitle=1mm, title=\textcolor{quarto-callout-tip-color}{\faLightbulb}\hspace{0.5em}{🔍 View Solution}, colframe=quarto-callout-tip-color-frame, rightrule=.15mm, opacityback=0, titlerule=0mm, colback=white, coltitle=black, bottomtitle=1mm, toprule=.15mm, leftrule=.75mm, bottomrule=.15mm, colbacktitle=quarto-callout-tip-color!10!white, left=2mm, opacitybacktitle=0.6]

\begin{enumerate}
\def\labelenumi{\alph{enumi})}
\item
  \(113.1\,\text{cm}^{3}\)
\item
  \(1077.6\,\text{cm}^{3}\)
\item
  \(367.6\,\text{cm}^{3}\)
\item
  \(1470.3\,\text{cm}^{3}\)
\item
  \(188.5\,\text{cm}^{3}\)
\item
  \(75.4\,\text{cm}^{3}\)
\end{enumerate}

\end{tcolorbox}

{Cylinder formula} {Squaring numbers} {Use of pi}

\end{exercise}

∑ × π ÷ √ ◆ ∞ ± ∫ ≠ Δ

\subsection{🟡 Medium}

\begin{exercise}[Volume of Trapezoidal
Prism]\protect\hypertarget{exr-vol-4}{}\label{exr-vol-4}

Find the volume of each trapezoidal prism.

\begin{enumerate}
\def\labelenumi{\alph{enumi})}
\item
  \(\text{Trapezoidal prism: }a=2,\ b=6,\ h=3,\ l=8\text{ cm}\)
\item
  \(\text{Trapezoidal prism: }a=2,\ b=4,\ h=4,\ l=10\text{ cm}\)
\item
  \(\text{Trapezoidal prism: }a=5,\ b=8,\ h=4,\ l=7\text{ cm}\)
\item
  \(\text{Trapezoidal prism: }a=3,\ b=8,\ h=2,\ l=11\text{ cm}\)
\item
  \(\text{Trapezoidal prism: }a=4,\ b=7,\ h=2,\ l=8\text{ cm}\)
\item
  \(\text{Trapezoidal prism: }a=5,\ b=8,\ h=3,\ l=9\text{ cm}\)
\end{enumerate}

\begin{tcolorbox}[enhanced jigsaw, arc=.35mm, breakable, toptitle=1mm, title=\textcolor{quarto-callout-tip-color}{\faLightbulb}\hspace{0.5em}{🔍 View Solution}, colframe=quarto-callout-tip-color-frame, rightrule=.15mm, opacityback=0, titlerule=0mm, colback=white, coltitle=black, bottomtitle=1mm, toprule=.15mm, leftrule=.75mm, bottomrule=.15mm, colbacktitle=quarto-callout-tip-color!10!white, left=2mm, opacitybacktitle=0.6]

\begin{enumerate}
\def\labelenumi{\alph{enumi})}
\item
  \(96.0\,\text{cm}^{3}\)
\item
  \(120.0\,\text{cm}^{3}\)
\item
  \(182.0\,\text{cm}^{3}\)
\item
  \(121.0\,\text{cm}^{3}\)
\item
  \(88.0\,\text{cm}^{3}\)
\item
  \(175.5\,\text{cm}^{3}\)
\end{enumerate}

\end{tcolorbox}

{Area trapezoid} {Volume formula prism} {Multi step calculation}

\end{exercise}

\begin{exercise}[Volume of Parallelogram
Prism]\protect\hypertarget{exr-vol-5}{}\label{exr-vol-5}

Find the volume of each parallelogram prism.

\begin{enumerate}
\def\labelenumi{\alph{enumi})}
\item
  \(\text{Parallelogram prism: }b=4,\ h=8,\ l=10\text{ cm}\)
\item
  \(\text{Parallelogram prism: }b=9,\ h=8,\ l=11\text{ cm}\)
\item
  \(\text{Parallelogram prism: }b=7,\ h=7,\ l=5\text{ cm}\)
\item
  \(\text{Parallelogram prism: }b=10,\ h=8,\ l=7\text{ cm}\)
\item
  \(\text{Parallelogram prism: }b=4,\ h=6,\ l=13\text{ cm}\)
\item
  \(\text{Parallelogram prism: }b=6,\ h=3,\ l=6\text{ cm}\)
\end{enumerate}

\begin{tcolorbox}[enhanced jigsaw, arc=.35mm, breakable, toptitle=1mm, title=\textcolor{quarto-callout-tip-color}{\faLightbulb}\hspace{0.5em}{🔍 View Solution}, colframe=quarto-callout-tip-color-frame, rightrule=.15mm, opacityback=0, titlerule=0mm, colback=white, coltitle=black, bottomtitle=1mm, toprule=.15mm, leftrule=.75mm, bottomrule=.15mm, colbacktitle=quarto-callout-tip-color!10!white, left=2mm, opacitybacktitle=0.6]

\begin{enumerate}
\def\labelenumi{\alph{enumi})}
\item
  \(320\,\text{cm}^{3}\)
\item
  \(792\,\text{cm}^{3}\)
\item
  \(245\,\text{cm}^{3}\)
\item
  \(560\,\text{cm}^{3}\)
\item
  \(312\,\text{cm}^{3}\)
\item
  \(108\,\text{cm}^{3}\)
\end{enumerate}

\end{tcolorbox}

{Area parallelogram} {Volume formula prism}

\end{exercise}

\begin{exercise}[Volume of Cylinder
(diameter)]\protect\hypertarget{exr-vol-7}{}\label{exr-vol-7}

Find the volume. Round to 1 decimal place.

\begin{enumerate}
\def\labelenumi{\alph{enumi})}
\item
  \(\text{Cylinder: }d=7\text{ cm},\ h=11\text{ cm}\)
\item
  \(\text{Cylinder: }d=12\text{ cm},\ h=14\text{ cm}\)
\item
  \(\text{Cylinder: }d=4\text{ cm},\ h=5\text{ cm}\)
\item
  \(\text{Cylinder: }d=4\text{ cm},\ h=13\text{ cm}\)
\item
  \(\text{Cylinder: }d=15\text{ cm},\ h=11\text{ cm}\)
\item
  \(\text{Cylinder: }d=18\text{ cm},\ h=5\text{ cm}\)
\end{enumerate}

\begin{tcolorbox}[enhanced jigsaw, arc=.35mm, breakable, toptitle=1mm, title=\textcolor{quarto-callout-tip-color}{\faLightbulb}\hspace{0.5em}{🔍 View Solution}, colframe=quarto-callout-tip-color-frame, rightrule=.15mm, opacityback=0, titlerule=0mm, colback=white, coltitle=black, bottomtitle=1mm, toprule=.15mm, leftrule=.75mm, bottomrule=.15mm, colbacktitle=quarto-callout-tip-color!10!white, left=2mm, opacitybacktitle=0.6]

\begin{enumerate}
\def\labelenumi{\alph{enumi})}
\item
  \(423.3\,\text{cm}^{3}\)
\item
  \(1583.4\,\text{cm}^{3}\)
\item
  \(62.8\,\text{cm}^{3}\)
\item
  \(163.4\,\text{cm}^{3}\)
\item
  \(1943.9\,\text{cm}^{3}\)
\item
  \(1272.3\,\text{cm}^{3}\)
\end{enumerate}

\end{tcolorbox}

{Cylinder formula} {Diameter to radius} {Use of pi}

\end{exercise}

\begin{exercise}[Volume of
Half-Cylinder]\protect\hypertarget{exr-vol-8}{}\label{exr-vol-8}

Find the volume. Round to 1 decimal place.

\begin{enumerate}
\def\labelenumi{\alph{enumi})}
\item
  \(\text{Half-cylinder: }r=6\text{ cm},\ l=6\text{ cm}\)
\item
  \(\text{Half-cylinder: }r=8\text{ cm},\ l=8\text{ cm}\)
\item
  \(\text{Half-cylinder: }r=3\text{ cm},\ l=10\text{ cm}\)
\item
  \(\text{Half-cylinder: }r=3\text{ cm},\ l=7\text{ cm}\)
\item
  \(\text{Half-cylinder: }r=5\text{ cm},\ l=7\text{ cm}\)
\item
  \(\text{Half-cylinder: }r=5\text{ cm},\ l=5\text{ cm}\)
\end{enumerate}

\begin{tcolorbox}[enhanced jigsaw, arc=.35mm, breakable, toptitle=1mm, title=\textcolor{quarto-callout-tip-color}{\faLightbulb}\hspace{0.5em}{🔍 View Solution}, colframe=quarto-callout-tip-color-frame, rightrule=.15mm, opacityback=0, titlerule=0mm, colback=white, coltitle=black, bottomtitle=1mm, toprule=.15mm, leftrule=.75mm, bottomrule=.15mm, colbacktitle=quarto-callout-tip-color!10!white, left=2mm, opacitybacktitle=0.6]

\begin{enumerate}
\def\labelenumi{\alph{enumi})}
\item
  \(339.3\,\text{cm}^{3}\)
\item
  \(804.2\,\text{cm}^{3}\)
\item
  \(141.4\,\text{cm}^{3}\)
\item
  \(99.0\,\text{cm}^{3}\)
\item
  \(274.9\,\text{cm}^{3}\)
\item
  \(196.3\,\text{cm}^{3}\)
\end{enumerate}

\end{tcolorbox}

{Cylinder formula} {Fraction of solid} {Use of pi}

\end{exercise}

\begin{exercise}[Volume of
Quarter-Cylinder]\protect\hypertarget{exr-vol-9}{}\label{exr-vol-9}

Find the volume. Round to 1 decimal place.

\begin{enumerate}
\def\labelenumi{\alph{enumi})}
\item
  \(\text{Quarter-cylinder: }r=6\text{ cm},\ l=6\text{ cm}\)
\item
  \(\text{Quarter-cylinder: }r=8\text{ cm},\ l=6\text{ cm}\)
\item
  \(\text{Quarter-cylinder: }r=3\text{ cm},\ l=11\text{ cm}\)
\item
  \(\text{Quarter-cylinder: }r=9\text{ cm},\ l=4\text{ cm}\)
\item
  \(\text{Quarter-cylinder: }r=4\text{ cm},\ l=7\text{ cm}\)
\item
  \(\text{Quarter-cylinder: }r=6\text{ cm},\ l=10\text{ cm}\)
\end{enumerate}

\begin{tcolorbox}[enhanced jigsaw, arc=.35mm, breakable, toptitle=1mm, title=\textcolor{quarto-callout-tip-color}{\faLightbulb}\hspace{0.5em}{🔍 View Solution}, colframe=quarto-callout-tip-color-frame, rightrule=.15mm, opacityback=0, titlerule=0mm, colback=white, coltitle=black, bottomtitle=1mm, toprule=.15mm, leftrule=.75mm, bottomrule=.15mm, colbacktitle=quarto-callout-tip-color!10!white, left=2mm, opacitybacktitle=0.6]

\begin{enumerate}
\def\labelenumi{\alph{enumi})}
\item
  \(169.6\,\text{cm}^{3}\)
\item
  \(301.6\,\text{cm}^{3}\)
\item
  \(77.8\,\text{cm}^{3}\)
\item
  \(254.5\,\text{cm}^{3}\)
\item
  \(88.0\,\text{cm}^{3}\)
\item
  \(282.7\,\text{cm}^{3}\)
\end{enumerate}

\end{tcolorbox}

{Cylinder formula} {Fraction of solid} {Use of pi}

\end{exercise}

∑ × π ÷ √ ◆ ∞ ± ∫ ≠ Δ

\subsection{🔴 Challenging}

\begin{exercise}[Volume of Cylindrical
Sector]\protect\hypertarget{exr-vol-10}{}\label{exr-vol-10}

Find the volume. Round to 1 decimal place.

\begin{enumerate}
\def\labelenumi{\alph{enumi})}
\item
  \(\text{Sector cylinder: }r=4,\ l=8\text{ cm},\ \theta=90^\circ\)
\item
  \(\text{Sector cylinder: }r=8,\ l=13\text{ cm},\ \theta=120^\circ\)
\item
  \(\text{Sector cylinder: }r=5,\ l=14\text{ cm},\ \theta=270^\circ\)
\item
  \(\text{Sector cylinder: }r=4,\ l=4\text{ cm},\ \theta=120^\circ\)
\item
  \(\text{Sector cylinder: }r=3,\ l=10\text{ cm},\ \theta=120^\circ\)
\item
  \(\text{Sector cylinder: }r=7,\ l=10\text{ cm},\ \theta=300^\circ\)
\end{enumerate}

\begin{tcolorbox}[enhanced jigsaw, arc=.35mm, breakable, toptitle=1mm, title=\textcolor{quarto-callout-tip-color}{\faLightbulb}\hspace{0.5em}{🔍 View Solution}, colframe=quarto-callout-tip-color-frame, rightrule=.15mm, opacityback=0, titlerule=0mm, colback=white, coltitle=black, bottomtitle=1mm, toprule=.15mm, leftrule=.75mm, bottomrule=.15mm, colbacktitle=quarto-callout-tip-color!10!white, left=2mm, opacitybacktitle=0.6]

\begin{enumerate}
\def\labelenumi{\alph{enumi})}
\item
  \(100.5\,\text{cm}^{3}\)
\item
  \(871.3\,\text{cm}^{3}\)
\item
  \(824.7\,\text{cm}^{3}\)
\item
  \(67.0\,\text{cm}^{3}\)
\item
  \(94.2\,\text{cm}^{3}\)
\item
  \(1282.8\,\text{cm}^{3}\)
\end{enumerate}

\end{tcolorbox}

{Cylinder formula} {Sector fraction} {Use of pi} {Proportional
reasoning}

\end{exercise}

∑ × π ÷ √ ◆ ∞ ± ∫ ≠ Δ

\begin{center}\rule{0.5\linewidth}{0.5pt}\end{center}

\section{🧩 Part C --- Composite Solids \& Missing
Dimensions}\label{sec-composite}

\subsection{🟡 Medium}

\begin{exercise}[Composite: L-Shaped
Prism]\protect\hypertarget{exr-comp-3}{}\label{exr-comp-3}

Find the volume of the L-shaped prism.

\begin{enumerate}
\def\labelenumi{\alph{enumi})}
\item
  \[\text{L-prism: }(4\times4)\text{ and }(2\times4)\text{ cm},\ l=11\text{ cm}\]
\item
  \[\text{L-prism: }(4\times3)\text{ and }(3\times4)\text{ cm},\ l=14\text{ cm}\]
\item
  \[\text{L-prism: }(6\times4)\text{ and }(4\times4)\text{ cm},\ l=8\text{ cm}\]
\item
  \[\text{L-prism: }(7\times5)\text{ and }(4\times3)\text{ cm},\ l=11\text{ cm}\]
\item
  \[\text{L-prism: }(8\times5)\text{ and }(2\times2)\text{ cm},\ l=7\text{ cm}\]
\item
  \[\text{L-prism: }(7\times6)\text{ and }(3\times3)\text{ cm},\ l=10\text{ cm}\]
\end{enumerate}

\begin{tcolorbox}[enhanced jigsaw, arc=.35mm, breakable, toptitle=1mm, title=\textcolor{quarto-callout-tip-color}{\faLightbulb}\hspace{0.5em}{🔍 View Solution}, colframe=quarto-callout-tip-color-frame, rightrule=.15mm, opacityback=0, titlerule=0mm, colback=white, coltitle=black, bottomtitle=1mm, toprule=.15mm, leftrule=.75mm, bottomrule=.15mm, colbacktitle=quarto-callout-tip-color!10!white, left=2mm, opacitybacktitle=0.6]

\begin{enumerate}
\def\labelenumi{\alph{enumi})}
\item
  \[264\,\text{cm}^{3}\]
\item
  \[336\,\text{cm}^{3}\]
\item
  \[320\,\text{cm}^{3}\]
\item
  \[517\,\text{cm}^{3}\]
\item
  \[308\,\text{cm}^{3}\]
\item
  \[510\,\text{cm}^{3}\]
\end{enumerate}

\end{tcolorbox}

{Composite volume} {Area l shape} {Multi step calculation}

\end{exercise}

\begin{exercise}[Find Missing Prism
Dimension]\protect\hypertarget{exr-comp-7}{}\label{exr-comp-7}

Find the missing dimension.

\begin{enumerate}
\def\labelenumi{\alph{enumi})}
\item
  \[\text{Cuboid: }V=80\,\text{cm}^{3},\ w=4,\ h=2\text{ cm. Find }l.\]
\item
  \[\text{Cuboid: }V=160\,\text{cm}^{3},\ w=8,\ h=5\text{ cm. Find }l.\]
\item
  \[\text{Cuboid: }V=200\,\text{cm}^{3},\ w=4,\ h=5\text{ cm. Find }l.\]
\item
  \[\text{Cuboid: }V=210\,\text{cm}^{3},\ w=6,\ h=7\text{ cm. Find }l.\]
\item
  \[\text{Cuboid: }V=48\,\text{cm}^{3},\ w=2,\ h=3\text{ cm. Find }l.\]
\item
  \[\text{Cuboid: }V=100\,\text{cm}^{3},\ w=5,\ h=5\text{ cm. Find }l.\]
\end{enumerate}

\begin{tcolorbox}[enhanced jigsaw, arc=.35mm, breakable, toptitle=1mm, title=\textcolor{quarto-callout-tip-color}{\faLightbulb}\hspace{0.5em}{🔍 View Solution}, colframe=quarto-callout-tip-color-frame, rightrule=.15mm, opacityback=0, titlerule=0mm, colback=white, coltitle=black, bottomtitle=1mm, toprule=.15mm, leftrule=.75mm, bottomrule=.15mm, colbacktitle=quarto-callout-tip-color!10!white, left=2mm, opacitybacktitle=0.6]

\begin{enumerate}
\def\labelenumi{\alph{enumi})}
\item
  \[l=10\text{ cm}\]
\item
  \[l=4\text{ cm}\]
\item
  \[l=10\text{ cm}\]
\item
  \[l=5\text{ cm}\]
\item
  \[l=8\text{ cm}\]
\item
  \[l=4\text{ cm}\]
\end{enumerate}

\end{tcolorbox}

{Rearranging formula} {Division of integers} {Algebraic reasoning}

\end{exercise}

\begin{exercise}[Find Missing Height --- Triangular
Prism]\protect\hypertarget{exr-comp-9}{}\label{exr-comp-9}

Find the value of x.

\begin{enumerate}
\def\labelenumi{\alph{enumi})}
\item
  \[\text{Triangular prism: }V=125.0\,\text{cm}^{3},\ b=5\text{ cm},\ l=10\text{ cm. Find }x.\]
\item
  \[\text{Triangular prism: }V=100.0\,\text{cm}^{3},\ b=5\text{ cm},\ l=10\text{ cm. Find }x.\]
\item
  \[\text{Triangular prism: }V=192.5\,\text{cm}^{3},\ b=5\text{ cm},\ l=11\text{ cm. Find }x.\]
\item
  \[\text{Triangular prism: }V=192.0\,\text{cm}^{3},\ b=8\text{ cm},\ l=12\text{ cm. Find }x.\]
\item
  \[\text{Triangular prism: }V=67.5\,\text{cm}^{3},\ b=5\text{ cm},\ l=9\text{ cm. Find }x.\]
\item
  \[\text{Triangular prism: }V=140.0\,\text{cm}^{3},\ b=8\text{ cm},\ l=7\text{ cm. Find }x.\]
\end{enumerate}

\begin{tcolorbox}[enhanced jigsaw, arc=.35mm, breakable, toptitle=1mm, title=\textcolor{quarto-callout-tip-color}{\faLightbulb}\hspace{0.5em}{🔍 View Solution}, colframe=quarto-callout-tip-color-frame, rightrule=.15mm, opacityback=0, titlerule=0mm, colback=white, coltitle=black, bottomtitle=1mm, toprule=.15mm, leftrule=.75mm, bottomrule=.15mm, colbacktitle=quarto-callout-tip-color!10!white, left=2mm, opacitybacktitle=0.6]

\begin{enumerate}
\def\labelenumi{\alph{enumi})}
\item
  \[x=5\text{ cm}\]
\item
  \[x=4\text{ cm}\]
\item
  \[x=7\text{ cm}\]
\item
  \[x=4\text{ cm}\]
\item
  \[x=3\text{ cm}\]
\item
  \[x=5\text{ cm}\]
\end{enumerate}

\end{tcolorbox}

{Rearranging formula} {Division of integers}

\end{exercise}

\begin{exercise}[Real-World: Water
Tank]\protect\hypertarget{exr-comp-10}{}\label{exr-comp-10}

Solve the word problem. Show all working.

\begin{enumerate}
\def\labelenumi{\alph{enumi})}
\item
  \[\text{Cylindrical tank: }r=1.8\text{ m},\ h=1.6\text{ m.}\ \text{How many litres?}\]
\item
  \[\text{Cylindrical tank: }r=0.8\text{ m},\ h=2.7\text{ m.}\ \text{How many litres?}\]
\item
  \[\text{Cylindrical tank: }r=1.1\text{ m},\ h=3.9\text{ m.}\ \text{How many litres?}\]
\item
  \[\text{Cylindrical tank: }r=1.3\text{ m},\ h=2.2\text{ m.}\ \text{How many litres?}\]
\item
  \[\text{Cylindrical tank: }r=2.1\text{ m},\ h=4.3\text{ m.}\ \text{How many litres?}\]
\item
  \[\text{Cylindrical tank: }r=2.0\text{ m},\ h=1.5\text{ m.}\ \text{How many litres?}\]
\end{enumerate}

\begin{tcolorbox}[enhanced jigsaw, arc=.35mm, breakable, toptitle=1mm, title=\textcolor{quarto-callout-tip-color}{\faLightbulb}\hspace{0.5em}{🔍 View Solution}, colframe=quarto-callout-tip-color-frame, rightrule=.15mm, opacityback=0, titlerule=0mm, colback=white, coltitle=black, bottomtitle=1mm, toprule=.15mm, leftrule=.75mm, bottomrule=.15mm, colbacktitle=quarto-callout-tip-color!10!white, left=2mm, opacitybacktitle=0.6]

\begin{enumerate}
\def\labelenumi{\alph{enumi})}
\item
  \[16{,}290\text{ L}\]
\item
  \[5430\text{ L}\]
\item
  \[14{,}830\text{ L}\]
\item
  \[11{,}680\text{ L}\]
\item
  \[59{,}570\text{ L}\]
\item
  \[18{,}850\text{ L}\]
\end{enumerate}

\end{tcolorbox}

{Cylinder formula} {Unit conversion m3 to litres} {Real world
application}

\end{exercise}

\begin{exercise}[Real-World: Shed
Roof]\protect\hypertarget{exr-comp-11}{}\label{exr-comp-11}

Solve the word problem. Show all working.

\begin{enumerate}
\def\labelenumi{\alph{enumi})}
\item
  \[\text{Triangular prism roof: }b=5\text{ m},\ h=2\text{ m},\ l=12\text{ m. Find volume.}\]
\item
  \[\text{Triangular prism roof: }b=6\text{ m},\ h=2\text{ m},\ l=13\text{ m. Find volume.}\]
\item
  \[\text{Triangular prism roof: }b=3\text{ m},\ h=3\text{ m},\ l=6\text{ m. Find volume.}\]
\item
  \[\text{Triangular prism roof: }b=7\text{ m},\ h=5\text{ m},\ l=9\text{ m. Find volume.}\]
\item
  \[\text{Triangular prism roof: }b=7\text{ m},\ h=2\text{ m},\ l=8\text{ m. Find volume.}\]
\item
  \[\text{Triangular prism roof: }b=3\text{ m},\ h=5\text{ m},\ l=10\text{ m. Find volume.}\]
\end{enumerate}

\begin{tcolorbox}[enhanced jigsaw, arc=.35mm, breakable, toptitle=1mm, title=\textcolor{quarto-callout-tip-color}{\faLightbulb}\hspace{0.5em}{🔍 View Solution}, colframe=quarto-callout-tip-color-frame, rightrule=.15mm, opacityback=0, titlerule=0mm, colback=white, coltitle=black, bottomtitle=1mm, toprule=.15mm, leftrule=.75mm, bottomrule=.15mm, colbacktitle=quarto-callout-tip-color!10!white, left=2mm, opacitybacktitle=0.6]

\begin{enumerate}
\def\labelenumi{\alph{enumi})}
\item
  \[60.0\,\text{m}^{3}\]
\item
  \[78.0\,\text{m}^{3}\]
\item
  \[27.0\,\text{m}^{3}\]
\item
  \[157.5\,\text{m}^{3}\]
\item
  \[56.0\,\text{m}^{3}\]
\item
  \[75.0\,\text{m}^{3}\]
\end{enumerate}

\end{tcolorbox}

{Prism formula} {Real world application}

\end{exercise}

∑ × π ÷ √ ◆ ∞ ± ∫ ≠ Δ

\subsection{🔴 Challenging}

\begin{exercise}[Composite: Prism +
Half-Cylinder]\protect\hypertarget{exr-comp-1}{}\label{exr-comp-1}

Find the total volume.

\begin{enumerate}
\def\labelenumi{\alph{enumi})}
\item
  \[\text{Prism+half-cyl: }l=14,w=6,h=7\text{ cm},\ r=3.0\text{ cm}\]
\item
  \[\text{Prism+half-cyl: }l=6,w=8,h=8\text{ cm},\ r=4.0\text{ cm}\]
\item
  \[\text{Prism+half-cyl: }l=9,w=6,h=3\text{ cm},\ r=3.0\text{ cm}\]
\item
  \[\text{Prism+half-cyl: }l=12,w=7,h=8\text{ cm},\ r=3.5\text{ cm}\]
\item
  \[\text{Prism+half-cyl: }l=8,w=4,h=6\text{ cm},\ r=2.0\text{ cm}\]
\item
  \[\text{Prism+half-cyl: }l=6,w=8,h=4\text{ cm},\ r=4.0\text{ cm}\]
\end{enumerate}

\begin{tcolorbox}[enhanced jigsaw, arc=.35mm, breakable, toptitle=1mm, title=\textcolor{quarto-callout-tip-color}{\faLightbulb}\hspace{0.5em}{🔍 View Solution}, colframe=quarto-callout-tip-color-frame, rightrule=.15mm, opacityback=0, titlerule=0mm, colback=white, coltitle=black, bottomtitle=1mm, toprule=.15mm, leftrule=.75mm, bottomrule=.15mm, colbacktitle=quarto-callout-tip-color!10!white, left=2mm, opacitybacktitle=0.6]

\begin{enumerate}
\def\labelenumi{\alph{enumi})}
\item
  \[785.9\,\text{cm}^{3}\]
\item
  \[534.8\,\text{cm}^{3}\]
\item
  \[289.2\,\text{cm}^{3}\]
\item
  \[902.9\,\text{cm}^{3}\]
\item
  \[242.3\,\text{cm}^{3}\]
\item
  \[342.8\,\text{cm}^{3}\]
\end{enumerate}

\end{tcolorbox}

{Composite volume} {Multi step calculation} {Combining shapes}

\end{exercise}

\begin{exercise}[Composite:
L-Prism]\protect\hypertarget{exr-comp-2}{}\label{exr-comp-2}

Find the total volume.

\begin{enumerate}
\def\labelenumi{\alph{enumi})}
\item
  \[\text{L-prism: }(4\times3)\text{ and }(2\times4)\text{ cm},\ l=14\text{ cm}\]
\item
  \[\text{L-prism: }(7\times5)\text{ and }(4\times2)\text{ cm},\ l=12\text{ cm}\]
\item
  \[\text{L-prism: }(4\times4)\text{ and }(4\times4)\text{ cm},\ l=10\text{ cm}\]
\item
  \[\text{L-prism: }(10\times6)\text{ and }(4\times2)\text{ cm},\ l=10\text{ cm}\]
\item
  \[\text{L-prism: }(9\times3)\text{ and }(4\times3)\text{ cm},\ l=10\text{ cm}\]
\item
  \[\text{L-prism: }(5\times5)\text{ and }(3\times3)\text{ cm},\ l=14\text{ cm}\]
\end{enumerate}

\begin{tcolorbox}[enhanced jigsaw, arc=.35mm, breakable, toptitle=1mm, title=\textcolor{quarto-callout-tip-color}{\faLightbulb}\hspace{0.5em}{🔍 View Solution}, colframe=quarto-callout-tip-color-frame, rightrule=.15mm, opacityback=0, titlerule=0mm, colback=white, coltitle=black, bottomtitle=1mm, toprule=.15mm, leftrule=.75mm, bottomrule=.15mm, colbacktitle=quarto-callout-tip-color!10!white, left=2mm, opacitybacktitle=0.6]

\begin{enumerate}
\def\labelenumi{\alph{enumi})}
\item
  \[280\,\text{cm}^{3}\]
\item
  \[516\,\text{cm}^{3}\]
\item
  \[320\,\text{cm}^{3}\]
\item
  \[680\,\text{cm}^{3}\]
\item
  \[390\,\text{cm}^{3}\]
\item
  \[476\,\text{cm}^{3}\]
\end{enumerate}

\end{tcolorbox}

{Composite volume} {Decomposing shapes} {Multi step calculation}

\end{exercise}

\begin{exercise}[Composite: Subtract
Cylinder]\protect\hypertarget{exr-comp-4}{}\label{exr-comp-4}

Find the volume of the solid.

\begin{enumerate}
\def\labelenumi{\alph{enumi})}
\item
  \[\text{Box}(l=10,w=9,h=8\text{ cm)}\text{, hole }r=2\text{ cm}\]
\item
  \[\text{Box}(l=8,w=10,h=4\text{ cm)}\text{, hole }r=2\text{ cm}\]
\item
  \[\text{Box}(l=12,w=9,h=4\text{ cm)}\text{, hole }r=1\text{ cm}\]
\item
  \[\text{Box}(l=13,w=6,h=6\text{ cm)}\text{, hole }r=3\text{ cm}\]
\item
  \[\text{Box}(l=10,w=8,h=6\text{ cm)}\text{, hole }r=3\text{ cm}\]
\item
  \[\text{Box}(l=12,w=6,h=8\text{ cm)}\text{, hole }r=3\text{ cm}\]
\end{enumerate}

\begin{tcolorbox}[enhanced jigsaw, arc=.35mm, breakable, toptitle=1mm, title=\textcolor{quarto-callout-tip-color}{\faLightbulb}\hspace{0.5em}{🔍 View Solution}, colframe=quarto-callout-tip-color-frame, rightrule=.15mm, opacityback=0, titlerule=0mm, colback=white, coltitle=black, bottomtitle=1mm, toprule=.15mm, leftrule=.75mm, bottomrule=.15mm, colbacktitle=quarto-callout-tip-color!10!white, left=2mm, opacitybacktitle=0.6]

\begin{enumerate}
\def\labelenumi{\alph{enumi})}
\item
  \[594.3\,\text{cm}^{3}\]
\item
  \[219.5\,\text{cm}^{3}\]
\item
  \[394.3\,\text{cm}^{3}\]
\item
  \[100.4\,\text{cm}^{3}\]
\item
  \[197.3\,\text{cm}^{3}\]
\item
  \[236.7\,\text{cm}^{3}\]
\end{enumerate}

\end{tcolorbox}

{Composite volume} {Subtraction of volumes} {Use of pi}

\end{exercise}

\begin{exercise}[Composite: Hollow
Cylinder]\protect\hypertarget{exr-comp-5}{}\label{exr-comp-5}

Find the volume of material.

\begin{enumerate}
\def\labelenumi{\alph{enumi})}
\item
  \[\text{Hollow pipe: }R=7,\ r=3,\ l=10\text{ cm}\]
\item
  \[\text{Hollow pipe: }R=8,\ r=2,\ l=17\text{ cm}\]
\item
  \[\text{Hollow pipe: }R=5,\ r=2,\ l=14\text{ cm}\]
\item
  \[\text{Hollow pipe: }R=6,\ r=1,\ l=17\text{ cm}\]
\item
  \[\text{Hollow pipe: }R=5,\ r=1,\ l=18\text{ cm}\]
\item
  \[\text{Hollow pipe: }R=6,\ r=1,\ l=15\text{ cm}\]
\end{enumerate}

\begin{tcolorbox}[enhanced jigsaw, arc=.35mm, breakable, toptitle=1mm, title=\textcolor{quarto-callout-tip-color}{\faLightbulb}\hspace{0.5em}{🔍 View Solution}, colframe=quarto-callout-tip-color-frame, rightrule=.15mm, opacityback=0, titlerule=0mm, colback=white, coltitle=black, bottomtitle=1mm, toprule=.15mm, leftrule=.75mm, bottomrule=.15mm, colbacktitle=quarto-callout-tip-color!10!white, left=2mm, opacitybacktitle=0.6]

\begin{enumerate}
\def\labelenumi{\alph{enumi})}
\item
  \[1256.7\,\text{cm}^{3}\]
\item
  \[3204.5\,\text{cm}^{3}\]
\item
  \[923.7\,\text{cm}^{3}\]
\item
  \[1869.3\,\text{cm}^{3}\]
\item
  \[1357.2\,\text{cm}^{3}\]
\item
  \[1649.4\,\text{cm}^{3}\]
\end{enumerate}

\end{tcolorbox}

{Composite volume} {Subtraction of volumes} {Use of pi}

\end{exercise}

\begin{exercise}[Composite: T-Shape
Cuboids]\protect\hypertarget{exr-comp-6}{}\label{exr-comp-6}

Find the total volume of the compound solid.

\begin{enumerate}
\def\labelenumi{\alph{enumi})}
\item
  \[\text{T-shape: base }(9\times7\times2)\text{ cm, top }(7\times5\times9)\text{ cm.}\ \text{Find total volume.}\]
\item
  \[\text{T-shape: base }(5\times5\times5)\text{ cm, top }(4\times3\times6)\text{ cm.}\ \text{Find total volume.}\]
\item
  \[\text{T-shape: base }(12\times8\times3)\text{ cm, top }(4\times5\times8)\text{ cm.}\ \text{Find total volume.}\]
\item
  \[\text{T-shape: base }(6\times9\times4)\text{ cm, top }(2\times3\times7)\text{ cm.}\ \text{Find total volume.}\]
\item
  \[\text{T-shape: base }(12\times4\times3)\text{ cm, top }(8\times2\times7)\text{ cm.}\ \text{Find total volume.}\]
\item
  \[\text{T-shape: base }(8\times6\times4)\text{ cm, top }(6\times2\times6)\text{ cm.}\ \text{Find total volume.}\]
\end{enumerate}

\begin{tcolorbox}[enhanced jigsaw, arc=.35mm, breakable, toptitle=1mm, title=\textcolor{quarto-callout-tip-color}{\faLightbulb}\hspace{0.5em}{🔍 View Solution}, colframe=quarto-callout-tip-color-frame, rightrule=.15mm, opacityback=0, titlerule=0mm, colback=white, coltitle=black, bottomtitle=1mm, toprule=.15mm, leftrule=.75mm, bottomrule=.15mm, colbacktitle=quarto-callout-tip-color!10!white, left=2mm, opacitybacktitle=0.6]

\begin{enumerate}
\def\labelenumi{\alph{enumi})}
\item
  \[441\,\text{cm}^{3}\]
\item
  \[197\,\text{cm}^{3}\]
\item
  \[448\,\text{cm}^{3}\]
\item
  \[258\,\text{cm}^{3}\]
\item
  \[256\,\text{cm}^{3}\]
\item
  \[264\,\text{cm}^{3}\]
\end{enumerate}

\end{tcolorbox}

{Composite volume} {Decomposing shapes} {Multi step calculation}

\end{exercise}

\begin{exercise}[Find Missing Cylinder
Dimension]\protect\hypertarget{exr-comp-8}{}\label{exr-comp-8}

Find the missing dimension. Round to 1 d.p.

\begin{enumerate}
\def\labelenumi{\alph{enumi})}
\item
  \[\text{Cylinder: }V=169.6\,\text{cm}^{3},\ h=6\text{ cm. Find }r.\]
\item
  \[\text{Cylinder: }V=615.8\,\text{cm}^{3},\ r=7\text{ cm. Find }h.\]
\item
  \[\text{Cylinder: }V=314.2\,\text{cm}^{3},\ r=5\text{ cm. Find }h.\]
\item
  \[\text{Cylinder: }V=311.0\,\text{cm}^{3},\ h=11\text{ cm. Find }r.\]
\item
  \[\text{Cylinder: }V=150.8\,\text{cm}^{3},\ r=2\text{ cm. Find }h.\]
\item
  \[\text{Cylinder: }V=1608.5\,\text{cm}^{3},\ h=8\text{ cm. Find }r.\]
\end{enumerate}

\begin{tcolorbox}[enhanced jigsaw, arc=.35mm, breakable, toptitle=1mm, title=\textcolor{quarto-callout-tip-color}{\faLightbulb}\hspace{0.5em}{🔍 View Solution}, colframe=quarto-callout-tip-color-frame, rightrule=.15mm, opacityback=0, titlerule=0mm, colback=white, coltitle=black, bottomtitle=1mm, toprule=.15mm, leftrule=.75mm, bottomrule=.15mm, colbacktitle=quarto-callout-tip-color!10!white, left=2mm, opacitybacktitle=0.6]

\begin{enumerate}
\def\labelenumi{\alph{enumi})}
\item
  \[r=3\text{ cm}\]
\item
  \[h=4\text{ cm}\]
\item
  \[h=4\text{ cm}\]
\item
  \[r=3\text{ cm}\]
\item
  \[h=12\text{ cm}\]
\item
  \[r=8\text{ cm}\]
\end{enumerate}

\end{tcolorbox}

{Rearranging formula} {Square root} {Use of pi} {Algebraic reasoning}

\end{exercise}

\begin{exercise}[Similar Cylinders --- Find
Volume]\protect\hypertarget{exr-comp-12}{}\label{exr-comp-12}

Use the scale factor to find the volume.

\begin{enumerate}
\def\labelenumi{\alph{enumi})}
\item
  \[\text{Similar cylinders A and B: }l_A=3\text{ cm},\ l_B=8\text{ cm}.\ V_A=58\,\text{cm}^{3}.\ \text{Find }V_B.\]
\item
  \[\text{Similar cylinders A and B: }l_A=5\text{ cm},\ l_B=7\text{ cm}.\ V_A=171\,\text{cm}^{3}.\ \text{Find }V_B.\]
\item
  \[\text{Similar cylinders A and B: }l_A=6\text{ cm},\ l_B=9\text{ cm}.\ V_A=163\,\text{cm}^{3}.\ \text{Find }V_B.\]
\item
  \[\text{Similar cylinders A and B: }l_A=8\text{ cm},\ l_B=12\text{ cm}.\ V_A=144\,\text{cm}^{3}.\ \text{Find }V_B.\]
\item
  \[\text{Similar cylinders A and B: }l_A=3\text{ cm},\ l_B=5\text{ cm}.\ V_A=57\,\text{cm}^{3}.\ \text{Find }V_B.\]
\item
  \[\text{Similar cylinders A and B: }l_A=4\text{ cm},\ l_B=5\text{ cm}.\ V_A=194\,\text{cm}^{3}.\ \text{Find }V_B.\]
\end{enumerate}

\begin{tcolorbox}[enhanced jigsaw, arc=.35mm, breakable, toptitle=1mm, title=\textcolor{quarto-callout-tip-color}{\faLightbulb}\hspace{0.5em}{🔍 View Solution}, colframe=quarto-callout-tip-color-frame, rightrule=.15mm, opacityback=0, titlerule=0mm, colback=white, coltitle=black, bottomtitle=1mm, toprule=.15mm, leftrule=.75mm, bottomrule=.15mm, colbacktitle=quarto-callout-tip-color!10!white, left=2mm, opacitybacktitle=0.6]

\begin{enumerate}
\def\labelenumi{\alph{enumi})}
\item
  \[1099.9\,\text{cm}^{3}\]
\item
  \[469.2\,\text{cm}^{3}\]
\item
  \[550.1\,\text{cm}^{3}\]
\item
  \[486.0\,\text{cm}^{3}\]
\item
  \[263.9\,\text{cm}^{3}\]
\item
  \[378.9\,\text{cm}^{3}\]
\end{enumerate}

\end{tcolorbox}

{Similar solids} {Scale factor} {Cubing scale factor}

\end{exercise}

∑ × π ÷ √ ◆ ∞ ± ∫ ≠ Δ

\begin{center}\rule{0.5\linewidth}{0.5pt}\end{center}

\section{🎨 Part D --- Surface Area}\label{sec-sa}

\subsection{🟢 Easy}

\begin{exercise}[Surface Area of a
Cuboid]\protect\hypertarget{exr-sa-1}{}\label{exr-sa-1}

Find the total surface area.

\begin{enumerate}
\def\labelenumi{\alph{enumi})}
\item
  \(\text{Cuboid: }l=10,\ w=6,\ h=6\text{ cm}\)
\item
  \(\text{Cuboid: }l=3,\ w=3,\ h=4\text{ cm}\)
\item
  \(\text{Cuboid: }l=6,\ w=3,\ h=4\text{ cm}\)
\item
  \(\text{Cuboid: }l=7,\ w=7,\ h=5\text{ cm}\)
\item
  \(\text{Cuboid: }l=3,\ w=6,\ h=7\text{ cm}\)
\item
  \(\text{Cuboid: }l=5,\ w=2,\ h=7\text{ cm}\)
\end{enumerate}

\begin{tcolorbox}[enhanced jigsaw, arc=.35mm, breakable, toptitle=1mm, title=\textcolor{quarto-callout-tip-color}{\faLightbulb}\hspace{0.5em}{🔍 View Solution}, colframe=quarto-callout-tip-color-frame, rightrule=.15mm, opacityback=0, titlerule=0mm, colback=white, coltitle=black, bottomtitle=1mm, toprule=.15mm, leftrule=.75mm, bottomrule=.15mm, colbacktitle=quarto-callout-tip-color!10!white, left=2mm, opacitybacktitle=0.6]

\begin{enumerate}
\def\labelenumi{\alph{enumi})}
\item
  \(312\,\text{cm}^{2}\)
\item
  \(66\,\text{cm}^{2}\)
\item
  \(108\,\text{cm}^{2}\)
\item
  \(238\,\text{cm}^{2}\)
\item
  \(162\,\text{cm}^{2}\)
\item
  \(118\,\text{cm}^{2}\)
\end{enumerate}

\end{tcolorbox}

{Sa formula cuboid} {Counting faces} {Multiplication of integers}

\end{exercise}

\begin{exercise}[Surface Area of Open-Top
Box]\protect\hypertarget{exr-sa-2}{}\label{exr-sa-2}

Find the surface area (no lid).

\begin{enumerate}
\def\labelenumi{\alph{enumi})}
\item
  \(\text{Open box: }l=7,\ w=7,\ h=5\text{ cm}\)
\item
  \(\text{Open box: }l=4,\ w=4,\ h=6\text{ cm}\)
\item
  \(\text{Open box: }l=9,\ w=3,\ h=4\text{ cm}\)
\item
  \(\text{Open box: }l=9,\ w=4,\ h=6\text{ cm}\)
\item
  \(\text{Open box: }l=12,\ w=3,\ h=5\text{ cm}\)
\item
  \(\text{Open box: }l=4,\ w=5,\ h=4\text{ cm}\)
\end{enumerate}

\begin{tcolorbox}[enhanced jigsaw, arc=.35mm, breakable, toptitle=1mm, title=\textcolor{quarto-callout-tip-color}{\faLightbulb}\hspace{0.5em}{🔍 View Solution}, colframe=quarto-callout-tip-color-frame, rightrule=.15mm, opacityback=0, titlerule=0mm, colback=white, coltitle=black, bottomtitle=1mm, toprule=.15mm, leftrule=.75mm, bottomrule=.15mm, colbacktitle=quarto-callout-tip-color!10!white, left=2mm, opacitybacktitle=0.6]

\begin{enumerate}
\def\labelenumi{\alph{enumi})}
\item
  \(189\,\text{cm}^{2}\)
\item
  \(112\,\text{cm}^{2}\)
\item
  \(123\,\text{cm}^{2}\)
\item
  \(192\,\text{cm}^{2}\)
\item
  \(186\,\text{cm}^{2}\)
\item
  \(92\,\text{cm}^{2}\)
\end{enumerate}

\end{tcolorbox}

{Sa formula cuboid} {Counting faces} {Omit one face}

\end{exercise}

\begin{exercise}[Surface Area: Open
Tube]\protect\hypertarget{exr-sa-8}{}\label{exr-sa-8}

Find the curved surface area. Round to 1 d.p.

\begin{enumerate}
\def\labelenumi{\alph{enumi})}
\item
  \(\text{Open tube: }r=6\text{ cm},\ l=5\text{ cm}\)
\item
  \(\text{Open tube: }r=2\text{ cm},\ l=16\text{ cm}\)
\item
  \(\text{Open tube: }r=8\text{ cm},\ l=11\text{ cm}\)
\item
  \(\text{Open tube: }r=4\text{ cm},\ l=18\text{ cm}\)
\item
  \(\text{Open tube: }r=5\text{ cm},\ l=6\text{ cm}\)
\item
  \(\text{Open tube: }r=7\text{ cm},\ l=10\text{ cm}\)
\end{enumerate}

\begin{tcolorbox}[enhanced jigsaw, arc=.35mm, breakable, toptitle=1mm, title=\textcolor{quarto-callout-tip-color}{\faLightbulb}\hspace{0.5em}{🔍 View Solution}, colframe=quarto-callout-tip-color-frame, rightrule=.15mm, opacityback=0, titlerule=0mm, colback=white, coltitle=black, bottomtitle=1mm, toprule=.15mm, leftrule=.75mm, bottomrule=.15mm, colbacktitle=quarto-callout-tip-color!10!white, left=2mm, opacitybacktitle=0.6]

\begin{enumerate}
\def\labelenumi{\alph{enumi})}
\item
  \(188.5\,\text{cm}^{2}\)
\item
  \(201.1\,\text{cm}^{2}\)
\item
  \(552.9\,\text{cm}^{2}\)
\item
  \(452.4\,\text{cm}^{2}\)
\item
  \(188.5\,\text{cm}^{2}\)
\item
  \(439.8\,\text{cm}^{2}\)
\end{enumerate}

\end{tcolorbox}

{Curved surface area} {Use of pi}

\end{exercise}

∑ × π ÷ √ ◆ ∞ ± ∫ ≠ Δ

\subsection{🟡 Medium}

\begin{exercise}[Surface Area of a Triangular
Prism]\protect\hypertarget{exr-sa-3}{}\label{exr-sa-3}

Find the total surface area.

\begin{enumerate}
\def\labelenumi{\alph{enumi})}
\item
  \(\text{Tri-prism: }b=4,\ h=4,\ s_1=6,\ s_2=5,\ l=13\text{ cm}\)
\item
  \(\text{Tri-prism: }b=5,\ h=5,\ s_1=7,\ s_2=5,\ l=14\text{ cm}\)
\item
  \(\text{Tri-prism: }b=4,\ h=6,\ s_1=7,\ s_2=8,\ l=6\text{ cm}\)
\item
  \(\text{Tri-prism: }b=4,\ h=7,\ s_1=5,\ s_2=8,\ l=14\text{ cm}\)
\item
  \(\text{Tri-prism: }b=3,\ h=3,\ s_1=7,\ s_2=8,\ l=9\text{ cm}\)
\item
  \(\text{Tri-prism: }b=3,\ h=4,\ s_1=5,\ s_2=7,\ l=12\text{ cm}\)
\end{enumerate}

\begin{tcolorbox}[enhanced jigsaw, arc=.35mm, breakable, toptitle=1mm, title=\textcolor{quarto-callout-tip-color}{\faLightbulb}\hspace{0.5em}{🔍 View Solution}, colframe=quarto-callout-tip-color-frame, rightrule=.15mm, opacityback=0, titlerule=0mm, colback=white, coltitle=black, bottomtitle=1mm, toprule=.15mm, leftrule=.75mm, bottomrule=.15mm, colbacktitle=quarto-callout-tip-color!10!white, left=2mm, opacitybacktitle=0.6]

\begin{enumerate}
\def\labelenumi{\alph{enumi})}
\item
  \(211.0\,\text{cm}^{2}\)
\item
  \(263.0\,\text{cm}^{2}\)
\item
  \(138.0\,\text{cm}^{2}\)
\item
  \(266.0\,\text{cm}^{2}\)
\item
  \(171.0\,\text{cm}^{2}\)
\item
  \(192.0\,\text{cm}^{2}\)
\end{enumerate}

\end{tcolorbox}

{Area triangle} {Sa formula prism} {Counting faces}

\end{exercise}

\begin{exercise}[Surface Area of a Trapezoidal
Prism]\protect\hypertarget{exr-sa-4}{}\label{exr-sa-4}

Find the total surface area.

\begin{enumerate}
\def\labelenumi{\alph{enumi})}
\item
  \(\text{Trap-prism: }a=2,\ b=7,\ h=4,\ s_1=7,\ s_2=6,\ l=9\text{ cm}\)
\item
  \(\text{Trap-prism: }a=3,\ b=5,\ h=2,\ s_1=4,\ s_2=7,\ l=11\text{ cm}\)
\item
  \(\text{Trap-prism: }a=3,\ b=8,\ h=2,\ s_1=6,\ s_2=4,\ l=5\text{ cm}\)
\item
  \(\text{Trap-prism: }a=5,\ b=10,\ h=3,\ s_1=4,\ s_2=6,\ l=8\text{ cm}\)
\item
  \(\text{Trap-prism: }a=3,\ b=6,\ h=4,\ s_1=5,\ s_2=3,\ l=7\text{ cm}\)
\item
  \(\text{Trap-prism: }a=4,\ b=7,\ h=3,\ s_1=3,\ s_2=4,\ l=7\text{ cm}\)
\end{enumerate}

\begin{tcolorbox}[enhanced jigsaw, arc=.35mm, breakable, toptitle=1mm, title=\textcolor{quarto-callout-tip-color}{\faLightbulb}\hspace{0.5em}{🔍 View Solution}, colframe=quarto-callout-tip-color-frame, rightrule=.15mm, opacityback=0, titlerule=0mm, colback=white, coltitle=black, bottomtitle=1mm, toprule=.15mm, leftrule=.75mm, bottomrule=.15mm, colbacktitle=quarto-callout-tip-color!10!white, left=2mm, opacitybacktitle=0.6]

\begin{enumerate}
\def\labelenumi{\alph{enumi})}
\item
  \(234.0\,\text{cm}^{2}\)
\item
  \(225.0\,\text{cm}^{2}\)
\item
  \(127.0\,\text{cm}^{2}\)
\item
  \(245.0\,\text{cm}^{2}\)
\item
  \(155.0\,\text{cm}^{2}\)
\item
  \(159.0\,\text{cm}^{2}\)
\end{enumerate}

\end{tcolorbox}

{Area trapezoid} {Sa formula prism} {Counting faces}

\end{exercise}

\begin{exercise}[Surface Area of Closed
Cylinder]\protect\hypertarget{exr-sa-6}{}\label{exr-sa-6}

Find the total surface area. Round to 1 d.p.

\begin{enumerate}
\def\labelenumi{\alph{enumi})}
\item
  \(\text{Closed cylinder: }r=2\text{ cm},\ h=11\text{ cm}\)
\item
  \(\text{Closed cylinder: }r=10\text{ cm},\ h=7\text{ cm}\)
\item
  \(\text{Closed cylinder: }r=9\text{ cm},\ h=7\text{ cm}\)
\item
  \(\text{Closed cylinder: }r=10\text{ cm},\ h=3\text{ cm}\)
\item
  \(\text{Closed cylinder: }r=10\text{ cm},\ h=6\text{ cm}\)
\item
  \(\text{Closed cylinder: }r=2\text{ cm},\ h=12\text{ cm}\)
\end{enumerate}

\begin{tcolorbox}[enhanced jigsaw, arc=.35mm, breakable, toptitle=1mm, title=\textcolor{quarto-callout-tip-color}{\faLightbulb}\hspace{0.5em}{🔍 View Solution}, colframe=quarto-callout-tip-color-frame, rightrule=.15mm, opacityback=0, titlerule=0mm, colback=white, coltitle=black, bottomtitle=1mm, toprule=.15mm, leftrule=.75mm, bottomrule=.15mm, colbacktitle=quarto-callout-tip-color!10!white, left=2mm, opacitybacktitle=0.6]

\begin{enumerate}
\def\labelenumi{\alph{enumi})}
\item
  \(163.3\,\text{cm}^{2}\)
\item
  \(1068.1\,\text{cm}^{2}\)
\item
  \(904.7\,\text{cm}^{2}\)
\item
  \(816.8\,\text{cm}^{2}\)
\item
  \(1005.3\,\text{cm}^{2}\)
\item
  \(175.9\,\text{cm}^{2}\)
\end{enumerate}

\end{tcolorbox}

{Sa formula cylinder} {Use of pi} {Curved surface area}

\end{exercise}

\begin{exercise}[Surface Area: Open-Top
Cylinder]\protect\hypertarget{exr-sa-7}{}\label{exr-sa-7}

Find the surface area (no lid). Round to 1 d.p.

\begin{enumerate}
\def\labelenumi{\alph{enumi})}
\item
  \(\text{Open-top cylinder: }r=7\text{ cm},\ h=8\text{ cm}\)
\item
  \(\text{Open-top cylinder: }r=3\text{ cm},\ h=4\text{ cm}\)
\item
  \(\text{Open-top cylinder: }r=2\text{ cm},\ h=8\text{ cm}\)
\item
  \(\text{Open-top cylinder: }r=2\text{ cm},\ h=8\text{ cm}\)
\item
  \(\text{Open-top cylinder: }r=3\text{ cm},\ h=6\text{ cm}\)
\item
  \(\text{Open-top cylinder: }r=4\text{ cm},\ h=13\text{ cm}\)
\end{enumerate}

\begin{tcolorbox}[enhanced jigsaw, arc=.35mm, breakable, toptitle=1mm, title=\textcolor{quarto-callout-tip-color}{\faLightbulb}\hspace{0.5em}{🔍 View Solution}, colframe=quarto-callout-tip-color-frame, rightrule=.15mm, opacityback=0, titlerule=0mm, colback=white, coltitle=black, bottomtitle=1mm, toprule=.15mm, leftrule=.75mm, bottomrule=.15mm, colbacktitle=quarto-callout-tip-color!10!white, left=2mm, opacitybacktitle=0.6]

\begin{enumerate}
\def\labelenumi{\alph{enumi})}
\item
  \(505.8\,\text{cm}^{2}\)
\item
  \(103.7\,\text{cm}^{2}\)
\item
  \(113.1\,\text{cm}^{2}\)
\item
  \(113.1\,\text{cm}^{2}\)
\item
  \(141.4\,\text{cm}^{2}\)
\item
  \(377.0\,\text{cm}^{2}\)
\end{enumerate}

\end{tcolorbox}

{Sa formula cylinder} {Use of pi} {Omit one face}

\end{exercise}

∑ × π ÷ √ ◆ ∞ ± ∫ ≠ Δ

\subsection{🔴 Challenging}

\begin{exercise}[Surface Area of L-Shaped
Prism]\protect\hypertarget{exr-sa-5}{}\label{exr-sa-5}

Find the total surface area of the L-prism.

\begin{enumerate}
\def\labelenumi{\alph{enumi})}
\item
  \[\text{L-prism: outer }a=4,b=9,\text{ notch }c=3,d=2,\ l=12\text{ cm}\]
\item
  \[\text{L-prism: outer }a=9,b=8,\text{ notch }c=5,d=3,\ l=10\text{ cm}\]
\item
  \[\text{L-prism: outer }a=8,b=6,\text{ notch }c=2,d=4,\ l=12\text{ cm}\]
\item
  \[\text{L-prism: outer }a=10,b=5,\text{ notch }c=3,d=4,\ l=8\text{ cm}\]
\item
  \[\text{L-prism: outer }a=6,b=5,\text{ notch }c=4,d=3,\ l=12\text{ cm}\]
\item
  \[\text{L-prism: outer }a=9,b=10,\text{ notch }c=5,d=3,\ l=14\text{ cm}\]
\end{enumerate}

\begin{tcolorbox}[enhanced jigsaw, arc=.35mm, breakable, toptitle=1mm, title=\textcolor{quarto-callout-tip-color}{\faLightbulb}\hspace{0.5em}{🔍 View Solution}, colframe=quarto-callout-tip-color-frame, rightrule=.15mm, opacityback=0, titlerule=0mm, colback=white, coltitle=black, bottomtitle=1mm, toprule=.15mm, leftrule=.75mm, bottomrule=.15mm, colbacktitle=quarto-callout-tip-color!10!white, left=2mm, opacitybacktitle=0.6]

\begin{enumerate}
\def\labelenumi{\alph{enumi})}
\item
  \[370\,\text{cm}^{2}\]
\item
  \[444\,\text{cm}^{2}\]
\item
  \[408\,\text{cm}^{2}\]
\item
  \[326\,\text{cm}^{2}\]
\item
  \[316\,\text{cm}^{2}\]
\item
  \[656\,\text{cm}^{2}\]
\end{enumerate}

\end{tcolorbox}

{Area l shape} {Sa formula prism} {Multi step calculation}

\end{exercise}

\begin{exercise}[Surface Area of
Half-Cylinder]\protect\hypertarget{exr-sa-9}{}\label{exr-sa-9}

Find the total surface area. Round to 1 d.p.

\begin{enumerate}
\def\labelenumi{\alph{enumi})}
\item
  \(\text{Half-cylinder: }r=4\text{ cm},\ l=5\text{ cm}\)
\item
  \(\text{Half-cylinder: }r=5\text{ cm},\ l=6\text{ cm}\)
\item
  \(\text{Half-cylinder: }r=3\text{ cm},\ l=12\text{ cm}\)
\item
  \(\text{Half-cylinder: }r=5\text{ cm},\ l=9\text{ cm}\)
\item
  \(\text{Half-cylinder: }r=7\text{ cm},\ l=12\text{ cm}\)
\item
  \(\text{Half-cylinder: }r=3\text{ cm},\ l=12\text{ cm}\)
\end{enumerate}

\begin{tcolorbox}[enhanced jigsaw, arc=.35mm, breakable, toptitle=1mm, title=\textcolor{quarto-callout-tip-color}{\faLightbulb}\hspace{0.5em}{🔍 View Solution}, colframe=quarto-callout-tip-color-frame, rightrule=.15mm, opacityback=0, titlerule=0mm, colback=white, coltitle=black, bottomtitle=1mm, toprule=.15mm, leftrule=.75mm, bottomrule=.15mm, colbacktitle=quarto-callout-tip-color!10!white, left=2mm, opacitybacktitle=0.6]

\begin{enumerate}
\def\labelenumi{\alph{enumi})}
\item
  \(153.1\,\text{cm}^{2}\)
\item
  \(232.7\,\text{cm}^{2}\)
\item
  \(213.4\,\text{cm}^{2}\)
\item
  \(309.9\,\text{cm}^{2}\)
\item
  \(585.8\,\text{cm}^{2}\)
\item
  \(213.4\,\text{cm}^{2}\)
\end{enumerate}

\end{tcolorbox}

{Fraction of solid} {Curved surface area} {Area rectangle} {Area
semicircle}

\end{exercise}

\begin{exercise}[SA: Prism +
Half-Cylinder]\protect\hypertarget{exr-sa-10}{}\label{exr-sa-10}

Find the total surface area.

\begin{enumerate}
\def\labelenumi{\alph{enumi})}
\item
  \[\text{House prism: }l=12,\ w=6,\ h_{\text{wall}}=7,\ r=3.0\text{ cm}\]
\item
  \[\text{House prism: }l=13,\ w=7,\ h_{\text{wall}}=3,\ r=3.5\text{ cm}\]
\item
  \[\text{House prism: }l=9,\ w=7,\ h_{\text{wall}}=5,\ r=3.5\text{ cm}\]
\item
  \[\text{House prism: }l=8,\ w=7,\ h_{\text{wall}}=6,\ r=3.5\text{ cm}\]
\item
  \[\text{House prism: }l=6,\ w=8,\ h_{\text{wall}}=4,\ r=4.0\text{ cm}\]
\item
  \[\text{House prism: }l=6,\ w=4,\ h_{\text{wall}}=4,\ r=2.0\text{ cm}\]
\end{enumerate}

\begin{tcolorbox}[enhanced jigsaw, arc=.35mm, breakable, toptitle=1mm, title=\textcolor{quarto-callout-tip-color}{\faLightbulb}\hspace{0.5em}{🔍 View Solution}, colframe=quarto-callout-tip-color-frame, rightrule=.15mm, opacityback=0, titlerule=0mm, colback=white, coltitle=black, bottomtitle=1mm, toprule=.15mm, leftrule=.75mm, bottomrule=.15mm, colbacktitle=quarto-callout-tip-color!10!white, left=2mm, opacitybacktitle=0.6]

\begin{enumerate}
\def\labelenumi{\alph{enumi})}
\item
  \[465.3\,\text{cm}^{2}\]
\item
  \[392.3\,\text{cm}^{2}\]
\item
  \[360.4\,\text{cm}^{2}\]
\item
  \[362.4\,\text{cm}^{2}\]
\item
  \[285.6\,\text{cm}^{2}\]
\item
  \[154.3\,\text{cm}^{2}\]
\end{enumerate}

\end{tcolorbox}

{Composite sa} {Multi step calculation} {Combining shapes}

\end{exercise}

\begin{exercise}[SA: Cuboid with
Hole]\protect\hypertarget{exr-sa-11}{}\label{exr-sa-11}

Find the total outer surface area.

\begin{enumerate}
\def\labelenumi{\alph{enumi})}
\item
  \[\text{Cuboid}(l=11,w=8,h=5\text{ cm) with hole }r=1\text{ cm}\]
\item
  \[\text{Cuboid}(l=7,w=8,h=5\text{ cm) with hole }r=2\text{ cm}\]
\item
  \[\text{Cuboid}(l=10,w=6,h=3\text{ cm) with hole }r=1\text{ cm}\]
\item
  \[\text{Cuboid}(l=10,w=4,h=4\text{ cm) with hole }r=1\text{ cm}\]
\item
  \[\text{Cuboid}(l=10,w=7,h=7\text{ cm) with hole }r=1\text{ cm}\]
\item
  \[\text{Cuboid}(l=11,w=4,h=4\text{ cm) with hole }r=1\text{ cm}\]
\end{enumerate}

\begin{tcolorbox}[enhanced jigsaw, arc=.35mm, breakable, toptitle=1mm, title=\textcolor{quarto-callout-tip-color}{\faLightbulb}\hspace{0.5em}{🔍 View Solution}, colframe=quarto-callout-tip-color-frame, rightrule=.15mm, opacityback=0, titlerule=0mm, colback=white, coltitle=black, bottomtitle=1mm, toprule=.15mm, leftrule=.75mm, bottomrule=.15mm, colbacktitle=quarto-callout-tip-color!10!white, left=2mm, opacitybacktitle=0.6]

\begin{enumerate}
\def\labelenumi{\alph{enumi})}
\item
  \[362.9\,\text{cm}^{2}\]
\item
  \[249.4\,\text{cm}^{2}\]
\item
  \[212.9\,\text{cm}^{2}\]
\item
  \[188.9\,\text{cm}^{2}\]
\item
  \[374.9\,\text{cm}^{2}\]
\item
  \[204.9\,\text{cm}^{2}\]
\end{enumerate}

\end{tcolorbox}

{Composite sa} {Subtraction of areas} {Use of pi}

\end{exercise}

\begin{exercise}[SA: Cylinder +
Cone]\protect\hypertarget{exr-sa-12}{}\label{exr-sa-12}

Find the total surface area. Round to 1 d.p.

\begin{enumerate}
\def\labelenumi{\alph{enumi})}
\item
  \[\text{Cylinder+cone: }d=12\text{ cm},\ h_{\text{cyl}}=7,\ h_{\text{cone}}=4\text{ cm}\]
\item
  \[\text{Cylinder+cone: }d=6\text{ cm},\ h_{\text{cyl}}=10,\ h_{\text{cone}}=6\text{ cm}\]
\item
  \[\text{Cylinder+cone: }d=6\text{ cm},\ h_{\text{cyl}}=9,\ h_{\text{cone}}=7\text{ cm}\]
\item
  \[\text{Cylinder+cone: }d=10\text{ cm},\ h_{\text{cyl}}=14,\ h_{\text{cone}}=5\text{ cm}\]
\item
  \[\text{Cylinder+cone: }d=12\text{ cm},\ h_{\text{cyl}}=9,\ h_{\text{cone}}=8\text{ cm}\]
\item
  \[\text{Cylinder+cone: }d=6\text{ cm},\ h_{\text{cyl}}=12,\ h_{\text{cone}}=5\text{ cm}\]
\end{enumerate}

\begin{tcolorbox}[enhanced jigsaw, arc=.35mm, breakable, toptitle=1mm, title=\textcolor{quarto-callout-tip-color}{\faLightbulb}\hspace{0.5em}{🔍 View Solution}, colframe=quarto-callout-tip-color-frame, rightrule=.15mm, opacityback=0, titlerule=0mm, colback=white, coltitle=black, bottomtitle=1mm, toprule=.15mm, leftrule=.75mm, bottomrule=.15mm, colbacktitle=quarto-callout-tip-color!10!white, left=2mm, opacitybacktitle=0.6]

\begin{enumerate}
\def\labelenumi{\alph{enumi})}
\item
  \[512.9\,\text{cm}^{2}\]
\item
  \[280.0\,\text{cm}^{2}\]
\item
  \[269.7\,\text{cm}^{2}\]
\item
  \[629.4\,\text{cm}^{2}\]
\item
  \[640.9\,\text{cm}^{2}\]
\item
  \[309.4\,\text{cm}^{2}\]
\end{enumerate}

\end{tcolorbox}

{Composite sa} {Cone sa formula} {Curved surface area}

\end{exercise}

∑ × π ÷ √ ◆ ∞ ± ∫ ≠ Δ

\begin{center}\rule{0.5\linewidth}{0.5pt}\end{center}

\section{🔺 Part E --- Pyramids, Cones and Spheres}\label{sec-curved}

\subsection{🟢 Easy}

\begin{exercise}[Volume of Rectangular
Pyramid]\protect\hypertarget{exr-cur-1}{}\label{exr-cur-1}

Find the volume. Round to 1 d.p. where needed.

\begin{enumerate}
\def\labelenumi{\alph{enumi})}
\item
  \(\text{Rect pyramid: }l=10,\ w=7,\ h=9\text{ cm}\)
\item
  \(\text{Rect pyramid: }l=6,\ w=10,\ h=9\text{ cm}\)
\item
  \(\text{Rect pyramid: }l=4,\ w=8,\ h=7\text{ cm}\)
\item
  \(\text{Rect pyramid: }l=11,\ w=4,\ h=6\text{ cm}\)
\item
  \(\text{Rect pyramid: }l=8,\ w=4,\ h=9\text{ cm}\)
\item
  \(\text{Rect pyramid: }l=8,\ w=6,\ h=7\text{ cm}\)
\end{enumerate}

\begin{tcolorbox}[enhanced jigsaw, arc=.35mm, breakable, toptitle=1mm, title=\textcolor{quarto-callout-tip-color}{\faLightbulb}\hspace{0.5em}{🔍 View Solution}, colframe=quarto-callout-tip-color-frame, rightrule=.15mm, opacityback=0, titlerule=0mm, colback=white, coltitle=black, bottomtitle=1mm, toprule=.15mm, leftrule=.75mm, bottomrule=.15mm, colbacktitle=quarto-callout-tip-color!10!white, left=2mm, opacitybacktitle=0.6]

\begin{enumerate}
\def\labelenumi{\alph{enumi})}
\item
  \(210.0\,\text{cm}^{3}\)
\item
  \(180.0\,\text{cm}^{3}\)
\item
  \(74.7\,\text{cm}^{3}\)
\item
  \(88.0\,\text{cm}^{3}\)
\item
  \(96.0\,\text{cm}^{3}\)
\item
  \(112.0\,\text{cm}^{3}\)
\end{enumerate}

\end{tcolorbox}

{Pyramid formula} {One third factor}

\end{exercise}

\begin{exercise}[Volume of Square
Pyramid]\protect\hypertarget{exr-cur-2}{}\label{exr-cur-2}

Find the volume.

\begin{enumerate}
\def\labelenumi{\alph{enumi})}
\item
  \(\text{Square pyramid: }s=4\text{ cm},\ h=9\text{ cm}\)
\item
  \(\text{Square pyramid: }s=9\text{ cm},\ h=14\text{ cm}\)
\item
  \(\text{Square pyramid: }s=4\text{ cm},\ h=5\text{ cm}\)
\item
  \(\text{Square pyramid: }s=8\text{ cm},\ h=14\text{ cm}\)
\item
  \(\text{Square pyramid: }s=9\text{ cm},\ h=4\text{ cm}\)
\item
  \(\text{Square pyramid: }s=6\text{ cm},\ h=9\text{ cm}\)
\end{enumerate}

\begin{tcolorbox}[enhanced jigsaw, arc=.35mm, breakable, toptitle=1mm, title=\textcolor{quarto-callout-tip-color}{\faLightbulb}\hspace{0.5em}{🔍 View Solution}, colframe=quarto-callout-tip-color-frame, rightrule=.15mm, opacityback=0, titlerule=0mm, colback=white, coltitle=black, bottomtitle=1mm, toprule=.15mm, leftrule=.75mm, bottomrule=.15mm, colbacktitle=quarto-callout-tip-color!10!white, left=2mm, opacitybacktitle=0.6]

\begin{enumerate}
\def\labelenumi{\alph{enumi})}
\item
  \(48.0\,\text{cm}^{3}\)
\item
  \(378.0\,\text{cm}^{3}\)
\item
  \(26.7\,\text{cm}^{3}\)
\item
  \(298.7\,\text{cm}^{3}\)
\item
  \(108.0\,\text{cm}^{3}\)
\item
  \(108.0\,\text{cm}^{3}\)
\end{enumerate}

\end{tcolorbox}

{Pyramid formula} {Squaring numbers} {One third factor}

\end{exercise}

\begin{exercise}[Volume of Pyramid (area
given)]\protect\hypertarget{exr-cur-3}{}\label{exr-cur-3}

Find the volume.

\begin{enumerate}
\def\labelenumi{\alph{enumi})}
\item
  \(\text{Pyramid: }A=25\,\text{cm}^{2},\ h=10\text{ cm}\)
\item
  \(\text{Pyramid: }A=14\,\text{cm}^{2},\ h=7\text{ cm}\)
\item
  \(\text{Pyramid: }A=36\,\text{cm}^{2},\ h=8\text{ cm}\)
\item
  \(\text{Pyramid: }A=45\,\text{cm}^{2},\ h=11\text{ cm}\)
\item
  \(\text{Pyramid: }A=16\,\text{cm}^{2},\ h=8\text{ cm}\)
\item
  \(\text{Pyramid: }A=47\,\text{cm}^{2},\ h=7\text{ cm}\)
\end{enumerate}

\begin{tcolorbox}[enhanced jigsaw, arc=.35mm, breakable, toptitle=1mm, title=\textcolor{quarto-callout-tip-color}{\faLightbulb}\hspace{0.5em}{🔍 View Solution}, colframe=quarto-callout-tip-color-frame, rightrule=.15mm, opacityback=0, titlerule=0mm, colback=white, coltitle=black, bottomtitle=1mm, toprule=.15mm, leftrule=.75mm, bottomrule=.15mm, colbacktitle=quarto-callout-tip-color!10!white, left=2mm, opacitybacktitle=0.6]

\begin{enumerate}
\def\labelenumi{\alph{enumi})}
\item
  \(83.3\,\text{cm}^{3}\)
\item
  \(32.7\,\text{cm}^{3}\)
\item
  \(96.0\,\text{cm}^{3}\)
\item
  \(165.0\,\text{cm}^{3}\)
\item
  \(42.7\,\text{cm}^{3}\)
\item
  \(109.7\,\text{cm}^{3}\)
\end{enumerate}

\end{tcolorbox}

{Pyramid formula} {One third factor}

\end{exercise}

∑ × π ÷ √ ◆ ∞ ± ∫ ≠ Δ

\subsection{🟡 Medium}

\begin{exercise}[Find Missing Pyramid
Dimension]\protect\hypertarget{exr-cur-4}{}\label{exr-cur-4}

Find the missing dimension.

\begin{enumerate}
\def\labelenumi{\alph{enumi})}
\item
  \[\text{Pyramid: }V=50.0\,\text{cm}^{3},\ A=15\,\text{cm}^{2}\text{. Find }h.\]
\item
  \[\text{Pyramid: }V=60.0\,\text{cm}^{3},\ A=15\,\text{cm}^{2}\text{. Find }h.\]
\item
  \[\text{Pyramid: }V=130.0\,\text{cm}^{3},\ A=39\,\text{cm}^{2}\text{. Find }h.\]
\item
  \[\text{Pyramid: }V=58.7\,\text{cm}^{3},\ A=16\,\text{cm}^{2}\text{. Find }h.\]
\item
  \[\text{Pyramid: }V=45.0\,\text{cm}^{3},\ A=27\,\text{cm}^{2}\text{. Find }h.\]
\item
  \[\text{Pyramid: }V=124.0\,\text{cm}^{3},\ A=31\,\text{cm}^{2}\text{. Find }h.\]
\end{enumerate}

\begin{tcolorbox}[enhanced jigsaw, arc=.35mm, breakable, toptitle=1mm, title=\textcolor{quarto-callout-tip-color}{\faLightbulb}\hspace{0.5em}{🔍 View Solution}, colframe=quarto-callout-tip-color-frame, rightrule=.15mm, opacityback=0, titlerule=0mm, colback=white, coltitle=black, bottomtitle=1mm, toprule=.15mm, leftrule=.75mm, bottomrule=.15mm, colbacktitle=quarto-callout-tip-color!10!white, left=2mm, opacitybacktitle=0.6]

\begin{enumerate}
\def\labelenumi{\alph{enumi})}
\item
  \[h=10\text{ cm}\]
\item
  \[h=12\text{ cm}\]
\item
  \[h=10\text{ cm}\]
\item
  \[h=11\text{ cm}\]
\item
  \[h=5\text{ cm}\]
\item
  \[h=12\text{ cm}\]
\end{enumerate}

\end{tcolorbox}

{Rearranging formula} {Pyramid formula}

\end{exercise}

\begin{exercise}[Volume of a
Cone]\protect\hypertarget{exr-cur-5}{}\label{exr-cur-5}

Find the volume. Round to 3 s.f.

\begin{enumerate}
\def\labelenumi{\alph{enumi})}
\item
  \(\text{Cone: }r=4\text{ cm},\ h=4\text{ cm}\)
\item
  \(\text{Cone: }r=7\text{ cm},\ h=4\text{ cm}\)
\item
  \(\text{Cone: }r=8\text{ cm},\ h=11\text{ cm}\)
\item
  \(\text{Cone: }r=10\text{ cm},\ h=16\text{ cm}\)
\item
  \(\text{Cone: }r=5\text{ cm},\ h=6\text{ cm}\)
\item
  \(\text{Cone: }r=3\text{ cm},\ h=6\text{ cm}\)
\end{enumerate}

\begin{tcolorbox}[enhanced jigsaw, arc=.35mm, breakable, toptitle=1mm, title=\textcolor{quarto-callout-tip-color}{\faLightbulb}\hspace{0.5em}{🔍 View Solution}, colframe=quarto-callout-tip-color-frame, rightrule=.15mm, opacityback=0, titlerule=0mm, colback=white, coltitle=black, bottomtitle=1mm, toprule=.15mm, leftrule=.75mm, bottomrule=.15mm, colbacktitle=quarto-callout-tip-color!10!white, left=2mm, opacitybacktitle=0.6]

\begin{enumerate}
\def\labelenumi{\alph{enumi})}
\item
  \(67.0\,\text{cm}^{3}\)
\item
  \(205.3\,\text{cm}^{3}\)
\item
  \(737.2\,\text{cm}^{3}\)
\item
  \(1675.5\,\text{cm}^{3}\)
\item
  \(157.1\,\text{cm}^{3}\)
\item
  \(56.5\,\text{cm}^{3}\)
\end{enumerate}

\end{tcolorbox}

{Cone formula} {Use of pi} {One third factor}

\end{exercise}

\begin{exercise}[Volume of a
Sphere]\protect\hypertarget{exr-cur-7}{}\label{exr-cur-7}

Find the volume. Round to 3 s.f.

\begin{enumerate}
\def\labelenumi{\alph{enumi})}
\item
  \(\text{Sphere: }r=6\text{ cm}\)
\item
  \(\text{Sphere: }r=10\text{ cm}\)
\item
  \(\text{Sphere: }r=8\text{ cm}\)
\item
  \(\text{Sphere: }r=9\text{ cm}\)
\item
  \(\text{Sphere: }r=8\text{ cm}\)
\item
  \(\text{Sphere: }r=8\text{ cm}\)
\end{enumerate}

\begin{tcolorbox}[enhanced jigsaw, arc=.35mm, breakable, toptitle=1mm, title=\textcolor{quarto-callout-tip-color}{\faLightbulb}\hspace{0.5em}{🔍 View Solution}, colframe=quarto-callout-tip-color-frame, rightrule=.15mm, opacityback=0, titlerule=0mm, colback=white, coltitle=black, bottomtitle=1mm, toprule=.15mm, leftrule=.75mm, bottomrule=.15mm, colbacktitle=quarto-callout-tip-color!10!white, left=2mm, opacitybacktitle=0.6]

\begin{enumerate}
\def\labelenumi{\alph{enumi})}
\item
  \(904.8\,\text{cm}^{3}\)
\item
  \(4188.8\,\text{cm}^{3}\)
\item
  \(2144.7\,\text{cm}^{3}\)
\item
  \(3053.6\,\text{cm}^{3}\)
\item
  \(2144.7\,\text{cm}^{3}\)
\item
  \(2144.7\,\text{cm}^{3}\)
\end{enumerate}

\end{tcolorbox}

{Sphere formula} {Use of pi} {Cubing numbers}

\end{exercise}

\begin{exercise}[Sphere Volume (diameter
given)]\protect\hypertarget{exr-cur-8}{}\label{exr-cur-8}

Find the volume. Round to 1 d.p.

\begin{enumerate}
\def\labelenumi{\alph{enumi})}
\item
  \(\text{Sphere: }d=11\text{ cm}\)
\item
  \(\text{Sphere: }d=5\text{ cm}\)
\item
  \(\text{Sphere: }d=10\text{ cm}\)
\item
  \(\text{Sphere: }d=5\text{ cm}\)
\item
  \(\text{Sphere: }d=9\text{ cm}\)
\item
  \(\text{Sphere: }d=4\text{ cm}\)
\end{enumerate}

\begin{tcolorbox}[enhanced jigsaw, arc=.35mm, breakable, toptitle=1mm, title=\textcolor{quarto-callout-tip-color}{\faLightbulb}\hspace{0.5em}{🔍 View Solution}, colframe=quarto-callout-tip-color-frame, rightrule=.15mm, opacityback=0, titlerule=0mm, colback=white, coltitle=black, bottomtitle=1mm, toprule=.15mm, leftrule=.75mm, bottomrule=.15mm, colbacktitle=quarto-callout-tip-color!10!white, left=2mm, opacitybacktitle=0.6]

\begin{enumerate}
\def\labelenumi{\alph{enumi})}
\item
  \(696.9\,\text{cm}^{3}\)
\item
  \(65.4\,\text{cm}^{3}\)
\item
  \(523.6\,\text{cm}^{3}\)
\item
  \(65.4\,\text{cm}^{3}\)
\item
  \(381.7\,\text{cm}^{3}\)
\item
  \(33.5\,\text{cm}^{3}\)
\end{enumerate}

\end{tcolorbox}

{Sphere formula} {Diameter to radius}

\end{exercise}

\begin{exercise}[Volume of a
Hemisphere]\protect\hypertarget{exr-cur-9}{}\label{exr-cur-9}

Find the volume. Round to 1 d.p.

\begin{enumerate}
\def\labelenumi{\alph{enumi})}
\item
  \(\text{Hemisphere: }r=3\text{ cm}\)
\item
  \(\text{Hemisphere: }r=4\text{ cm}\)
\item
  \(\text{Hemisphere: }r=10\text{ cm}\)
\item
  \(\text{Hemisphere: }r=3\text{ cm}\)
\item
  \(\text{Hemisphere: }r=7\text{ cm}\)
\item
  \(\text{Hemisphere: }r=8\text{ cm}\)
\end{enumerate}

\begin{tcolorbox}[enhanced jigsaw, arc=.35mm, breakable, toptitle=1mm, title=\textcolor{quarto-callout-tip-color}{\faLightbulb}\hspace{0.5em}{🔍 View Solution}, colframe=quarto-callout-tip-color-frame, rightrule=.15mm, opacityback=0, titlerule=0mm, colback=white, coltitle=black, bottomtitle=1mm, toprule=.15mm, leftrule=.75mm, bottomrule=.15mm, colbacktitle=quarto-callout-tip-color!10!white, left=2mm, opacitybacktitle=0.6]

\begin{enumerate}
\def\labelenumi{\alph{enumi})}
\item
  \(56.5\,\text{cm}^{3}\)
\item
  \(134.0\,\text{cm}^{3}\)
\item
  \(2094.4\,\text{cm}^{3}\)
\item
  \(56.5\,\text{cm}^{3}\)
\item
  \(718.4\,\text{cm}^{3}\)
\item
  \(1072.3\,\text{cm}^{3}\)
\end{enumerate}

\end{tcolorbox}

{Sphere formula} {Fraction of solid}

\end{exercise}

\begin{exercise}[Surface Area of a
Sphere]\protect\hypertarget{exr-cur-11}{}\label{exr-cur-11}

Find the surface area. Round to 1 d.p.

\begin{enumerate}
\def\labelenumi{\alph{enumi})}
\item
  \(\text{Sphere: }r=4\text{ cm. Find surface area.}\)
\item
  \(\text{Sphere: }r=4\text{ cm. Find surface area.}\)
\item
  \(\text{Sphere: }r=2\text{ cm. Find surface area.}\)
\item
  \(\text{Sphere: }r=4\text{ cm. Find surface area.}\)
\item
  \(\text{Sphere: }r=2\text{ cm. Find surface area.}\)
\item
  \(\text{Sphere: }r=6\text{ cm. Find surface area.}\)
\end{enumerate}

\begin{tcolorbox}[enhanced jigsaw, arc=.35mm, breakable, toptitle=1mm, title=\textcolor{quarto-callout-tip-color}{\faLightbulb}\hspace{0.5em}{🔍 View Solution}, colframe=quarto-callout-tip-color-frame, rightrule=.15mm, opacityback=0, titlerule=0mm, colback=white, coltitle=black, bottomtitle=1mm, toprule=.15mm, leftrule=.75mm, bottomrule=.15mm, colbacktitle=quarto-callout-tip-color!10!white, left=2mm, opacitybacktitle=0.6]

\begin{enumerate}
\def\labelenumi{\alph{enumi})}
\item
  \(201.1\,\text{cm}^{2}\)
\item
  \(201.1\,\text{cm}^{2}\)
\item
  \(50.3\,\text{cm}^{2}\)
\item
  \(201.1\,\text{cm}^{2}\)
\item
  \(50.3\,\text{cm}^{2}\)
\item
  \(452.4\,\text{cm}^{2}\)
\end{enumerate}

\end{tcolorbox}

{Sphere sa formula} {Use of pi}

\end{exercise}

\begin{exercise}[Surface Area of Sphere (diameter
given)]\protect\hypertarget{exr-cur-12}{}\label{exr-cur-12}

Find the surface area. Round to 1 d.p.

\begin{enumerate}
\def\labelenumi{\alph{enumi})}
\item
  \(\text{Sphere: diameter }d=15\text{ cm. Find surface area.}\)
\item
  \(\text{Sphere: diameter }d=4\text{ cm. Find surface area.}\)
\item
  \(\text{Sphere: diameter }d=9\text{ cm. Find surface area.}\)
\item
  \(\text{Sphere: diameter }d=20\text{ cm. Find surface area.}\)
\item
  \(\text{Sphere: diameter }d=5\text{ cm. Find surface area.}\)
\item
  \(\text{Sphere: diameter }d=11\text{ cm. Find surface area.}\)
\end{enumerate}

\begin{tcolorbox}[enhanced jigsaw, arc=.35mm, breakable, toptitle=1mm, title=\textcolor{quarto-callout-tip-color}{\faLightbulb}\hspace{0.5em}{🔍 View Solution}, colframe=quarto-callout-tip-color-frame, rightrule=.15mm, opacityback=0, titlerule=0mm, colback=white, coltitle=black, bottomtitle=1mm, toprule=.15mm, leftrule=.75mm, bottomrule=.15mm, colbacktitle=quarto-callout-tip-color!10!white, left=2mm, opacitybacktitle=0.6]

\begin{enumerate}
\def\labelenumi{\alph{enumi})}
\item
  \(706.9\,\text{cm}^{2}\)
\item
  \(50.3\,\text{cm}^{2}\)
\item
  \(254.5\,\text{cm}^{2}\)
\item
  \(1256.6\,\text{cm}^{2}\)
\item
  \(78.5\,\text{cm}^{2}\)
\item
  \(380.1\,\text{cm}^{2}\)
\end{enumerate}

\end{tcolorbox}

{Sphere sa formula} {Diameter to radius} {Use of pi}

\end{exercise}

\begin{exercise}[Surface Area of a
Hemisphere]\protect\hypertarget{exr-cur-13}{}\label{exr-cur-13}

Find the total surface area. Round to 1 d.p.

\begin{enumerate}
\def\labelenumi{\alph{enumi})}
\item
  \(\text{Hemisphere: }r=12\text{ cm. Find total surface area.}\)
\item
  \(\text{Hemisphere: }r=7\text{ cm. Find total surface area.}\)
\item
  \(\text{Hemisphere: }r=12\text{ cm. Find total surface area.}\)
\item
  \(\text{Hemisphere: }r=3\text{ cm. Find total surface area.}\)
\item
  \(\text{Hemisphere: }r=5\text{ cm. Find total surface area.}\)
\item
  \(\text{Hemisphere: }r=10\text{ cm. Find total surface area.}\)
\end{enumerate}

\begin{tcolorbox}[enhanced jigsaw, arc=.35mm, breakable, toptitle=1mm, title=\textcolor{quarto-callout-tip-color}{\faLightbulb}\hspace{0.5em}{🔍 View Solution}, colframe=quarto-callout-tip-color-frame, rightrule=.15mm, opacityback=0, titlerule=0mm, colback=white, coltitle=black, bottomtitle=1mm, toprule=.15mm, leftrule=.75mm, bottomrule=.15mm, colbacktitle=quarto-callout-tip-color!10!white, left=2mm, opacitybacktitle=0.6]

\begin{enumerate}
\def\labelenumi{\alph{enumi})}
\item
  \(1357.2\,\text{cm}^{2}\)
\item
  \(461.8\,\text{cm}^{2}\)
\item
  \(1357.2\,\text{cm}^{2}\)
\item
  \(84.8\,\text{cm}^{2}\)
\item
  \(235.6\,\text{cm}^{2}\)
\item
  \(942.5\,\text{cm}^{2}\)
\end{enumerate}

\end{tcolorbox}

{Sphere sa formula} {Fraction of solid} {Area circle}

\end{exercise}

\begin{exercise}[SA of Cone (slant height
given)]\protect\hypertarget{exr-cur-15}{}\label{exr-cur-15}

Find the total surface area. Round to 1 d.p.

\begin{enumerate}
\def\labelenumi{\alph{enumi})}
\item
  \[\text{Cone: base radius }r=5\text{ cm},\ \text{slant height }l=11\text{ cm.}\ \text{Find total SA.}\]
\item
  \[\text{Cone: base radius }r=5\text{ cm},\ \text{slant height }l=6\text{ cm.}\ \text{Find total SA.}\]
\item
  \[\text{Cone: base radius }r=4\text{ cm},\ \text{slant height }l=13\text{ cm.}\ \text{Find total SA.}\]
\item
  \[\text{Cone: base radius }r=9\text{ cm},\ \text{slant height }l=12\text{ cm.}\ \text{Find total SA.}\]
\item
  \[\text{Cone: base radius }r=5\text{ cm},\ \text{slant height }l=7\text{ cm.}\ \text{Find total SA.}\]
\item
  \[\text{Cone: base radius }r=9\text{ cm},\ \text{slant height }l=10\text{ cm.}\ \text{Find total SA.}\]
\end{enumerate}

\begin{tcolorbox}[enhanced jigsaw, arc=.35mm, breakable, toptitle=1mm, title=\textcolor{quarto-callout-tip-color}{\faLightbulb}\hspace{0.5em}{🔍 View Solution}, colframe=quarto-callout-tip-color-frame, rightrule=.15mm, opacityback=0, titlerule=0mm, colback=white, coltitle=black, bottomtitle=1mm, toprule=.15mm, leftrule=.75mm, bottomrule=.15mm, colbacktitle=quarto-callout-tip-color!10!white, left=2mm, opacitybacktitle=0.6]

\begin{enumerate}
\def\labelenumi{\alph{enumi})}
\item
  \[251.3\,\text{cm}^{2}\]
\item
  \[172.8\,\text{cm}^{2}\]
\item
  \[213.6\,\text{cm}^{2}\]
\item
  \[593.8\,\text{cm}^{2}\]
\item
  \[188.5\,\text{cm}^{2}\]
\item
  \[537.2\,\text{cm}^{2}\]
\end{enumerate}

\end{tcolorbox}

{Cone sa formula} {Use of pi}

\end{exercise}

∑ × π ÷ √ ◆ ∞ ± ∫ ≠ Δ

\subsection{🔴 Challenging}

\begin{exercise}[Cone Volume using
Pythagoras]\protect\hypertarget{exr-cur-6}{}\label{exr-cur-6}

Find the volume. Round to 3 s.f.

\begin{enumerate}
\def\labelenumi{\alph{enumi})}
\item
  \[\text{Cone: }r=5\text{ cm},\ l=11.18\text{ cm. Find volume.}\]
\item
  \[\text{Cone: }r=4\text{ cm},\ l=9.849\text{ cm. Find volume.}\]
\item
  \[\text{Cone: }r=3\text{ cm},\ l=14.318\text{ cm. Find volume.}\]
\item
  \[\text{Cone: }r=7\text{ cm},\ l=8.062\text{ cm. Find volume.}\]
\item
  \[\text{Cone: }r=3\text{ cm},\ l=14.318\text{ cm. Find volume.}\]
\item
  \[\text{Cone: }r=4\text{ cm},\ l=14.56\text{ cm. Find volume.}\]
\end{enumerate}

\begin{tcolorbox}[enhanced jigsaw, arc=.35mm, breakable, toptitle=1mm, title=\textcolor{quarto-callout-tip-color}{\faLightbulb}\hspace{0.5em}{🔍 View Solution}, colframe=quarto-callout-tip-color-frame, rightrule=.15mm, opacityback=0, titlerule=0mm, colback=white, coltitle=black, bottomtitle=1mm, toprule=.15mm, leftrule=.75mm, bottomrule=.15mm, colbacktitle=quarto-callout-tip-color!10!white, left=2mm, opacitybacktitle=0.6]

\begin{enumerate}
\def\labelenumi{\alph{enumi})}
\item
  \[261.8\,\text{cm}^{3}\]
\item
  \[150.8\,\text{cm}^{3}\]
\item
  \[131.9\,\text{cm}^{3}\]
\item
  \[205.3\,\text{cm}^{3}\]
\item
  \[131.9\,\text{cm}^{3}\]
\item
  \[234.6\,\text{cm}^{3}\]
\end{enumerate}

\end{tcolorbox}

{Cone formula} {Pythagoras theorem}

\end{exercise}

\begin{exercise}[Composite: Cone +
Hemisphere]\protect\hypertarget{exr-cur-10}{}\label{exr-cur-10}

Find the total volume. Round to 3 s.f.

\begin{enumerate}
\def\labelenumi{\alph{enumi})}
\item
  \[\text{Hemisphere}(r=5)+\text{Cone}(r=5,h=13)\text{ cm}\]
\item
  \[\text{Hemisphere}(r=7)+\text{Cone}(r=7,h=12)\text{ cm}\]
\item
  \[\text{Hemisphere}(r=3)+\text{Cone}(r=3,h=9)\text{ cm}\]
\item
  \[\text{Hemisphere}(r=4)+\text{Cone}(r=4,h=11)\text{ cm}\]
\item
  \[\text{Hemisphere}(r=4)+\text{Cone}(r=4,h=6)\text{ cm}\]
\item
  \[\text{Hemisphere}(r=3)+\text{Cone}(r=3,h=5)\text{ cm}\]
\end{enumerate}

\begin{tcolorbox}[enhanced jigsaw, arc=.35mm, breakable, toptitle=1mm, title=\textcolor{quarto-callout-tip-color}{\faLightbulb}\hspace{0.5em}{🔍 View Solution}, colframe=quarto-callout-tip-color-frame, rightrule=.15mm, opacityback=0, titlerule=0mm, colback=white, coltitle=black, bottomtitle=1mm, toprule=.15mm, leftrule=.75mm, bottomrule=.15mm, colbacktitle=quarto-callout-tip-color!10!white, left=2mm, opacitybacktitle=0.6]

\begin{enumerate}
\def\labelenumi{\alph{enumi})}
\item
  \[602.1\,\text{cm}^{3}\]
\item
  \[1334.2\,\text{cm}^{3}\]
\item
  \[141.3\,\text{cm}^{3}\]
\item
  \[318.3\,\text{cm}^{3}\]
\item
  \[234.5\,\text{cm}^{3}\]
\item
  \[103.6\,\text{cm}^{3}\]
\end{enumerate}

\end{tcolorbox}

{Composite volume} {Cone formula} {Sphere formula}

\end{exercise}

\begin{exercise}[Surface Area of a
Cone]\protect\hypertarget{exr-cur-14}{}\label{exr-cur-14}

Find the total surface area. Round to 1 d.p.

\begin{enumerate}
\def\labelenumi{\alph{enumi})}
\item
  \[\text{Cone: }r=6\text{ cm},\ h=10\text{ cm.}\ \text{Find total SA.}\]
\item
  \[\text{Cone: }r=5\text{ cm},\ h=12\text{ cm.}\ \text{Find total SA.}\]
\item
  \[\text{Cone: }r=4\text{ cm},\ h=7\text{ cm.}\ \text{Find total SA.}\]
\item
  \[\text{Cone: }r=9\text{ cm},\ h=4\text{ cm.}\ \text{Find total SA.}\]
\item
  \[\text{Cone: }r=7\text{ cm},\ h=13\text{ cm.}\ \text{Find total SA.}\]
\item
  \[\text{Cone: }r=5\text{ cm},\ h=14\text{ cm.}\ \text{Find total SA.}\]
\end{enumerate}

\begin{tcolorbox}[enhanced jigsaw, arc=.35mm, breakable, toptitle=1mm, title=\textcolor{quarto-callout-tip-color}{\faLightbulb}\hspace{0.5em}{🔍 View Solution}, colframe=quarto-callout-tip-color-frame, rightrule=.15mm, opacityback=0, titlerule=0mm, colback=white, coltitle=black, bottomtitle=1mm, toprule=.15mm, leftrule=.75mm, bottomrule=.15mm, colbacktitle=quarto-callout-tip-color!10!white, left=2mm, opacitybacktitle=0.6]

\begin{enumerate}
\def\labelenumi{\alph{enumi})}
\item
  \[332.9\,\text{cm}^{2}\]
\item
  \[282.7\,\text{cm}^{2}\]
\item
  \[151.6\,\text{cm}^{2}\]
\item
  \[533.0\,\text{cm}^{2}\]
\item
  \[478.5\,\text{cm}^{2}\]
\item
  \[312.1\,\text{cm}^{2}\]
\end{enumerate}

\end{tcolorbox}

{Cone sa formula} {Pythagoras theorem}

\end{exercise}

∑ × π ÷ √ ◆ ∞ ± ∫ ≠ Δ

\begin{center}\rule{0.5\linewidth}{0.5pt}\end{center}

\section{🔣 Part F --- Algebraic Volume}\label{sec-algebra}

\begin{tcolorbox}[enhanced jigsaw, arc=.35mm, breakable, toptitle=1mm, title={📌 What to expect}, colframe=quarto-callout-note-color-frame, rightrule=.15mm, opacityback=0, titlerule=0mm, colback=white, coltitle=black, bottomtitle=1mm, toprule=.15mm, leftrule=.75mm, bottomrule=.15mm, colbacktitle=quarto-callout-note-color!10!white, left=2mm, opacitybacktitle=0.6]

These questions give dimensions as algebraic expressions (e.g.~\(2x\),
\(x+3\)). Apply the volume formula, expand brackets, and simplify using
index laws.

\end{tcolorbox}

\subsection{🔴 Challenging}

\begin{exercise}[Algebraic Volume --- Rectangular
Prism]\protect\hypertarget{exr-alg-1}{}\label{exr-alg-1}

Express V in terms of x.

\begin{enumerate}
\def\labelenumi{\alph{enumi})}
\item
  \[\text{Rect prism: }l=4x,\ w=3x,\ h=7\text{ cm.}\ \text{Express }V.\]
\item
  \[\text{Rect prism: }l=3x,\ w=3x,\ h=4\text{ cm.}\ \text{Express }V.\]
\item
  \[\text{Rect prism: }l=5x,\ w=3x,\ h=5\text{ cm.}\ \text{Express }V.\]
\item
  \[\text{Rect prism: }l=2x,\ w=2x,\ h=8\text{ cm.}\ \text{Express }V.\]
\item
  \[\text{Rect prism: }l=3x,\ w=2x,\ h=6\text{ cm.}\ \text{Express }V.\]
\item
  \[\text{Rect prism: }l=2x,\ w=2x,\ h=8\text{ cm.}\ \text{Express }V.\]
\end{enumerate}

\begin{tcolorbox}[enhanced jigsaw, arc=.35mm, breakable, toptitle=1mm, title=\textcolor{quarto-callout-tip-color}{\faLightbulb}\hspace{0.5em}{🔍 View Solution}, colframe=quarto-callout-tip-color-frame, rightrule=.15mm, opacityback=0, titlerule=0mm, colback=white, coltitle=black, bottomtitle=1mm, toprule=.15mm, leftrule=.75mm, bottomrule=.15mm, colbacktitle=quarto-callout-tip-color!10!white, left=2mm, opacitybacktitle=0.6]

\begin{enumerate}
\def\labelenumi{\alph{enumi})}
\item
  \[V=84x^2\,\text{cm}^{3}\]
\item
  \[V=36x^2\,\text{cm}^{3}\]
\item
  \[V=75x^2\,\text{cm}^{3}\]
\item
  \[V=32x^2\,\text{cm}^{3}\]
\item
  \[V=36x^2\,\text{cm}^{3}\]
\item
  \[V=32x^2\,\text{cm}^{3}\]
\end{enumerate}

\end{tcolorbox}

{Algebraic expressions} {Index laws}

\end{exercise}

\begin{exercise}[Algebraic Volume --- Triangular
Prism]\protect\hypertarget{exr-alg-2}{}\label{exr-alg-2}

Find V in terms of x. Simplify.

\begin{enumerate}
\def\labelenumi{\alph{enumi})}
\item
  \[\text{Tri-prism: base}=2x,\ h=x,\ \text{length}=(x+8)\text{ cm.}\ \text{Find }V.\]
\item
  \[\text{Tri-prism: base}=2x,\ h=x,\ \text{length}=(x+2)\text{ cm.}\ \text{Find }V.\]
\item
  \[\text{Tri-prism: base}=2x,\ h=x,\ \text{length}=(x+5)\text{ cm.}\ \text{Find }V.\]
\item
  \[\text{Tri-prism: base}=2x,\ h=x,\ \text{length}=(x+8)\text{ cm.}\ \text{Find }V.\]
\item
  \[\text{Tri-prism: base}=2x,\ h=x,\ \text{length}=(x+7)\text{ cm.}\ \text{Find }V.\]
\item
  \[\text{Tri-prism: base}=2x,\ h=x,\ \text{length}=(x+3)\text{ cm.}\ \text{Find }V.\]
\end{enumerate}

\begin{tcolorbox}[enhanced jigsaw, arc=.35mm, breakable, toptitle=1mm, title=\textcolor{quarto-callout-tip-color}{\faLightbulb}\hspace{0.5em}{🔍 View Solution}, colframe=quarto-callout-tip-color-frame, rightrule=.15mm, opacityback=0, titlerule=0mm, colback=white, coltitle=black, bottomtitle=1mm, toprule=.15mm, leftrule=.75mm, bottomrule=.15mm, colbacktitle=quarto-callout-tip-color!10!white, left=2mm, opacitybacktitle=0.6]

\begin{enumerate}
\def\labelenumi{\alph{enumi})}
\item
  \[V=x^3+8x^2\]
\item
  \[V=x^3+2x^2\]
\item
  \[V=x^3+5x^2\]
\item
  \[V=x^3+8x^2\]
\item
  \[V=x^3+7x^2\]
\item
  \[V=x^3+3x^2\]
\end{enumerate}

\end{tcolorbox}

{Expanding brackets} {Volume formula prism}

\end{exercise}

\begin{exercise}[Algebraic Volume ---
Cylinder]\protect\hypertarget{exr-alg-3}{}\label{exr-alg-3}

Express V in terms of x. Leave π in answer.

\begin{enumerate}
\def\labelenumi{\alph{enumi})}
\item
  \[\text{Cylinder: }r=1x\text{ cm},\ h=4x\text{ cm.}\ \text{Express }V.\]
\item
  \[\text{Cylinder: }r=3x\text{ cm},\ h=3x\text{ cm.}\ \text{Express }V.\]
\item
  \[\text{Cylinder: }r=3x\text{ cm},\ h=2x\text{ cm.}\ \text{Express }V.\]
\item
  \[\text{Cylinder: }r=3x\text{ cm},\ h=2x\text{ cm.}\ \text{Express }V.\]
\item
  \[\text{Cylinder: }r=2x\text{ cm},\ h=3x\text{ cm.}\ \text{Express }V.\]
\item
  \[\text{Cylinder: }r=1x\text{ cm},\ h=3x\text{ cm.}\ \text{Express }V.\]
\end{enumerate}

\begin{tcolorbox}[enhanced jigsaw, arc=.35mm, breakable, toptitle=1mm, title=\textcolor{quarto-callout-tip-color}{\faLightbulb}\hspace{0.5em}{🔍 View Solution}, colframe=quarto-callout-tip-color-frame, rightrule=.15mm, opacityback=0, titlerule=0mm, colback=white, coltitle=black, bottomtitle=1mm, toprule=.15mm, leftrule=.75mm, bottomrule=.15mm, colbacktitle=quarto-callout-tip-color!10!white, left=2mm, opacitybacktitle=0.6]

\begin{enumerate}
\def\labelenumi{\alph{enumi})}
\item
  \[V=4\pi x^3\,\text{cm}^{3}\]
\item
  \[V=27\pi x^3\,\text{cm}^{3}\]
\item
  \[V=18\pi x^3\,\text{cm}^{3}\]
\item
  \[V=18\pi x^3\,\text{cm}^{3}\]
\item
  \[V=12\pi x^3\,\text{cm}^{3}\]
\item
  \[V=3\pi x^3\,\text{cm}^{3}\]
\end{enumerate}

\end{tcolorbox}

{Algebraic expressions} {Cylinder formula}

\end{exercise}

\begin{exercise}[Solve for x given
Volume]\protect\hypertarget{exr-alg-4}{}\label{exr-alg-4}

Find the value of x.

\begin{enumerate}
\def\labelenumi{\alph{enumi})}
\item
  \[\text{Prism: }V=192\,\text{cm}^{3},\ l=x,\ w=x,\ h=3x.\ \text{Find }x.\]
\item
  \[\text{Prism: }V=128\,\text{cm}^{3},\ l=x,\ w=x,\ h=2x.\ \text{Find }x.\]
\item
  \[\text{Prism: }V=81\,\text{cm}^{3},\ l=x,\ w=x,\ h=3x.\ \text{Find }x.\]
\item
  \[\text{Prism: }V=192\,\text{cm}^{3},\ l=x,\ w=x,\ h=3x.\ \text{Find }x.\]
\item
  \[\text{Prism: }V=162\,\text{cm}^{3},\ l=x,\ w=x,\ h=6x.\ \text{Find }x.\]
\item
  \[\text{Prism: }V=48\,\text{cm}^{3},\ l=x,\ w=x,\ h=6x.\ \text{Find }x.\]
\end{enumerate}

\begin{tcolorbox}[enhanced jigsaw, arc=.35mm, breakable, toptitle=1mm, title=\textcolor{quarto-callout-tip-color}{\faLightbulb}\hspace{0.5em}{🔍 View Solution}, colframe=quarto-callout-tip-color-frame, rightrule=.15mm, opacityback=0, titlerule=0mm, colback=white, coltitle=black, bottomtitle=1mm, toprule=.15mm, leftrule=.75mm, bottomrule=.15mm, colbacktitle=quarto-callout-tip-color!10!white, left=2mm, opacitybacktitle=0.6]

\begin{enumerate}
\def\labelenumi{\alph{enumi})}
\item
  \[x=4\text{ cm}\]
\item
  \[x=4\text{ cm}\]
\item
  \[x=3\text{ cm}\]
\item
  \[x=4\text{ cm}\]
\item
  \[x=3\text{ cm}\]
\item
  \[x=2\text{ cm}\]
\end{enumerate}

\end{tcolorbox}

{Cube root} {Rearranging formula}

\end{exercise}

∑ × π ÷ √ ◆ ∞ ± ∫ ≠ Δ

\begin{center}\rule{0.5\linewidth}{0.5pt}\end{center}

\section{⚡ Test Your Knowledge}\label{sec-mensuration-worksheet}

\chapter{Unit 4: Statistics}\label{unit-4-statistics}

\section{🎯 Learning Targets}\label{learning-targets-3}

By the end of this unit, you should be able to:

\textbf{LT1 --- Averages and Range from a List}

\begin{itemize}
\tightlist
\item
  Calculate the \textbf{mean}, \textbf{median}, \textbf{mode} and
  \textbf{range} from a list of values
\item
  Find a \textbf{missing value} given the mean of a dataset
\item
  Find the \textbf{new mean} when a value is added to a dataset
\item
  Find two unknown values given their \textbf{ratio} and the overall
  mean
\item
  Read and interpret a \textbf{stem and leaf diagram} to find averages
  and range
\end{itemize}

\textbf{LT2 --- Frequency Tables (Ungrouped)}

\begin{itemize}
\tightlist
\item
  Calculate the \textbf{mean} from a frequency table using
  \(\bar{x} = \dfrac{\sum fx}{\sum f}\)
\item
  Find the \textbf{median} from a frequency table using cumulative
  frequencies
\item
  Identify the \textbf{mode} from a frequency table
\item
  Find the largest and smallest possible value of an unknown frequency
  \(x\) given a target \textbf{median}
\end{itemize}

\textbf{LT3 --- Grouped Frequency Tables}

\begin{itemize}
\tightlist
\item
  Identify the \textbf{modal class} from a grouped frequency table
\item
  Identify the \textbf{median class} from a grouped frequency table
\item
  \textbf{Estimate the mean} using midpoints:
  \(\bar{x} \approx \dfrac{\sum fx}{\sum f}\)
\item
  Calculate a \textbf{combined mean} when two groups with different
  sizes are merged
\item
  Find \textbf{missing frequencies} \(a\) and \(b\) given total
  frequency and estimated mean
\end{itemize}

\textbf{LT4 --- Cumulative Frequency}

\begin{itemize}
\tightlist
\item
  Build a \textbf{cumulative frequency table} from a frequency table
\item
  Read values from a cumulative frequency table (how many below / above
  a value)
\item
  Use a \textbf{cumulative frequency graph} to estimate \(Q_1\), \(Q_2\)
  (median) and \(Q_3\)
\item
  Estimate the \textbf{IQR} from a cumulative frequency graph
\item
  Estimate the number of items \textbf{above a given value}
\item
  Find missing values \textbf{A} and \textbf{B} in a cumulative
  frequency table
\end{itemize}

\textbf{LT5 --- Box Plots and IQR}

\begin{itemize}
\tightlist
\item
  Calculate the \textbf{IQR} from raw data: \(\text{IQR} = Q_3 - Q_1\)
\item
  Draw a \textbf{box-and-whisker plot} from a five-number summary
\item
  Identify \textbf{outliers} using the rule: value
  \(> Q_3 + 1.5 \times \text{IQR}\) or \(< Q_1 - 1.5 \times \text{IQR}\)
\item
  Find unknown values in an ordered list given conditions on median,
  mode and range
\end{itemize}

\begin{center}\rule{0.5\linewidth}{0.5pt}\end{center}

\section{🗂️ Formula Cards}\label{formula-cards-2}

\subsection*{Statistics --- Key Formulas \&
Rules}\label{statistics-key-formulas-rules}
\addcontentsline{toc}{subsection}{Statistics --- Key Formulas \& Rules}

\subsubsection*{Averages and Range}\label{averages-and-range}
\addcontentsline{toc}{subsubsection}{Averages and Range}

\textbf{💡 Pro Tip} Always \textbf{sort your data first} before finding
the median or quartiles.

\textbf{Mean}:
\[\bar{x} = \frac{\sum x}{n} = \frac{\text{sum of all values}}{\text{number of values}}\]

\textbf{Median}: Middle value when data is ordered.\\
If \(n\) is even: mean of the two middle values.

\textbf{Mode}: Most frequently occurring value.

\textbf{Range}: \[\text{Range} = \text{maximum} - \text{minimum}\]

\textbf{Example}: Data: \(3, 7, 7, 9, 14\)
\[\bar{x} = \frac{3+7+7+9+14}{5} = 8 \qquad \text{Median} = 7 \qquad \text{Mode} = 7 \qquad \text{Range} = 11\]

\subsubsection*{Mean from a Frequency
Table}\label{mean-from-a-frequency-table}
\addcontentsline{toc}{subsubsection}{Mean from a Frequency Table}

\textbf{🎯 Memory Trick} Multiply each value by its frequency first ---
then divide by the \textbf{total frequency}, not the number of rows!

\[\bar{x} = \frac{\sum fx}{\sum f}\]

\begin{longtable}[]{@{}ccc@{}}
\toprule\noalign{}
\(x\) & \(f\) & \(fx\) \\
\midrule\noalign{}
\endhead
\bottomrule\noalign{}
\endlastfoot
2 & 3 & 6 \\
4 & 7 & 28 \\
6 & 5 & 30 \\
\textbf{Total} & \textbf{15} & \textbf{64} \\
\end{longtable}

\[\bar{x} = \frac{64}{15} \approx 4.3\]

\subsubsection*{Grouped Frequency
Tables}\label{grouped-frequency-tables}
\addcontentsline{toc}{subsubsection}{Grouped Frequency Tables}

\textbf{⚠️ Common Mistake} Use the \textbf{midpoint} of each class ---
not the lower or upper boundary!

\textbf{Estimated mean}:
\[\bar{x} \approx \frac{\sum fx}{\sum f} \quad \text{where } x = \text{midpoint of class}\]

\textbf{Modal class}: Class with the \textbf{highest frequency}.

\textbf{Median class}: Class containing the \(\dfrac{n}{2}\)th value
(use cumulative frequency).

\textbf{Example}: Class \(10 \leq x < 20\), frequency \(= 8\) → midpoint
\(x = 15\), so \(fx = 120\)

\subsubsection*{Cumulative Frequency}\label{cumulative-frequency}
\addcontentsline{toc}{subsubsection}{Cumulative Frequency}

\textbf{📝 Quick Method} Read across from the CF axis to the curve, then
\textbf{down} to the \(x\)-axis.

\textbf{Building CF table}: Each row = previous cumulative total \(+\)
current frequency.

\textbf{Reading quartiles from a CF graph} (total = \(n\)):

\begin{longtable}[]{@{}lc@{}}
\toprule\noalign{}
Measure & Read at CF value \\
\midrule\noalign{}
\endhead
\bottomrule\noalign{}
\endlastfoot
Lower quartile \(Q_1\) & \(\dfrac{n}{4}\) \\
Median \(Q_2\) & \(\dfrac{n}{2}\) \\
Upper quartile \(Q_3\) & \(\dfrac{3n}{4}\) \\
\end{longtable}

\[\text{IQR} = Q_3 - Q_1\]

\subsubsection*{Box Plots and Outliers}\label{box-plots-and-outliers}
\addcontentsline{toc}{subsubsection}{Box Plots and Outliers}

\textbf{Five-number summary}: Min, \(Q_1\), Median, \(Q_3\), Max

\textbf{Box plot}: Box spans \(Q_1\) to \(Q_3\). Median line inside.
Whiskers extend to min and max.

\textbf{Outlier rule}:
\[\text{Outlier if: } x > Q_3 + 1.5 \times \text{IQR} \quad \text{or} \quad x < Q_1 - 1.5 \times \text{IQR}\]

\textbf{Example}: \(Q_1 = 35,\ Q_3 = 55,\ \text{IQR} = 20\)\\
Upper fence \(= 55 + 1.5 \times 20 = 85\)\\
Any value \(> 85\) is an outlier.

\subsubsection*{Combined Mean}\label{combined-mean}
\addcontentsline{toc}{subsubsection}{Combined Mean}

\textbf{💡 Pro Tip} You \textbf{cannot} just average two means --- you
must use the totals!

\[\bar{x}_{\text{combined}} = \frac{n_1 \bar{x}_1 + n_2 \bar{x}_2}{n_1 + n_2}\]

\textbf{Example}: Group 1: \(n=100\), mean \(= 18.2\). Group 2:
\(n=50\), mean \(= 17.3\)
\[\bar{x} = \frac{100 \times 18.2 + 50 \times 17.3}{150} = \frac{1820 + 865}{150} = \frac{2685}{150} = 17.9\]

\begin{center}\rule{0.5\linewidth}{0.5pt}\end{center}

\section{📊 Part A --- Averages and Range}\label{sec-averages}

\begin{tcolorbox}[enhanced jigsaw, arc=.35mm, breakable, toptitle=1mm, title={🔑 Key Formulas}, colframe=quarto-callout-note-color-frame, rightrule=.15mm, opacityback=0, titlerule=0mm, colback=white, coltitle=black, bottomtitle=1mm, toprule=.15mm, leftrule=.75mm, bottomrule=.15mm, colbacktitle=quarto-callout-note-color!10!white, left=2mm, opacitybacktitle=0.6]

\[\bar{x} = \frac{\sum x}{n} \qquad \text{Median} = \text{middle value (sorted)} \qquad \text{Range} = \text{max} - \text{min}\]

\end{tcolorbox}

\subsection{🟢 Easy}

\begin{exercise}[Mean from a
List]\protect\hypertarget{exr-avg-1}{}\label{exr-avg-1}

Calculate the mean.

\begin{enumerate}
\def\labelenumi{\alph{enumi})}
\tightlist
\item
\end{enumerate}

\[15,\ 17,\ 20,\ 2,\ 8\]

\begin{enumerate}
\def\labelenumi{\alph{enumi})}
\setcounter{enumi}{1}
\tightlist
\item
\end{enumerate}

\[17,\ 10,\ 7,\ 3,\ 18,\ 17,\ 12,\ 4\]

\begin{enumerate}
\def\labelenumi{\alph{enumi})}
\setcounter{enumi}{2}
\tightlist
\item
\end{enumerate}

\[13,\ 3,\ 15,\ 6,\ 13,\ 14\]

\begin{enumerate}
\def\labelenumi{\alph{enumi})}
\setcounter{enumi}{3}
\tightlist
\item
\end{enumerate}

\[11,\ 10,\ 16,\ 7,\ 11,\ 13,\ 6,\ 16\]

\begin{tcolorbox}[enhanced jigsaw, arc=.35mm, breakable, toptitle=1mm, title=\textcolor{quarto-callout-tip-color}{\faLightbulb}\hspace{0.5em}{🔍 View Solution}, colframe=quarto-callout-tip-color-frame, rightrule=.15mm, opacityback=0, titlerule=0mm, colback=white, coltitle=black, bottomtitle=1mm, toprule=.15mm, leftrule=.75mm, bottomrule=.15mm, colbacktitle=quarto-callout-tip-color!10!white, left=2mm, opacitybacktitle=0.6]

\textbf{a)} \[\bar{x} = 12.4\]

\[\quad \dfrac{15+17+20+2+8}{5} = 12.4\]

\textbf{b)} \[\bar{x} = 11\]

\[\quad \dfrac{17+10+7+3+18+17+12+4}{8} = 11\]

\textbf{c)} \[\bar{x} = 10.67\]

\[\quad \dfrac{13+3+15+6+13+14}{6} = 10.67\]

\textbf{d)} \[\bar{x} = 11.25\]

\[\quad \dfrac{11+10+16+7+11+13+6+16}{8} = 11.25\]

\end{tcolorbox}

{Sum of values} {Division} {Mean formula}

\end{exercise}

\begin{exercise}[Median from a
List]\protect\hypertarget{exr-avg-2}{}\label{exr-avg-2}

Find the median.

\begin{enumerate}
\def\labelenumi{\alph{enumi})}
\tightlist
\item
\end{enumerate}

\[1,\ 2,\ 13,\ 15,\ 19,\ 20\]

\begin{enumerate}
\def\labelenumi{\alph{enumi})}
\setcounter{enumi}{1}
\tightlist
\item
\end{enumerate}

\[5,\ 7,\ 8,\ 10,\ 12,\ 18\]

\begin{enumerate}
\def\labelenumi{\alph{enumi})}
\setcounter{enumi}{2}
\tightlist
\item
\end{enumerate}

\[3,\ 11,\ 14,\ 15,\ 16,\ 18,\ 19\]

\begin{enumerate}
\def\labelenumi{\alph{enumi})}
\setcounter{enumi}{3}
\tightlist
\item
\end{enumerate}

\[2,\ 2,\ 6,\ 8,\ 8,\ 10,\ 14,\ 16,\ 20\]

\begin{tcolorbox}[enhanced jigsaw, arc=.35mm, breakable, toptitle=1mm, title=\textcolor{quarto-callout-tip-color}{\faLightbulb}\hspace{0.5em}{🔍 View Solution}, colframe=quarto-callout-tip-color-frame, rightrule=.15mm, opacityback=0, titlerule=0mm, colback=white, coltitle=black, bottomtitle=1mm, toprule=.15mm, leftrule=.75mm, bottomrule=.15mm, colbacktitle=quarto-callout-tip-color!10!white, left=2mm, opacitybacktitle=0.6]

\textbf{a)} \[\text{Median} = 14\]

\[\quad \text{Mean of positions 3 and 4: }\frac{13+15}{2}=14\]

\textbf{b)} \[\text{Median} = 9\]

\[\quad \text{Mean of positions 3 and 4: }\frac{8+10}{2}=9\]

\textbf{c)} \[\text{Median} = 15\]

\[\quad \text{Ordered. Position }\frac{n+1}{2}=4\text{th value} = 15\]

\textbf{d)} \[\text{Median} = 8\]

\[\quad \text{Ordered. Position }\frac{n+1}{2}=5\text{th value} = 8\]

\end{tcolorbox}

{Ordering data} {Locating middle}

\end{exercise}

\begin{exercise}[Mode from a
List]\protect\hypertarget{exr-avg-3}{}\label{exr-avg-3}

Find the mode.

\begin{enumerate}
\def\labelenumi{\alph{enumi})}
\tightlist
\item
\end{enumerate}

\[2,\ 11,\ 2,\ 9,\ 2,\ 2\]

\begin{enumerate}
\def\labelenumi{\alph{enumi})}
\setcounter{enumi}{1}
\tightlist
\item
\end{enumerate}

\[2,\ 2,\ 2,\ 2,\ 3,\ 8\]

\begin{enumerate}
\def\labelenumi{\alph{enumi})}
\setcounter{enumi}{2}
\tightlist
\item
\end{enumerate}

\[5,\ 5,\ 5,\ 10,\ 3,\ 11\]

\begin{enumerate}
\def\labelenumi{\alph{enumi})}
\setcounter{enumi}{3}
\tightlist
\item
\end{enumerate}

\[11,\ 3,\ 11,\ 3,\ 3,\ 3\]

\begin{tcolorbox}[enhanced jigsaw, arc=.35mm, breakable, toptitle=1mm, title=\textcolor{quarto-callout-tip-color}{\faLightbulb}\hspace{0.5em}{🔍 View Solution}, colframe=quarto-callout-tip-color-frame, rightrule=.15mm, opacityback=0, titlerule=0mm, colback=white, coltitle=black, bottomtitle=1mm, toprule=.15mm, leftrule=.75mm, bottomrule=.15mm, colbacktitle=quarto-callout-tip-color!10!white, left=2mm, opacitybacktitle=0.6]

\textbf{a)} \[\text{Mode} = 2\]

\[\quad \text{Most frequent value(s): }2\]

\textbf{b)} \[\text{Mode} = 2\]

\[\quad \text{Most frequent value(s): }2\]

\textbf{c)} \[\text{Mode} = 5\]

\[\quad \text{Most frequent value(s): }5\]

\textbf{d)} \[\text{Mode} = 3\]

\[\quad \text{Most frequent value(s): }3\]

\end{tcolorbox}

{Frequency counting}

\end{exercise}

\begin{exercise}[Range from a
List]\protect\hypertarget{exr-avg-4}{}\label{exr-avg-4}

Find the range.

\begin{enumerate}
\def\labelenumi{\alph{enumi})}
\tightlist
\item
\end{enumerate}

\[24,\ 3,\ 12,\ 30,\ 12,\ 9,\ 4,\ 10,\ 16\]

\begin{enumerate}
\def\labelenumi{\alph{enumi})}
\setcounter{enumi}{1}
\tightlist
\item
\end{enumerate}

\[20,\ 7,\ 14,\ 29,\ 17,\ 23,\ 30,\ 9\]

\begin{enumerate}
\def\labelenumi{\alph{enumi})}
\setcounter{enumi}{2}
\tightlist
\item
\end{enumerate}

\[10,\ 18,\ 17,\ 21,\ 17,\ 4,\ 7,\ 17,\ 23\]

\begin{enumerate}
\def\labelenumi{\alph{enumi})}
\setcounter{enumi}{3}
\tightlist
\item
\end{enumerate}

\[30,\ 14,\ 6,\ 15,\ 19,\ 20,\ 13,\ 19\]

\begin{tcolorbox}[enhanced jigsaw, arc=.35mm, breakable, toptitle=1mm, title=\textcolor{quarto-callout-tip-color}{\faLightbulb}\hspace{0.5em}{🔍 View Solution}, colframe=quarto-callout-tip-color-frame, rightrule=.15mm, opacityback=0, titlerule=0mm, colback=white, coltitle=black, bottomtitle=1mm, toprule=.15mm, leftrule=.75mm, bottomrule=.15mm, colbacktitle=quarto-callout-tip-color!10!white, left=2mm, opacitybacktitle=0.6]

\textbf{a)} \[\text{Range} = 27\]

\[\quad 30 - 3 = 27\]

\textbf{b)} \[\text{Range} = 23\]

\[\quad 30 - 7 = 23\]

\textbf{c)} \[\text{Range} = 19\]

\[\quad 23 - 4 = 19\]

\textbf{d)} \[\text{Range} = 24\]

\[\quad 30 - 6 = 24\]

\end{tcolorbox}

{Max min} {Subtraction}

\end{exercise}

∑ × π ÷ √ ◆ ∞ ± ∫ ≠ Δ

\subsection{🟡 Medium}

\begin{exercise}[Mean, Median, Mode and
Range]\protect\hypertarget{exr-avg-5}{}\label{exr-avg-5}

Find the mean, median, mode and range.

\begin{enumerate}
\def\labelenumi{\alph{enumi})}
\tightlist
\item
\end{enumerate}

\[4,\ 8,\ 13,\ 3,\ 8,\ 9,\ 10,\ 8,\ 9\]

\begin{enumerate}
\def\labelenumi{\alph{enumi})}
\setcounter{enumi}{1}
\tightlist
\item
\end{enumerate}

\[3,\ 9,\ 15,\ 14,\ 15,\ 15,\ 11\]

\begin{enumerate}
\def\labelenumi{\alph{enumi})}
\setcounter{enumi}{2}
\tightlist
\item
\end{enumerate}

\[13,\ 12,\ 7,\ 10,\ 10,\ 7,\ 7,\ 8,\ 4\]

\begin{enumerate}
\def\labelenumi{\alph{enumi})}
\setcounter{enumi}{3}
\tightlist
\item
\end{enumerate}

\[5,\ 15,\ 6,\ 15,\ 8,\ 8,\ 13,\ 8\]

\begin{tcolorbox}[enhanced jigsaw, arc=.35mm, breakable, toptitle=1mm, title=\textcolor{quarto-callout-tip-color}{\faLightbulb}\hspace{0.5em}{🔍 View Solution}, colframe=quarto-callout-tip-color-frame, rightrule=.15mm, opacityback=0, titlerule=0mm, colback=white, coltitle=black, bottomtitle=1mm, toprule=.15mm, leftrule=.75mm, bottomrule=.15mm, colbacktitle=quarto-callout-tip-color!10!white, left=2mm, opacitybacktitle=0.6]

\textbf{a)}
\[\text{Mean}=8,\ \text{Median}=8,\ \text{Mode}=8,\ \text{Range}=10\]

\[\quad \bar{x}=\frac{72}{9}=8\quad \text{Median}=8\quad \text{Mode}=8\quad \text{Range}=10\]

\textbf{b)}
\[\text{Mean}=11.71,\ \text{Median}=14,\ \text{Mode}=15,\ \text{Range}=12\]

\[\quad \bar{x}=\frac{82}{7}=11.71\quad \text{Median}=14\quad \text{Mode}=15\quad \text{Range}=12\]

\textbf{c)}
\[\text{Mean}=8.67,\ \text{Median}=8,\ \text{Mode}=7,\ \text{Range}=9\]

\[\quad \bar{x}=\frac{78}{9}=8.67\quad \text{Median}=8\quad \text{Mode}=7\quad \text{Range}=9\]

\textbf{d)}
\[\text{Mean}=9.75,\ \text{Median}=8,\ \text{Mode}=8,\ \text{Range}=10\]

\[\quad \bar{x}=\frac{78}{8}=9.75\quad \text{Median}=8\quad \text{Mode}=8\quad \text{Range}=10\]

\end{tcolorbox}

{All averages} {Ordering data}

\end{exercise}

∑ × π ÷ √ ◆ ∞ ± ∫ ≠ Δ

\subsection{🔴 Challenging}

\begin{exercise}[Find Missing Value / Updated
Mean]\protect\hypertarget{exr-avg-6}{}\label{exr-avg-6}

Find the missing value or new mean.

\begin{enumerate}
\def\labelenumi{\alph{enumi})}
\tightlist
\item
\end{enumerate}

\[\text{Mean} = 14,\quad n=6\\[4pt]\text{Numbers: }14,\ 8,\ 13,\ 7,\ 4,\ x\]

\begin{enumerate}
\def\labelenumi{\alph{enumi})}
\setcounter{enumi}{1}
\tightlist
\item
\end{enumerate}

\[\text{Mean} = 26,\quad\text{three values: }30,\ 35,\ 30\\[4pt]\text{Other two in ratio }1:4\]

\begin{enumerate}
\def\labelenumi{\alph{enumi})}
\setcounter{enumi}{2}
\tightlist
\item
\end{enumerate}

\[\text{Mean} = 11,\quad n=5\\[4pt]\text{Numbers: }15,\ 14,\ 8,\ 4,\ x\]

\begin{enumerate}
\def\labelenumi{\alph{enumi})}
\setcounter{enumi}{3}
\tightlist
\item
\end{enumerate}

\[\text{Mean of }6\text{ numbers} = 3,\quad\text{new value} = 12\]

\begin{tcolorbox}[enhanced jigsaw, arc=.35mm, breakable, toptitle=1mm, title=\textcolor{quarto-callout-tip-color}{\faLightbulb}\hspace{0.5em}{🔍 View Solution}, colframe=quarto-callout-tip-color-frame, rightrule=.15mm, opacityback=0, titlerule=0mm, colback=white, coltitle=black, bottomtitle=1mm, toprule=.15mm, leftrule=.75mm, bottomrule=.15mm, colbacktitle=quarto-callout-tip-color!10!white, left=2mm, opacitybacktitle=0.6]

\textbf{a)} \[x = 38\]

\[\quad \text{Total} = 14\times6=84\quad x = 84-(14+8+13+7+4) = 38\]

\textbf{b)} \[7\text{ and }28\]

\[\quad \text{Step 1: Total}=26\times5=130\\[4pt]\text{Step 2: Remaining}=130-(30+35+30)=35\\[4pt]\text{Step 3: Ratio }1:4\text{ → }5\text{ equal parts, each }=\frac{35}{5}=7\\[4pt]\text{Step 4: Numbers are }7\text{ and }28\]

\textbf{c)} \[x = 14\]

\[\quad \text{Total} = 11\times5=55\quad x = 55-(15+14+8+4) = 14\]

\textbf{d)} \[\text{New mean} = 4.29\]

\[\quad \text{Old total} = 3\times6 = 18\quad \text{New total} = 18+12 = 30\quad \bar{x}_{\text{new}} = \frac{30}{7} = 4.29\]

\end{tcolorbox}

{Rearranging formula} {Sum from mean}

\end{exercise}

\begin{exercise}[Averages from Stem \&
Leaf]\protect\hypertarget{exr-avg-1}{}\label{exr-avg-1}

Use the stem and leaf diagram to find the median, range and modal class.

\begin{enumerate}
\def\labelenumi{\alph{enumi})}
\item
  \(\text{Find the median.}\)
\item
  \(\text{Find the range.}\)
\item
  \(\text{Find the mode.}\)
\end{enumerate}

\begin{tcolorbox}[enhanced jigsaw, arc=.35mm, breakable, toptitle=1mm, title=\textcolor{quarto-callout-tip-color}{\faLightbulb}\hspace{0.5em}{🔍 View Solution}, colframe=quarto-callout-tip-color-frame, rightrule=.15mm, opacityback=0, titlerule=0mm, colback=white, coltitle=black, bottomtitle=1mm, toprule=.15mm, leftrule=.75mm, bottomrule=.15mm, colbacktitle=quarto-callout-tip-color!10!white, left=2mm, opacitybacktitle=0.6]

\textbf{a)} \(\text{Median} = 31\)

\[\quad \text{Ordered: }27\text{ values, middle} = 31\]

\textbf{b)} \(\text{Range} = 47\)

\[\quad 59 - 12 = 47\]

\textbf{c)} \(\text{Mode} = 17\)

\[\quad \text{Most frequent value: }17\]

\end{tcolorbox}

{Reading stem leaf} {Ordering data}

\end{exercise}

\begin{exercise}[Averages from Stem \&
Leaf]\protect\hypertarget{exr-avg-2}{}\label{exr-avg-2}

Use the stem and leaf diagram to find the median, range and modal class.

\begin{enumerate}
\def\labelenumi{\alph{enumi})}
\item
  \(\text{Find the median.}\)
\item
  \(\text{Find the range.}\)
\item
  \(\text{Find the mode.}\)
\end{enumerate}

\begin{tcolorbox}[enhanced jigsaw, arc=.35mm, breakable, toptitle=1mm, title=\textcolor{quarto-callout-tip-color}{\faLightbulb}\hspace{0.5em}{🔍 View Solution}, colframe=quarto-callout-tip-color-frame, rightrule=.15mm, opacityback=0, titlerule=0mm, colback=white, coltitle=black, bottomtitle=1mm, toprule=.15mm, leftrule=.75mm, bottomrule=.15mm, colbacktitle=quarto-callout-tip-color!10!white, left=2mm, opacitybacktitle=0.6]

\textbf{a)} \(\text{Median} = 37\)

\[\quad \text{Ordered: }15\text{ values, middle} = 37\]

\textbf{b)} \(\text{Range} = 35\)

\[\quad 58 - 23 = 35\]

\textbf{c)} \(\text{Mode} = 23\)

\[\quad \text{Most frequent value: }23\]

\end{tcolorbox}

{Reading stem leaf} {Ordering data}

\end{exercise}

\begin{exercise}[Averages from Stem \&
Leaf]\protect\hypertarget{exr-avg-3}{}\label{exr-avg-3}

Use the stem and leaf diagram to find the median, range and modal class.

\begin{enumerate}
\def\labelenumi{\alph{enumi})}
\item
  \(\text{Find the median.}\)
\item
  \(\text{Find the range.}\)
\item
  \(\text{Find the mode.}\)
\end{enumerate}

\begin{tcolorbox}[enhanced jigsaw, arc=.35mm, breakable, toptitle=1mm, title=\textcolor{quarto-callout-tip-color}{\faLightbulb}\hspace{0.5em}{🔍 View Solution}, colframe=quarto-callout-tip-color-frame, rightrule=.15mm, opacityback=0, titlerule=0mm, colback=white, coltitle=black, bottomtitle=1mm, toprule=.15mm, leftrule=.75mm, bottomrule=.15mm, colbacktitle=quarto-callout-tip-color!10!white, left=2mm, opacitybacktitle=0.6]

\textbf{a)} \(\text{Median} = 60\)

\[\quad \text{Ordered: }11\text{ values, middle} = 60\]

\textbf{b)} \(\text{Range} = 36\)

\[\quad 77 - 41 = 36\]

\textbf{c)} \(\text{Mode} = 60\)

\[\quad \text{Most frequent value: }60\]

\end{tcolorbox}

{Reading stem leaf} {Ordering data}

\end{exercise}

∑ × π ÷ √ ◆ ∞ ± ∫ ≠ Δ

\begin{center}\rule{0.5\linewidth}{0.5pt}\end{center}

\section{📋 Part B --- Frequency Tables}\label{sec-freq}

\begin{tcolorbox}[enhanced jigsaw, arc=.35mm, breakable, toptitle=1mm, title={🔑 Key Formula}, colframe=quarto-callout-note-color-frame, rightrule=.15mm, opacityback=0, titlerule=0mm, colback=white, coltitle=black, bottomtitle=1mm, toprule=.15mm, leftrule=.75mm, bottomrule=.15mm, colbacktitle=quarto-callout-note-color!10!white, left=2mm, opacitybacktitle=0.6]

\[\bar{x} = \frac{\sum fx}{\sum f}\]

\end{tcolorbox}

\subsection{🟢 Easy}

\begin{exercise}[Mean from Frequency
Table]\protect\hypertarget{exr-freq-1}{}\label{exr-freq-1}

Calculate the mean from the frequency table.

\begin{enumerate}
\def\labelenumi{\alph{enumi})}
\tightlist
\item
\end{enumerate}

\[{\renewcommand{\arraystretch}{1.6}\begin{array}{lcccccc}
\hline x & 2 & 3 & 5 & 7 & 9 & 10 \\[4pt]
\hline f & 3 & 1 & 7 & 7 & 2 & 2 \\[4pt]
\hline
\end{array}}\]

\begin{enumerate}
\def\labelenumi{\alph{enumi})}
\setcounter{enumi}{1}
\tightlist
\item
\end{enumerate}

\[{\renewcommand{\arraystretch}{1.6}\begin{array}{lcccc}
\hline x & 4 & 6 & 8 & 9 \\[4pt]
\hline f & 4 & 4 & 6 & 6 \\[4pt]
\hline
\end{array}}\]

\begin{enumerate}
\def\labelenumi{\alph{enumi})}
\setcounter{enumi}{2}
\tightlist
\item
\end{enumerate}

\[{\renewcommand{\arraystretch}{1.6}\begin{array}{lccccc}
\hline x & 1 & 2 & 4 & 7 & 8 \\[4pt]
\hline f & 4 & 10 & 4 & 1 & 4 \\[4pt]
\hline
\end{array}}\]

\begin{enumerate}
\def\labelenumi{\alph{enumi})}
\setcounter{enumi}{3}
\tightlist
\item
\end{enumerate}

\[{\renewcommand{\arraystretch}{1.6}\begin{array}{lcccc}
\hline x & 2 & 4 & 5 & 6 \\[4pt]
\hline f & 5 & 5 & 5 & 3 \\[4pt]
\hline
\end{array}}\]

\begin{tcolorbox}[enhanced jigsaw, arc=.35mm, breakable, toptitle=1mm, title=\textcolor{quarto-callout-tip-color}{\faLightbulb}\hspace{0.5em}{🔍 View Solution}, colframe=quarto-callout-tip-color-frame, rightrule=.15mm, opacityback=0, titlerule=0mm, colback=white, coltitle=black, bottomtitle=1mm, toprule=.15mm, leftrule=.75mm, bottomrule=.15mm, colbacktitle=quarto-callout-tip-color!10!white, left=2mm, opacitybacktitle=0.6]

\textbf{a)} \[\bar{x} = 5.95\]

\[\quad \frac{2\times3+3\times1+5\times7+7\times7+9\times2+10\times2}{22} = \frac{131}{22} = 5.95\]

\textbf{b)} \[\bar{x} = 7.1\]

\[\quad \frac{4\times4+6\times4+8\times6+9\times6}{20} = \frac{142}{20} = 7.1\]

\textbf{c)} \[\bar{x} = 3.43\]

\[\quad \frac{1\times4+2\times10+4\times4+7\times1+8\times4}{23} = \frac{79}{23} = 3.43\]

\textbf{d)} \[\bar{x} = 4.06\]

\[\quad \frac{2\times5+4\times5+5\times5+6\times3}{18} = \frac{73}{18} = 4.06\]

\end{tcolorbox}

{Fx column} {Sum of frequencies}

\end{exercise}

\begin{exercise}[Mode from Frequency
Table]\protect\hypertarget{exr-freq-3}{}\label{exr-freq-3}

Find the mode from the frequency table.

\begin{enumerate}
\def\labelenumi{\alph{enumi})}
\tightlist
\item
\end{enumerate}

\[{\renewcommand{\arraystretch}{1.6}\begin{array}{lcccccc}
\hline x & 1 & 2 & 3 & 4 & 5 & 9 \\[4pt]
\hline f & 1 & 2 & 10 & 2 & 2 & 6 \\[4pt]
\hline
\end{array}}\]

\begin{enumerate}
\def\labelenumi{\alph{enumi})}
\setcounter{enumi}{1}
\tightlist
\item
\end{enumerate}

\[{\renewcommand{\arraystretch}{1.6}\begin{array}{lccccc}
\hline x & 1 & 2 & 5 & 7 & 8 \\[4pt]
\hline f & 4 & 4 & 9 & 9 & 1 \\[4pt]
\hline
\end{array}}\]

\begin{enumerate}
\def\labelenumi{\alph{enumi})}
\setcounter{enumi}{2}
\tightlist
\item
\end{enumerate}

\[{\renewcommand{\arraystretch}{1.6}\begin{array}{lcccc}
\hline x & 1 & 3 & 8 & 9 \\[4pt]
\hline f & 3 & 7 & 5 & 2 \\[4pt]
\hline
\end{array}}\]

\begin{enumerate}
\def\labelenumi{\alph{enumi})}
\setcounter{enumi}{3}
\tightlist
\item
\end{enumerate}

\[{\renewcommand{\arraystretch}{1.6}\begin{array}{lccccc}
\hline x & 2 & 4 & 7 & 9 & 10 \\[4pt]
\hline f & 2 & 5 & 3 & 9 & 7 \\[4pt]
\hline
\end{array}}\]

\begin{tcolorbox}[enhanced jigsaw, arc=.35mm, breakable, toptitle=1mm, title=\textcolor{quarto-callout-tip-color}{\faLightbulb}\hspace{0.5em}{🔍 View Solution}, colframe=quarto-callout-tip-color-frame, rightrule=.15mm, opacityback=0, titlerule=0mm, colback=white, coltitle=black, bottomtitle=1mm, toprule=.15mm, leftrule=.75mm, bottomrule=.15mm, colbacktitle=quarto-callout-tip-color!10!white, left=2mm, opacitybacktitle=0.6]

\textbf{a)} \[\text{Mode} = 3\]

\[\quad \text{Highest frequency is }10\text{, giving mode(s): }3\]

\textbf{b)} \[\text{Mode} = 5, 7\]

\[\quad \text{Highest frequency is }9\text{, giving mode(s): }5, 7\]

\textbf{c)} \[\text{Mode} = 3\]

\[\quad \text{Highest frequency is }7\text{, giving mode(s): }3\]

\textbf{d)} \[\text{Mode} = 9\]

\[\quad \text{Highest frequency is }9\text{, giving mode(s): }9\]

\end{tcolorbox}

{Reading frequency table}

\end{exercise}

∑ × π ÷ √ ◆ ∞ ± ∫ ≠ Δ

\subsection{🟡 Medium}

\begin{exercise}[Median from Frequency
Table]\protect\hypertarget{exr-freq-2}{}\label{exr-freq-2}

Find the median from the frequency table.

\begin{enumerate}
\def\labelenumi{\alph{enumi})}
\tightlist
\item
\end{enumerate}

\[{\renewcommand{\arraystretch}{1.6}\begin{array}{lcccccc}
\hline x & 1 & 2 & 5 & 6 & 7 & 10 \\[4pt]
\hline f & 10 & 4 & 1 & 1 & 10 & 5 \\[4pt]
\hline
\end{array}}\]

\begin{enumerate}
\def\labelenumi{\alph{enumi})}
\setcounter{enumi}{1}
\tightlist
\item
\end{enumerate}

\[{\renewcommand{\arraystretch}{1.6}\begin{array}{lccccc}
\hline x & 2 & 3 & 6 & 8 & 9 \\[4pt]
\hline f & 5 & 9 & 7 & 3 & 3 \\[4pt]
\hline
\end{array}}\]

\begin{enumerate}
\def\labelenumi{\alph{enumi})}
\setcounter{enumi}{2}
\tightlist
\item
\end{enumerate}

\[{\renewcommand{\arraystretch}{1.6}\begin{array}{lccccc}
\hline x & 2 & 3 & 7 & 9 & 10 \\[4pt]
\hline f & 2 & 10 & 1 & 9 & 4 \\[4pt]
\hline
\end{array}}\]

\begin{enumerate}
\def\labelenumi{\alph{enumi})}
\setcounter{enumi}{3}
\tightlist
\item
\end{enumerate}

\[{\renewcommand{\arraystretch}{1.6}\begin{array}{lccccc}
\hline x & 1 & 2 & 3 & 9 & 10 \\[4pt]
\hline f & 1 & 6 & 7 & 8 & 9 \\[4pt]
\hline
\end{array}}\]

\begin{tcolorbox}[enhanced jigsaw, arc=.35mm, breakable, toptitle=1mm, title=\textcolor{quarto-callout-tip-color}{\faLightbulb}\hspace{0.5em}{🔍 View Solution}, colframe=quarto-callout-tip-color-frame, rightrule=.15mm, opacityback=0, titlerule=0mm, colback=white, coltitle=black, bottomtitle=1mm, toprule=.15mm, leftrule=.75mm, bottomrule=.15mm, colbacktitle=quarto-callout-tip-color!10!white, left=2mm, opacitybacktitle=0.6]

\textbf{a)} \[\text{Median} = 6\]

\[\quad n=31\quad \text{Middle position }\frac{n+1}{2}=16.0\quad \text{Median}=6\]

\textbf{b)} \[\text{Median} = 3\]

\[\quad n=27\quad \text{Middle position }\frac{n+1}{2}=14.0\quad \text{Median}=3\]

\textbf{c)} \[\text{Median} = 8\]

\[\quad n=26\quad \text{Middle position }\frac{n+1}{2}=13.5\quad \text{Median}=8\]

\textbf{d)} \[\text{Median} = 9\]

\[\quad n=31\quad \text{Middle position }\frac{n+1}{2}=16.0\quad \text{Median}=9\]

\end{tcolorbox}

{Cumulative frequency} {Locating middle}

\end{exercise}

\begin{exercise}[Mean, Median \& Mode from
Table]\protect\hypertarget{exr-freq-4}{}\label{exr-freq-4}

Find the mean, median and mode from the frequency table.

\begin{enumerate}
\def\labelenumi{\alph{enumi})}
\tightlist
\item
\end{enumerate}

\[\text{Calculate the mean.}\]

\begin{enumerate}
\def\labelenumi{\alph{enumi})}
\setcounter{enumi}{1}
\tightlist
\item
\end{enumerate}

\[\text{Find the median.}\]

\begin{enumerate}
\def\labelenumi{\alph{enumi})}
\setcounter{enumi}{2}
\tightlist
\item
\end{enumerate}

\[\text{(c) Find the mode.}\]

\begin{tcolorbox}[enhanced jigsaw, arc=.35mm, breakable, toptitle=1mm, title=\textcolor{quarto-callout-tip-color}{\faLightbulb}\hspace{0.5em}{🔍 View Solution}, colframe=quarto-callout-tip-color-frame, rightrule=.15mm, opacityback=0, titlerule=0mm, colback=white, coltitle=black, bottomtitle=1mm, toprule=.15mm, leftrule=.75mm, bottomrule=.15mm, colbacktitle=quarto-callout-tip-color!10!white, left=2mm, opacitybacktitle=0.6]

\textbf{a)} \[\bar{x} = 4.58\]

\[\quad \bar{x}=\dfrac{\sum fx}{\sum f}=\dfrac{1\times3+2\times4+4\times5+8\times7}{19}=\dfrac{87}{19}=4.58\]

\textbf{b)} \[\text{Median} = 4\]

\[\quad n=19\quad \text{Position }\frac{n+1}{2}=10.0\quad \text{Build CF: }1: CF=3,\ 2: CF=7,\ 4: CF=12,\ 8: CF=19,\\quad \text{Median}=4\]

\textbf{c)} \[\text{Mode} = 8\]

\[\quad \text{Highest frequency = }7\text{ at }x=8\]

\end{tcolorbox}

{Fx column} {Cumulative frequency} {Reading frequency table}

\end{exercise}

∑ × π ÷ √ ◆ ∞ ± ∫ ≠ Δ

\subsection{🔴 Challenging}

\begin{exercise}[Find x Given Median (Frequency
Table)]\protect\hypertarget{exr-freq-5}{}\label{exr-freq-5}

Find the largest and smallest possible values of x such that the median
is as given.

\begin{enumerate}
\def\labelenumi{\alph{enumi})}
\tightlist
\item
\end{enumerate}

\[\text{Median} = 6\\[6pt]{\renewcommand{\arraystretch}{1.6}\begin{array}{lcccc}
\hline \text{Score} & 5 & 6 & 7 & 8 \\[4pt]
\hline \text{Freq} & 3 & 7 & x & 5 \\[4pt]
\hline
\end{array}}\]

\begin{enumerate}
\def\labelenumi{\alph{enumi})}
\setcounter{enumi}{1}
\tightlist
\item
\end{enumerate}

\[\text{Median} = 6\\[6pt]{\renewcommand{\arraystretch}{1.6}\begin{array}{lcccc}
\hline \text{Score} & 5 & 6 & 7 & 8 \\[4pt]
\hline \text{Freq} & 5 & 9 & x & 6 \\[4pt]
\hline
\end{array}}\]

\begin{enumerate}
\def\labelenumi{\alph{enumi})}
\setcounter{enumi}{2}
\tightlist
\item
\end{enumerate}

\[\text{Median} = 6\\[6pt]{\renewcommand{\arraystretch}{1.6}\begin{array}{lcccc}
\hline \text{Score} & 5 & 6 & 7 & 8 \\[4pt]
\hline \text{Freq} & 4 & 5 & x & 4 \\[4pt]
\hline
\end{array}}\]

\begin{enumerate}
\def\labelenumi{\alph{enumi})}
\setcounter{enumi}{3}
\tightlist
\item
\end{enumerate}

\[\text{Median} = 6\\[6pt]{\renewcommand{\arraystretch}{1.6}\begin{array}{lcccc}
\hline \text{Score} & 5 & 6 & 7 & 8 \\[4pt]
\hline \text{Freq} & 5 & 5 & x & 7 \\[4pt]
\hline
\end{array}}\]

\begin{tcolorbox}[enhanced jigsaw, arc=.35mm, breakable, toptitle=1mm, title=\textcolor{quarto-callout-tip-color}{\faLightbulb}\hspace{0.5em}{🔍 View Solution}, colframe=quarto-callout-tip-color-frame, rightrule=.15mm, opacityback=0, titlerule=0mm, colback=white, coltitle=black, bottomtitle=1mm, toprule=.15mm, leftrule=.75mm, bottomrule=.15mm, colbacktitle=quarto-callout-tip-color!10!white, left=2mm, opacitybacktitle=0.6]

\textbf{a)} \[x_{\min} = 0,\ x_{\max} = 4\]

\[\quad \text{For median = 6: cumulative freq up to 6} \geq \frac{n+1}{2}\quad \text{Largest: }3+7-5-1=4\quad \text{Smallest: }0\]

\textbf{b)} \[x_{\min} = 0,\ x_{\max} = 7\]

\[\quad \text{For median = 6: cumulative freq up to 6} \geq \frac{n+1}{2}\quad \text{Largest: }5+9-6-1=7\quad \text{Smallest: }0\]

\textbf{c)} \[x_{\min} = 0,\ x_{\max} = 4\]

\[\quad \text{For median = 6: cumulative freq up to 6} \geq \frac{n+1}{2}\quad \text{Largest: }4+5-4-1=4\quad \text{Smallest: }0\]

\textbf{d)} \[x_{\min} = 0,\ x_{\max} = 2\]

\[\quad \text{For median = 6: cumulative freq up to 6} \geq \frac{n+1}{2}\quad \text{Largest: }5+5-7-1=2\quad \text{Smallest: }0\]

\end{tcolorbox}

{Cumulative frequency} {Inequality reasoning}

\end{exercise}

∑ × π ÷ √ ◆ ∞ ± ∫ ≠ Δ

\begin{center}\rule{0.5\linewidth}{0.5pt}\end{center}

\section{📦 Part C --- Grouped Frequency Tables}\label{sec-grouped}

\begin{tcolorbox}[enhanced jigsaw, arc=.35mm, breakable, toptitle=1mm, title={🔑 Key Points}, colframe=quarto-callout-note-color-frame, rightrule=.15mm, opacityback=0, titlerule=0mm, colback=white, coltitle=black, bottomtitle=1mm, toprule=.15mm, leftrule=.75mm, bottomrule=.15mm, colbacktitle=quarto-callout-note-color!10!white, left=2mm, opacitybacktitle=0.6]

Use the \textbf{midpoint} of each class interval for estimated mean:
\(\bar{x} \approx \dfrac{\sum fx}{\sum f}\)

\textbf{Modal class} = class with highest frequency. \textbf{Median
class} = class containing the \(\frac{n}{2}\)th value.

\end{tcolorbox}

\subsection{🟢 Easy}

\begin{exercise}[Estimated Mean
(Grouped)]\protect\hypertarget{exr-grp-1}{}\label{exr-grp-1}

Estimate the mean from the grouped frequency table.

\begin{enumerate}
\def\labelenumi{\alph{enumi})}
\tightlist
\item
\end{enumerate}

\[{\renewcommand{\arraystretch}{1.6}\begin{array}{lcc}
\hline \text{Class} & \text{Midpoint} & \text{Frequency} \\[4pt]
\hline 19 \leq x < 27  & 23 & 3 \\[4pt]
\hline 27 \leq x < 35  & 31 & 9 \\[4pt]
\hline 35 \leq x < 43  & 39 & 11 \\[4pt]
\hline 43 \leq x < 51  & 47 & 4 \\[4pt]
\hline 51 \leq x < 59  & 55 & 15 \\[4pt]
\hline 59 \leq x < 67  & 63 & 15 \\[4pt]
\hline \end{array}}\]

\begin{enumerate}
\def\labelenumi{\alph{enumi})}
\setcounter{enumi}{1}
\tightlist
\item
\end{enumerate}

\[{\renewcommand{\arraystretch}{1.6}\begin{array}{lcc}
\hline \text{Class} & \text{Midpoint} & \text{Frequency} \\[4pt]
\hline 14 \leq x < 18  & 16 & 17 \\[4pt]
\hline 18 \leq x < 22  & 20 & 3 \\[4pt]
\hline 22 \leq x < 26  & 24 & 19 \\[4pt]
\hline 26 \leq x < 30  & 28 & 10 \\[4pt]
\hline 30 \leq x < 34  & 32 & 3 \\[4pt]
\hline 34 \leq x < 38  & 36 & 5 \\[4pt]
\hline \end{array}}\]

\begin{enumerate}
\def\labelenumi{\alph{enumi})}
\setcounter{enumi}{2}
\tightlist
\item
\end{enumerate}

\[{\renewcommand{\arraystretch}{1.6}\begin{array}{lcc}
\hline \text{Class} & \text{Midpoint} & \text{Frequency} \\[4pt]
\hline 29 \leq x < 37  & 33 & 15 \\[4pt]
\hline 37 \leq x < 45  & 41 & 14 \\[4pt]
\hline 45 \leq x < 53  & 49 & 20 \\[4pt]
\hline 53 \leq x < 61  & 57 & 5 \\[4pt]
\hline \end{array}}\]

\begin{enumerate}
\def\labelenumi{\alph{enumi})}
\setcounter{enumi}{3}
\tightlist
\item
\end{enumerate}

\[{\renewcommand{\arraystretch}{1.6}\begin{array}{lcc}
\hline \text{Class} & \text{Midpoint} & \text{Frequency} \\[4pt]
\hline 10 \leq x < 15  & 12.5 & 11 \\[4pt]
\hline 15 \leq x < 20  & 17.5 & 16 \\[4pt]
\hline 20 \leq x < 25  & 22.5 & 6 \\[4pt]
\hline 25 \leq x < 30  & 27.5 & 11 \\[4pt]
\hline 30 \leq x < 35  & 32.5 & 20 \\[4pt]
\hline \end{array}}\]

\begin{tcolorbox}[enhanced jigsaw, arc=.35mm, breakable, toptitle=1mm, title=\textcolor{quarto-callout-tip-color}{\faLightbulb}\hspace{0.5em}{🔍 View Solution}, colframe=quarto-callout-tip-color-frame, rightrule=.15mm, opacityback=0, titlerule=0mm, colback=white, coltitle=black, bottomtitle=1mm, toprule=.15mm, leftrule=.75mm, bottomrule=.15mm, colbacktitle=quarto-callout-tip-color!10!white, left=2mm, opacitybacktitle=0.6]

\textbf{a)} \[\text{Estimated mean} \approx 48\]

\[\quad \frac{23\times3+31\times9+39\times11+47\times4+55\times15+63\times15}{57}= \frac{2735}{57} \approx 48\]

\textbf{b)} \[\text{Estimated mean} \approx 23.6\]

\[\quad \frac{16\times17+20\times3+24\times19+28\times10+32\times3+36\times5}{57}= \frac{1344}{57} \approx 23.6\]

\textbf{c)} \[\text{Estimated mean} \approx 43.2\]

\[\quad \frac{33\times15+41\times14+49\times20+57\times5}{54}= \frac{2334}{54} \approx 43.2\]

\textbf{d)} \[\text{Estimated mean} \approx 23.5\]

\[\quad \frac{12.5\times11+17.5\times16+22.5\times6+27.5\times11+32.5\times20}{64}= \frac{1505}{64} \approx 23.5\]

\end{tcolorbox}

{Midpoints} {Fx column}

\end{exercise}

\begin{exercise}[Modal Class
(Grouped)]\protect\hypertarget{exr-grp-2}{}\label{exr-grp-2}

Write down the modal class.

\begin{enumerate}
\def\labelenumi{\alph{enumi})}
\tightlist
\item
\end{enumerate}

\[{\renewcommand{\arraystretch}{1.6}\begin{array}{lc}
\hline \text{Class} & \text{Frequency} \\[4pt]
\hline 20 \leq x < 27  & 7 \\[4pt]
\hline 27 \leq x < 34  & 12 \\[4pt]
\hline 34 \leq x < 41  & 12 \\[4pt]
\hline 41 \leq x < 48  & 11 \\[4pt]
\hline \end{array}}\]

\begin{enumerate}
\def\labelenumi{\alph{enumi})}
\setcounter{enumi}{1}
\tightlist
\item
\end{enumerate}

\[{\renewcommand{\arraystretch}{1.6}\begin{array}{lc}
\hline \text{Class} & \text{Frequency} \\[4pt]
\hline 7 \leq x < 13  & 10 \\[4pt]
\hline 13 \leq x < 19  & 13 \\[4pt]
\hline 19 \leq x < 25  & 6 \\[4pt]
\hline 25 \leq x < 31  & 3 \\[4pt]
\hline 31 \leq x < 37  & 16 \\[4pt]
\hline 37 \leq x < 43  & 19 \\[4pt]
\hline \end{array}}\]

\begin{enumerate}
\def\labelenumi{\alph{enumi})}
\setcounter{enumi}{2}
\tightlist
\item
\end{enumerate}

\[{\renewcommand{\arraystretch}{1.6}\begin{array}{lc}
\hline \text{Class} & \text{Frequency} \\[4pt]
\hline 10 \leq x < 19  & 8 \\[4pt]
\hline 19 \leq x < 28  & 18 \\[4pt]
\hline 28 \leq x < 37  & 13 \\[4pt]
\hline 37 \leq x < 46  & 19 \\[4pt]
\hline 46 \leq x < 55  & 17 \\[4pt]
\hline \end{array}}\]

\begin{enumerate}
\def\labelenumi{\alph{enumi})}
\setcounter{enumi}{3}
\tightlist
\item
\end{enumerate}

\[{\renewcommand{\arraystretch}{1.6}\begin{array}{lc}
\hline \text{Class} & \text{Frequency} \\[4pt]
\hline 15 \leq x < 24  & 13 \\[4pt]
\hline 24 \leq x < 33  & 5 \\[4pt]
\hline 33 \leq x < 42  & 3 \\[4pt]
\hline 42 \leq x < 51  & 6 \\[4pt]
\hline 51 \leq x < 60  & 18 \\[4pt]
\hline 60 \leq x < 69  & 18 \\[4pt]
\hline \end{array}}\]

\begin{tcolorbox}[enhanced jigsaw, arc=.35mm, breakable, toptitle=1mm, title=\textcolor{quarto-callout-tip-color}{\faLightbulb}\hspace{0.5em}{🔍 View Solution}, colframe=quarto-callout-tip-color-frame, rightrule=.15mm, opacityback=0, titlerule=0mm, colback=white, coltitle=black, bottomtitle=1mm, toprule=.15mm, leftrule=.75mm, bottomrule=.15mm, colbacktitle=quarto-callout-tip-color!10!white, left=2mm, opacitybacktitle=0.6]

\textbf{a)} \[\text{Modal class: }27 \leq x < 34\]

\[\quad \text{Highest frequency = }12\text{, class: }27 \leq x < 34\]

\textbf{b)} \[\text{Modal class: }37 \leq x < 43\]

\[\quad \text{Highest frequency = }19\text{, class: }37 \leq x < 43\]

\textbf{c)} \[\text{Modal class: }37 \leq x < 46\]

\[\quad \text{Highest frequency = }19\text{, class: }37 \leq x < 46\]

\textbf{d)} \[\text{Modal class: }51 \leq x < 60\]

\[\quad \text{Highest frequency = }18\text{, class: }51 \leq x < 60\]

\end{tcolorbox}

{Reading grouped table}

\end{exercise}

∑ × π ÷ √ ◆ ∞ ± ∫ ≠ Δ

\subsection{🟡 Medium}

\begin{exercise}[Median Class
(Grouped)]\protect\hypertarget{exr-grp-3}{}\label{exr-grp-3}

Identify the median class.

\begin{enumerate}
\def\labelenumi{\alph{enumi})}
\tightlist
\item
\end{enumerate}

\[{\renewcommand{\arraystretch}{1.6}\begin{array}{lc}
\hline \text{Class} & \text{Frequency} \\[4pt]
\hline 20 \leq x < 27  & 11 \\[4pt]
\hline 27 \leq x < 34  & 8 \\[4pt]
\hline 34 \leq x < 41  & 8 \\[4pt]
\hline 41 \leq x < 48  & 8 \\[4pt]
\hline 48 \leq x < 55  & 19 \\[4pt]
\hline \end{array}}\]

\begin{enumerate}
\def\labelenumi{\alph{enumi})}
\setcounter{enumi}{1}
\tightlist
\item
\end{enumerate}

\[{\renewcommand{\arraystretch}{1.6}\begin{array}{lc}
\hline \text{Class} & \text{Frequency} \\[4pt]
\hline 18 \leq x < 28  & 10 \\[4pt]
\hline 28 \leq x < 38  & 9 \\[4pt]
\hline 38 \leq x < 48  & 18 \\[4pt]
\hline 48 \leq x < 58  & 14 \\[4pt]
\hline 58 \leq x < 68  & 16 \\[4pt]
\hline \end{array}}\]

\begin{enumerate}
\def\labelenumi{\alph{enumi})}
\setcounter{enumi}{2}
\tightlist
\item
\end{enumerate}

\[{\renewcommand{\arraystretch}{1.6}\begin{array}{lc}
\hline \text{Class} & \text{Frequency} \\[4pt]
\hline 14 \leq x < 23  & 15 \\[4pt]
\hline 23 \leq x < 32  & 15 \\[4pt]
\hline 32 \leq x < 41  & 12 \\[4pt]
\hline 41 \leq x < 50  & 18 \\[4pt]
\hline 50 \leq x < 59  & 14 \\[4pt]
\hline 59 \leq x < 68  & 12 \\[4pt]
\hline \end{array}}\]

\begin{enumerate}
\def\labelenumi{\alph{enumi})}
\setcounter{enumi}{3}
\tightlist
\item
\end{enumerate}

\[{\renewcommand{\arraystretch}{1.6}\begin{array}{lc}
\hline \text{Class} & \text{Frequency} \\[4pt]
\hline 9 \leq x < 19  & 8 \\[4pt]
\hline 19 \leq x < 29  & 12 \\[4pt]
\hline 29 \leq x < 39  & 19 \\[4pt]
\hline 39 \leq x < 49  & 13 \\[4pt]
\hline 49 \leq x < 59  & 3 \\[4pt]
\hline 59 \leq x < 69  & 9 \\[4pt]
\hline \end{array}}\]

\begin{tcolorbox}[enhanced jigsaw, arc=.35mm, breakable, toptitle=1mm, title=\textcolor{quarto-callout-tip-color}{\faLightbulb}\hspace{0.5em}{🔍 View Solution}, colframe=quarto-callout-tip-color-frame, rightrule=.15mm, opacityback=0, titlerule=0mm, colback=white, coltitle=black, bottomtitle=1mm, toprule=.15mm, leftrule=.75mm, bottomrule=.15mm, colbacktitle=quarto-callout-tip-color!10!white, left=2mm, opacitybacktitle=0.6]

\textbf{a)} \[\text{Median class: }34 \leq x < 41\]

\[\quad n=54\quad \frac{n}{2}=27\quad \text{Median class: }34 \leq x < 41\]

\textbf{b)} \[\text{Median class: }38 \leq x < 48\]

\[\quad n=67\quad \frac{n}{2}=33.5\quad \text{Median class: }38 \leq x < 48\]

\textbf{c)} \[\text{Median class: }41 \leq x < 50\]

\[\quad n=86\quad \frac{n}{2}=43\quad \text{Median class: }41 \leq x < 50\]

\textbf{d)} \[\text{Median class: }29 \leq x < 39\]

\[\quad n=64\quad \frac{n}{2}=32\quad \text{Median class: }29 \leq x < 39\]

\end{tcolorbox}

{Cumulative frequency} {Locating middle}

\end{exercise}

∑ × π ÷ √ ◆ ∞ ± ∫ ≠ Δ

\subsection{🔴 Challenging}

\begin{exercise}[Modal, Median \& Mean
(Grouped)]\protect\hypertarget{exr-grp-4}{}\label{exr-grp-4}

Find the modal class, median class, and estimated mean.

\begin{enumerate}
\def\labelenumi{\alph{enumi})}
\tightlist
\item
\end{enumerate}

\[\text{{Write down the modal class.}}\]

\begin{enumerate}
\def\labelenumi{\alph{enumi})}
\setcounter{enumi}{1}
\tightlist
\item
\end{enumerate}

\[\text{{Identify the median class.}}\]

\begin{enumerate}
\def\labelenumi{\alph{enumi})}
\setcounter{enumi}{2}
\tightlist
\item
\end{enumerate}

\[\text{{Estimate the mean.}}\]

\begin{tcolorbox}[enhanced jigsaw, arc=.35mm, breakable, toptitle=1mm, title=\textcolor{quarto-callout-tip-color}{\faLightbulb}\hspace{0.5em}{🔍 View Solution}, colframe=quarto-callout-tip-color-frame, rightrule=.15mm, opacityback=0, titlerule=0mm, colback=white, coltitle=black, bottomtitle=1mm, toprule=.15mm, leftrule=.75mm, bottomrule=.15mm, colbacktitle=quarto-callout-tip-color!10!white, left=2mm, opacitybacktitle=0.6]

\textbf{a)} \[\text{Modal class: }24 \leq x < 33\]

\[\quad \text{Highest frequency = }20\text{ → modal class: }24 \leq x < 33\]

\textbf{b)} \[\text{Median class: }33 \leq x < 42\]

\[\quad n=50\quad \frac{n}{2}=25\quad \text{Build CF to find class containing position }25\quad \text{Median class: }33 \leq x < 42\]

\textbf{c)} \[\bar{x} \approx 35.2\]

\[\quad \bar{x}=\dfrac{\sum fx}{\sum f}=\dfrac{19.5\times4+28.5\times20+37.5\times11+46.5\times15}{50}=\dfrac{1758}{50}\approx35.2\]

\end{tcolorbox}

{Midpoints} {Cumulative frequency} {Reading grouped table}

\end{exercise}

\begin{exercise}[Combined Mean (Two
Groups)]\protect\hypertarget{exr-grp-5}{}\label{exr-grp-5}

Find the combined mean of the two groups.

\begin{enumerate}
\def\labelenumi{\alph{enumi})}
\tightlist
\item
\end{enumerate}

\[n_1=50,\ \bar{x}_1=17.8\\[4pt]n_2=50,\ \bar{x}_2=10.1\]

\begin{enumerate}
\def\labelenumi{\alph{enumi})}
\setcounter{enumi}{1}
\tightlist
\item
\end{enumerate}

\[n_1=120,\ \bar{x}_1=23.1\\[4pt]n_2=50,\ \bar{x}_2=13.7\]

\begin{enumerate}
\def\labelenumi{\alph{enumi})}
\setcounter{enumi}{2}
\tightlist
\item
\end{enumerate}

\[n_1=80,\ \bar{x}_1=23.3\\[4pt]n_2=20,\ \bar{x}_2=15.5\]

\begin{enumerate}
\def\labelenumi{\alph{enumi})}
\setcounter{enumi}{3}
\tightlist
\item
\end{enumerate}

\[n_1=50,\ \bar{x}_1=23.9\\[4pt]n_2=50,\ \bar{x}_2=21.7\]

\begin{tcolorbox}[enhanced jigsaw, arc=.35mm, breakable, toptitle=1mm, title=\textcolor{quarto-callout-tip-color}{\faLightbulb}\hspace{0.5em}{🔍 View Solution}, colframe=quarto-callout-tip-color-frame, rightrule=.15mm, opacityback=0, titlerule=0mm, colback=white, coltitle=black, bottomtitle=1mm, toprule=.15mm, leftrule=.75mm, bottomrule=.15mm, colbacktitle=quarto-callout-tip-color!10!white, left=2mm, opacitybacktitle=0.6]

\textbf{a)} \[\text{Combined mean} = 13.95\]

\[\quad \frac{50\times17.8+50\times10.1}{50+50}=\frac{890.0+505.0}{100}=\frac{1395.0}{100}=13.95\]

\textbf{b)} \[\text{Combined mean} = 20.34\]

\[\quad \frac{120\times23.1+50\times13.7}{120+50}=\frac{2772.0+685.0}{170}=\frac{3457.0}{170}=20.34\]

\textbf{c)} \[\text{Combined mean} = 21.74\]

\[\quad \frac{80\times23.3+20\times15.5}{80+20}=\frac{1864.0+310.0}{100}=\frac{2174.0}{100}=21.74\]

\textbf{d)} \[\text{Combined mean} = 22.8\]

\[\quad \frac{50\times23.9+50\times21.7}{50+50}=\frac{1195.0+1085.0}{100}=\frac{2280.0}{100}=22.8\]

\end{tcolorbox}

{Weighted mean} {Total from mean}

\end{exercise}

\begin{exercise}[Find Missing Frequencies
(Grouped)]\protect\hypertarget{exr-grp-6}{}\label{exr-grp-6}

Find the values of a and b.

\begin{enumerate}
\def\labelenumi{\alph{enumi})}
\tightlist
\item
\end{enumerate}

\[n = 100,\quad \bar{x} \approx 43.5\\[6pt]{\renewcommand{\arraystretch}{1.6}\begin{array}{lc}
\hline \text{Class} & \text{Frequency} \\[4pt]
\hline 20 \leq x < 30 & 15 \\[4pt]
\hline 30 \leq x < 40 & 25 \\[4pt]
\hline 40 \leq x < 50 & 35 \\[4pt]
\hline 50 \leq x < 60 & a \\[4pt]
\hline 60 \leq x < 70 & b \\[4pt]
\hline \end{array}}\]

\begin{enumerate}
\def\labelenumi{\alph{enumi})}
\setcounter{enumi}{1}
\tightlist
\item
\end{enumerate}

\[n = 100,\quad \bar{x} \approx 43.5\\[6pt]{\renewcommand{\arraystretch}{1.6}\begin{array}{lc}
\hline \text{Class} & \text{Frequency} \\[4pt]
\hline 20 \leq x < 30 & 15 \\[4pt]
\hline 30 \leq x < 40 & 25 \\[4pt]
\hline 40 \leq x < 50 & 35 \\[4pt]
\hline 50 \leq x < 60 & a \\[4pt]
\hline 60 \leq x < 70 & b \\[4pt]
\hline \end{array}}\]

\begin{enumerate}
\def\labelenumi{\alph{enumi})}
\setcounter{enumi}{2}
\tightlist
\item
\end{enumerate}

\[n = 100,\quad \bar{x} \approx 43.5\\[6pt]{\renewcommand{\arraystretch}{1.6}\begin{array}{lc}
\hline \text{Class} & \text{Frequency} \\[4pt]
\hline 20 \leq x < 30 & 15 \\[4pt]
\hline 30 \leq x < 40 & 25 \\[4pt]
\hline 40 \leq x < 50 & 35 \\[4pt]
\hline 50 \leq x < 60 & a \\[4pt]
\hline 60 \leq x < 70 & b \\[4pt]
\hline \end{array}}\]

\begin{enumerate}
\def\labelenumi{\alph{enumi})}
\setcounter{enumi}{3}
\tightlist
\item
\end{enumerate}

\[n = 150,\quad \bar{x} \approx 48\\[6pt]{\renewcommand{\arraystretch}{1.6}\begin{array}{lc}
\hline \text{Class} & \text{Frequency} \\[4pt]
\hline 20 \leq x < 30 & 13 \\[4pt]
\hline 30 \leq x < 40 & 15 \\[4pt]
\hline 40 \leq x < 50 & 40 \\[4pt]
\hline 50 \leq x < 60 & a \\[4pt]
\hline 60 \leq x < 70 & b \\[4pt]
\hline \end{array}}\]

\begin{tcolorbox}[enhanced jigsaw, arc=.35mm, breakable, toptitle=1mm, title=\textcolor{quarto-callout-tip-color}{\faLightbulb}\hspace{0.5em}{🔍 View Solution}, colframe=quarto-callout-tip-color-frame, rightrule=.15mm, opacityback=0, titlerule=0mm, colback=white, coltitle=black, bottomtitle=1mm, toprule=.15mm, leftrule=.75mm, bottomrule=.15mm, colbacktitle=quarto-callout-tip-color!10!white, left=2mm, opacitybacktitle=0.6]

\textbf{a)} \[a = 10,\ b = 15\]

\[\quad \text{Step 1: }a+b=100-75=25\\[4pt]\text{Step 2: }\Sigma fx=43.5\times100=4350\\[4pt]\text{Known }\Sigma fx=2825\quad\Rightarrow 55a+65b=1525\\[4pt]\text{Step 3: Solve simultaneously: }a=10,\ b=15\]

\textbf{b)} \[a = 10,\ b = 15\]

\[\quad \text{Step 1: }a+b=100-75=25\\[4pt]\text{Step 2: }\Sigma fx=43.5\times100=4350\\[4pt]\text{Known }\Sigma fx=2825\quad\Rightarrow 55a+65b=1525\\[4pt]\text{Step 3: Solve simultaneously: }a=10,\ b=15\]

\textbf{c)} \[a = 10,\ b = 15\]

\[\quad \text{Step 1: }a+b=100-75=25\\[4pt]\text{Step 2: }\Sigma fx=43.5\times100=4350\\[4pt]\text{Known }\Sigma fx=2825\quad\Rightarrow 55a+65b=1525\\[4pt]\text{Step 3: Solve simultaneously: }a=10,\ b=15\]

\textbf{d)} \[a = 78,\ b = 4\]

\[\quad \text{Step 1: }a+b=150-68=82\\[4pt]\text{Step 2: }\Sigma fx=48\times150=7200\\[4pt]\text{Known }\Sigma fx=2650\quad\Rightarrow 55a+65b=4550\\[4pt]\text{Step 3: Solve simultaneously: }a=78,\ b=4\]

\end{tcolorbox}

{Simultaneous equations} {Mean formula}

\end{exercise}

∑ × π ÷ √ ◆ ∞ ± ∫ ≠ Δ

\begin{center}\rule{0.5\linewidth}{0.5pt}\end{center}

\section{📈 Part D --- Cumulative Frequency}\label{sec-cumfreq}

\begin{tcolorbox}[enhanced jigsaw, arc=.35mm, breakable, toptitle=1mm, title={🔑 Key Points}, colframe=quarto-callout-note-color-frame, rightrule=.15mm, opacityback=0, titlerule=0mm, colback=white, coltitle=black, bottomtitle=1mm, toprule=.15mm, leftrule=.75mm, bottomrule=.15mm, colbacktitle=quarto-callout-note-color!10!white, left=2mm, opacitybacktitle=0.6]

\begin{longtable}[]{@{}ll@{}}
\toprule\noalign{}
Measure & Position on CF axis \\
\midrule\noalign{}
\endhead
\bottomrule\noalign{}
\endlastfoot
Median (\(Q_2\)) & \(\frac{n}{2}\) \\
Lower quartile (\(Q_1\)) & \(\frac{n}{4}\) \\
Upper quartile (\(Q_3\)) & \(\frac{3n}{4}\) \\
\end{longtable}

\(\text{IQR} = Q_3 - Q_1\)

\end{tcolorbox}

\subsection{🟢 Easy}

\begin{exercise}[Build Cumulative Frequency
Table]\protect\hypertarget{exr-cf-1}{}\label{exr-cf-1}

Complete the cumulative frequency table.

\begin{enumerate}
\def\labelenumi{\alph{enumi})}
\tightlist
\item
\end{enumerate}

\[{\renewcommand{\arraystretch}{1.6}\begin{array}{lc}
\hline \text{Class} & f \\[4pt]
\hline 28 \leq x < 37 & 5 \\[4pt]
\hline 37 \leq x < 46 & 11 \\[4pt]
\hline 46 \leq x < 55 & 13 \\[4pt]
\hline 55 \leq x < 64 & 25 \\[4pt]
\hline 64 \leq x < 73 & 24 \\[4pt]
\hline 73 \leq x < 82 & 25 \\[4pt]
\hline
\end{array}}\]

\begin{tcolorbox}[enhanced jigsaw, arc=.35mm, breakable, toptitle=1mm, title=\textcolor{quarto-callout-tip-color}{\faLightbulb}\hspace{0.5em}{🔍 View Solution}, colframe=quarto-callout-tip-color-frame, rightrule=.15mm, opacityback=0, titlerule=0mm, colback=white, coltitle=black, bottomtitle=1mm, toprule=.15mm, leftrule=.75mm, bottomrule=.15mm, colbacktitle=quarto-callout-tip-color!10!white, left=2mm, opacitybacktitle=0.6]

\textbf{a)} \[\begin{array}{lc}
\hline x & \text{Cum. Freq} \\[4pt]
\hline x < 37 & 5 \\[4pt]
\hline x < 46 & 16 \\[4pt]
\hline x < 55 & 29 \\[4pt]
\hline x < 64 & 54 \\[4pt]
\hline x < 73 & 78 \\[4pt]
\hline x < 82 & 103 \\[4pt]
\hline
\end{array}\]

\[\quad \text{Running totals: }5, 16, 29, 54, 78, 103\]

\end{tcolorbox}

{Running totals}

\end{exercise}

\begin{exercise}[Read Cumulative Frequency
Table]\protect\hypertarget{exr-cf-2}{}\label{exr-cf-2}

Use the table to answer the questions.

\begin{enumerate}
\def\labelenumi{\alph{enumi})}
\tightlist
\item
\end{enumerate}

\[{\renewcommand{\arraystretch}{1.6}\begin{array}{lc}
\hline x & \text{Cum. Freq} \\[4pt]
\hline x < 39 & 12 \\[4pt]
\hline x < 54 & 30 \\[4pt]
\hline x < 69 & 34 \\[4pt]
\hline x < 84 & 54 \\[4pt]
\hline x < 99 & 62 \\[4pt]
\hline
\end{array}}\\[8pt]\text{(a) Items} \leq 54\text{?}\\[6pt]\text{(b) Items} > 54\text{?}\]

\begin{tcolorbox}[enhanced jigsaw, arc=.35mm, breakable, toptitle=1mm, title=\textcolor{quarto-callout-tip-color}{\faLightbulb}\hspace{0.5em}{🔍 View Solution}, colframe=quarto-callout-tip-color-frame, rightrule=.15mm, opacityback=0, titlerule=0mm, colback=white, coltitle=black, bottomtitle=1mm, toprule=.15mm, leftrule=.75mm, bottomrule=.15mm, colbacktitle=quarto-callout-tip-color!10!white, left=2mm, opacitybacktitle=0.6]

\textbf{a)} \[\text{(a) }30\quad\text{(b) }32\]

\[\quad \text{(a) CF at }54=30\quad\text{(b) }62-30=32\]

\end{tcolorbox}

{Reading cf table} {Subtraction}

\end{exercise}

∑ × π ÷ √ ◆ ∞ ± ∫ ≠ Δ

\subsection{🟡 Medium}

\begin{exercise}[IQR from Cumulative
Frequency]\protect\hypertarget{exr-cf-3}{}\label{exr-cf-3}

Use the cumulative frequency graph to estimate the IQR.

\begin{enumerate}
\def\labelenumi{\alph{enumi})}
\tightlist
\item
\end{enumerate}

\[\text{{Complete the CF table.}}\]

\begin{enumerate}
\def\labelenumi{\alph{enumi})}
\setcounter{enumi}{1}
\tightlist
\item
\end{enumerate}

\[\text{{Estimate the median from the graph.}}\]

\begin{enumerate}
\def\labelenumi{\alph{enumi})}
\setcounter{enumi}{2}
\tightlist
\item
\end{enumerate}

\[\text{{Estimate the IQR.}}\]

\begin{enumerate}
\def\labelenumi{\alph{enumi})}
\setcounter{enumi}{3}
\tightlist
\item
\end{enumerate}

\[\text{How many items are above }95?\]

\begin{tcolorbox}[enhanced jigsaw, arc=.35mm, breakable, toptitle=1mm, title=\textcolor{quarto-callout-tip-color}{\faLightbulb}\hspace{0.5em}{🔍 View Solution}, colframe=quarto-callout-tip-color-frame, rightrule=.15mm, opacityback=0, titlerule=0mm, colback=white, coltitle=black, bottomtitle=1mm, toprule=.15mm, leftrule=.75mm, bottomrule=.15mm, colbacktitle=quarto-callout-tip-color!10!white, left=2mm, opacitybacktitle=0.6]

\textbf{a)} \[{\renewcommand{\arraystretch}{1.6}\begin{array}{lc}
\hline x & \text{Cum. Freq} \\[4pt]
\hline x < 35 & 18 \\[4pt]
\hline x < 54 & 23 \\[4pt]
\hline x < 73 & 32 \\[4pt]
\hline x < 92 & 55 \\[4pt]
\hline x < 111 & 69 \\[4pt]
\hline x < 130 & 76 \\[4pt]
\hline x < 149 & 88 \\[4pt]
\hline
\end{array}}\]

\[\quad \text{Running totals: }18, 23, 32, 55, 69, 76, 88\]

\textbf{b)} \[\text{Median} \approx 82.9\]

\[\quad \text{Read at }\frac{n}{2}=\frac{88}{2}=44\text{ on CF axis}\quad \text{Median}\approx82.9\]

\textbf{c)} \[\text{IQR} = 56.7\]

\[\quad Q_1\approx50.2\text{ (read at }22.0\text{)}\quad Q_3\approx106.9\text{ (read at }66.0\text{)}\quad \text{IQR}=106.9-50.2=56.7\]

\textbf{d)} \[\approx 31\text{ items}\]

\[\quad \text{Read CF at }95\approx57\quad \text{Above: }88-57=31\]

\end{tcolorbox}

{Interpolation} {Quartile location}

\end{exercise}

∑ × π ÷ √ ◆ ∞ ± ∫ ≠ Δ

\subsection{🔴 Challenging}

\begin{exercise}[Estimate Frequency Above a
Value]\protect\hypertarget{exr-cf-4}{}\label{exr-cf-4}

Estimate the number of items above the given value.

\begin{enumerate}
\def\labelenumi{\alph{enumi})}
\tightlist
\item
\end{enumerate}

\[{\renewcommand{\arraystretch}{1.6}\begin{array}{lc}
\hline x & \text{CF} \\[4pt]
\hline x < 32 & 25 \\[4pt]
\hline x < 50 & 45 \\[4pt]
\hline x < 68 & 52 \\[4pt]
\hline x < 86 & 69 \\[4pt]
\hline x < 104 & 74 \\[4pt]
\hline x < 122 & 88 \\[4pt]
\hline
\end{array}}\\[8pt]\text{How many items above }93?\]

\begin{tcolorbox}[enhanced jigsaw, arc=.35mm, breakable, toptitle=1mm, title=\textcolor{quarto-callout-tip-color}{\faLightbulb}\hspace{0.5em}{🔍 View Solution}, colframe=quarto-callout-tip-color-frame, rightrule=.15mm, opacityback=0, titlerule=0mm, colback=white, coltitle=black, bottomtitle=1mm, toprule=.15mm, leftrule=.75mm, bottomrule=.15mm, colbacktitle=quarto-callout-tip-color!10!white, left=2mm, opacitybacktitle=0.6]

\textbf{a)} \[\approx 17\]

\[\quad \text{Read CF at }93\approx71\quad 88-71=17\]

\end{tcolorbox}

{Interpolation} {Reading cf graph}

\end{exercise}

\begin{exercise}[Find Missing CF Values A and
B]\protect\hypertarget{exr-cf-5}{}\label{exr-cf-5}

Find the values of A and B in the cumulative frequency table.

\begin{enumerate}
\def\labelenumi{\alph{enumi})}
\tightlist
\item
\end{enumerate}

\[\text{Find A and B.}\\[6pt]{\renewcommand{\arraystretch}{1.6}\begin{array}{lc}
\hline x & \text{Cum. Freq} \\[4pt]
\hline x < 21 & 9 \\[4pt]
\hline x < 28 & A \\[4pt]
\hline x < 35 & 42 \\[4pt]
\hline x < 42 & B \\[4pt]
\hline x < 49 & 73 \\[4pt]
\hline \end{array}}\]

\begin{tcolorbox}[enhanced jigsaw, arc=.35mm, breakable, toptitle=1mm, title=\textcolor{quarto-callout-tip-color}{\faLightbulb}\hspace{0.5em}{🔍 View Solution}, colframe=quarto-callout-tip-color-frame, rightrule=.15mm, opacityback=0, titlerule=0mm, colback=white, coltitle=black, bottomtitle=1mm, toprule=.15mm, leftrule=.75mm, bottomrule=.15mm, colbacktitle=quarto-callout-tip-color!10!white, left=2mm, opacitybacktitle=0.6]

\textbf{a)} \[A = 29,\ B = 55\]

\[\quad A = 9+20=29\quad B = 42+13=55\]

\end{tcolorbox}

{Running totals} {Algebraic reasoning}

\end{exercise}

\begin{exercise}[Cumulative Frequency --- Full
Analysis]\protect\hypertarget{exr-cf-6}{}\label{exr-cf-6}

Use the frequency table and cumulative frequency graph to answer the
questions.

\begin{enumerate}
\def\labelenumi{\alph{enumi})}
\tightlist
\item
\end{enumerate}

\[\text{Complete the cumulative frequency table.}\]

\begin{enumerate}
\def\labelenumi{\alph{enumi})}
\setcounter{enumi}{1}
\item
  \(\text{Estimate the median from the graph.}\)
\item
  \(\text{Estimate }Q_1\text{ and }Q_3.\)
\item
  \(\text{Calculate the IQR.}\)
\item
  \(\text{Estimate how many items are above }57.\)
\end{enumerate}

\begin{tcolorbox}[enhanced jigsaw, arc=.35mm, breakable, toptitle=1mm, title=\textcolor{quarto-callout-tip-color}{\faLightbulb}\hspace{0.5em}{🔍 View Solution}, colframe=quarto-callout-tip-color-frame, rightrule=.15mm, opacityback=0, titlerule=0mm, colback=white, coltitle=black, bottomtitle=1mm, toprule=.15mm, leftrule=.75mm, bottomrule=.15mm, colbacktitle=quarto-callout-tip-color!10!white, left=2mm, opacitybacktitle=0.6]

\textbf{a)} \[\begin{array}{lc}
\hline x & \text{CF} \\[4pt]
\hline x < 39 & 19 \\[4pt]
\hline x < 50 & 36 \\[4pt]
\hline x < 61 & 56 \\[4pt]
\hline x < 72 & 62 \\[4pt]
\hline x < 83 & 71 \\[4pt]
\hline
\end{array}\]

\[\quad \text{Running totals: }19, 36, 56, 62, 71\]

\textbf{b)} \(\text{Median} \approx 49.7\)

\[\quad \text{Read at }\frac{n}{2}=\frac{71}{2}=35\quad \Rightarrow \text{Median}\approx49.7\]

\textbf{c)} \(Q_1 \approx 38.3,\quad Q_3 \approx 59.5\)

\[\quad Q_1\text{: read at }\frac{n}{4}=17.8\Rightarrow Q_1\approx38.3\\[4pt]Q_3\text{: read at }\frac{3n}{4}=53.2\Rightarrow Q_3\approx59.5\]

\textbf{d)} \(\text{IQR} = 21.2\)

\[\quad \text{IQR} = Q_3 - Q_1 = 59.5 - 38.3 = 21.2\]

\textbf{e)} \(\approx 22\text{ items}\)

\[\quad \text{Read CF at }57\approx49\quad 71-49=22\]

\end{tcolorbox}

{Cumulative frequency} {Interpolation} {Quartile location} {Iqr
calculation} {Reading cf graph}

\end{exercise}

∑ × π ÷ √ ◆ ∞ ± ∫ ≠ Δ

\begin{center}\rule{0.5\linewidth}{0.5pt}\end{center}

\section{🎯 Part E --- Box Plots and IQR}\label{sec-boxplot}

\begin{tcolorbox}[enhanced jigsaw, arc=.35mm, breakable, toptitle=1mm, title={🔑 Key Points}, colframe=quarto-callout-note-color-frame, rightrule=.15mm, opacityback=0, titlerule=0mm, colback=white, coltitle=black, bottomtitle=1mm, toprule=.15mm, leftrule=.75mm, bottomrule=.15mm, colbacktitle=quarto-callout-note-color!10!white, left=2mm, opacitybacktitle=0.6]

\[\text{IQR} = Q_3 - Q_1\]

\textbf{Outlier rule:} A value is an outlier if it is
\(> Q_3 + 1.5 \times \text{IQR}\) or \(< Q_1 - 1.5 \times \text{IQR}\)

\end{tcolorbox}

\subsection{🟢 Easy}

\begin{exercise}[IQR from Raw
Data]\protect\hypertarget{exr-bp-1}{}\label{exr-bp-1}

Find the interquartile range.

\begin{enumerate}
\def\labelenumi{\alph{enumi})}
\tightlist
\item
  \(4,\ 9,\ 9,\ 10,\ 12,\ 12,\ 13,\ 17,\ 19,\ 22,\ 23,\ 24,\ 26,\ 29,\ 30\)
\end{enumerate}

\begin{tcolorbox}[enhanced jigsaw, arc=.35mm, breakable, toptitle=1mm, title=\textcolor{quarto-callout-tip-color}{\faLightbulb}\hspace{0.5em}{🔍 View Solution}, colframe=quarto-callout-tip-color-frame, rightrule=.15mm, opacityback=0, titlerule=0mm, colback=white, coltitle=black, bottomtitle=1mm, toprule=.15mm, leftrule=.75mm, bottomrule=.15mm, colbacktitle=quarto-callout-tip-color!10!white, left=2mm, opacitybacktitle=0.6]

\textbf{a)} \(\text{IQR} = 14\)

\[\quad Q_1=10,\ Q_3=24\quad \text{IQR}=24-10=14\]

\end{tcolorbox}

{Quartile calculation} {Subtraction}

\end{exercise}

∑ × π ÷ √ ◆ ∞ ± ∫ ≠ Δ

\subsection{🟡 Medium}

\begin{exercise}[Draw Box Plot from 5-Number
Summary]\protect\hypertarget{exr-bp-2}{}\label{exr-bp-2}

Find the IQR and draw a box plot.

\begin{enumerate}
\def\labelenumi{\alph{enumi})}
\tightlist
\item
\end{enumerate}

\[\min=10,\ Q_1=13,\ \text{med}=16,\ Q_3=20,\ \max=24\]

\begin{tcolorbox}[enhanced jigsaw, arc=.35mm, breakable, toptitle=1mm, title=\textcolor{quarto-callout-tip-color}{\faLightbulb}\hspace{0.5em}{🔍 View Solution}, colframe=quarto-callout-tip-color-frame, rightrule=.15mm, opacityback=0, titlerule=0mm, colback=white, coltitle=black, bottomtitle=1mm, toprule=.15mm, leftrule=.75mm, bottomrule=.15mm, colbacktitle=quarto-callout-tip-color!10!white, left=2mm, opacitybacktitle=0.6]

\textbf{a)} \[\text{IQR}=7\]

\[\quad \text{IQR}=Q_3-Q_1=20-13=7\]

\end{tcolorbox}

{Five number summary} {Box plot construction}

\end{exercise}

\begin{exercise}[Find Unknown Values from Box Plot
Conditions]\protect\hypertarget{exr-bp-4}{}\label{exr-bp-4}

Find the unknown values using the given conditions.

\begin{enumerate}
\def\labelenumi{\alph{enumi})}
\tightlist
\item
\end{enumerate}

\[h,\ 5,\ 5,\ 7,\ j,\ 15,\ k,\ k\\[6pt]\text{Median}=8,\ \text{Mode}=20,\ \text{Range}=17\]

\begin{tcolorbox}[enhanced jigsaw, arc=.35mm, breakable, toptitle=1mm, title=\textcolor{quarto-callout-tip-color}{\faLightbulb}\hspace{0.5em}{🔍 View Solution}, colframe=quarto-callout-tip-color-frame, rightrule=.15mm, opacityback=0, titlerule=0mm, colback=white, coltitle=black, bottomtitle=1mm, toprule=.15mm, leftrule=.75mm, bottomrule=.15mm, colbacktitle=quarto-callout-tip-color!10!white, left=2mm, opacitybacktitle=0.6]

\textbf{a)} \[h=3,\ j=9,\ k=20\]

\[\quad k=20\text{ (mode, appears twice)}\quad h=k-17=3\text{ (range)}\quad j=2\times8-7=9\text{ (median)}\]

\end{tcolorbox}

{Simultaneous conditions} {Median rule} {Mode rule} {Range rule}

\end{exercise}

\begin{exercise}[Read a Box
Plot]\protect\hypertarget{exr-bp-5}{}\label{exr-bp-5}

Use the box plot to answer the questions.

\begin{enumerate}
\def\labelenumi{\alph{enumi})}
\item
  \(\text{Write down the minimum value.}\)
\item
  \(\text{Write down }Q_1.\)
\item
  \(\text{Write down the median.}\)
\item
  \(\text{Write down }Q_3.\)
\item
  \(\text{Calculate the IQR.}\)
\end{enumerate}

\begin{tcolorbox}[enhanced jigsaw, arc=.35mm, breakable, toptitle=1mm, title=\textcolor{quarto-callout-tip-color}{\faLightbulb}\hspace{0.5em}{🔍 View Solution}, colframe=quarto-callout-tip-color-frame, rightrule=.15mm, opacityback=0, titlerule=0mm, colback=white, coltitle=black, bottomtitle=1mm, toprule=.15mm, leftrule=.75mm, bottomrule=.15mm, colbacktitle=quarto-callout-tip-color!10!white, left=2mm, opacitybacktitle=0.6]

\textbf{a)} \(\text{Min} = 15\)

\[\quad \text{Read the left whisker end: }15\]

\textbf{b)} \(Q_1 = 21\)

\[\quad \text{Read the left edge of the box: }21\]

\textbf{c)} \(\text{Median} = 27\)

\[\quad \text{Read the line inside the box: }27\]

\textbf{d)} \(Q_3 = 35\)

\[\quad \text{Read the right edge of the box: }35\]

\textbf{e)} \(\text{IQR} = 14\)

\[\quad \text{IQR} = Q_3 - Q_1 = 35 - 21 = 14\]

\end{tcolorbox}

{Reading box plot} {Quartile identification} {Iqr calculation}

\end{exercise}

∑ × π ÷ √ ◆ ∞ ± ∫ ≠ Δ

\subsection{🔴 Challenging}

\begin{exercise}[Identify Outliers (1.5 × IQR
Rule)]\protect\hypertarget{exr-bp-3}{}\label{exr-bp-3}

Identify outliers using the 1.5 × IQR rule and draw a box plot.

\begin{enumerate}
\def\labelenumi{\alph{enumi})}
\tightlist
\item
\end{enumerate}

\[14,\ 22,\ 24,\ 30,\ 31,\ 37,\ 48,\ 48,\ 65,\ 67,\ 107\]

\begin{tcolorbox}[enhanced jigsaw, arc=.35mm, breakable, toptitle=1mm, title=\textcolor{quarto-callout-tip-color}{\faLightbulb}\hspace{0.5em}{🔍 View Solution}, colframe=quarto-callout-tip-color-frame, rightrule=.15mm, opacityback=0, titlerule=0mm, colback=white, coltitle=black, bottomtitle=1mm, toprule=.15mm, leftrule=.75mm, bottomrule=.15mm, colbacktitle=quarto-callout-tip-color!10!white, left=2mm, opacitybacktitle=0.6]

\textbf{a)}
\[Q_1=24,\ Q_3=65,\ \text{IQR}=41\\\text{Upper fence}=126.5\\107>126.5\Rightarrow\text{outlier}\]

\[\quad \text{Step 1: }Q_1=24,\ Q_3=65\\[4pt]\text{Step 2: IQR}=Q_3-Q_1=65-24=41\\[4pt]\text{Step 3: Upper fence}=Q_3+1.5\times\text{IQR}=65+1.5\times41=126.5\\[4pt]\text{Step 4: }107>126.5\therefore\text{ outlier }\checkmark\]

\end{tcolorbox}

{Outlier rule} {Quartile calculation}

\end{exercise}

∑ × π ÷ √ ◆ ∞ ± ∫ ≠ Δ

\begin{center}\rule{0.5\linewidth}{0.5pt}\end{center}

\section{⚡ Test Your Knowledge}\label{sec-worksheet}

\chapter{Unit 5: Linear Equations}\label{unit-5-linear-equations}

\section{🎯 Learning Targets}\label{learning-targets-4}

\begin{longtable}[]{@{}
  >{\raggedright\arraybackslash}p{(\linewidth - 4\tabcolsep) * \real{0.1364}}
  >{\raggedright\arraybackslash}p{(\linewidth - 4\tabcolsep) * \real{0.3636}}
  >{\raggedright\arraybackslash}p{(\linewidth - 4\tabcolsep) * \real{0.5000}}@{}}
\toprule\noalign{}
\begin{minipage}[b]{\linewidth}\raggedright
\#
\end{minipage} & \begin{minipage}[b]{\linewidth}\raggedright
I can\ldots{}
\end{minipage} & \begin{minipage}[b]{\linewidth}\raggedright
Key Skill
\end{minipage} \\
\midrule\noalign{}
\endhead
\bottomrule\noalign{}
\endlastfoot
LT1 & Solve one-step equations using inverse operations & Addition,
subtraction, multiplication, division \\
LT2 & Solve two-step equations, including brackets and fractions & Two
inverse operations in sequence \\
LT3 & Solve equations with unknowns on both sides & Collecting like
terms \\
LT4 & Expand brackets and solve the resulting equation & Single and
double brackets \\
LT5 & Solve equations involving algebraic fractions & Clearing
fractions, cross-multiplication \\
LT6 & Form and solve equations from word problems & Number riddles,
ages, geometry, formulas \\
\end{longtable}

\section{🗂️ Formula Card}\label{formula-card-1}

\subsection*{Linear Equations --- Solving
Strategies}\label{linear-equations-solving-strategies}
\addcontentsline{toc}{subsection}{Linear Equations --- Solving
Strategies}

\subsubsection*{The Balance Method}\label{the-balance-method}
\addcontentsline{toc}{subsubsection}{The Balance Method}

\textbf{💡 Golden Rule} Whatever you do to one side, you \textbf{must}
do to the other side.

\textbf{One-step} \[x + 5 = 12 \;\Rightarrow\; x = 12 - 5 = 7\]
\[3x = 21 \;\Rightarrow\; x = 21 \div 3 = 7\]

\textbf{Two-step} --- undo addition/subtraction \emph{first}, then
multiplication/division
\[2x + 3 = 11 \;\Rightarrow\; 2x = 8 \;\Rightarrow\; x = 4\]

\subsubsection*{Unknowns on Both Sides}\label{unknowns-on-both-sides}
\addcontentsline{toc}{subsubsection}{Unknowns on Both Sides}

\textbf{🎯 Strategy} Move all \(x\) terms to one side, all numbers to
the other.

\textbf{Method}: \[5x + 2 = 2x + 8\] \[5x - 2x = 8 - 2\]
\[3x = 6 \;\Rightarrow\; x = 2\]

\subsubsection*{Expanding Brackets}\label{expanding-brackets}
\addcontentsline{toc}{subsubsection}{Expanding Brackets}

\textbf{⚠️ Remember} The number outside multiplies \textbf{every} term
inside the bracket.

\textbf{Single bracket}:
\[3(2x + 5) = 27 \;\Rightarrow\; 6x + 15 = 27 \;\Rightarrow\; x = 2\]

\textbf{Two brackets}: \[2(x + 3) + 4(x - 1) = 18\]
\[2x + 6 + 4x - 4 = 18\]
\[6x + 2 = 18 \;\Rightarrow\; x = \dfrac{8}{3}\]

\subsubsection*{Algebraic Fractions}\label{algebraic-fractions}
\addcontentsline{toc}{subsubsection}{Algebraic Fractions}

\textbf{📝 Key Move} Multiply every term by the denominator (or LCM) to
clear fractions first.

\textbf{Single fraction}:
\[\dfrac{2x + 1}{3} = 5 \;\Rightarrow\; 2x + 1 = 15 \;\Rightarrow\; x = 7\]

\textbf{Fractions on both sides --- cross multiply}:
\[\dfrac{x + 1}{2} = \dfrac{x + 3}{4} \;\Rightarrow\; 4(x+1) = 2(x+3) \;\Rightarrow\; x = 1\]

\textbf{Adding fractions --- use LCM}:
\[\dfrac{x}{2} + \dfrac{x}{3} = 5 \;\xrightarrow{\times 6}\; 3x + 2x = 30 \;\Rightarrow\; x = 6\]

\subsubsection*{Forming Equations from
Words}\label{forming-equations-from-words}
\addcontentsline{toc}{subsubsection}{Forming Equations from Words}

\textbf{Step-by-step}:

\begin{enumerate}
\def\labelenumi{\arabic{enumi}.}
\tightlist
\item
  Define your variable --- \emph{let \(x = \ldots\)}
\item
  Translate each condition into algebra
\item
  Form the equation
\item
  Solve and check your answer makes sense
\end{enumerate}

\textbf{Example}: Three consecutive integers sum to 51.
\[x + (x+1) + (x+2) = 51 \;\Rightarrow\; 3x + 3 = 51 \;\Rightarrow\; x = 16\]
Numbers: 16, 17, 18. ✓

\section{📝 Practice Problems}\label{practice-problems-2}

\subsection{🟢 Easy}

\begin{exercise}[]\protect\hypertarget{exr-u5-1}{}\label{exr-u5-1}

~

\textbf{a)} \(\displaystyle t - 9 = -14\)

\textbf{b)} \(\displaystyle t - 4 = 0\)

\textbf{c)} \(\displaystyle y - 8 = -9\)

\textbf{d)} \(\displaystyle t + 11 = 19\)

\textbf{e)} \(\displaystyle x + 6 = 12\)

\textbf{f)} \(\displaystyle x - 12 = -10\)

\textbf{g)} \(\displaystyle y - 5 = 17\)

\textbf{h)} \(\displaystyle p - 12 = 20\)

\textbf{i)} \(\displaystyle y + 10 = 19\)

\textbf{j)} \(\displaystyle m - 5 = 13\)

\textbf{k)} \(\displaystyle p - 4 = -15\)

\textbf{l)} \(\displaystyle y + 10 = 12\)

\begin{tcolorbox}[enhanced jigsaw, arc=.35mm, breakable, toptitle=1mm, title=\textcolor{quarto-callout-tip-color}{\faLightbulb}\hspace{0.5em}{🔍 View Solutions}, colframe=quarto-callout-tip-color-frame, rightrule=.15mm, opacityback=0, titlerule=0mm, colback=white, coltitle=black, bottomtitle=1mm, toprule=.15mm, leftrule=.75mm, bottomrule=.15mm, colbacktitle=quarto-callout-tip-color!10!white, left=2mm, opacitybacktitle=0.6]

\textbf{a)} \(\displaystyle t - 9 = -14\)

\begin{align*}
t = -5
\end{align*}

\textbf{b)} \(\displaystyle t - 4 = 0\)

\begin{align*}
t = 4
\end{align*}

\textbf{c)} \(\displaystyle y - 8 = -9\)

\begin{align*}
y = -1
\end{align*}

\textbf{d)} \(\displaystyle t + 11 = 19\)

\begin{align*}
t = 8
\end{align*}

\textbf{e)} \(\displaystyle x + 6 = 12\)

\begin{align*}
x = 6
\end{align*}

\textbf{f)} \(\displaystyle x - 12 = -10\)

\begin{align*}
x = 2
\end{align*}

\textbf{g)} \(\displaystyle y - 5 = 17\)

\begin{align*}
y = 22
\end{align*}

\textbf{h)} \(\displaystyle p - 12 = 20\)

\begin{align*}
p = 32
\end{align*}

\textbf{i)} \(\displaystyle y + 10 = 19\)

\begin{align*}
y = 9
\end{align*}

\textbf{j)} \(\displaystyle m - 5 = 13\)

\begin{align*}
m = 18
\end{align*}

\textbf{k)} \(\displaystyle p - 4 = -15\)

\begin{align*}
p = -11
\end{align*}

\textbf{l)} \(\displaystyle y + 10 = 12\)

\begin{align*}
y = 2
\end{align*}

\end{tcolorbox}

\end{exercise}

\begin{exercise}[]\protect\hypertarget{exr-u5-2}{}\label{exr-u5-2}

~

\textbf{a)} \(\displaystyle 6n = -36\)

\textbf{b)} \(\displaystyle 7y = -49\)

\textbf{c)} \(\displaystyle 8x = -56\)

\textbf{d)} \(\displaystyle 7n = 63\)

\textbf{e)} \(\displaystyle 2n = 26\)

\textbf{f)} \(\displaystyle 10m = -70\)

\textbf{g)} \(\displaystyle 3m = 21\)

\textbf{h)} \(\displaystyle 12n = 108\)

\textbf{i)} \(\displaystyle 11n = -44\)

\textbf{j)} \(\displaystyle 3t = -27\)

\textbf{k)} \(\displaystyle 5t = 70\)

\textbf{l)} \(\displaystyle 3n = -9\)

\begin{tcolorbox}[enhanced jigsaw, arc=.35mm, breakable, toptitle=1mm, title=\textcolor{quarto-callout-tip-color}{\faLightbulb}\hspace{0.5em}{🔍 View Solutions}, colframe=quarto-callout-tip-color-frame, rightrule=.15mm, opacityback=0, titlerule=0mm, colback=white, coltitle=black, bottomtitle=1mm, toprule=.15mm, leftrule=.75mm, bottomrule=.15mm, colbacktitle=quarto-callout-tip-color!10!white, left=2mm, opacitybacktitle=0.6]

\textbf{a)} \(\displaystyle 6n = -36\)

\begin{align*}
n = -6
\end{align*}

\textbf{b)} \(\displaystyle 7y = -49\)

\begin{align*}
y = -7
\end{align*}

\textbf{c)} \(\displaystyle 8x = -56\)

\begin{align*}
x = -7
\end{align*}

\textbf{d)} \(\displaystyle 7n = 63\)

\begin{align*}
n = 9
\end{align*}

\textbf{e)} \(\displaystyle 2n = 26\)

\begin{align*}
n = 13
\end{align*}

\textbf{f)} \(\displaystyle 10m = -70\)

\begin{align*}
m = -7
\end{align*}

\textbf{g)} \(\displaystyle 3m = 21\)

\begin{align*}
m = 7
\end{align*}

\textbf{h)} \(\displaystyle 12n = 108\)

\begin{align*}
n = 9
\end{align*}

\textbf{i)} \(\displaystyle 11n = -44\)

\begin{align*}
n = -4
\end{align*}

\textbf{j)} \(\displaystyle 3t = -27\)

\begin{align*}
t = -9
\end{align*}

\textbf{k)} \(\displaystyle 5t = 70\)

\begin{align*}
t = 14
\end{align*}

\textbf{l)} \(\displaystyle 3n = -9\)

\begin{align*}
n = -3
\end{align*}

\end{tcolorbox}

\end{exercise}

\begin{exercise}[]\protect\hypertarget{exr-u5-3}{}\label{exr-u5-3}

~

\textbf{a)} \(\displaystyle \dfrac{x}{8} = -2\)

\textbf{b)} \(\displaystyle \dfrac{m}{7} = -5\)

\textbf{c)} \(\displaystyle \dfrac{n}{7} = -4\)

\textbf{d)} \(\displaystyle \dfrac{t}{6} = 12\)

\textbf{e)} \(\displaystyle \dfrac{t}{3} = 9\)

\textbf{f)} \(\displaystyle \dfrac{t}{4} = 7\)

\textbf{g)} \(\displaystyle \dfrac{t}{5} = -5\)

\textbf{h)} \(\displaystyle \dfrac{m}{8} = -2\)

\begin{tcolorbox}[enhanced jigsaw, arc=.35mm, breakable, toptitle=1mm, title=\textcolor{quarto-callout-tip-color}{\faLightbulb}\hspace{0.5em}{🔍 View Solutions}, colframe=quarto-callout-tip-color-frame, rightrule=.15mm, opacityback=0, titlerule=0mm, colback=white, coltitle=black, bottomtitle=1mm, toprule=.15mm, leftrule=.75mm, bottomrule=.15mm, colbacktitle=quarto-callout-tip-color!10!white, left=2mm, opacitybacktitle=0.6]

\textbf{a)} \(\displaystyle \dfrac{x}{8} = -2\)

\begin{align*}
x = -16
\end{align*}

\textbf{b)} \(\displaystyle \dfrac{m}{7} = -5\)

\begin{align*}
m = -35
\end{align*}

\textbf{c)} \(\displaystyle \dfrac{n}{7} = -4\)

\begin{align*}
n = -28
\end{align*}

\textbf{d)} \(\displaystyle \dfrac{t}{6} = 12\)

\begin{align*}
t = 72
\end{align*}

\textbf{e)} \(\displaystyle \dfrac{t}{3} = 9\)

\begin{align*}
t = 27
\end{align*}

\textbf{f)} \(\displaystyle \dfrac{t}{4} = 7\)

\begin{align*}
t = 28
\end{align*}

\textbf{g)} \(\displaystyle \dfrac{t}{5} = -5\)

\begin{align*}
t = -25
\end{align*}

\textbf{h)} \(\displaystyle \dfrac{m}{8} = -2\)

\begin{align*}
m = -16
\end{align*}

\end{tcolorbox}

\end{exercise}

\begin{exercise}[]\protect\hypertarget{exr-u5-4}{}\label{exr-u5-4}

~

\textbf{a)} \(\displaystyle 10t - 3 = -73\)

\textbf{b)} \(\displaystyle 2y + 2 = -10\)

\textbf{c)} \(\displaystyle 7y - 4 = 10\)

\textbf{d)} \(\displaystyle 9t - 6 = -42\)

\textbf{e)} \(\displaystyle 10y + 7 = 107\)

\textbf{f)} \(\displaystyle 8m + 1 = -31\)

\textbf{g)} \(\displaystyle 9p - 8 = -53\)

\textbf{h)} \(\displaystyle 4y + 3 = 19\)

\begin{tcolorbox}[enhanced jigsaw, arc=.35mm, breakable, toptitle=1mm, title=\textcolor{quarto-callout-tip-color}{\faLightbulb}\hspace{0.5em}{🔍 View Solutions}, colframe=quarto-callout-tip-color-frame, rightrule=.15mm, opacityback=0, titlerule=0mm, colback=white, coltitle=black, bottomtitle=1mm, toprule=.15mm, leftrule=.75mm, bottomrule=.15mm, colbacktitle=quarto-callout-tip-color!10!white, left=2mm, opacitybacktitle=0.6]

\textbf{a)} \(\displaystyle 10t - 3 = -73\)

\begin{align*}
10t = -70 \\ t = -7
\end{align*}

\textbf{b)} \(\displaystyle 2y + 2 = -10\)

\begin{align*}
2y = -12 \\ y = -6
\end{align*}

\textbf{c)} \(\displaystyle 7y - 4 = 10\)

\begin{align*}
7y = 14 \\ y = 2
\end{align*}

\textbf{d)} \(\displaystyle 9t - 6 = -42\)

\begin{align*}
9t = -36 \\ t = -4
\end{align*}

\textbf{e)} \(\displaystyle 10y + 7 = 107\)

\begin{align*}
10y = 100 \\ y = 10
\end{align*}

\textbf{f)} \(\displaystyle 8m + 1 = -31\)

\begin{align*}
8m = -32 \\ m = -4
\end{align*}

\textbf{g)} \(\displaystyle 9p - 8 = -53\)

\begin{align*}
9p = -45 \\ p = -5
\end{align*}

\textbf{h)} \(\displaystyle 4y + 3 = 19\)

\begin{align*}
4y = 16 \\ y = 4
\end{align*}

\end{tcolorbox}

\end{exercise}

∑ × π ÷ √ ◆ ∞ ± ∫ ≠ Δ

\subsection{🟡 Medium}

\begin{exercise}[]\protect\hypertarget{exr-u5-5}{}\label{exr-u5-5}

~

\textbf{a)} \(\displaystyle \dfrac{m + 6}{6} = 2\)

\textbf{b)} \(\displaystyle \dfrac{n - 8}{6} = 1\)

\textbf{c)} \(\displaystyle \dfrac{n + 2}{8} = -1\)

\textbf{d)} \(\displaystyle \dfrac{t + 8}{7} = 2\)

\textbf{e)} \(\displaystyle \dfrac{x + 1}{8} = 1\)

\textbf{f)} \(\displaystyle \dfrac{p + 2}{2} = -4\)

\begin{tcolorbox}[enhanced jigsaw, arc=.35mm, breakable, toptitle=1mm, title=\textcolor{quarto-callout-tip-color}{\faLightbulb}\hspace{0.5em}{🔍 View Solutions}, colframe=quarto-callout-tip-color-frame, rightrule=.15mm, opacityback=0, titlerule=0mm, colback=white, coltitle=black, bottomtitle=1mm, toprule=.15mm, leftrule=.75mm, bottomrule=.15mm, colbacktitle=quarto-callout-tip-color!10!white, left=2mm, opacitybacktitle=0.6]

\textbf{a)} \(\displaystyle \dfrac{m + 6}{6} = 2\)

\begin{align*}
m + 6 = 12 \\ m = 6
\end{align*}

\textbf{b)} \(\displaystyle \dfrac{n - 8}{6} = 1\)

\begin{align*}
n - 8 = 6 \\ n = 14
\end{align*}

\textbf{c)} \(\displaystyle \dfrac{n + 2}{8} = -1\)

\begin{align*}
n + 2 = -8 \\ n = -10
\end{align*}

\textbf{d)} \(\displaystyle \dfrac{t + 8}{7} = 2\)

\begin{align*}
t + 8 = 14 \\ t = 6
\end{align*}

\textbf{e)} \(\displaystyle \dfrac{x + 1}{8} = 1\)

\begin{align*}
x + 1 = 8 \\ x = 7
\end{align*}

\textbf{f)} \(\displaystyle \dfrac{p + 2}{2} = -4\)

\begin{align*}
p + 2 = -8 \\ p = -10
\end{align*}

\end{tcolorbox}

\end{exercise}

\begin{exercise}[]\protect\hypertarget{exr-u5-6}{}\label{exr-u5-6}

~

\textbf{a)} \(\displaystyle 4(x + 1) = 0\)

\textbf{b)} \(\displaystyle 3(x - 6) = -36\)

\textbf{c)} \(\displaystyle 5(t - 6) = 15\)

\textbf{d)} \(\displaystyle 3(y + 7) = 48\)

\textbf{e)} \(\displaystyle 4(y + 8) = 76\)

\textbf{f)} \(\displaystyle 3(m - 2) = -3\)

\begin{tcolorbox}[enhanced jigsaw, arc=.35mm, breakable, toptitle=1mm, title=\textcolor{quarto-callout-tip-color}{\faLightbulb}\hspace{0.5em}{🔍 View Solutions}, colframe=quarto-callout-tip-color-frame, rightrule=.15mm, opacityback=0, titlerule=0mm, colback=white, coltitle=black, bottomtitle=1mm, toprule=.15mm, leftrule=.75mm, bottomrule=.15mm, colbacktitle=quarto-callout-tip-color!10!white, left=2mm, opacitybacktitle=0.6]

\textbf{a)} \(\displaystyle 4(x + 1) = 0\)

\begin{align*}
x + 1 = 0 \\ x = -1
\end{align*}

\textbf{b)} \(\displaystyle 3(x - 6) = -36\)

\begin{align*}
x - 6 = -12 \\ x = -6
\end{align*}

\textbf{c)} \(\displaystyle 5(t - 6) = 15\)

\begin{align*}
t - 6 = 3 \\ t = 9
\end{align*}

\textbf{d)} \(\displaystyle 3(y + 7) = 48\)

\begin{align*}
y + 7 = 16 \\ y = 9
\end{align*}

\textbf{e)} \(\displaystyle 4(y + 8) = 76\)

\begin{align*}
y + 8 = 19 \\ y = 11
\end{align*}

\textbf{f)} \(\displaystyle 3(m - 2) = -3\)

\begin{align*}
m - 2 = -1 \\ m = 1
\end{align*}

\end{tcolorbox}

\end{exercise}

\begin{exercise}[]\protect\hypertarget{exr-u5-7}{}\label{exr-u5-7}

~

\textbf{a)} \(\displaystyle 7m + 6 = 8m\)

\textbf{b)} \(\displaystyle 5x - 8 = 4x - 6\)

\textbf{c)} \(\displaystyle 10x + 8 = 4x + 2\)

\textbf{d)} \(\displaystyle 3x - 3 = x - 15\)

\textbf{e)} \(\displaystyle 7x + 6 = 2x + 1\)

\textbf{f)} \(\displaystyle 9n + 7 = 4n - 13\)

\begin{tcolorbox}[enhanced jigsaw, arc=.35mm, breakable, toptitle=1mm, title=\textcolor{quarto-callout-tip-color}{\faLightbulb}\hspace{0.5em}{🔍 View Solutions}, colframe=quarto-callout-tip-color-frame, rightrule=.15mm, opacityback=0, titlerule=0mm, colback=white, coltitle=black, bottomtitle=1mm, toprule=.15mm, leftrule=.75mm, bottomrule=.15mm, colbacktitle=quarto-callout-tip-color!10!white, left=2mm, opacitybacktitle=0.6]

\textbf{a)} \(\displaystyle 7m + 6 = 8m\)

\begin{align*}
-m + 6 = 0 \\ -m = -6 \\ m = 6
\end{align*}

\textbf{b)} \(\displaystyle 5x - 8 = 4x - 6\)

\begin{align*}
x - 8 = -6 \\ x = 2
\end{align*}

\textbf{c)} \(\displaystyle 10x + 8 = 4x + 2\)

\begin{align*}
6x + 8 = 2 \\ 6x = -6 \\ x = -1
\end{align*}

\textbf{d)} \(\displaystyle 3x - 3 = x - 15\)

\begin{align*}
2x - 3 = -15 \\ 2x = -12 \\ x = -6
\end{align*}

\textbf{e)} \(\displaystyle 7x + 6 = 2x + 1\)

\begin{align*}
5x + 6 = 1 \\ 5x = -5 \\ x = -1
\end{align*}

\textbf{f)} \(\displaystyle 9n + 7 = 4n - 13\)

\begin{align*}
5n + 7 = -13 \\ 5n = -20 \\ n = -4
\end{align*}

\end{tcolorbox}

\end{exercise}

\begin{exercise}[]\protect\hypertarget{exr-u5-8}{}\label{exr-u5-8}

~

\textbf{a)} \(\displaystyle 5(t + 6) = 3(t + 12)\)

\textbf{b)} \(\displaystyle 8(m - 6) = 7(m - 7)\)

\textbf{c)} \(\displaystyle 3(y + 1) = 6(y - 2)\)

\textbf{d)} \(\displaystyle 4(y - 3) = 5(y - 1)\)

\begin{tcolorbox}[enhanced jigsaw, arc=.35mm, breakable, toptitle=1mm, title=\textcolor{quarto-callout-tip-color}{\faLightbulb}\hspace{0.5em}{🔍 View Solutions}, colframe=quarto-callout-tip-color-frame, rightrule=.15mm, opacityback=0, titlerule=0mm, colback=white, coltitle=black, bottomtitle=1mm, toprule=.15mm, leftrule=.75mm, bottomrule=.15mm, colbacktitle=quarto-callout-tip-color!10!white, left=2mm, opacitybacktitle=0.6]

\textbf{a)} \(\displaystyle 5(t + 6) = 3(t + 12)\)

\begin{align*}
2t + 30 = 36 \\ 2t = 6 \\ t = 3
\end{align*}

\textbf{b)} \(\displaystyle 8(m - 6) = 7(m - 7)\)

\begin{align*}
m - 48 = -49 \\ m = -1
\end{align*}

\textbf{c)} \(\displaystyle 3(y + 1) = 6(y - 2)\)

\begin{align*}
-3y + 3 = -12 \\ -3y = -15 \\ y = 5
\end{align*}

\textbf{d)} \(\displaystyle 4(y - 3) = 5(y - 1)\)

\begin{align*}
-y - 12 = -5 \\ -y = 7 \\ y = -7
\end{align*}

\end{tcolorbox}

\end{exercise}

\begin{exercise}[]\protect\hypertarget{exr-u5-9}{}\label{exr-u5-9}

~

\textbf{a)} \(\displaystyle 6(2t - 6) = -24\)

\textbf{b)} \(\displaystyle 5(5y + 7) = 35\)

\textbf{c)} \(\displaystyle 2(3m + 4) = -28\)

\textbf{d)} \(\displaystyle 4(5m + 1) = 144\)

\textbf{e)} \(\displaystyle 6(5t - 4) = -24\)

\textbf{f)} \(\displaystyle 3(2n - 7) = 3\)

\begin{tcolorbox}[enhanced jigsaw, arc=.35mm, breakable, toptitle=1mm, title=\textcolor{quarto-callout-tip-color}{\faLightbulb}\hspace{0.5em}{🔍 View Solutions}, colframe=quarto-callout-tip-color-frame, rightrule=.15mm, opacityback=0, titlerule=0mm, colback=white, coltitle=black, bottomtitle=1mm, toprule=.15mm, leftrule=.75mm, bottomrule=.15mm, colbacktitle=quarto-callout-tip-color!10!white, left=2mm, opacitybacktitle=0.6]

\textbf{a)} \(\displaystyle 6(2t - 6) = -24\)

\begin{align*}
12t - 36 = -24 \\ 12t = 12 \\ t = 1
\end{align*}

\textbf{b)} \(\displaystyle 5(5y + 7) = 35\)

\begin{align*}
25y + 35 = 35 \\ 25y = 0 \\ y = 0
\end{align*}

\textbf{c)} \(\displaystyle 2(3m + 4) = -28\)

\begin{align*}
6m + 8 = -28 \\ 6m = -36 \\ m = -6
\end{align*}

\textbf{d)} \(\displaystyle 4(5m + 1) = 144\)

\begin{align*}
20m + 4 = 144 \\ 20m = 140 \\ m = 7
\end{align*}

\textbf{e)} \(\displaystyle 6(5t - 4) = -24\)

\begin{align*}
30t - 24 = -24 \\ 30t = 0 \\ t = 0
\end{align*}

\textbf{f)} \(\displaystyle 3(2n - 7) = 3\)

\begin{align*}
6n - 21 = 3 \\ 6n = 24 \\ n = 4
\end{align*}

\end{tcolorbox}

\end{exercise}

\begin{exercise}[]\protect\hypertarget{exr-u5-13}{}\label{exr-u5-13}

~

\textbf{a)} \(\displaystyle \dfrac{5p + 7}{2} = 16\)

\textbf{b)} \(\displaystyle \dfrac{4x + 4}{4} = 1\)

\textbf{c)} \(\displaystyle \dfrac{4m + 2}{5} = 2\)

\textbf{d)} \(\displaystyle \dfrac{4m - 6}{5} = -6\)

\textbf{e)} \(\displaystyle \dfrac{6x - 8}{4} = 7\)

\textbf{f)} \(\displaystyle \dfrac{2p + 8}{3} = 0\)

\begin{tcolorbox}[enhanced jigsaw, arc=.35mm, breakable, toptitle=1mm, title=\textcolor{quarto-callout-tip-color}{\faLightbulb}\hspace{0.5em}{🔍 View Solutions}, colframe=quarto-callout-tip-color-frame, rightrule=.15mm, opacityback=0, titlerule=0mm, colback=white, coltitle=black, bottomtitle=1mm, toprule=.15mm, leftrule=.75mm, bottomrule=.15mm, colbacktitle=quarto-callout-tip-color!10!white, left=2mm, opacitybacktitle=0.6]

\textbf{a)} \(\displaystyle \dfrac{5p + 7}{2} = 16\)

\begin{align*}
5p + 7 = 32 \\ 5p = 25 \\ p = 5
\end{align*}

\textbf{b)} \(\displaystyle \dfrac{4x + 4}{4} = 1\)

\begin{align*}
4x + 4 = 4 \\ 4x = 0 \\ x = 0
\end{align*}

\textbf{c)} \(\displaystyle \dfrac{4m + 2}{5} = 2\)

\begin{align*}
4m + 2 = 10 \\ 4m = 8 \\ m = 2
\end{align*}

\textbf{d)} \(\displaystyle \dfrac{4m - 6}{5} = -6\)

\begin{align*}
4m - 6 = -30 \\ 4m = -24 \\ m = -6
\end{align*}

\textbf{e)} \(\displaystyle \dfrac{6x - 8}{4} = 7\)

\begin{align*}
6x - 8 = 28 \\ 6x = 36 \\ x = 6
\end{align*}

\textbf{f)} \(\displaystyle \dfrac{2p + 8}{3} = 0\)

\begin{align*}
2p + 8 = 0 \\ 2p = -8 \\ p = -4
\end{align*}

\end{tcolorbox}

\end{exercise}

\begin{exercise}[]\protect\hypertarget{exr-u5-17-1}{}\label{exr-u5-17-1}

~

Sara buys tickets. Adult tickets cost \$7 each and child tickets cost
\$4 each. The total cost is \$29. Sara bought 3 adult tickets.

\textbf{(i)} Write a formula for the total cost \(T\).

\textbf{(ii)} How many child tickets were bought?

\begin{tcolorbox}[enhanced jigsaw, arc=.35mm, breakable, toptitle=1mm, title=\textcolor{quarto-callout-tip-color}{\faLightbulb}\hspace{0.5em}{🔍 View Solution}, colframe=quarto-callout-tip-color-frame, rightrule=.15mm, opacityback=0, titlerule=0mm, colback=white, coltitle=black, bottomtitle=1mm, toprule=.15mm, leftrule=.75mm, bottomrule=.15mm, colbacktitle=quarto-callout-tip-color!10!white, left=2mm, opacitybacktitle=0.6]

\begin{align*}
\text{(i) } T = 7a + 4c \\ \text{(ii) } 21 + 4c = 29 \\ 4c = 8 \\ c = 2
\end{align*}

\textbf{Answer:} \(T = 7a + 4c,\; c = 2\)

\end{tcolorbox}

\end{exercise}

\begin{exercise}[]\protect\hypertarget{exr-u5-17-2}{}\label{exr-u5-17-2}

~

The sum of four consecutive integers is 70.

\textbf{(i)} Let the first integer be \(x\). Write an equation.

\textbf{(ii)} Find the numbers.

\begin{tcolorbox}[enhanced jigsaw, arc=.35mm, breakable, toptitle=1mm, title=\textcolor{quarto-callout-tip-color}{\faLightbulb}\hspace{0.5em}{🔍 View Solution}, colframe=quarto-callout-tip-color-frame, rightrule=.15mm, opacityback=0, titlerule=0mm, colback=white, coltitle=black, bottomtitle=1mm, toprule=.15mm, leftrule=.75mm, bottomrule=.15mm, colbacktitle=quarto-callout-tip-color!10!white, left=2mm, opacitybacktitle=0.6]

\begin{align*}
\text{(i) } x + (x + 1) + (x + 2) + (x + 3) = 70 \\ \text{(ii) } 4x + 6 = 70 \\ x = 16 \\ \text{Numbers: } 16, 17, 18, 19
\end{align*}

\textbf{Answer:} \(x = 16,\; \text{numbers: } 16, 17, 18, 19\)

\end{tcolorbox}

\end{exercise}

\begin{exercise}[]\protect\hypertarget{exr-u5-17-3}{}\label{exr-u5-17-3}

~

Ibrahim is \(x\) years old. Yusuf is 6 years older than Ibrahim. The sum
of their ages is 54.

\textbf{(i)} Form an equation in \(x\).

\textbf{(ii)} Find the ages of Ibrahim and Yusuf.

\begin{tcolorbox}[enhanced jigsaw, arc=.35mm, breakable, toptitle=1mm, title=\textcolor{quarto-callout-tip-color}{\faLightbulb}\hspace{0.5em}{🔍 View Solution}, colframe=quarto-callout-tip-color-frame, rightrule=.15mm, opacityback=0, titlerule=0mm, colback=white, coltitle=black, bottomtitle=1mm, toprule=.15mm, leftrule=.75mm, bottomrule=.15mm, colbacktitle=quarto-callout-tip-color!10!white, left=2mm, opacitybacktitle=0.6]

\begin{align*}
\text{(i) } x + (x + 6) = 54 \\ \text{(ii) } 2x + 6 = 54 \\ 2x = 48 \\ x = 24 \\ \text{Ibrahim = 24, Yusuf = 30}
\end{align*}

\textbf{Answer:} \(x = 24,\; Ibrahim=24,\; Yusuf=30\)

\end{tcolorbox}

\end{exercise}

\begin{exercise}[]\protect\hypertarget{exr-u5-17-4}{}\label{exr-u5-17-4}

~

The angles of a triangle are \((x + 13)^\circ\), \((2x + 14)^\circ\) and
\((3x + 45)^\circ\).

\begin{itemize}
\tightlist
\item
  Angles in a triangle sum to \(180^\circ\).
\end{itemize}

\textbf{(i)} Form an equation in \(x\).

\textbf{(ii)} Find \(x\) and each angle.

\begin{tcolorbox}[enhanced jigsaw, arc=.35mm, breakable, toptitle=1mm, title=\textcolor{quarto-callout-tip-color}{\faLightbulb}\hspace{0.5em}{🔍 View Solution}, colframe=quarto-callout-tip-color-frame, rightrule=.15mm, opacityback=0, titlerule=0mm, colback=white, coltitle=black, bottomtitle=1mm, toprule=.15mm, leftrule=.75mm, bottomrule=.15mm, colbacktitle=quarto-callout-tip-color!10!white, left=2mm, opacitybacktitle=0.6]

\begin{align*}
\text{(i) } 6x + 72 = 180 \\ \text{(ii) } x = 18 \\ \text{Angles: }31^\circ,\; 50^\circ,\; 99^\circ
\end{align*}

\textbf{Answer:} \(x = 18,\; 31^\circ,\; 50^\circ,\; 99^\circ\)

\end{tcolorbox}

\end{exercise}

\begin{exercise}[]\protect\hypertarget{exr-u5-18-1}{}\label{exr-u5-18-1}

~

The lengths, in cm, of the sides of the triangle are \(3(x + 2)\),
\(2x - 2\) and \(2x + 7\).

\textbf{(i)} Write an expression for the perimeter of the triangle.

\textbf{(ii)} The perimeter is 32 cm. Find \(x\).

\begin{tcolorbox}[enhanced jigsaw, arc=.35mm, breakable, toptitle=1mm, title=\textcolor{quarto-callout-tip-color}{\faLightbulb}\hspace{0.5em}{🔍 View Solution}, colframe=quarto-callout-tip-color-frame, rightrule=.15mm, opacityback=0, titlerule=0mm, colback=white, coltitle=black, bottomtitle=1mm, toprule=.15mm, leftrule=.75mm, bottomrule=.15mm, colbacktitle=quarto-callout-tip-color!10!white, left=2mm, opacitybacktitle=0.6]

\begin{align*}
\text{(i) } (3(x + 2)) + (2x - 2) + (2x + 7) = 7x + 11 \text{ cm} \\ \text{(ii) } 7x + 11 = 32 \\ 7x = 21 \\ x = 3
\end{align*}

\textbf{Answer:} \((i)\; 7x + 11 \text{ cm},\quad (ii)\; x = 3\)

\end{tcolorbox}

\end{exercise}

\begin{exercise}[]\protect\hypertarget{exr-u5-18-2}{}\label{exr-u5-18-2}

~

4(2x + 3)4x - 34x + 1Diagram NOT accurately drawn

The lengths, in cm, of the sides of the triangle are \(4(2x + 3)\),
\(4x - 3\) and \(4x + 1\).

\textbf{(i)} Write an expression for the perimeter of the triangle.

\textbf{(ii)} The perimeter is 58 cm. Find \(x\).

\begin{tcolorbox}[enhanced jigsaw, arc=.35mm, breakable, toptitle=1mm, title=\textcolor{quarto-callout-tip-color}{\faLightbulb}\hspace{0.5em}{🔍 View Solution}, colframe=quarto-callout-tip-color-frame, rightrule=.15mm, opacityback=0, titlerule=0mm, colback=white, coltitle=black, bottomtitle=1mm, toprule=.15mm, leftrule=.75mm, bottomrule=.15mm, colbacktitle=quarto-callout-tip-color!10!white, left=2mm, opacitybacktitle=0.6]

\begin{align*}
\text{(i) } (4(2x + 3)) + (4x - 3) + (4x + 1) = 16x + 10 \text{ cm} \\ \text{(ii) } 16x + 10 = 58 \\ 16x = 48 \\ x = 3
\end{align*}

\textbf{Answer:} \((i)\; 16x + 10 \text{ cm},\quad (ii)\; x = 3\)

\end{tcolorbox}

\end{exercise}

\begin{exercise}[]\protect\hypertarget{exr-u5-19-1}{}\label{exr-u5-19-1}

~

(2x + 16)°(3x + 13)°(3x + 47)°Diagram NOT accurately drawn

The angles of a triangle are \((2x + 16)^\circ\), \((3x + 13)^\circ\)
and \((3x + 47)^\circ\).

\begin{itemize}
\tightlist
\item
  Angles in a triangle sum to \(180^\circ\).
\end{itemize}

\textbf{(i)} Form an equation in \(x\).

\textbf{(ii)} Find \(x\) and each angle.

\begin{tcolorbox}[enhanced jigsaw, arc=.35mm, breakable, toptitle=1mm, title=\textcolor{quarto-callout-tip-color}{\faLightbulb}\hspace{0.5em}{🔍 View Solution}, colframe=quarto-callout-tip-color-frame, rightrule=.15mm, opacityback=0, titlerule=0mm, colback=white, coltitle=black, bottomtitle=1mm, toprule=.15mm, leftrule=.75mm, bottomrule=.15mm, colbacktitle=quarto-callout-tip-color!10!white, left=2mm, opacitybacktitle=0.6]

\begin{align*}
\text{(i) } 8x + 76 = 180 \\ \text{(ii) } x = 13 \\ \text{Angles: }42^\circ,\; 52^\circ,\; 86^\circ
\end{align*}

\textbf{Answer:} \(x=13,\; 42^\circ,\; 52^\circ,\; 86^\circ\)

\end{tcolorbox}

\end{exercise}

\begin{exercise}[]\protect\hypertarget{exr-u5-19-2}{}\label{exr-u5-19-2}

~

The angles of a triangle are \((3x + 22)^\circ\), \((2x + 24)^\circ\)
and \((x + 74)^\circ\).

\begin{itemize}
\tightlist
\item
  Angles in a triangle sum to \(180^\circ\).
\end{itemize}

\textbf{(i)} Form an equation in \(x\).

\textbf{(ii)} Find \(x\) and each angle.

\begin{tcolorbox}[enhanced jigsaw, arc=.35mm, breakable, toptitle=1mm, title=\textcolor{quarto-callout-tip-color}{\faLightbulb}\hspace{0.5em}{🔍 View Solution}, colframe=quarto-callout-tip-color-frame, rightrule=.15mm, opacityback=0, titlerule=0mm, colback=white, coltitle=black, bottomtitle=1mm, toprule=.15mm, leftrule=.75mm, bottomrule=.15mm, colbacktitle=quarto-callout-tip-color!10!white, left=2mm, opacitybacktitle=0.6]

\begin{align*}
\text{(i) } 6x + 120 = 180 \\ \text{(ii) } x = 10 \\ \text{Angles: }52^\circ,\; 44^\circ,\; 84^\circ
\end{align*}

\textbf{Answer:} \(x=10,\; 52^\circ,\; 44^\circ,\; 84^\circ\)

\end{tcolorbox}

\end{exercise}

\begin{exercise}[]\protect\hypertarget{exr-u5-21-1}{}\label{exr-u5-21-1}

~

4x + 6 cmx + 4 cmDiagram NOT accurately drawn

A rectangle has length \((4x + 6)\) cm and width \((x + 4)\) cm.

\textbf{(i)} Write an expression for the perimeter.

\textbf{(ii)} The perimeter is 90 cm. Find \(x\).

\begin{tcolorbox}[enhanced jigsaw, arc=.35mm, breakable, toptitle=1mm, title=\textcolor{quarto-callout-tip-color}{\faLightbulb}\hspace{0.5em}{🔍 View Solution}, colframe=quarto-callout-tip-color-frame, rightrule=.15mm, opacityback=0, titlerule=0mm, colback=white, coltitle=black, bottomtitle=1mm, toprule=.15mm, leftrule=.75mm, bottomrule=.15mm, colbacktitle=quarto-callout-tip-color!10!white, left=2mm, opacitybacktitle=0.6]

\begin{align*}
\text{(i) } 2(4x + 6 + x + 4) = 10x + 20 \text{ cm} \\ \text{(ii) } 10x + 20 = 90 \\ 10x = 70 \\ x = 7
\end{align*}

\textbf{Answer:} \((i)\; 10x + 20 \text{ cm},\quad (ii)\; x = 7\)

\end{tcolorbox}

\end{exercise}

\begin{exercise}[]\protect\hypertarget{exr-u5-21-2}{}\label{exr-u5-21-2}

~

A rectangle has length \((2x + 8)\) cm and width \((x + 6)\) cm.

\textbf{(i)} Write an expression for the perimeter.

\textbf{(ii)} The perimeter is 64 cm. Find \(x\).

\begin{tcolorbox}[enhanced jigsaw, arc=.35mm, breakable, toptitle=1mm, title=\textcolor{quarto-callout-tip-color}{\faLightbulb}\hspace{0.5em}{🔍 View Solution}, colframe=quarto-callout-tip-color-frame, rightrule=.15mm, opacityback=0, titlerule=0mm, colback=white, coltitle=black, bottomtitle=1mm, toprule=.15mm, leftrule=.75mm, bottomrule=.15mm, colbacktitle=quarto-callout-tip-color!10!white, left=2mm, opacitybacktitle=0.6]

\begin{align*}
\text{(i) } 2(2x + 8 + x + 6) = 6x + 28 \text{ cm} \\ \text{(ii) } 6x + 28 = 64 \\ 6x = 36 \\ x = 6
\end{align*}

\textbf{Answer:} \((i)\; 6x + 28 \text{ cm},\quad (ii)\; x = 6\)

\end{tcolorbox}

\end{exercise}

\begin{exercise}[]\protect\hypertarget{exr-u5-22-1}{}\label{exr-u5-22-1}

~

2x - 4x + 2x + 5Diagram NOT accurately drawn

An isosceles triangle has equal sides \((2x - 4)\) cm and \((x + 2)\)
cm. The base is \((x + 5)\) cm.

\textbf{(i)} Find \(x\).

\textbf{(ii)} Find the perimeter of the triangle.

\begin{tcolorbox}[enhanced jigsaw, arc=.35mm, breakable, toptitle=1mm, title=\textcolor{quarto-callout-tip-color}{\faLightbulb}\hspace{0.5em}{🔍 View Solution}, colframe=quarto-callout-tip-color-frame, rightrule=.15mm, opacityback=0, titlerule=0mm, colback=white, coltitle=black, bottomtitle=1mm, toprule=.15mm, leftrule=.75mm, bottomrule=.15mm, colbacktitle=quarto-callout-tip-color!10!white, left=2mm, opacitybacktitle=0.6]

\begin{align*}
\text{(i) } 2x - 4 = x + 2 \\ x = 6 \\ \text{(ii) Perimeter} = 27 \text{ cm}
\end{align*}

\textbf{Answer:} \(x = 6,\; P = 27 \text{ cm}\)

\end{tcolorbox}

\end{exercise}

\begin{exercise}[]\protect\hypertarget{exr-u5-22-2}{}\label{exr-u5-22-2}

~

An isosceles triangle has equal sides \((3x + 5)\) cm and \((x + 17)\)
cm. The base is \((2x + 2)\) cm.

\textbf{(i)} Find \(x\).

\textbf{(ii)} Find the perimeter of the triangle.

\begin{tcolorbox}[enhanced jigsaw, arc=.35mm, breakable, toptitle=1mm, title=\textcolor{quarto-callout-tip-color}{\faLightbulb}\hspace{0.5em}{🔍 View Solution}, colframe=quarto-callout-tip-color-frame, rightrule=.15mm, opacityback=0, titlerule=0mm, colback=white, coltitle=black, bottomtitle=1mm, toprule=.15mm, leftrule=.75mm, bottomrule=.15mm, colbacktitle=quarto-callout-tip-color!10!white, left=2mm, opacitybacktitle=0.6]

\begin{align*}
\text{(i) } 3x + 5 = x + 17 \\ 2x = 12 \\ x = 6 \\ \text{(ii) Perimeter} = 60 \text{ cm}
\end{align*}

\textbf{Answer:} \(x = 6,\; P = 60 \text{ cm}\)

\end{tcolorbox}

\end{exercise}

∑ × π ÷ √ ◆ ∞ ± ∫ ≠ Δ

\subsection{🔴 Challenging}

\begin{exercise}[]\protect\hypertarget{exr-u5-10}{}\label{exr-u5-10}

~

\textbf{a)} \(\displaystyle 2(4x + 2) + 3(x + 2) = -34\)

\textbf{b)} \(\displaystyle 2(y + 4) + 3(4y - 5) = 49\)

\textbf{c)} \(\displaystyle 2(y + 4) + 4(3y - 3) = 66\)

\textbf{d)} \(\displaystyle 4(2t - 2) + 5(2t + 4) = 102\)

\begin{tcolorbox}[enhanced jigsaw, arc=.35mm, breakable, toptitle=1mm, title=\textcolor{quarto-callout-tip-color}{\faLightbulb}\hspace{0.5em}{🔍 View Solutions}, colframe=quarto-callout-tip-color-frame, rightrule=.15mm, opacityback=0, titlerule=0mm, colback=white, coltitle=black, bottomtitle=1mm, toprule=.15mm, leftrule=.75mm, bottomrule=.15mm, colbacktitle=quarto-callout-tip-color!10!white, left=2mm, opacitybacktitle=0.6]

\textbf{a)} \(\displaystyle 2(4x + 2) + 3(x + 2) = -34\)

\begin{align*}
8x + 4 + 3x + 6 = -34 \\ 11x + 10 = -34 \\ 11x = -44 \\ x = -4
\end{align*}

\textbf{b)} \(\displaystyle 2(y + 4) + 3(4y - 5) = 49\)

\begin{align*}
2y + 8 + 12y - 15 = 49 \\ 14y - 7 = 49 \\ 14y = 56 \\ y = 4
\end{align*}

\textbf{c)} \(\displaystyle 2(y + 4) + 4(3y - 3) = 66\)

\begin{align*}
2y + 8 + 12y - 12 = 66 \\ 14y - 4 = 66 \\ 14y = 70 \\ y = 5
\end{align*}

\textbf{d)} \(\displaystyle 4(2t - 2) + 5(2t + 4) = 102\)

\begin{align*}
8t - 8 + 10t + 20 = 102 \\ 18t + 12 = 102 \\ 18t = 90 \\ t = 5
\end{align*}

\end{tcolorbox}

\end{exercise}

\begin{exercise}[]\protect\hypertarget{exr-u5-11}{}\label{exr-u5-11}

~

\textbf{a)} \(\displaystyle (n + 3)^{2} = (n + 4)(n + 1)\)

\textbf{b)} \(\displaystyle (3y + 4)^{2} = (3y + 5)(3y + 2)\)

\textbf{c)} \(\displaystyle (2y + 4)^{2} = (2y + 6)(2y + 3)\)

\textbf{d)} \(\displaystyle (2m + 2)^{2} = (2m + 1)(2m + 4)\)

\begin{tcolorbox}[enhanced jigsaw, arc=.35mm, breakable, toptitle=1mm, title=\textcolor{quarto-callout-tip-color}{\faLightbulb}\hspace{0.5em}{🔍 View Solutions}, colframe=quarto-callout-tip-color-frame, rightrule=.15mm, opacityback=0, titlerule=0mm, colback=white, coltitle=black, bottomtitle=1mm, toprule=.15mm, leftrule=.75mm, bottomrule=.15mm, colbacktitle=quarto-callout-tip-color!10!white, left=2mm, opacitybacktitle=0.6]

\textbf{a)} \(\displaystyle (n + 3)^{2} = (n + 4)(n + 1)\)

\begin{align*}
\text{Expand both sides, } 1n^2 \text{ cancels:} \\ 6n + 9 = 5n + 4 \\ 1n = -5 \\ n = -5
\end{align*}

\textbf{b)} \(\displaystyle (3y + 4)^{2} = (3y + 5)(3y + 2)\)

\begin{align*}
\text{Expand both sides, } 9y^2 \text{ cancels:} \\ 24y + 16 = 21y + 10 \\ 3y = -6 \\ y = -2
\end{align*}

\textbf{c)} \(\displaystyle (2y + 4)^{2} = (2y + 6)(2y + 3)\)

\begin{align*}
\text{Expand both sides, } 4y^2 \text{ cancels:} \\ 16y + 16 = 18y + 18 \\ -2y = 2 \\ y = -1
\end{align*}

\textbf{d)} \(\displaystyle (2m + 2)^{2} = (2m + 1)(2m + 4)\)

\begin{align*}
\text{Expand both sides, } 4m^2 \text{ cancels:} \\ 8m + 4 = 10m + 4 \\ -2m = 0 \\ m = 0
\end{align*}

\end{tcolorbox}

\end{exercise}

\begin{exercise}[]\protect\hypertarget{exr-u5-12}{}\label{exr-u5-12}

~

\textbf{a)} \(\displaystyle 5(3y + 2) + 4(3y - 1) = -2(3y + 5)\)

\textbf{b)} \(\displaystyle 5(3p + 1) + 2(3p + 1) = -2(2p + 6)\)

\textbf{c)} \(\displaystyle 4(3y + 1) + 4(y - 2) = -2(y + 5)\)

\textbf{d)} \(\displaystyle 3(2n - 1) + 3(3n - 5) = -1(n - 3)\)

\begin{tcolorbox}[enhanced jigsaw, arc=.35mm, breakable, toptitle=1mm, title=\textcolor{quarto-callout-tip-color}{\faLightbulb}\hspace{0.5em}{🔍 View Solutions}, colframe=quarto-callout-tip-color-frame, rightrule=.15mm, opacityback=0, titlerule=0mm, colback=white, coltitle=black, bottomtitle=1mm, toprule=.15mm, leftrule=.75mm, bottomrule=.15mm, colbacktitle=quarto-callout-tip-color!10!white, left=2mm, opacitybacktitle=0.6]

\textbf{a)} \(\displaystyle 5(3y + 2) + 4(3y - 1) = -2(3y + 5)\)

\begin{align*}
15y + 10 + 12y - 4 = -6y - 10 \\ 27y + 6 = -6y - 10 \\ 33y = -16 \\ y = -\dfrac{16}{33}
\end{align*}

\textbf{b)} \(\displaystyle 5(3p + 1) + 2(3p + 1) = -2(2p + 6)\)

\begin{align*}
15p + 5 + 6p + 2 = -4p - 12 \\ 21p + 7 = -4p - 12 \\ 25p = -19 \\ p = -\dfrac{19}{25}
\end{align*}

\textbf{c)} \(\displaystyle 4(3y + 1) + 4(y - 2) = -2(y + 5)\)

\begin{align*}
12y + 4 + 4y - 8 = -2y - 10 \\ 16y - 4 = -2y - 10 \\ 18y = -6 \\ y = -\dfrac{1}{3}
\end{align*}

\textbf{d)} \(\displaystyle 3(2n - 1) + 3(3n - 5) = -1(n - 3)\)

\begin{align*}
6n - 3 + 9n - 15 = -n + 3 \\ 15n - 18 = -n + 3 \\ 16n = 21 \\ n = \dfrac{21}{16}
\end{align*}

\end{tcolorbox}

\end{exercise}

\begin{exercise}[]\protect\hypertarget{exr-u5-14}{}\label{exr-u5-14}

~

\textbf{a)} \(\displaystyle \dfrac{5y + 2}{3} = \dfrac{2y + 26}{4}\)

\textbf{b)} \(\displaystyle \dfrac{4p + 2}{3} = \dfrac{3p + 1}{3}\)

\textbf{c)} \(\displaystyle \dfrac{4n - 2}{2} = \dfrac{n - 4}{4}\)

\textbf{d)} \(\displaystyle \dfrac{3x + 2}{2} = \dfrac{2x}{4}\)

\begin{tcolorbox}[enhanced jigsaw, arc=.35mm, breakable, toptitle=1mm, title=\textcolor{quarto-callout-tip-color}{\faLightbulb}\hspace{0.5em}{🔍 View Solutions}, colframe=quarto-callout-tip-color-frame, rightrule=.15mm, opacityback=0, titlerule=0mm, colback=white, coltitle=black, bottomtitle=1mm, toprule=.15mm, leftrule=.75mm, bottomrule=.15mm, colbacktitle=quarto-callout-tip-color!10!white, left=2mm, opacitybacktitle=0.6]

\textbf{a)} \(\displaystyle \dfrac{5y + 2}{3} = \dfrac{2y + 26}{4}\)

\begin{align*}
\text{Cross multiply: } 4(5y + 2) = 3(2y + 26) \\ 20y + 8 = 6y + 78 \\ 14y = 70 \\ y = 5
\end{align*}

\textbf{b)} \(\displaystyle \dfrac{4p + 2}{3} = \dfrac{3p + 1}{3}\)

\begin{align*}
\text{Cross multiply: } 3(4p + 2) = 3(3p + 1) \\ 12p + 6 = 9p + 3 \\ 3p = -3 \\ p = -1
\end{align*}

\textbf{c)} \(\displaystyle \dfrac{4n - 2}{2} = \dfrac{n - 4}{4}\)

\begin{align*}
\text{Cross multiply: } 4(4n - 2) = 2(n - 4) \\ 16n - 8 = 2n - 8 \\ 14n = 0 \\ n = 0
\end{align*}

\textbf{d)} \(\displaystyle \dfrac{3x + 2}{2} = \dfrac{2x}{4}\)

\begin{align*}
\text{Cross multiply: } 4(3x + 2) = 2(2x) \\ 12x + 8 = 4x \\ 8x = -8 \\ x = -1
\end{align*}

\end{tcolorbox}

\end{exercise}

\begin{exercise}[]\protect\hypertarget{exr-u5-15}{}\label{exr-u5-15}

~

\textbf{a)}
\(\displaystyle \dfrac{2p - 1}{2} + \dfrac{2p + 5}{2} = -10\)

\textbf{b)} \(\displaystyle \dfrac{2n + 5}{2} + \dfrac{n - 1}{2} = 17\)

\textbf{c)} \(\displaystyle \dfrac{m - 1}{2} + \dfrac{3m}{5} = -6\)

\textbf{d)} \(\displaystyle \dfrac{p - 1}{5} + \dfrac{2p + 3}{5} = -2\)

\begin{tcolorbox}[enhanced jigsaw, arc=.35mm, breakable, toptitle=1mm, title=\textcolor{quarto-callout-tip-color}{\faLightbulb}\hspace{0.5em}{🔍 View Solutions}, colframe=quarto-callout-tip-color-frame, rightrule=.15mm, opacityback=0, titlerule=0mm, colback=white, coltitle=black, bottomtitle=1mm, toprule=.15mm, leftrule=.75mm, bottomrule=.15mm, colbacktitle=quarto-callout-tip-color!10!white, left=2mm, opacitybacktitle=0.6]

\textbf{a)}
\(\displaystyle \dfrac{2p - 1}{2} + \dfrac{2p + 5}{2} = -10\)

\begin{align*}
\text{Multiply by }2: 2p - 1 + 2p + 5 = -20 \\ 4p + 4 = -20 \\ 4p = -24 \\ p = -6
\end{align*}

\textbf{b)} \(\displaystyle \dfrac{2n + 5}{2} + \dfrac{n - 1}{2} = 17\)

\begin{align*}
\text{Multiply by }2: 2n + 5 + n - 1 = 34 \\ 3n + 4 = 34 \\ 3n = 30 \\ n = 10
\end{align*}

\textbf{c)} \(\displaystyle \dfrac{m - 1}{2} + \dfrac{3m}{5} = -6\)

\begin{align*}
\text{Multiply by }10: 5m - 5 + 6m = -60 \\ 11m - 5 = -60 \\ 11m = -55 \\ m = -5
\end{align*}

\textbf{d)} \(\displaystyle \dfrac{p - 1}{5} + \dfrac{2p + 3}{5} = -2\)

\begin{align*}
\text{Multiply by }5: p - 1 + 2p + 3 = -10 \\ 3p + 2 = -10 \\ 3p = -12 \\ p = -4
\end{align*}

\end{tcolorbox}

\end{exercise}

\begin{exercise}[]\protect\hypertarget{exr-u5-16}{}\label{exr-u5-16}

~

\textbf{a)}
\(\displaystyle \dfrac{t - 2}{3} - \dfrac{2t + 2}{5} = \dfrac{3t - 1}{3} + \dfrac{t - 19}{3}\)

\textbf{b)}
\(\displaystyle \dfrac{2p}{4} - \dfrac{2p - 4}{3} = \dfrac{3p}{2} + \dfrac{p - 32}{3}\)

\textbf{c)}
\(\displaystyle \dfrac{y - 2}{4} - \dfrac{3y - 3}{4} = \dfrac{2y - 4}{2} + \dfrac{y + 16}{4}\)

\begin{tcolorbox}[enhanced jigsaw, arc=.35mm, breakable, toptitle=1mm, title=\textcolor{quarto-callout-tip-color}{\faLightbulb}\hspace{0.5em}{🔍 View Solutions}, colframe=quarto-callout-tip-color-frame, rightrule=.15mm, opacityback=0, titlerule=0mm, colback=white, coltitle=black, bottomtitle=1mm, toprule=.15mm, leftrule=.75mm, bottomrule=.15mm, colbacktitle=quarto-callout-tip-color!10!white, left=2mm, opacitybacktitle=0.6]

\textbf{a)}
\(\displaystyle \dfrac{t - 2}{3} - \dfrac{2t + 2}{5} = \dfrac{3t - 1}{3} + \dfrac{t - 19}{3}\)

\begin{align*}
\text{Multiply by }15: 5t - 10 - 6t + 6 = 15t - 5 + 5t - 95 \\ -21t = -84 \\ t = 4
\end{align*}

\textbf{b)}
\(\displaystyle \dfrac{2p}{4} - \dfrac{2p - 4}{3} = \dfrac{3p}{2} + \dfrac{p - 32}{3}\)

\begin{align*}
\text{Multiply by }12: 6p - 8p - 16 = 18p + 4p - 128 \\ -24p = -144 \\ p = 6
\end{align*}

\textbf{c)}
\(\displaystyle \dfrac{y - 2}{4} - \dfrac{3y - 3}{4} = \dfrac{2y - 4}{2} + \dfrac{y + 16}{4}\)

\begin{align*}
\text{Multiply by }4: y - 2 - 3y - 3 = 4y - 8 + y + 16 \\ -7y = 7 \\ y = -1
\end{align*}

\end{tcolorbox}

\end{exercise}

\begin{exercise}[]\protect\hypertarget{exr-u5-20-1}{}\label{exr-u5-20-1}

~

ABCD(6x - 23)°(3x + 59)°(3x + 13)°

ABCD is a trapezium.

\begin{itemize}
\tightlist
\item
  BC is parallel to AD.
\item
  Angle A = \((6x - 23)°\), Angle B = \((3x + 59)°\), Angle D =
  \((3x + 13)°\).
\end{itemize}

\textbf{(i)} Find \(x\).

\textbf{(ii)} Find the size of the largest angle in the trapezium.

\begin{tcolorbox}[enhanced jigsaw, arc=.35mm, breakable, toptitle=1mm, title=\textcolor{quarto-callout-tip-color}{\faLightbulb}\hspace{0.5em}{🔍 View Solution}, colframe=quarto-callout-tip-color-frame, rightrule=.15mm, opacityback=0, titlerule=0mm, colback=white, coltitle=black, bottomtitle=1mm, toprule=.15mm, leftrule=.75mm, bottomrule=.15mm, colbacktitle=quarto-callout-tip-color!10!white, left=2mm, opacitybacktitle=0.6]

\begin{align*}
\text{(i) Co-interior: } (6x-23) + (3x+59) = 180 \\ 9x + 36 = 180 \\ x = 16 \\ \text{(ii) } \text{A=}73^\circ,\; \text{B=}107^\circ,\; \text{C=}119^\circ,\; \text{D=}61^\circ \\ \text{Largest = }119^\circ
\end{align*}

\textbf{Answer:} \(x=16,\; \text{largest} = 119^\circ\)

\end{tcolorbox}

\end{exercise}

\begin{exercise}[]\protect\hypertarget{exr-u5-20-2}{}\label{exr-u5-20-2}

~

ABCD(5x - 26)°(3x + 54)°(2x + 16)°

ABCD is a trapezium.

\begin{itemize}
\tightlist
\item
  BC is parallel to AD.
\item
  Angle A = \((5x - 26)°\), Angle B = \((3x + 54)°\), Angle D =
  \((2x + 16)°\).
\end{itemize}

\textbf{(i)} Find \(x\).

\textbf{(ii)} Find the size of the largest angle in the trapezium.

\begin{tcolorbox}[enhanced jigsaw, arc=.35mm, breakable, toptitle=1mm, title=\textcolor{quarto-callout-tip-color}{\faLightbulb}\hspace{0.5em}{🔍 View Solution}, colframe=quarto-callout-tip-color-frame, rightrule=.15mm, opacityback=0, titlerule=0mm, colback=white, coltitle=black, bottomtitle=1mm, toprule=.15mm, leftrule=.75mm, bottomrule=.15mm, colbacktitle=quarto-callout-tip-color!10!white, left=2mm, opacitybacktitle=0.6]

\begin{align*}
\text{(i) Co-interior: } (5x-26) + (3x+54) = 180 \\ 8x + 28 = 180 \\ x = 19 \\ \text{(ii) } \text{A=}69^\circ,\; \text{B=}111^\circ,\; \text{C=}126^\circ,\; \text{D=}54^\circ \\ \text{Largest = }126^\circ
\end{align*}

\textbf{Answer:} \(x=19,\; \text{largest} = 126^\circ\)

\end{tcolorbox}

\end{exercise}

∑ × π ÷ √ ◆ ∞ ± ∫ ≠ Δ

\section{⚡ Test Your Knowledge}\label{test-your-knowledge}

\chapter{Unit 6: Quadratic Equations}\label{unit-6-quadratic-equations}

\section{🎯 Learning Targets}\label{learning-targets-5}

\begin{longtable}[]{@{}
  >{\raggedright\arraybackslash}p{(\linewidth - 4\tabcolsep) * \real{0.3333}}
  >{\raggedright\arraybackslash}p{(\linewidth - 4\tabcolsep) * \real{0.3333}}
  >{\raggedright\arraybackslash}p{(\linewidth - 4\tabcolsep) * \real{0.3333}}@{}}
\toprule\noalign{}
\begin{minipage}[b]{\linewidth}\raggedright
\#
\end{minipage} & \begin{minipage}[b]{\linewidth}\raggedright
Learning Target
\end{minipage} & \begin{minipage}[b]{\linewidth}\raggedright
I can\ldots{}
\end{minipage} \\
\midrule\noalign{}
\endhead
\bottomrule\noalign{}
\endlastfoot
\textbf{LT1} & Factorising --- Simple & Factorise \(x^2 + bx + c\) using
the sum-product method \\
\textbf{LT2} & Factorising --- Harder & Factorise \(ax^2 + bx + c\)
(\(a \neq 1\)) using the \(ac\)-method \\
\textbf{LT3} & Special Forms & Recognise and factorise difference of two
squares and perfect squares \\
\textbf{LT4} & Solve by Factorising & Solve quadratic equations by
factorising \\
\textbf{LT5} & Quadratic Formula \& Discriminant & Apply
\(x = \dfrac{-b \pm \sqrt{b^2-4ac}}{2a}\); use \(\Delta = b^2-4ac\) to
determine nature of roots and find unknown parameters \\
\textbf{LT6} & Completing the Square & Write \(ax^2+bx+c\) in completed
square form and use it to solve \\
\textbf{LT7} & Word Problems & Form and solve quadratic equations from
real-world contexts (area, integers: consecutive, even, odd) \\
\textbf{LT8} & Geometry & Apply quadratic equations to geometric
problems (right triangles, rectangles, L-shapes, cuboids) \\
\textbf{LT9} & Roots \& Vieta's Formulas & Form equations from roots;
use \(\alpha+\beta = -b/a\) and \(\alpha\beta = c/a\) to find symmetric
expressions \\
\textbf{LT10} & Algebraic Fraction Equations & Solve equations with
algebraic fractions by multiplying through to form a quadratic \\
\textbf{LT11} & Disguised Quadratics & Recognise and solve
\(x^4 + bx^2 + c = 0\) and \((x+a)^2 + b(x+a) + c = 0\) using
substitution \\
\textbf{LT12} & Surd Equations & Solve equations involving square roots
by squaring both sides; identify and reject extraneous roots \\
\textbf{LT13} & Indices → Quadratic & Solve exponential equations like
\((b^{x+p})^{x+q} = b^k\) by equating exponents to form a quadratic \\
\end{longtable}

\section{🗂️ Formula Card}\label{formula-card-2}

\subsection*{Quadratic Equations --- Formula
Card}\label{quadratic-equations-formula-card}
\addcontentsline{toc}{subsection}{Quadratic Equations --- Formula Card}

\subsubsection*{Factorising (a = 1)}\label{factorising-a-1}
\addcontentsline{toc}{subsubsection}{Factorising (a = 1)}

\textbf{💡 Sum-Product Method} Find two numbers that \textbf{multiply}
to \(c\) and \textbf{add} to \(b\).

\textbf{Standard form}: \(x^2 + bx + c = (x + p)(x + q)\)

where \(p \times q = c\) and \(p + q = b\)

\textbf{Examples}:
\[x^2 + 5x + 6 = (x+2)(x+3) \quad [2\times3=6,\;2+3=5]\]
\[x^2 - x - 12 = (x+3)(x-4) \quad [3\times(-4)=-12,\;3+(-4)=-1]\]

\subsubsection*{Factorising (a ≠ 1) --- ac
Method}\label{factorising-a-1-ac-method}
\addcontentsline{toc}{subsubsection}{Factorising (a ≠ 1) --- ac Method}

\textbf{🎯 ac Method} Multiply \(a \times c\), then find two numbers
that multiply to \(ac\) and add to \(b\).

\textbf{Steps}: For \(ax^2 + bx + c\):

\begin{enumerate}
\def\labelenumi{\arabic{enumi}.}
\tightlist
\item
  Find \(p \times q = ac\) and \(p + q = b\)
\item
  Split the middle term: \(ax^2 + px + qx + c\)
\item
  Group and factorise
\end{enumerate}

\textbf{Example}: \(2x^2 + 7x + 3\), \(\;ac = 6\), find
\(p,q: p\times q=6,\; p+q=7 \Rightarrow 1,6\)
\[2x^2 + x + 6x + 3 = x(2x+1) + 3(2x+1) = (x+3)(2x+1)\]

\subsubsection*{Special Forms}\label{special-forms}
\addcontentsline{toc}{subsubsection}{Special Forms}

\textbf{⚠️ Spot These First!} Always check for special forms before
using other methods.

\textbf{Difference of Two Squares}: \[a^2 - b^2 = (a+b)(a-b)\]

\textbf{Perfect Square}:
\[a^2 + 2ab + b^2 = (a+b)^2 \qquad a^2 - 2ab + b^2 = (a-b)^2\]

\textbf{Examples}: \[x^2 - 25 = (x+5)(x-5)\] \[x^2 + 8x + 16 = (x+4)^2\]

\subsubsection*{Quadratic Formula}\label{quadratic-formula}
\addcontentsline{toc}{subsubsection}{Quadratic Formula}

\textbf{📝 Always works!} Use when factorising is not possible.

For \(ax^2 + bx + c = 0\): \[x = \dfrac{-b \pm \sqrt{b^2 - 4ac}}{2a}\]

\textbf{The Discriminant} \(\Delta = b^2 - 4ac\):

\begin{longtable}[]{@{}ll@{}}
\toprule\noalign{}
\(\Delta > 0\) & Two distinct real roots \\
\midrule\noalign{}
\endhead
\bottomrule\noalign{}
\endlastfoot
\(\Delta = 0\) & One repeated root \\
\(\Delta < 0\) & No real roots \\
\end{longtable}

\textbf{Example}: \(x^2 - 5x + 3 = 0\)
\[x = \dfrac{5 \pm \sqrt{25-12}}{2} = \dfrac{5 \pm \sqrt{13}}{2}\]

\subsubsection*{Completing the Square}\label{completing-the-square}
\addcontentsline{toc}{subsubsection}{Completing the Square}

\textbf{Monic} (\(a = 1\)):
\(x^2 + bx + c = \left(x + \dfrac{b}{2}\right)^2 - \left(\dfrac{b}{2}\right)^2 + c\)

\textbf{Non-monic} (\(a \neq 1\)):
\(ax^2 + bx + c = a\left(x + \dfrac{b}{2a}\right)^2 - \dfrac{b^2}{4a} + c\)

\textbf{To solve}: Write in form \((x+p)^2 = q\), then
\(x = -p \pm \sqrt{q}\)

\textbf{Example}: \(x^2 + 6x + 2 = 0\)
\[( x+3)^2 - 9 + 2 = 0 \implies (x+3)^2 = 7 \implies x = -3 \pm \sqrt{7}\]

\section{📝 Practice Problems}\label{practice-problems-3}

\subsection{🟢 Easy}

\begin{exercise}[]\protect\hypertarget{exr-u6-1}{}\label{exr-u6-1}

\emph{Factorise each expression.}

\textbf{a)} \(\displaystyle n^2 + 3n + 2\)

\textbf{b)} \(\displaystyle y^2 + 8y + 16\)

\textbf{c)} \(\displaystyle x^2 + 11x + 18\)

\textbf{d)} \(\displaystyle x^2 + 17x + 70\)

\textbf{e)} \(\displaystyle y^2 + 3y + 2\)

\textbf{f)} \(\displaystyle x^2 + 13x + 36\)

\textbf{g)} \(\displaystyle m^2 + 13m + 36\)

\textbf{h)} \(\displaystyle n^2 + 12n + 32\)

\begin{tcolorbox}[enhanced jigsaw, arc=.35mm, breakable, toptitle=1mm, title=\textcolor{quarto-callout-tip-color}{\faLightbulb}\hspace{0.5em}{🔍 View Solutions}, colframe=quarto-callout-tip-color-frame, rightrule=.15mm, opacityback=0, titlerule=0mm, colback=white, coltitle=black, bottomtitle=1mm, toprule=.15mm, leftrule=.75mm, bottomrule=.15mm, colbacktitle=quarto-callout-tip-color!10!white, left=2mm, opacitybacktitle=0.6]

\textbf{a)} \(\displaystyle n^2 + 3n + 2\)

\begin{align*}
\text{Find } p, q: p \times q = 2,\; p + q = 3 \\ p = 2,\; q = 1 \\ n^2 + 3n + 2 = (n + 2)(n + 1)
\end{align*}

\textbf{b)} \(\displaystyle y^2 + 8y + 16\)

\begin{align*}
\text{Find } p, q: p \times q = 16,\; p + q = 8 \\ p = 4,\; q = 4 \\ y^2 + 8y + 16 = (y + 4)(y + 4)
\end{align*}

\textbf{c)} \(\displaystyle x^2 + 11x + 18\)

\begin{align*}
\text{Find } p, q: p \times q = 18,\; p + q = 11 \\ p = 2,\; q = 9 \\ x^2 + 11x + 18 = (x + 2)(x + 9)
\end{align*}

\textbf{d)} \(\displaystyle x^2 + 17x + 70\)

\begin{align*}
\text{Find } p, q: p \times q = 70,\; p + q = 17 \\ p = 10,\; q = 7 \\ x^2 + 17x + 70 = (x + 10)(x + 7)
\end{align*}

\textbf{e)} \(\displaystyle y^2 + 3y + 2\)

\begin{align*}
\text{Find } p, q: p \times q = 2,\; p + q = 3 \\ p = 1,\; q = 2 \\ y^2 + 3y + 2 = (y + 1)(y + 2)
\end{align*}

\textbf{f)} \(\displaystyle x^2 + 13x + 36\)

\begin{align*}
\text{Find } p, q: p \times q = 36,\; p + q = 13 \\ p = 4,\; q = 9 \\ x^2 + 13x + 36 = (x + 4)(x + 9)
\end{align*}

\textbf{g)} \(\displaystyle m^2 + 13m + 36\)

\begin{align*}
\text{Find } p, q: p \times q = 36,\; p + q = 13 \\ p = 9,\; q = 4 \\ m^2 + 13m + 36 = (m + 9)(m + 4)
\end{align*}

\textbf{h)} \(\displaystyle n^2 + 12n + 32\)

\begin{align*}
\text{Find } p, q: p \times q = 32,\; p + q = 12 \\ p = 4,\; q = 8 \\ n^2 + 12n + 32 = (n + 4)(n + 8)
\end{align*}

\end{tcolorbox}

\end{exercise}

\begin{exercise}[]\protect\hypertarget{exr-u6-2}{}\label{exr-u6-2}

\emph{Factorise each expression.}

\textbf{a)} \(\displaystyle m^2 - 2m - 3\)

\textbf{b)} \(\displaystyle y^2 + y - 30\)

\textbf{c)} \(\displaystyle x^2 - 2x - 24\)

\textbf{d)} \(\displaystyle x^2 - 5x - 14\)

\textbf{e)} \(\displaystyle x^2 + 5x - 50\)

\textbf{f)} \(\displaystyle x^2 - x - 72\)

\textbf{g)} \(\displaystyle n^2 + 5n - 14\)

\textbf{h)} \(\displaystyle y^2 + 4y - 60\)

\begin{tcolorbox}[enhanced jigsaw, arc=.35mm, breakable, toptitle=1mm, title=\textcolor{quarto-callout-tip-color}{\faLightbulb}\hspace{0.5em}{🔍 View Solutions}, colframe=quarto-callout-tip-color-frame, rightrule=.15mm, opacityback=0, titlerule=0mm, colback=white, coltitle=black, bottomtitle=1mm, toprule=.15mm, leftrule=.75mm, bottomrule=.15mm, colbacktitle=quarto-callout-tip-color!10!white, left=2mm, opacitybacktitle=0.6]

\textbf{a)} \(\displaystyle m^2 - 2m - 3\)

\begin{align*}
\text{Find } p, q: p \times q = -3,\; p + q = -2 \\ p = 1,\; q = -3 \\ m^2 - 2m - 3 = (m + 1)(m - 3)
\end{align*}

\textbf{b)} \(\displaystyle y^2 + y - 30\)

\begin{align*}
\text{Find } p, q: p \times q = -30,\; p + q = 1 \\ p = 6,\; q = -5 \\ y^2 + y - 30 = (y + 6)(y - 5)
\end{align*}

\textbf{c)} \(\displaystyle x^2 - 2x - 24\)

\begin{align*}
\text{Find } p, q: p \times q = -24,\; p + q = -2 \\ p = 4,\; q = -6 \\ x^2 - 2x - 24 = (x + 4)(x - 6)
\end{align*}

\textbf{d)} \(\displaystyle x^2 - 5x - 14\)

\begin{align*}
\text{Find } p, q: p \times q = -14,\; p + q = -5 \\ p = 2,\; q = -7 \\ x^2 - 5x - 14 = (x + 2)(x - 7)
\end{align*}

\textbf{e)} \(\displaystyle x^2 + 5x - 50\)

\begin{align*}
\text{Find } p, q: p \times q = -50,\; p + q = 5 \\ p = 10,\; q = -5 \\ x^2 + 5x - 50 = (x + 10)(x - 5)
\end{align*}

\textbf{f)} \(\displaystyle x^2 - x - 72\)

\begin{align*}
\text{Find } p, q: p \times q = -72,\; p + q = -1 \\ p = 8,\; q = -9 \\ x^2 - x - 72 = (x + 8)(x - 9)
\end{align*}

\textbf{g)} \(\displaystyle n^2 + 5n - 14\)

\begin{align*}
\text{Find } p, q: p \times q = -14,\; p + q = 5 \\ p = 7,\; q = -2 \\ n^2 + 5n - 14 = (n + 7)(n - 2)
\end{align*}

\textbf{h)} \(\displaystyle y^2 + 4y - 60\)

\begin{align*}
\text{Find } p, q: p \times q = -60,\; p + q = 4 \\ p = 10,\; q = -6 \\ y^2 + 4y - 60 = (y + 10)(y - 6)
\end{align*}

\end{tcolorbox}

\end{exercise}

\begin{exercise}[]\protect\hypertarget{exr-u6-6}{}\label{exr-u6-6}

\emph{Factorise using the difference of two squares.}

\textbf{a)} \(\displaystyle x^2 - 1\)

\textbf{b)} \(\displaystyle y^2 - 121\)

\textbf{c)} \(\displaystyle 4x^2 - 100\)

\textbf{d)} \(\displaystyle 4m^2 - 49\)

\textbf{e)} \(\displaystyle 9x^2 - 25\)

\textbf{f)} \(\displaystyle n^2 - 121\)

\textbf{g)} \(\displaystyle n^2 - 4\)

\textbf{h)} \(\displaystyle 4m^2 - 9\)

\begin{tcolorbox}[enhanced jigsaw, arc=.35mm, breakable, toptitle=1mm, title=\textcolor{quarto-callout-tip-color}{\faLightbulb}\hspace{0.5em}{🔍 View Solutions}, colframe=quarto-callout-tip-color-frame, rightrule=.15mm, opacityback=0, titlerule=0mm, colback=white, coltitle=black, bottomtitle=1mm, toprule=.15mm, leftrule=.75mm, bottomrule=.15mm, colbacktitle=quarto-callout-tip-color!10!white, left=2mm, opacitybacktitle=0.6]

\textbf{a)} \(\displaystyle x^2 - 1\)

\begin{align*}
x^2 - 1^2 = (x + 1)(x - 1)
\end{align*}

\textbf{b)} \(\displaystyle y^2 - 121\)

\begin{align*}
y^2 - 11^2 = (y + 11)(y - 11)
\end{align*}

\textbf{c)} \(\displaystyle 4x^2 - 100\)

\begin{align*}
(2x)^2 - (10)^2 = (2x + 10)(2x - 10)
\end{align*}

\textbf{d)} \(\displaystyle 4m^2 - 49\)

\begin{align*}
(2m)^2 - (7)^2 = (2m + 7)(2m - 7)
\end{align*}

\textbf{e)} \(\displaystyle 9x^2 - 25\)

\begin{align*}
(3x)^2 - (5)^2 = (3x + 5)(3x - 5)
\end{align*}

\textbf{f)} \(\displaystyle n^2 - 121\)

\begin{align*}
n^2 - 11^2 = (n + 11)(n - 11)
\end{align*}

\textbf{g)} \(\displaystyle n^2 - 4\)

\begin{align*}
n^2 - 2^2 = (n + 2)(n - 2)
\end{align*}

\textbf{h)} \(\displaystyle 4m^2 - 9\)

\begin{align*}
(2m)^2 - (3)^2 = (2m + 3)(2m - 3)
\end{align*}

\end{tcolorbox}

\end{exercise}

\begin{exercise}[]\protect\hypertarget{exr-u6-7}{}\label{exr-u6-7}

\emph{Recognise and factorise each perfect square.}

\textbf{a)} \(\displaystyle y^2 - 2y + 1\)

\textbf{b)} \(\displaystyle n^2 + 18n + 81\)

\textbf{c)} \(\displaystyle y^2 + 18y + 81\)

\textbf{d)} \(\displaystyle x^2 + 12x + 36\)

\textbf{e)} \(\displaystyle m^2 - 20m + 100\)

\textbf{f)} \(\displaystyle n^2 + 2n + 1\)

\begin{tcolorbox}[enhanced jigsaw, arc=.35mm, breakable, toptitle=1mm, title=\textcolor{quarto-callout-tip-color}{\faLightbulb}\hspace{0.5em}{🔍 View Solutions}, colframe=quarto-callout-tip-color-frame, rightrule=.15mm, opacityback=0, titlerule=0mm, colback=white, coltitle=black, bottomtitle=1mm, toprule=.15mm, leftrule=.75mm, bottomrule=.15mm, colbacktitle=quarto-callout-tip-color!10!white, left=2mm, opacitybacktitle=0.6]

\textbf{a)} \(\displaystyle y^2 - 2y + 1\)

\begin{align*}
(y - 1)^2 \quad \text{[perfect square]}
\end{align*}

\textbf{b)} \(\displaystyle n^2 + 18n + 81\)

\begin{align*}
(n + 9)^2 \quad \text{[perfect square]}
\end{align*}

\textbf{c)} \(\displaystyle y^2 + 18y + 81\)

\begin{align*}
(y + 9)^2 \quad \text{[perfect square]}
\end{align*}

\textbf{d)} \(\displaystyle x^2 + 12x + 36\)

\begin{align*}
(x + 6)^2 \quad \text{[perfect square]}
\end{align*}

\textbf{e)} \(\displaystyle m^2 - 20m + 100\)

\begin{align*}
(m - 10)^2 \quad \text{[perfect square]}
\end{align*}

\textbf{f)} \(\displaystyle n^2 + 2n + 1\)

\begin{align*}
(n + 1)^2 \quad \text{[perfect square]}
\end{align*}

\end{tcolorbox}

\end{exercise}

\begin{exercise}[]\protect\hypertarget{exr-u6-8}{}\label{exr-u6-8}

\emph{Solve each quadratic equation.}

\textbf{a)} \(\displaystyle x^2 - 1 = 0\)

\textbf{b)} \(\displaystyle x^2 - 7x + 6 = 0\)

\textbf{c)} \(\displaystyle y^2 + y - 42 = 0\)

\textbf{d)} \(\displaystyle y^2 + 3y - 28 = 0\)

\textbf{e)} \(\displaystyle y^2 + 3y - 10 = 0\)

\textbf{f)} \(\displaystyle n^2 + 5n + 4 = 0\)

\textbf{g)} \(\displaystyle m^2 + 14m + 48 = 0\)

\textbf{h)} \(\displaystyle y^2 - 6y + 5 = 0\)

\begin{tcolorbox}[enhanced jigsaw, arc=.35mm, breakable, toptitle=1mm, title=\textcolor{quarto-callout-tip-color}{\faLightbulb}\hspace{0.5em}{🔍 View Solutions}, colframe=quarto-callout-tip-color-frame, rightrule=.15mm, opacityback=0, titlerule=0mm, colback=white, coltitle=black, bottomtitle=1mm, toprule=.15mm, leftrule=.75mm, bottomrule=.15mm, colbacktitle=quarto-callout-tip-color!10!white, left=2mm, opacitybacktitle=0.6]

\textbf{a)} \(\displaystyle x^2 - 1 = 0\)

\begin{align*}
\text{Find } p, q: p \times q = -1,\; p + q = 0 \\ p = 1,\; q = -1 \\ (x + 1)(x - 1) = 0 \\ x = -1 \text{ or } x = 1
\end{align*}

\textbf{b)} \(\displaystyle x^2 - 7x + 6 = 0\)

\begin{align*}
\text{Find } p, q: p \times q = 6,\; p + q = -7 \\ p = -1,\; q = -6 \\ (x - 1)(x - 6) = 0 \\ x = 1 \text{ or } x = 6
\end{align*}

\textbf{c)} \(\displaystyle y^2 + y - 42 = 0\)

\begin{align*}
\text{Find } p, q: p \times q = -42,\; p + q = 1 \\ p = 7,\; q = -6 \\ (y + 7)(y - 6) = 0 \\ y = -7 \text{ or } y = 6
\end{align*}

\textbf{d)} \(\displaystyle y^2 + 3y - 28 = 0\)

\begin{align*}
\text{Find } p, q: p \times q = -28,\; p + q = 3 \\ p = -4,\; q = 7 \\ (y - 4)(y + 7) = 0 \\ y = 4 \text{ or } y = -7
\end{align*}

\textbf{e)} \(\displaystyle y^2 + 3y - 10 = 0\)

\begin{align*}
\text{Find } p, q: p \times q = -10,\; p + q = 3 \\ p = 5,\; q = -2 \\ (y + 5)(y - 2) = 0 \\ y = -5 \text{ or } y = 2
\end{align*}

\textbf{f)} \(\displaystyle n^2 + 5n + 4 = 0\)

\begin{align*}
\text{Find } p, q: p \times q = 4,\; p + q = 5 \\ p = 1,\; q = 4 \\ (n + 1)(n + 4) = 0 \\ n = -1 \text{ or } n = -4
\end{align*}

\textbf{g)} \(\displaystyle m^2 + 14m + 48 = 0\)

\begin{align*}
\text{Find } p, q: p \times q = 48,\; p + q = 14 \\ p = 6,\; q = 8 \\ (m + 6)(m + 8) = 0 \\ m = -6 \text{ or } m = -8
\end{align*}

\textbf{h)} \(\displaystyle y^2 - 6y + 5 = 0\)

\begin{align*}
\text{Find } p, q: p \times q = 5,\; p + q = -6 \\ p = -5,\; q = -1 \\ (y - 5)(y - 1) = 0 \\ y = 5 \text{ or } y = 1
\end{align*}

\end{tcolorbox}

\end{exercise}

∑ × π ÷ √ ◆ ∞ ± ∫ ≠ Δ

\subsection{🟡 Medium}

\begin{exercise}[]\protect\hypertarget{exr-u6-3}{}\label{exr-u6-3}

\emph{Factorise each expression.}

\textbf{a)} \(\displaystyle y^2 - 3y + 2\)

\textbf{b)} \(\displaystyle y^2 - 7y + 10\)

\textbf{c)} \(\displaystyle n^2 - 9n + 14\)

\textbf{d)} \(\displaystyle y^2 - 14y + 48\)

\textbf{e)} \(\displaystyle y^2 - 12y + 36\)

\textbf{f)} \(\displaystyle y^2 - 7y + 10\)

\begin{tcolorbox}[enhanced jigsaw, arc=.35mm, breakable, toptitle=1mm, title=\textcolor{quarto-callout-tip-color}{\faLightbulb}\hspace{0.5em}{🔍 View Solutions}, colframe=quarto-callout-tip-color-frame, rightrule=.15mm, opacityback=0, titlerule=0mm, colback=white, coltitle=black, bottomtitle=1mm, toprule=.15mm, leftrule=.75mm, bottomrule=.15mm, colbacktitle=quarto-callout-tip-color!10!white, left=2mm, opacitybacktitle=0.6]

\textbf{a)} \(\displaystyle y^2 - 3y + 2\)

\begin{align*}
\text{Find } p, q: p \times q = 2,\; p + q = -3 \\ p = -2,\; q = -1 \\ y^2 - 3y + 2 = (y - 2)(y - 1)
\end{align*}

\textbf{b)} \(\displaystyle y^2 - 7y + 10\)

\begin{align*}
\text{Find } p, q: p \times q = 10,\; p + q = -7 \\ p = -5,\; q = -2 \\ y^2 - 7y + 10 = (y - 5)(y - 2)
\end{align*}

\textbf{c)} \(\displaystyle n^2 - 9n + 14\)

\begin{align*}
\text{Find } p, q: p \times q = 14,\; p + q = -9 \\ p = -2,\; q = -7 \\ n^2 - 9n + 14 = (n - 2)(n - 7)
\end{align*}

\textbf{d)} \(\displaystyle y^2 - 14y + 48\)

\begin{align*}
\text{Find } p, q: p \times q = 48,\; p + q = -14 \\ p = -8,\; q = -6 \\ y^2 - 14y + 48 = (y - 8)(y - 6)
\end{align*}

\textbf{e)} \(\displaystyle y^2 - 12y + 36\)

\begin{align*}
\text{Find } p, q: p \times q = 36,\; p + q = -12 \\ p = -6,\; q = -6 \\ y^2 - 12y + 36 = (y - 6)(y - 6)
\end{align*}

\textbf{f)} \(\displaystyle y^2 - 7y + 10\)

\begin{align*}
\text{Find } p, q: p \times q = 10,\; p + q = -7 \\ p = -5,\; q = -2 \\ y^2 - 7y + 10 = (y - 5)(y - 2)
\end{align*}

\end{tcolorbox}

\end{exercise}

\begin{exercise}[]\protect\hypertarget{exr-u6-4}{}\label{exr-u6-4}

\emph{Factorise each expression.}

\textbf{a)} \(\displaystyle 3m^2 + 14m + 8\)

\textbf{b)} \(\displaystyle 4y^2 + 30y + 36\)

\textbf{c)} \(\displaystyle 4x^2 + 9x + 2\)

\textbf{d)} \(\displaystyle 4x^2 + 16x + 12\)

\textbf{e)} \(\displaystyle 3n^2 + 23n + 30\)

\textbf{f)} \(\displaystyle 3m^2 + 18m + 24\)

\begin{tcolorbox}[enhanced jigsaw, arc=.35mm, breakable, toptitle=1mm, title=\textcolor{quarto-callout-tip-color}{\faLightbulb}\hspace{0.5em}{🔍 View Solutions}, colframe=quarto-callout-tip-color-frame, rightrule=.15mm, opacityback=0, titlerule=0mm, colback=white, coltitle=black, bottomtitle=1mm, toprule=.15mm, leftrule=.75mm, bottomrule=.15mm, colbacktitle=quarto-callout-tip-color!10!white, left=2mm, opacitybacktitle=0.6]

\textbf{a)} \(\displaystyle 3m^2 + 14m + 8\)

\begin{align*}
ac = 3 \times 8 = 24 \\ \text{Find } p, q: p \times q = 24,\; p + q = 14 \\ p = 2,\; q = 12 \\ 3m^2 + 2m + 12m + 8 \\ m(3m + 2) + \text{group} \\ 3m^2 + 14m + 8 = (3m + 2)(m + 4)
\end{align*}

\textbf{b)} \(\displaystyle 4y^2 + 30y + 36\)

\begin{align*}
ac = 4 \times 36 = 144 \\ \text{Find } p, q: p \times q = 144,\; p + q = 30 \\ p = 6,\; q = 24 \\ 4y^2 + 6y + 24y + 36 \\ y(4y + 3) + \text{group} \\ 4y^2 + 30y + 36 = (4y + 6)(y + 6)
\end{align*}

\textbf{c)} \(\displaystyle 4x^2 + 9x + 2\)

\begin{align*}
ac = 4 \times 2 = 8 \\ \text{Find } p, q: p \times q = 8,\; p + q = 9 \\ p = 1,\; q = 8 \\ 4x^2 + x + 8x + 2 \\ x(4x + 1) + \text{group} \\ 4x^2 + 9x + 2 = (4x + 1)(x + 2)
\end{align*}

\textbf{d)} \(\displaystyle 4x^2 + 16x + 12\)

\begin{align*}
ac = 4 \times 12 = 48 \\ \text{Find } p, q: p \times q = 48,\; p + q = 16 \\ p = 4,\; q = 12 \\ 4x^2 + 4x + 12x + 12 \\ x(4x + 1) + \text{group} \\ 4x^2 + 16x + 12 = (4x + 4)(x + 3)
\end{align*}

\textbf{e)} \(\displaystyle 3n^2 + 23n + 30\)

\begin{align*}
ac = 3 \times 30 = 90 \\ \text{Find } p, q: p \times q = 90,\; p + q = 23 \\ p = 5,\; q = 18 \\ 3n^2 + 5n + 18n + 30 \\ n(3n + 5) + \text{group} \\ 3n^2 + 23n + 30 = (3n + 5)(n + 6)
\end{align*}

\textbf{f)} \(\displaystyle 3m^2 + 18m + 24\)

\begin{align*}
ac = 3 \times 24 = 72 \\ \text{Find } p, q: p \times q = 72,\; p + q = 18 \\ p = 6,\; q = 12 \\ 3m^2 + 6m + 12m + 24 \\ m(3m + 2) + \text{group} \\ 3m^2 + 18m + 24 = (3m + 6)(m + 4)
\end{align*}

\end{tcolorbox}

\end{exercise}

\begin{exercise}[]\protect\hypertarget{exr-u6-9}{}\label{exr-u6-9}

\emph{Solve each quadratic equation.}

\textbf{a)} \(\displaystyle 2y^2 - 13y + 6 = 0\)

\textbf{b)} \(\displaystyle 4x^2 - 27x + 18 = 0\)

\textbf{c)} \(\displaystyle 4y^2 - 20y - 24 = 0\)

\textbf{d)} \(\displaystyle 2x^2 - 8x + 6 = 0\)

\textbf{e)} \(\displaystyle 4m^2 - 11m + 6 = 0\)

\textbf{f)} \(\displaystyle 2m^2 - 6m - 8 = 0\)

\begin{tcolorbox}[enhanced jigsaw, arc=.35mm, breakable, toptitle=1mm, title=\textcolor{quarto-callout-tip-color}{\faLightbulb}\hspace{0.5em}{🔍 View Solutions}, colframe=quarto-callout-tip-color-frame, rightrule=.15mm, opacityback=0, titlerule=0mm, colback=white, coltitle=black, bottomtitle=1mm, toprule=.15mm, leftrule=.75mm, bottomrule=.15mm, colbacktitle=quarto-callout-tip-color!10!white, left=2mm, opacitybacktitle=0.6]

\textbf{a)} \(\displaystyle 2y^2 - 13y + 6 = 0\)

\begin{align*}
ac = 2 \times 6 = 12 \\ p = -12,\; q = -1 \\ (2y - 1)(y - 6) = 0 \\ y = \dfrac{1}{2} \text{ or } y = 6
\end{align*}

\textbf{b)} \(\displaystyle 4x^2 - 27x + 18 = 0\)

\begin{align*}
ac = 4 \times 18 = 72 \\ p = -24,\; q = -3 \\ (4x - 3)(x - 6) = 0 \\ x = \dfrac{3}{4} \text{ or } x = 6
\end{align*}

\textbf{c)} \(\displaystyle 4y^2 - 20y - 24 = 0\)

\begin{align*}
ac = 4 \times -24 = -96 \\ p = -24,\; q = 4 \\ (4y + 4)(y - 6) = 0 \\ y = -1 \text{ or } y = 6
\end{align*}

\textbf{d)} \(\displaystyle 2x^2 - 8x + 6 = 0\)

\begin{align*}
ac = 2 \times 6 = 12 \\ p = -6,\; q = -2 \\ (2x - 6)(x - 1) = 0 \\ x = 3 \text{ or } x = 1
\end{align*}

\textbf{e)} \(\displaystyle 4m^2 - 11m + 6 = 0\)

\begin{align*}
ac = 4 \times 6 = 24 \\ p = -8,\; q = -3 \\ (4m - 3)(m - 2) = 0 \\ m = \dfrac{3}{4} \text{ or } m = 2
\end{align*}

\textbf{f)} \(\displaystyle 2m^2 - 6m - 8 = 0\)

\begin{align*}
ac = 2 \times -8 = -16 \\ p = -8,\; q = 2 \\ (2m + 2)(m - 4) = 0 \\ m = -1 \text{ or } m = 4
\end{align*}

\end{tcolorbox}

\end{exercise}

\begin{exercise}[]\protect\hypertarget{exr-u6-10}{}\label{exr-u6-10}

\emph{Solve each equation.}

\textbf{a)} \(\displaystyle m^2 - 64 = 0\)

\textbf{b)} \(\displaystyle 4x^2 - 64 = 0\)

\textbf{c)} \(\displaystyle m^2 + 16m + 64 = 0\)

\textbf{d)} \(\displaystyle 4x^2 - 64 = 0\)

\textbf{e)} \(\displaystyle y^2 - 25 = 0\)

\textbf{f)} \(\displaystyle y^2 + 14y + 49 = 0\)

\begin{tcolorbox}[enhanced jigsaw, arc=.35mm, breakable, toptitle=1mm, title=\textcolor{quarto-callout-tip-color}{\faLightbulb}\hspace{0.5em}{🔍 View Solutions}, colframe=quarto-callout-tip-color-frame, rightrule=.15mm, opacityback=0, titlerule=0mm, colback=white, coltitle=black, bottomtitle=1mm, toprule=.15mm, leftrule=.75mm, bottomrule=.15mm, colbacktitle=quarto-callout-tip-color!10!white, left=2mm, opacitybacktitle=0.6]

\textbf{a)} \(\displaystyle m^2 - 64 = 0\)

\begin{align*}
(m + 8)(m - 8) = 0 \\ m = 8 \text{ or } m = -8
\end{align*}

\textbf{b)} \(\displaystyle 4x^2 - 64 = 0\)

\begin{align*}
(2x + 8)(2x - 8) = 0 \\ x = 4 \text{ or } x = -4
\end{align*}

\textbf{c)} \(\displaystyle m^2 + 16m + 64 = 0\)

\begin{align*}
(m + 8)^2 = 0 \\ m = -8 \text{ (repeated root)}
\end{align*}

\textbf{d)} \(\displaystyle 4x^2 - 64 = 0\)

\begin{align*}
(2x + 8)(2x - 8) = 0 \\ x = 4 \text{ or } x = -4
\end{align*}

\textbf{e)} \(\displaystyle y^2 - 25 = 0\)

\begin{align*}
(y + 5)(y - 5) = 0 \\ y = 5 \text{ or } y = -5
\end{align*}

\textbf{f)} \(\displaystyle y^2 + 14y + 49 = 0\)

\begin{align*}
(y + 7)^2 = 0 \\ y = -7 \text{ (repeated root)}
\end{align*}

\end{tcolorbox}

\end{exercise}

\begin{exercise}[]\protect\hypertarget{exr-u6-11}{}\label{exr-u6-11}

\emph{Solve using the quadratic formula.}

\textbf{a)} \(\displaystyle x^2 + 7x + 6 = 0\)

\textbf{b)} \(\displaystyle y^2 + 15y + 56 = 0\)

\textbf{c)} \(\displaystyle m^2 - 2m - 15 = 0\)

\textbf{d)} \(\displaystyle y^2 - 2y - 8 = 0\)

\textbf{e)} \(\displaystyle n^2 + 4n - 32 = 0\)

\textbf{f)} \(\displaystyle m^2 - 7m + 6 = 0\)

\begin{tcolorbox}[enhanced jigsaw, arc=.35mm, breakable, toptitle=1mm, title=\textcolor{quarto-callout-tip-color}{\faLightbulb}\hspace{0.5em}{🔍 View Solutions}, colframe=quarto-callout-tip-color-frame, rightrule=.15mm, opacityback=0, titlerule=0mm, colback=white, coltitle=black, bottomtitle=1mm, toprule=.15mm, leftrule=.75mm, bottomrule=.15mm, colbacktitle=quarto-callout-tip-color!10!white, left=2mm, opacitybacktitle=0.6]

\textbf{a)} \(\displaystyle x^2 + 7x + 6 = 0\)

\begin{align*}
x = \dfrac{-(7) \pm \sqrt{(7)^2 - 4(1)(6)}}{2} \\ x = \dfrac{-7 \pm 5}{2} \\ x = -1 \text{ or } x = -6
\end{align*}

\textbf{b)} \(\displaystyle y^2 + 15y + 56 = 0\)

\begin{align*}
y = \dfrac{-(15) \pm \sqrt{(15)^2 - 4(1)(56)}}{2} \\ y = \dfrac{-15 \pm 1}{2} \\ y = -7 \text{ or } y = -8
\end{align*}

\textbf{c)} \(\displaystyle m^2 - 2m - 15 = 0\)

\begin{align*}
m = \dfrac{-(-2) \pm \sqrt{(-2)^2 - 4(1)(-15)}}{2} \\ m = \dfrac{2 \pm 8}{2} \\ m = 5 \text{ or } m = -3
\end{align*}

\textbf{d)} \(\displaystyle y^2 - 2y - 8 = 0\)

\begin{align*}
y = \dfrac{-(-2) \pm \sqrt{(-2)^2 - 4(1)(-8)}}{2} \\ y = \dfrac{2 \pm 6}{2} \\ y = 4 \text{ or } y = -2
\end{align*}

\textbf{e)} \(\displaystyle n^2 + 4n - 32 = 0\)

\begin{align*}
n = \dfrac{-(4) \pm \sqrt{(4)^2 - 4(1)(-32)}}{2} \\ n = \dfrac{-4 \pm 12}{2} \\ n = 4 \text{ or } n = -8
\end{align*}

\textbf{f)} \(\displaystyle m^2 - 7m + 6 = 0\)

\begin{align*}
m = \dfrac{-(-7) \pm \sqrt{(-7)^2 - 4(1)(6)}}{2} \\ m = \dfrac{7 \pm 5}{2} \\ m = 6 \text{ or } m = 1
\end{align*}

\end{tcolorbox}

\end{exercise}

\begin{exercise}[]\protect\hypertarget{exr-u6-14}{}\label{exr-u6-14}

\emph{Write each expression in completed square form.}

\textbf{a)} \(\displaystyle n^2 + 2n - 6\)

\textbf{b)} \(\displaystyle x^2 + 4x\)

\textbf{c)} \(\displaystyle x^2 - 10x + 4\)

\textbf{d)} \(\displaystyle y^2 - 8y + 7\)

\textbf{e)} \(\displaystyle y^2 + 6y - 2\)

\textbf{f)} \(\displaystyle y^2 + y - 8\)

\begin{tcolorbox}[enhanced jigsaw, arc=.35mm, breakable, toptitle=1mm, title=\textcolor{quarto-callout-tip-color}{\faLightbulb}\hspace{0.5em}{🔍 View Solutions}, colframe=quarto-callout-tip-color-frame, rightrule=.15mm, opacityback=0, titlerule=0mm, colback=white, coltitle=black, bottomtitle=1mm, toprule=.15mm, leftrule=.75mm, bottomrule=.15mm, colbacktitle=quarto-callout-tip-color!10!white, left=2mm, opacitybacktitle=0.6]

\textbf{a)} \(\displaystyle n^2 + 2n - 6\)

\begin{align*}
(n + 1)^2 - 7
\end{align*}

\textbf{b)} \(\displaystyle x^2 + 4x\)

\begin{align*}
(x + 2)^2 - 4
\end{align*}

\textbf{c)} \(\displaystyle x^2 - 10x + 4\)

\begin{align*}
\left(x - \dfrac{10}{2}\right)^2 - \left(\dfrac{10}{2}\right)^2 + 4 \\ (x - 5)^2 - 25 + 4 \\ (x - 5)^2 - 21
\end{align*}

\textbf{d)} \(\displaystyle y^2 - 8y + 7\)

\begin{align*}
\left(y - \dfrac{8}{2}\right)^2 - \left(\dfrac{8}{2}\right)^2 + 7 \\ (y - 4)^2 - 16 + 7 \\ (y - 4)^2 - 9
\end{align*}

\textbf{e)} \(\displaystyle y^2 + 6y - 2\)

\begin{align*}
(y + 3)^2 - 11
\end{align*}

\textbf{f)} \(\displaystyle y^2 + y - 8\)

\begin{align*}
(y + \dfrac{1}{2})^2 - \dfrac{33}{4}
\end{align*}

\end{tcolorbox}

\end{exercise}

\begin{exercise}[]\protect\hypertarget{exr-u6-18-1}{}\label{exr-u6-18-1}

~

A rectangle has length \((x + 3)\) cm and width \((x + 2)\) cm. The area
is 42 cm².

\textbf{(i)} Form a quadratic equation in \(x\).

\textbf{(ii)} Solve the equation, giving exact answers.

\textbf{(iii)} Find the dimensions of the rectangle.

\begin{tcolorbox}[enhanced jigsaw, arc=.35mm, breakable, toptitle=1mm, title=\textcolor{quarto-callout-tip-color}{\faLightbulb}\hspace{0.5em}{🔍 View Solution}, colframe=quarto-callout-tip-color-frame, rightrule=.15mm, opacityback=0, titlerule=0mm, colback=white, coltitle=black, bottomtitle=1mm, toprule=.15mm, leftrule=.75mm, bottomrule=.15mm, colbacktitle=quarto-callout-tip-color!10!white, left=2mm, opacitybacktitle=0.6]

\begin{align*}
\text{(i) } (x + 3)(x + 2) = 42 \\ x^2 + 5x + 6 = 42 \\ x^2 +5x - 36 = 0 \\ \text{(ii) } x = \dfrac{-(5) \pm \sqrt{(5)^2 - 4(1)(-36)}}{2} \\ x = \dfrac{-5 \pm 13}{2} \\ x = 4 \text{ or } x = -9 \\ \text{(iii) Taking positive root: } x = 4 \\ \text{Length } = 4 + 3 = 7 \text{ cm, }\text{ Width } = 4 + 2 = 6 \text{ cm}
\end{align*}

\textbf{Answer:} \(x = 4\)

\end{tcolorbox}

\end{exercise}

\begin{exercise}[]\protect\hypertarget{exr-u6-18-2}{}\label{exr-u6-18-2}

~

A rectangle has length \((x + 3)\) cm and width \((x + 4)\) cm. The area
is 30 cm².

\textbf{(i)} Form a quadratic equation in \(x\).

\textbf{(ii)} Solve the equation, giving exact answers.

\textbf{(iii)} Find the dimensions of the rectangle.

\begin{tcolorbox}[enhanced jigsaw, arc=.35mm, breakable, toptitle=1mm, title=\textcolor{quarto-callout-tip-color}{\faLightbulb}\hspace{0.5em}{🔍 View Solution}, colframe=quarto-callout-tip-color-frame, rightrule=.15mm, opacityback=0, titlerule=0mm, colback=white, coltitle=black, bottomtitle=1mm, toprule=.15mm, leftrule=.75mm, bottomrule=.15mm, colbacktitle=quarto-callout-tip-color!10!white, left=2mm, opacitybacktitle=0.6]

\begin{align*}
\text{(i) } (x + 3)(x + 4) = 30 \\ x^2 + 7x + 12 = 30 \\ x^2 +7x - 18 = 0 \\ \text{(ii) } x = \dfrac{-(7) \pm \sqrt{(7)^2 - 4(1)(-18)}}{2} \\ x = \dfrac{-7 \pm 11}{2} \\ x = 2 \text{ or } x = -9 \\ \text{(iii) Taking positive root: } x = 2 \\ \text{Length } = 2 + 3 = 5 \text{ cm, }\text{ Width } = 2 + 4 = 6 \text{ cm}
\end{align*}

\textbf{Answer:} \(x = 2\)

\end{tcolorbox}

\end{exercise}

\begin{exercise}[]\protect\hypertarget{exr-u6-18-3}{}\label{exr-u6-18-3}

~

A rectangle has length \((x + 3)\) cm and width \((x + 6)\) cm. The area
is 130 cm².

\textbf{(i)} Form a quadratic equation in \(x\).

\textbf{(ii)} Solve the equation, giving exact answers.

\textbf{(iii)} Find the dimensions of the rectangle.

\begin{tcolorbox}[enhanced jigsaw, arc=.35mm, breakable, toptitle=1mm, title=\textcolor{quarto-callout-tip-color}{\faLightbulb}\hspace{0.5em}{🔍 View Solution}, colframe=quarto-callout-tip-color-frame, rightrule=.15mm, opacityback=0, titlerule=0mm, colback=white, coltitle=black, bottomtitle=1mm, toprule=.15mm, leftrule=.75mm, bottomrule=.15mm, colbacktitle=quarto-callout-tip-color!10!white, left=2mm, opacitybacktitle=0.6]

\begin{align*}
\text{(i) } (x + 3)(x + 6) = 130 \\ x^2 + 9x + 18 = 130 \\ x^2 +9x - 112 = 0 \\ \text{(ii) } x = \dfrac{-(9) \pm \sqrt{(9)^2 - 4(1)(-112)}}{2} \\ x = \dfrac{-9 \pm 23}{2} \\ x = 7 \text{ or } x = -16 \\ \text{(iii) Taking positive root: } x = 7 \\ \text{Length } = 7 + 3 = 10 \text{ cm, }\text{ Width } = 7 + 6 = 13 \text{ cm}
\end{align*}

\textbf{Answer:} \(x = 7\)

\end{tcolorbox}

\end{exercise}

\begin{exercise}[]\protect\hypertarget{exr-u6-19-1}{}\label{exr-u6-19-1}

~

The product of two consecutive positive integers is 240.

\textbf{(i)} Form a quadratic equation in \(n\).

\textbf{(ii)} Solve the equation exactly.

\textbf{(iii)} Find the two integers.

\begin{tcolorbox}[enhanced jigsaw, arc=.35mm, breakable, toptitle=1mm, title=\textcolor{quarto-callout-tip-color}{\faLightbulb}\hspace{0.5em}{🔍 View Solution}, colframe=quarto-callout-tip-color-frame, rightrule=.15mm, opacityback=0, titlerule=0mm, colback=white, coltitle=black, bottomtitle=1mm, toprule=.15mm, leftrule=.75mm, bottomrule=.15mm, colbacktitle=quarto-callout-tip-color!10!white, left=2mm, opacitybacktitle=0.6]

\begin{align*}
\text{(i) } n(n+1) = 240 \\ n^2 + n - 240 = 0 \\ \text{(ii) } n = \dfrac{-(1) \pm \sqrt{(1)^2 - 4(1)(-240)}}{2} \\ n = \dfrac{-1 \pm 31}{2} \\ n = 15 \text{ or } n = -16 \\ \text{(iii) Since integers are positive, } n = 15 \\ \text{Integers: } 15 \text{ and } 16
\end{align*}

\textbf{Answer:} \(Integers: 15 and 16\)

\end{tcolorbox}

\end{exercise}

\begin{exercise}[]\protect\hypertarget{exr-u6-19-2}{}\label{exr-u6-19-2}

~

The product of two consecutive positive integers is 90.

\textbf{(i)} Form a quadratic equation in \(n\).

\textbf{(ii)} Solve the equation exactly.

\textbf{(iii)} Find the two integers.

\begin{tcolorbox}[enhanced jigsaw, arc=.35mm, breakable, toptitle=1mm, title=\textcolor{quarto-callout-tip-color}{\faLightbulb}\hspace{0.5em}{🔍 View Solution}, colframe=quarto-callout-tip-color-frame, rightrule=.15mm, opacityback=0, titlerule=0mm, colback=white, coltitle=black, bottomtitle=1mm, toprule=.15mm, leftrule=.75mm, bottomrule=.15mm, colbacktitle=quarto-callout-tip-color!10!white, left=2mm, opacitybacktitle=0.6]

\begin{align*}
\text{(i) } n(n+1) = 90 \\ n^2 + n - 90 = 0 \\ \text{(ii) } n = \dfrac{-(1) \pm \sqrt{(1)^2 - 4(1)(-90)}}{2} \\ n = \dfrac{-1 \pm 19}{2} \\ n = 9 \text{ or } n = -10 \\ \text{(iii) Since integers are positive, } n = 9 \\ \text{Integers: } 9 \text{ and } 10
\end{align*}

\textbf{Answer:} \(Integers: 9 and 10\)

\end{tcolorbox}

\end{exercise}

\begin{exercise}[]\protect\hypertarget{exr-u6-19-3}{}\label{exr-u6-19-3}

~

The product of two consecutive positive integers is 240.

\textbf{(i)} Form a quadratic equation in \(n\).

\textbf{(ii)} Solve the equation exactly.

\textbf{(iii)} Find the two integers.

\begin{tcolorbox}[enhanced jigsaw, arc=.35mm, breakable, toptitle=1mm, title=\textcolor{quarto-callout-tip-color}{\faLightbulb}\hspace{0.5em}{🔍 View Solution}, colframe=quarto-callout-tip-color-frame, rightrule=.15mm, opacityback=0, titlerule=0mm, colback=white, coltitle=black, bottomtitle=1mm, toprule=.15mm, leftrule=.75mm, bottomrule=.15mm, colbacktitle=quarto-callout-tip-color!10!white, left=2mm, opacitybacktitle=0.6]

\begin{align*}
\text{(i) } n(n+1) = 240 \\ n^2 + n - 240 = 0 \\ \text{(ii) } n = \dfrac{-(1) \pm \sqrt{(1)^2 - 4(1)(-240)}}{2} \\ n = \dfrac{-1 \pm 31}{2} \\ n = 15 \text{ or } n = -16 \\ \text{(iii) Since integers are positive, } n = 15 \\ \text{Integers: } 15 \text{ and } 16
\end{align*}

\textbf{Answer:} \(Integers: 15 and 16\)

\end{tcolorbox}

\end{exercise}

\begin{exercise}[]\protect\hypertarget{exr-u6-21-1}{}\label{exr-u6-21-1}

~

A rectangle has perimeter 32 cm and area 28 cm². Let the length be \(x\)
cm, so the width is \((16 - x)\) cm.

\textbf{(i)} Show that \(x^2 - 16x + 28 = 0\).

\textbf{(ii)} Solve the equation exactly.

\textbf{(iii)} Find the dimensions of the rectangle.

\begin{tcolorbox}[enhanced jigsaw, arc=.35mm, breakable, toptitle=1mm, title=\textcolor{quarto-callout-tip-color}{\faLightbulb}\hspace{0.5em}{🔍 View Solution}, colframe=quarto-callout-tip-color-frame, rightrule=.15mm, opacityback=0, titlerule=0mm, colback=white, coltitle=black, bottomtitle=1mm, toprule=.15mm, leftrule=.75mm, bottomrule=.15mm, colbacktitle=quarto-callout-tip-color!10!white, left=2mm, opacitybacktitle=0.6]

\begin{align*}
\text{(i) } x(16 - x) = 28 \\ x^2 - 16x + 28 = 0 \\ \text{(ii) } x = \dfrac{-(-16) \pm \sqrt{(-16)^2 - 4(1)(28)}}{2} \\ x = \dfrac{16 \pm 12}{2} \\ x = 14 \text{ or } x = 2 \\ \text{(iii) } x = 14 \text{ or } x = 2 \\ \text{Length } = 14 \text{ cm, Width } = 2 \text{ cm}
\end{align*}

\textbf{Answer:} \(Length = 14 cm, Width = 2 cm\)

\end{tcolorbox}

\end{exercise}

\begin{exercise}[]\protect\hypertarget{exr-u6-21-2}{}\label{exr-u6-21-2}

~

A rectangle has perimeter 32 cm and area 55 cm². Let the length be \(x\)
cm, so the width is \((16 - x)\) cm.

\textbf{(i)} Show that \(x^2 - 16x + 55 = 0\).

\textbf{(ii)} Solve the equation exactly.

\textbf{(iii)} Find the dimensions of the rectangle.

\begin{tcolorbox}[enhanced jigsaw, arc=.35mm, breakable, toptitle=1mm, title=\textcolor{quarto-callout-tip-color}{\faLightbulb}\hspace{0.5em}{🔍 View Solution}, colframe=quarto-callout-tip-color-frame, rightrule=.15mm, opacityback=0, titlerule=0mm, colback=white, coltitle=black, bottomtitle=1mm, toprule=.15mm, leftrule=.75mm, bottomrule=.15mm, colbacktitle=quarto-callout-tip-color!10!white, left=2mm, opacitybacktitle=0.6]

\begin{align*}
\text{(i) } x(16 - x) = 55 \\ x^2 - 16x + 55 = 0 \\ \text{(ii) } x = \dfrac{-(-16) \pm \sqrt{(-16)^2 - 4(1)(55)}}{2} \\ x = \dfrac{16 \pm 6}{2} \\ x = 11 \text{ or } x = 5 \\ \text{(iii) } x = 11 \text{ or } x = 5 \\ \text{Length } = 11 \text{ cm, Width } = 5 \text{ cm}
\end{align*}

\textbf{Answer:} \(Length = 11 cm, Width = 5 cm\)

\end{tcolorbox}

\end{exercise}

\begin{exercise}[]\protect\hypertarget{exr-u6-23}{}\label{exr-u6-23}

\emph{Find the discriminant and state the nature of the roots.}

\textbf{a)} \(\displaystyle 3y^2 - 5y = 0\)

\textbf{b)} \(\displaystyle m^2 - 7m + 5 = 0\)

\textbf{c)} \(\displaystyle y^2 + 8y + 6 = 0\)

\textbf{d)} \(\displaystyle n^2 + 5n - 5 = 0\)

\textbf{e)} \(\displaystyle 3x^2 + 2x + 4 = 0\)

\textbf{f)} \(\displaystyle m^2 + 8m - 2 = 0\)

\begin{tcolorbox}[enhanced jigsaw, arc=.35mm, breakable, toptitle=1mm, title=\textcolor{quarto-callout-tip-color}{\faLightbulb}\hspace{0.5em}{🔍 View Solutions}, colframe=quarto-callout-tip-color-frame, rightrule=.15mm, opacityback=0, titlerule=0mm, colback=white, coltitle=black, bottomtitle=1mm, toprule=.15mm, leftrule=.75mm, bottomrule=.15mm, colbacktitle=quarto-callout-tip-color!10!white, left=2mm, opacitybacktitle=0.6]

\textbf{a)} \(\displaystyle 3y^2 - 5y = 0\)

\begin{align*}
\Delta = b^2 - 4ac = (-5)^2 - 4(3)(0) \\ \Delta = 25 - 0 = 25 \\ \Delta > 0 \Rightarrow \text{Two distinct real roots}
\end{align*}

\textbf{b)} \(\displaystyle m^2 - 7m + 5 = 0\)

\begin{align*}
\Delta = b^2 - 4ac = (-7)^2 - 4(1)(5) \\ \Delta = 49 - 20 = 29 \\ \Delta > 0 \Rightarrow \text{Two distinct real roots}
\end{align*}

\textbf{c)} \(\displaystyle y^2 + 8y + 6 = 0\)

\begin{align*}
\Delta = b^2 - 4ac = (8)^2 - 4(1)(6) \\ \Delta = 64 - 24 = 40 \\ \Delta > 0 \Rightarrow \text{Two distinct real roots}
\end{align*}

\textbf{d)} \(\displaystyle n^2 + 5n - 5 = 0\)

\begin{align*}
\Delta = b^2 - 4ac = (5)^2 - 4(1)(-5) \\ \Delta = 25 + 20 = 45 \\ \Delta > 0 \Rightarrow \text{Two distinct real roots}
\end{align*}

\textbf{e)} \(\displaystyle 3x^2 + 2x + 4 = 0\)

\begin{align*}
\Delta = b^2 - 4ac = (2)^2 - 4(3)(4) \\ \Delta = 4 - 48 = -44 \\ \Delta < 0 \Rightarrow \text{No real roots}
\end{align*}

\textbf{f)} \(\displaystyle m^2 + 8m - 2 = 0\)

\begin{align*}
\Delta = b^2 - 4ac = (8)^2 - 4(1)(-2) \\ \Delta = 64 + 8 = 72 \\ \Delta > 0 \Rightarrow \text{Two distinct real roots}
\end{align*}

\end{tcolorbox}

\end{exercise}

\begin{exercise}[]\protect\hypertarget{exr-u6-25-1}{}\label{exr-u6-25-1}

~

The product of two consecutive positive even integers is 624.

\textbf{(i)} Form a quadratic equation in \(n\).

\textbf{(ii)} Solve the equation exactly.

\textbf{(iii)} Find the two integers.

\begin{tcolorbox}[enhanced jigsaw, arc=.35mm, breakable, toptitle=1mm, title=\textcolor{quarto-callout-tip-color}{\faLightbulb}\hspace{0.5em}{🔍 View Solution}, colframe=quarto-callout-tip-color-frame, rightrule=.15mm, opacityback=0, titlerule=0mm, colback=white, coltitle=black, bottomtitle=1mm, toprule=.15mm, leftrule=.75mm, bottomrule=.15mm, colbacktitle=quarto-callout-tip-color!10!white, left=2mm, opacitybacktitle=0.6]

\begin{align*}
\text{(i) } n(n+2) = 624 \\ n^2 + 2n - 624 = 0 \\ \text{(ii) } n = \dfrac{-(2) \pm \sqrt{(2)^2 - 4(1)(-624)}}{2} \\ n = \dfrac{-2 \pm 50}{2} \\ n = 24 \text{ or } n = -26 \\ \text{Since integers are positive even, } n = 24 \\ \text{Integers: } 24 \text{ and } 26
\end{align*}

\textbf{Answer:} \(Integers: 24 and 26\)

\end{tcolorbox}

\end{exercise}

\begin{exercise}[]\protect\hypertarget{exr-u6-25-2}{}\label{exr-u6-25-2}

~

The product of two consecutive positive even integers is 48.

\textbf{(i)} Form a quadratic equation in \(n\).

\textbf{(ii)} Solve the equation exactly.

\textbf{(iii)} Find the two integers.

\begin{tcolorbox}[enhanced jigsaw, arc=.35mm, breakable, toptitle=1mm, title=\textcolor{quarto-callout-tip-color}{\faLightbulb}\hspace{0.5em}{🔍 View Solution}, colframe=quarto-callout-tip-color-frame, rightrule=.15mm, opacityback=0, titlerule=0mm, colback=white, coltitle=black, bottomtitle=1mm, toprule=.15mm, leftrule=.75mm, bottomrule=.15mm, colbacktitle=quarto-callout-tip-color!10!white, left=2mm, opacitybacktitle=0.6]

\begin{align*}
\text{(i) } n(n+2) = 48 \\ n^2 + 2n - 48 = 0 \\ \text{(ii) } n = \dfrac{-(2) \pm \sqrt{(2)^2 - 4(1)(-48)}}{2} \\ n = \dfrac{-2 \pm 14}{2} \\ n = 6 \text{ or } n = -8 \\ \text{Since integers are positive even, } n = 6 \\ \text{Integers: } 6 \text{ and } 8
\end{align*}

\textbf{Answer:} \(Integers: 6 and 8\)

\end{tcolorbox}

\end{exercise}

\begin{exercise}[]\protect\hypertarget{exr-u6-25-3}{}\label{exr-u6-25-3}

~

The product of two consecutive positive even integers is 528.

\textbf{(i)} Form a quadratic equation in \(n\).

\textbf{(ii)} Solve the equation exactly.

\textbf{(iii)} Find the two integers.

\begin{tcolorbox}[enhanced jigsaw, arc=.35mm, breakable, toptitle=1mm, title=\textcolor{quarto-callout-tip-color}{\faLightbulb}\hspace{0.5em}{🔍 View Solution}, colframe=quarto-callout-tip-color-frame, rightrule=.15mm, opacityback=0, titlerule=0mm, colback=white, coltitle=black, bottomtitle=1mm, toprule=.15mm, leftrule=.75mm, bottomrule=.15mm, colbacktitle=quarto-callout-tip-color!10!white, left=2mm, opacitybacktitle=0.6]

\begin{align*}
\text{(i) } n(n+2) = 528 \\ n^2 + 2n - 528 = 0 \\ \text{(ii) } n = \dfrac{-(2) \pm \sqrt{(2)^2 - 4(1)(-528)}}{2} \\ n = \dfrac{-2 \pm 46}{2} \\ n = 22 \text{ or } n = -24 \\ \text{Since integers are positive even, } n = 22 \\ \text{Integers: } 22 \text{ and } 24
\end{align*}

\textbf{Answer:} \(Integers: 22 and 24\)

\end{tcolorbox}

\end{exercise}

\begin{exercise}[]\protect\hypertarget{exr-u6-27}{}\label{exr-u6-27}

\emph{Form a quadratic equation with roots \(\alpha\) and \(\beta\).
Give your answer in the form \(y^2 + py + q = 0\).}

\textbf{a)} \(\displaystyle \alpha = -1,\; \beta = -2\)

\textbf{b)} \(\displaystyle \alpha = -2,\; \beta = 5\)

\textbf{c)} \(\displaystyle \alpha = 2,\; \beta = 3\)

\textbf{d)} \(\displaystyle \alpha = 1,\; \beta = -4\)

\begin{tcolorbox}[enhanced jigsaw, arc=.35mm, breakable, toptitle=1mm, title=\textcolor{quarto-callout-tip-color}{\faLightbulb}\hspace{0.5em}{🔍 View Solutions}, colframe=quarto-callout-tip-color-frame, rightrule=.15mm, opacityback=0, titlerule=0mm, colback=white, coltitle=black, bottomtitle=1mm, toprule=.15mm, leftrule=.75mm, bottomrule=.15mm, colbacktitle=quarto-callout-tip-color!10!white, left=2mm, opacitybacktitle=0.6]

\textbf{a)} \(\displaystyle \alpha = -1,\; \beta = -2\)

\begin{align*}
\text{Sum of roots: } \alpha + \beta = -1 + (-2) = -3 \\ \text{Product of roots: } \alpha\beta = (-1) \times (-2) = 2 \\ \text{Equation: } y^2 - (\alpha+\beta)y + \alpha\beta = 0 \\ y^2 + 3y + 2 = 0
\end{align*}

\textbf{b)} \(\displaystyle \alpha = -2,\; \beta = 5\)

\begin{align*}
\text{Sum of roots: } \alpha + \beta = -2 + (5) = 3 \\ \text{Product of roots: } \alpha\beta = (-2) \times (5) = -10 \\ \text{Equation: } y^2 - (\alpha+\beta)y + \alpha\beta = 0 \\ y^2 - 3y - 10 = 0
\end{align*}

\textbf{c)} \(\displaystyle \alpha = 2,\; \beta = 3\)

\begin{align*}
\text{Sum of roots: } \alpha + \beta = 2 + (3) = 5 \\ \text{Product of roots: } \alpha\beta = (2) \times (3) = 6 \\ \text{Equation: } x^2 - (\alpha+\beta)x + \alpha\beta = 0 \\ x^2 - 5x + 6 = 0
\end{align*}

\textbf{d)} \(\displaystyle \alpha = 1,\; \beta = -4\)

\begin{align*}
\text{Sum of roots: } \alpha + \beta = 1 + (-4) = -3 \\ \text{Product of roots: } \alpha\beta = (1) \times (-4) = -4 \\ \text{Equation: } m^2 - (\alpha+\beta)m + \alpha\beta = 0 \\ m^2 + 3m - 4 = 0
\end{align*}

\end{tcolorbox}

\end{exercise}

∑ × π ÷ √ ◆ ∞ ± ∫ ≠ Δ

\subsection{🔴 Challenging}

\begin{exercise}[]\protect\hypertarget{exr-u6-5}{}\label{exr-u6-5}

\emph{Factorise each expression.}

\textbf{a)} \(\displaystyle 4n^2 - 12n + 8\)

\textbf{b)} \(\displaystyle 2n^2 - 12n + 10\)

\textbf{c)} \(\displaystyle 4m^2 - 21m + 20\)

\textbf{d)} \(\displaystyle 3m^2 - 8m + 4\)

\begin{tcolorbox}[enhanced jigsaw, arc=.35mm, breakable, toptitle=1mm, title=\textcolor{quarto-callout-tip-color}{\faLightbulb}\hspace{0.5em}{🔍 View Solutions}, colframe=quarto-callout-tip-color-frame, rightrule=.15mm, opacityback=0, titlerule=0mm, colback=white, coltitle=black, bottomtitle=1mm, toprule=.15mm, leftrule=.75mm, bottomrule=.15mm, colbacktitle=quarto-callout-tip-color!10!white, left=2mm, opacitybacktitle=0.6]

\textbf{a)} \(\displaystyle 4n^2 - 12n + 8\)

\begin{align*}
ac = 4 \times 8 = 32 \\ \text{Find } p, q: p \times q = 32,\; p + q = -12 \\ p = -8,\; q = -4 \\ 4n^2 - 12n + 8 = (4n - 4)(n - 2)
\end{align*}

\textbf{b)} \(\displaystyle 2n^2 - 12n + 10\)

\begin{align*}
ac = 2 \times 10 = 20 \\ \text{Find } p, q: p \times q = 20,\; p + q = -12 \\ p = -10,\; q = -2 \\ 2n^2 - 12n + 10 = (2n - 2)(n - 5)
\end{align*}

\textbf{c)} \(\displaystyle 4m^2 - 21m + 20\)

\begin{align*}
ac = 4 \times 20 = 80 \\ \text{Find } p, q: p \times q = 80,\; p + q = -21 \\ p = -16,\; q = -5 \\ 4m^2 - 21m + 20 = (4m - 5)(m - 4)
\end{align*}

\textbf{d)} \(\displaystyle 3m^2 - 8m + 4\)

\begin{align*}
ac = 3 \times 4 = 12 \\ \text{Find } p, q: p \times q = 12,\; p + q = -8 \\ p = -6,\; q = -2 \\ 3m^2 - 8m + 4 = (3m - 2)(m - 2)
\end{align*}

\end{tcolorbox}

\end{exercise}

\begin{exercise}[]\protect\hypertarget{exr-u6-12}{}\label{exr-u6-12}

\emph{Solve using the quadratic formula. Leave answers in surd form.}

\textbf{a)} \(\displaystyle y^2 - 2y - 2 = 0\)

\textbf{b)} \(\displaystyle n^2 - 7n + 5 = 0\)

\textbf{c)} \(\displaystyle m^2 - 7m + 3 = 0\)

\textbf{d)} \(\displaystyle x^2 - 7x + 2 = 0\)

\begin{tcolorbox}[enhanced jigsaw, arc=.35mm, breakable, toptitle=1mm, title=\textcolor{quarto-callout-tip-color}{\faLightbulb}\hspace{0.5em}{🔍 View Solutions}, colframe=quarto-callout-tip-color-frame, rightrule=.15mm, opacityback=0, titlerule=0mm, colback=white, coltitle=black, bottomtitle=1mm, toprule=.15mm, leftrule=.75mm, bottomrule=.15mm, colbacktitle=quarto-callout-tip-color!10!white, left=2mm, opacitybacktitle=0.6]

\textbf{a)} \(\displaystyle y^2 - 2y - 2 = 0\)

\begin{align*}
y = \dfrac{-(-2) \pm \sqrt{(-2)^2 - 4(1)(-2)}}{2} \\ y = \dfrac{2 \pm \sqrt{12}}{2} \\ y = \dfrac{2 \pm 2\sqrt{3}}{2} \\ y = 1 \pm \sqrt{3}
\end{align*}

\textbf{b)} \(\displaystyle n^2 - 7n + 5 = 0\)

\begin{align*}
n = \dfrac{-(-7) \pm \sqrt{(-7)^2 - 4(1)(5)}}{2} \\ n = \dfrac{7 \pm \sqrt{29}}{2}
\end{align*}

\textbf{c)} \(\displaystyle m^2 - 7m + 3 = 0\)

\begin{align*}
m = \dfrac{-(-7) \pm \sqrt{(-7)^2 - 4(1)(3)}}{2} \\ m = \dfrac{7 \pm \sqrt{37}}{2}
\end{align*}

\textbf{d)} \(\displaystyle x^2 - 7x + 2 = 0\)

\begin{align*}
x = \dfrac{-(-7) \pm \sqrt{(-7)^2 - 4(1)(2)}}{2} \\ x = \dfrac{7 \pm \sqrt{41}}{2}
\end{align*}

\end{tcolorbox}

\end{exercise}

\begin{exercise}[]\protect\hypertarget{exr-u6-13}{}\label{exr-u6-13}

\emph{Use the discriminant to determine the nature of roots.}

\textbf{a)} \(\displaystyle 3x^2 - x + 7 = 0\)

\textbf{b)} \(\displaystyle 3x^2 - x + 1 = 0\)

\textbf{c)} \(\displaystyle 2n^2 + 4n + 6 = 0\)

\textbf{d)} \(\displaystyle n^2 + n + 4 = 0\)

\begin{tcolorbox}[enhanced jigsaw, arc=.35mm, breakable, toptitle=1mm, title=\textcolor{quarto-callout-tip-color}{\faLightbulb}\hspace{0.5em}{🔍 View Solutions}, colframe=quarto-callout-tip-color-frame, rightrule=.15mm, opacityback=0, titlerule=0mm, colback=white, coltitle=black, bottomtitle=1mm, toprule=.15mm, leftrule=.75mm, bottomrule=.15mm, colbacktitle=quarto-callout-tip-color!10!white, left=2mm, opacitybacktitle=0.6]

\textbf{a)} \(\displaystyle 3x^2 - x + 7 = 0\)

\begin{align*}
x = \dfrac{-(-1) \pm \sqrt{(-1)^2 - 4(3)(7)}}{6} \\ x = \dfrac{1 \pm 0}{6} \\ \text{Discriminant } = -83 < 0 \\ \text{No real solutions}
\end{align*}

\textbf{b)} \(\displaystyle 3x^2 - x + 1 = 0\)

\begin{align*}
x = \dfrac{-(-1) \pm \sqrt{(-1)^2 - 4(3)(1)}}{6} \\ x = \dfrac{1 \pm 0}{6} \\ \text{Discriminant } = -11 < 0 \\ \text{No real solutions}
\end{align*}

\textbf{c)} \(\displaystyle 2n^2 + 4n + 6 = 0\)

\begin{align*}
n = \dfrac{-(4) \pm \sqrt{(4)^2 - 4(2)(6)}}{4} \\ n = \dfrac{-4 \pm 0}{4} \\ \text{Discriminant } = -32 < 0 \\ \text{No real solutions}
\end{align*}

\textbf{d)} \(\displaystyle n^2 + n + 4 = 0\)

\begin{align*}
n = \dfrac{-(1) \pm \sqrt{(1)^2 - 4(1)(4)}}{2} \\ n = \dfrac{-1 \pm 0}{2} \\ \text{Discriminant } = -15 < 0 \\ \text{No real solutions}
\end{align*}

\end{tcolorbox}

\end{exercise}

\begin{exercise}[]\protect\hypertarget{exr-u6-15}{}\label{exr-u6-15}

\emph{Write in the form \(a(x+p)^2 + q\).}

\textbf{a)} \(\displaystyle 3m^2 + m - 4\)

\textbf{b)} \(\displaystyle -2x^2 + x + 3\)

\textbf{c)} \(\displaystyle -2x^2 + x + 4\)

\textbf{d)} \(\displaystyle 2y^2 - 5y + 5\)

\begin{tcolorbox}[enhanced jigsaw, arc=.35mm, breakable, toptitle=1mm, title=\textcolor{quarto-callout-tip-color}{\faLightbulb}\hspace{0.5em}{🔍 View Solutions}, colframe=quarto-callout-tip-color-frame, rightrule=.15mm, opacityback=0, titlerule=0mm, colback=white, coltitle=black, bottomtitle=1mm, toprule=.15mm, leftrule=.75mm, bottomrule=.15mm, colbacktitle=quarto-callout-tip-color!10!white, left=2mm, opacitybacktitle=0.6]

\textbf{a)} \(\displaystyle 3m^2 + m - 4\)

\begin{align*}
3\left[m^2 + \dfrac{1}{3}m\right] - 4 \\ 3\left(m + \dfrac{1}{6}\right)^2 - \dfrac{1}{12} - 4 \\ 3(m + \dfrac{1}{6})^2 - \dfrac{49}{12}
\end{align*}

\textbf{b)} \(\displaystyle -2x^2 + x + 3\)

\begin{align*}
-2\left[x^2 + \dfrac{1}{-2}x\right] + 3 \\ -2\left(x + \dfrac{1}{-4}\right)^2 + \dfrac{1}{8} + 3 \\ -2(x - \dfrac{1}{4})^2 + \dfrac{25}{8}
\end{align*}

\textbf{c)} \(\displaystyle -2x^2 + x + 4\)

\begin{align*}
-2\left[x^2 + \dfrac{1}{-2}x\right] + 4 \\ -2\left(x + \dfrac{1}{-4}\right)^2 + \dfrac{1}{8} + 4 \\ -2(x - \dfrac{1}{4})^2 + \dfrac{33}{8}
\end{align*}

\textbf{d)} \(\displaystyle 2y^2 - 5y + 5\)

\begin{align*}
2\left[y^2 + \dfrac{-5}{2}y\right] + 5 \\ 2\left(y + \dfrac{-5}{4}\right)^2 - \dfrac{25}{8} + 5 \\ 2(y - \dfrac{5}{4})^2 + \dfrac{15}{8}
\end{align*}

\end{tcolorbox}

\end{exercise}

\begin{exercise}[]\protect\hypertarget{exr-u6-16}{}\label{exr-u6-16}

\emph{Solve by completing the square.}

\textbf{a)} \(\displaystyle n^2 + 2n - 2 = 0\)

\textbf{b)} \(\displaystyle n^2 + 8n + 7 = 0\)

\textbf{c)} \(\displaystyle n^2 + 2n - 4 = 0\)

\textbf{d)} \(\displaystyle n^2 - 4n - 7 = 0\)

\begin{tcolorbox}[enhanced jigsaw, arc=.35mm, breakable, toptitle=1mm, title=\textcolor{quarto-callout-tip-color}{\faLightbulb}\hspace{0.5em}{🔍 View Solutions}, colframe=quarto-callout-tip-color-frame, rightrule=.15mm, opacityback=0, titlerule=0mm, colback=white, coltitle=black, bottomtitle=1mm, toprule=.15mm, leftrule=.75mm, bottomrule=.15mm, colbacktitle=quarto-callout-tip-color!10!white, left=2mm, opacitybacktitle=0.6]

\textbf{a)} \(\displaystyle n^2 + 2n - 2 = 0\)

\begin{align*}
n^2 + 2n - 2 = 0 \\ (n + 1)^2 - 3 = 0 \\ (n + 1)^2 = 3 \\ n + 1 = \pm \sqrt{3} \\ n = -1 + \sqrt{3} \text{ or } n = -1 - \sqrt{3}
\end{align*}

\textbf{b)} \(\displaystyle n^2 + 8n + 7 = 0\)

\begin{align*}
n^2 + 8n + 7 = 0 \\ (n + 4)^2 - 9 = 0 \\ (n + 4)^2 = 9 \\ n + 4 = \pm 3 \\ n = -1 \text{ or } n = -7
\end{align*}

\textbf{c)} \(\displaystyle n^2 + 2n - 4 = 0\)

\begin{align*}
n^2 + 2n - 4 = 0 \\ (n + 1)^2 - 5 = 0 \\ (n + 1)^2 = 5 \\ n + 1 = \pm \sqrt{5} \\ n = -1 + \sqrt{5} \text{ or } n = -1 - \sqrt{5}
\end{align*}

\textbf{d)} \(\displaystyle n^2 - 4n - 7 = 0\)

\begin{align*}
n^2 - 4n - 7 = 0 \\ (n - 2)^2 - 11 = 0 \\ (n - 2)^2 = 11 \\ n - 2 = \pm \sqrt{11} \\ n = 2 + \sqrt{11} \text{ or } n = 2 - \sqrt{11}
\end{align*}

\end{tcolorbox}

\end{exercise}

\begin{exercise}[]\protect\hypertarget{exr-u6-17}{}\label{exr-u6-17}

\emph{Solve by completing the square.}

\textbf{a)} \(\displaystyle 2y^2 - 4y = 0\)

\textbf{b)} \(\displaystyle 2x^2 + x - 1 = 0\)

\textbf{c)} \(\displaystyle 2m^2 - 5m - 1 = 0\)

\begin{tcolorbox}[enhanced jigsaw, arc=.35mm, breakable, toptitle=1mm, title=\textcolor{quarto-callout-tip-color}{\faLightbulb}\hspace{0.5em}{🔍 View Solutions}, colframe=quarto-callout-tip-color-frame, rightrule=.15mm, opacityback=0, titlerule=0mm, colback=white, coltitle=black, bottomtitle=1mm, toprule=.15mm, leftrule=.75mm, bottomrule=.15mm, colbacktitle=quarto-callout-tip-color!10!white, left=2mm, opacitybacktitle=0.6]

\textbf{a)} \(\displaystyle 2y^2 - 4y = 0\)

\begin{align*}
\text{Divide by } 2: y^2 -2y +0 = 0 \\ \text{Complete the square:} \\ \left(y - 1\right)^2 = 1 \\ y - 1 = \pm 1 \\ y = 2 \text{ or } y = 0
\end{align*}

\textbf{b)} \(\displaystyle 2x^2 + x - 1 = 0\)

\begin{align*}
\text{Divide by } 2: x^2 +\dfrac{1}{2}x -\dfrac{1}{2} = 0 \\ \text{Complete the square:} \\ \left(x + \dfrac{1}{4}\right)^2 = \dfrac{9}{16} \\ x + \dfrac{1}{4} = \pm \dfrac{3}{4} \\ x = \dfrac{1}{2} \text{ or } x = -1
\end{align*}

\textbf{c)} \(\displaystyle 2m^2 - 5m - 1 = 0\)

\begin{align*}
\text{Divide by } 2: m^2 -\dfrac{5}{2}m -\dfrac{1}{2} = 0 \\ \text{Complete the square:} \\ \left(m - \dfrac{5}{4}\right)^2 = \dfrac{33}{16} \\ m - \dfrac{5}{4} = \pm \dfrac{\sqrt{33}}{4} \\ m = \dfrac{5}{4} + \dfrac{\sqrt{33}}{4} \text{ or } m = \dfrac{5}{4} - \dfrac{\sqrt{33}}{4}
\end{align*}

\end{tcolorbox}

\end{exercise}

\begin{exercise}[]\protect\hypertarget{exr-u6-20-1}{}\label{exr-u6-20-1}

~

A ball is thrown upward with initial velocity 15 m/s. Its height (m)
after \(t\) seconds is \(h = -5t^2 + 15t + 15\).

\textbf{(i)} What does the constant 15 represent in this context?

\textbf{(ii)} Find when the ball hits the ground. Give exact answers in
surd form.

\textbf{(iii)} Find the time to 2 d.p. and interpret your answer.

\begin{tcolorbox}[enhanced jigsaw, arc=.35mm, breakable, toptitle=1mm, title=\textcolor{quarto-callout-tip-color}{\faLightbulb}\hspace{0.5em}{🔍 View Solution}, colframe=quarto-callout-tip-color-frame, rightrule=.15mm, opacityback=0, titlerule=0mm, colback=white, coltitle=black, bottomtitle=1mm, toprule=.15mm, leftrule=.75mm, bottomrule=.15mm, colbacktitle=quarto-callout-tip-color!10!white, left=2mm, opacitybacktitle=0.6]

\begin{align*}
\text{(i) The constant term 15 is the initial height (height at } t=0\text{).} \\ \text{(ii) Set } h=0: \\ -5t^2 + 15t + 15 = 0 \\ 5t^2 - 15t - 15 = 0 \\ t = \dfrac{-(-15) \pm \sqrt{(-15)^2 - 4(5)(-15)}}{10} \\ t = \dfrac{15 \pm \sqrt{525}}{10} \\ t = \dfrac{15 \pm 5\sqrt{21}}{10} \\ t = \dfrac{3 \pm \sqrt{21}}{2} \\ \text{(iii) Taking positive root: } t \approx 3.79 \text{ s} \\ \text{This is the time when the ball hits the ground.}
\end{align*}

\textbf{Answer:} \(t \approx 3.79 \text{ s (time to hit ground)}\)

\end{tcolorbox}

\end{exercise}

\begin{exercise}[]\protect\hypertarget{exr-u6-20-2}{}\label{exr-u6-20-2}

~

A ball is thrown upward with initial velocity 15 m/s. Its height (m)
after \(t\) seconds is \(h = -5t^2 + 15t + 5\).

\textbf{(i)} What does the constant 5 represent in this context?

\textbf{(ii)} Find when the ball hits the ground. Give exact answers in
surd form.

\textbf{(iii)} Find the time to 2 d.p. and interpret your answer.

\begin{tcolorbox}[enhanced jigsaw, arc=.35mm, breakable, toptitle=1mm, title=\textcolor{quarto-callout-tip-color}{\faLightbulb}\hspace{0.5em}{🔍 View Solution}, colframe=quarto-callout-tip-color-frame, rightrule=.15mm, opacityback=0, titlerule=0mm, colback=white, coltitle=black, bottomtitle=1mm, toprule=.15mm, leftrule=.75mm, bottomrule=.15mm, colbacktitle=quarto-callout-tip-color!10!white, left=2mm, opacitybacktitle=0.6]

\begin{align*}
\text{(i) The constant term 5 is the initial height (height at } t=0\text{).} \\ \text{(ii) Set } h=0: \\ -5t^2 + 15t + 5 = 0 \\ 5t^2 - 15t - 5 = 0 \\ t = \dfrac{-(-15) \pm \sqrt{(-15)^2 - 4(5)(-5)}}{10} \\ t = \dfrac{15 \pm \sqrt{325}}{10} \\ t = \dfrac{15 \pm 5\sqrt{13}}{10} \\ t = \dfrac{3 \pm \sqrt{13}}{2} \\ \text{(iii) Taking positive root: } t \approx 3.30 \text{ s} \\ \text{This is the time when the ball hits the ground.}
\end{align*}

\textbf{Answer:} \(t \approx 3.30 \text{ s (time to hit ground)}\)

\end{tcolorbox}

\end{exercise}

\begin{exercise}[]\protect\hypertarget{exr-u6-22-1}{}\label{exr-u6-22-1}

~

A right-angled triangle has legs \(x\) cm and \((x + 2)\) cm. The
hypotenuse is 10 cm.

\textbf{(i)} Use Pythagoras' theorem to form a quadratic equation.

\textbf{(ii)} Solve the equation exactly.

\textbf{(iii)} Find the lengths of the two legs.

\begin{tcolorbox}[enhanced jigsaw, arc=.35mm, breakable, toptitle=1mm, title=\textcolor{quarto-callout-tip-color}{\faLightbulb}\hspace{0.5em}{🔍 View Solution}, colframe=quarto-callout-tip-color-frame, rightrule=.15mm, opacityback=0, titlerule=0mm, colback=white, coltitle=black, bottomtitle=1mm, toprule=.15mm, leftrule=.75mm, bottomrule=.15mm, colbacktitle=quarto-callout-tip-color!10!white, left=2mm, opacitybacktitle=0.6]

\begin{align*}
\text{(i) Pythagoras: } x^2 + (x+2)^2 = 10^2 \\ x^2 + x^2 + 4x + 4 = 100 \\ 2x^2 + 4x - 96 = 0 \\ \text{(ii) } x = \dfrac{-(4) \pm \sqrt{(4)^2 - 4(2)(-96)}}{4} \\ x = \dfrac{-4 \pm 28}{4} \\ x = 6 \text{ or } x = -8 \\ \text{(iii) } x = 6 \text{ cm}
\end{align*}

\textbf{Answer:}
\(x = 6 \text{ cm, legs: } 6 \text{ and } 8 \text{ cm}\)

\end{tcolorbox}

\end{exercise}

\begin{exercise}[]\protect\hypertarget{exr-u6-22-2}{}\label{exr-u6-22-2}

~

A right-angled triangle has legs \(x\) cm and \((x + 4)\) cm. The
hypotenuse is 20 cm.

\textbf{(i)} Use Pythagoras' theorem to form a quadratic equation.

\textbf{(ii)} Solve the equation exactly.

\textbf{(iii)} Find the lengths of the two legs.

\begin{tcolorbox}[enhanced jigsaw, arc=.35mm, breakable, toptitle=1mm, title=\textcolor{quarto-callout-tip-color}{\faLightbulb}\hspace{0.5em}{🔍 View Solution}, colframe=quarto-callout-tip-color-frame, rightrule=.15mm, opacityback=0, titlerule=0mm, colback=white, coltitle=black, bottomtitle=1mm, toprule=.15mm, leftrule=.75mm, bottomrule=.15mm, colbacktitle=quarto-callout-tip-color!10!white, left=2mm, opacitybacktitle=0.6]

\begin{align*}
\text{(i) Pythagoras: } x^2 + (x+4)^2 = 20^2 \\ x^2 + x^2 + 8x + 16 = 400 \\ 2x^2 + 8x - 384 = 0 \\ \text{(ii) } x = \dfrac{-(8) \pm \sqrt{(8)^2 - 4(2)(-384)}}{4} \\ x = \dfrac{-8 \pm 56}{4} \\ x = 12 \text{ or } x = -16 \\ \text{(iii) } x = 12 \text{ cm}
\end{align*}

\textbf{Answer:}
\(x = 12 \text{ cm, legs: } 12 \text{ and } 16 \text{ cm}\)

\end{tcolorbox}

\end{exercise}

\begin{exercise}[]\protect\hypertarget{exr-u6-24}{}\label{exr-u6-24}

\emph{Find the value(s) of \(k\) for which the equation has real roots.}

\textbf{a)} \(\displaystyle n^2 - 6n + k = 0\)

\textbf{b)} \(\displaystyle m^2 + km + 9 = 0\)

\textbf{c)} \(\displaystyle x^2 + 4x + k = 0\)

\textbf{d)} \(\displaystyle y^2 + 6y + k = 0\)

\begin{tcolorbox}[enhanced jigsaw, arc=.35mm, breakable, toptitle=1mm, title=\textcolor{quarto-callout-tip-color}{\faLightbulb}\hspace{0.5em}{🔍 View Solutions}, colframe=quarto-callout-tip-color-frame, rightrule=.15mm, opacityback=0, titlerule=0mm, colback=white, coltitle=black, bottomtitle=1mm, toprule=.15mm, leftrule=.75mm, bottomrule=.15mm, colbacktitle=quarto-callout-tip-color!10!white, left=2mm, opacitybacktitle=0.6]

\textbf{a)} \(\displaystyle n^2 - 6n + k = 0\)

\begin{align*}
\text{Given: } n^2 - 6n + k = 0 \\ \text{For real roots: } \Delta \geq 0 \\ (-6)^2 - 4k \geq 0 \\ 36 \geq 4k \\ k \leq 9
\end{align*}

\textbf{b)} \(\displaystyle m^2 + km + 9 = 0\)

\begin{align*}
\text{Given: } m^2 + km + 9 = 0 \\ \text{For real roots: } \Delta \geq 0 \\ k^2 - 4(9) \geq 0 \\ k^2 \geq 36 \\ k \leq -6 \text{ or } k \geq 6
\end{align*}

\textbf{c)} \(\displaystyle x^2 + 4x + k = 0\)

\begin{align*}
\text{Given: } x^2 + 4x + k = 0 \\ \text{For real roots: } \Delta \geq 0 \\ (4)^2 - 4k \geq 0 \\ 16 \geq 4k \\ k \leq 4
\end{align*}

\textbf{d)} \(\displaystyle y^2 + 6y + k = 0\)

\begin{align*}
\text{Given: } y^2 + 6y + k = 0 \\ \text{For real roots: } \Delta \geq 0 \\ (6)^2 - 4k \geq 0 \\ 36 \geq 4k \\ k \leq 9
\end{align*}

\end{tcolorbox}

\end{exercise}

\begin{exercise}[]\protect\hypertarget{exr-u6-26-1}{}\label{exr-u6-26-1}

~

Two consecutive positive odd integers: the product of the larger and
twice the smaller is 30.

\textbf{(i)} Form a quadratic equation in \(n\).

\textbf{(ii)} Solve the equation exactly.

\textbf{(iii)} Find the two integers.

\begin{tcolorbox}[enhanced jigsaw, arc=.35mm, breakable, toptitle=1mm, title=\textcolor{quarto-callout-tip-color}{\faLightbulb}\hspace{0.5em}{🔍 View Solution}, colframe=quarto-callout-tip-color-frame, rightrule=.15mm, opacityback=0, titlerule=0mm, colback=white, coltitle=black, bottomtitle=1mm, toprule=.15mm, leftrule=.75mm, bottomrule=.15mm, colbacktitle=quarto-callout-tip-color!10!white, left=2mm, opacitybacktitle=0.6]

\begin{align*}
text{(i) } (n+2)(2n) = 30 \\ 2n^2 + 4n - 30 = 0 \\ text{(ii) } n = \dfrac{-(4) \pm \sqrt{(4)^2 - 4(2)(-30)}}{4} \\ n = \dfrac{-4 \pm 16}{4} \\ n = 3 \text{ or } n = -5 \\ text{Since integers are positive odd, } n = 3 \\ text{Integers: } 3 \\ text{ and } 5
\end{align*}

\textbf{Answer:} \(Integers: 3 and 5\)

\end{tcolorbox}

\end{exercise}

\begin{exercise}[]\protect\hypertarget{exr-u6-26-2}{}\label{exr-u6-26-2}

~

Two consecutive positive odd integers: the product of the larger and
twice the smaller is 646.

\textbf{(i)} Form a quadratic equation in \(n\).

\textbf{(ii)} Solve the equation exactly.

\textbf{(iii)} Find the two integers.

\begin{tcolorbox}[enhanced jigsaw, arc=.35mm, breakable, toptitle=1mm, title=\textcolor{quarto-callout-tip-color}{\faLightbulb}\hspace{0.5em}{🔍 View Solution}, colframe=quarto-callout-tip-color-frame, rightrule=.15mm, opacityback=0, titlerule=0mm, colback=white, coltitle=black, bottomtitle=1mm, toprule=.15mm, leftrule=.75mm, bottomrule=.15mm, colbacktitle=quarto-callout-tip-color!10!white, left=2mm, opacitybacktitle=0.6]

\begin{align*}
text{(i) } (n+2)(2n) = 646 \\ 2n^2 + 4n - 646 = 0 \\ text{(ii) } n = \dfrac{-(4) \pm \sqrt{(4)^2 - 4(2)(-646)}}{4} \\ n = \dfrac{-4 \pm 72}{4} \\ n = 17 \text{ or } n = -19 \\ text{Since integers are positive odd, } n = 17 \\ text{Integers: } 17 \\ text{ and } 19
\end{align*}

\textbf{Answer:} \(Integers: 17 and 19\)

\end{tcolorbox}

\end{exercise}

\begin{exercise}[]\protect\hypertarget{exr-u6-26-3}{}\label{exr-u6-26-3}

~

The product of two consecutive positive odd integers is 15.

\textbf{(i)} Form a quadratic equation in \(n\).

\textbf{(ii)} Solve the equation exactly.

\textbf{(iii)} Find the two integers.

\begin{tcolorbox}[enhanced jigsaw, arc=.35mm, breakable, toptitle=1mm, title=\textcolor{quarto-callout-tip-color}{\faLightbulb}\hspace{0.5em}{🔍 View Solution}, colframe=quarto-callout-tip-color-frame, rightrule=.15mm, opacityback=0, titlerule=0mm, colback=white, coltitle=black, bottomtitle=1mm, toprule=.15mm, leftrule=.75mm, bottomrule=.15mm, colbacktitle=quarto-callout-tip-color!10!white, left=2mm, opacitybacktitle=0.6]

\begin{align*}
text{(i) } n(n+2) = 15 \\ n^2 + 2n - 15 = 0 \\ text{(ii) } n = \dfrac{-(2) \pm \sqrt{(2)^2 - 4(1)(-15)}}{2} \\ n = \dfrac{-2 \pm 8}{2} \\ n = 3 \text{ or } n = -5 \\ text{Since integers are positive odd, } n = 3 \\ text{Integers: } 3 \\ text{ and } 5
\end{align*}

\textbf{Answer:} \(Integers: 3 and 5\)

\end{tcolorbox}

\end{exercise}

\begin{exercise}[]\protect\hypertarget{exr-u6-28}{}\label{exr-u6-28}

\emph{The equations below have roots \(\alpha\) and \(\beta\). Use
Vieta's formulas to find the given expression.}

\textbf{a)} \(\displaystyle 2m^2 + 6m - 5 = 0\)

\textbf{b)} \(\displaystyle x^2 + 6x + 6 = 0\)

\textbf{c)} \(\displaystyle x^2 + 2x - 3 = 0\)

\textbf{d)} \(\displaystyle y^2 + y - 2 = 0\)

\begin{tcolorbox}[enhanced jigsaw, arc=.35mm, breakable, toptitle=1mm, title=\textcolor{quarto-callout-tip-color}{\faLightbulb}\hspace{0.5em}{🔍 View Solutions}, colframe=quarto-callout-tip-color-frame, rightrule=.15mm, opacityback=0, titlerule=0mm, colback=white, coltitle=black, bottomtitle=1mm, toprule=.15mm, leftrule=.75mm, bottomrule=.15mm, colbacktitle=quarto-callout-tip-color!10!white, left=2mm, opacitybacktitle=0.6]

\textbf{a)} \(\displaystyle 2m^2 + 6m - 5 = 0\)

\begin{align*}
\text{From } 2m^2 + 6m - 5 = 0: \\ \alpha + \beta = \dfrac{-b}{a} = \dfrac{-6}{2} = -3 \\ \alpha\beta = \dfrac{c}{a} = \dfrac{-5}{2} = -\dfrac{5}{2} \\ \dfrac{1}{\alpha} + \dfrac{1}{\beta} = \dfrac{\alpha+\beta}{\alpha\beta} = \dfrac{-3}{-\dfrac{5}{2}} = \dfrac{6}{5}
\end{align*}

\textbf{b)} \(\displaystyle x^2 + 6x + 6 = 0\)

\begin{align*}
\text{From } x^2 + 6x + 6 = 0: \\ \alpha + \beta = \dfrac{-b}{a} = \dfrac{-6}{1} = -6 \\ \alpha\beta = \dfrac{c}{a} = \dfrac{6}{1} = 6 \\ \dfrac{1}{\alpha} + \dfrac{1}{\beta} = \dfrac{\alpha+\beta}{\alpha\beta} = \dfrac{-6}{6} = -1
\end{align*}

\textbf{c)} \(\displaystyle x^2 + 2x - 3 = 0\)

\begin{align*}
\text{From } x^2 + 2x - 3 = 0: \\ \alpha + \beta = \dfrac{-b}{a} = \dfrac{-2}{1} = -2 \\ \alpha\beta = \dfrac{c}{a} = \dfrac{-3}{1} = -3 \\ (\alpha - \beta)^2 = (\alpha+\beta)^2 - 4\alpha\beta = (-2)^2 - 4(-3) = 16
\end{align*}

\textbf{d)} \(\displaystyle y^2 + y - 2 = 0\)

\begin{align*}
\text{From } y^2 + y - 2 = 0: \\ \alpha + \beta = \dfrac{-b}{a} = \dfrac{-1}{1} = -1 \\ \alpha\beta = \dfrac{c}{a} = \dfrac{-2}{1} = -2 \\ \alpha^2 + \beta^2 = (\alpha+\beta)^2 - 2\alpha\beta = (-1)^2 - 2(-2) = 5
\end{align*}

\end{tcolorbox}

\end{exercise}

\begin{exercise}[]\protect\hypertarget{exr-u6-29-1}{}\label{exr-u6-29-1}

~

x + 5x - 22x - 1

Shown is a right-angled triangle. All sides are in cm.

\textbf{(a)} Show that \(x^2 - 5x - 14 = 0\).

\textbf{(b)} Find \(x\).

\textbf{(c)} Find the area of the triangle.

\begin{tcolorbox}[enhanced jigsaw, arc=.35mm, breakable, toptitle=1mm, title=\textcolor{quarto-callout-tip-color}{\faLightbulb}\hspace{0.5em}{🔍 View Solution}, colframe=quarto-callout-tip-color-frame, rightrule=.15mm, opacityback=0, titlerule=0mm, colback=white, coltitle=black, bottomtitle=1mm, toprule=.15mm, leftrule=.75mm, bottomrule=.15mm, colbacktitle=quarto-callout-tip-color!10!white, left=2mm, opacitybacktitle=0.6]

\begin{align*}
\text{(a) Pythagoras: } (x + 5)^2 + (x - 2)^2 = (2x - 1)^2 \\ x^2 -5x -14 = 0 \\ x^2 - 5x - 14 = 0 \\ \text{(b) } x = \dfrac{-(-5) \pm \sqrt{(-5)^2 - 4(1)(-14)}}{2} \\ x = \dfrac{5 \pm 9}{2} \\ x = 7 \text{ or } x = -2 \\ \text{Taking positive root: } x = 7 \\ \text{(c) Legs: } 12 \text{ cm and } 5 \text{ cm} \\ \text{Area } = \tfrac{1}{2} \times 12 \times 5 = 30 \text{ cm}^2
\end{align*}

\textbf{Answer:} \(x = 7,\; \text{Area} = 30 \text{ cm}^2\)

\end{tcolorbox}

\end{exercise}

\begin{exercise}[]\protect\hypertarget{exr-u6-29-2}{}\label{exr-u6-29-2}

~

x + 5x - 22x - 1

Shown is a right-angled triangle. All sides are in cm.

\textbf{(a)} Show that \(x^2 - 5x - 14 = 0\).

\textbf{(b)} Find \(x\).

\textbf{(c)} Find the area of the triangle.

\begin{tcolorbox}[enhanced jigsaw, arc=.35mm, breakable, toptitle=1mm, title=\textcolor{quarto-callout-tip-color}{\faLightbulb}\hspace{0.5em}{🔍 View Solution}, colframe=quarto-callout-tip-color-frame, rightrule=.15mm, opacityback=0, titlerule=0mm, colback=white, coltitle=black, bottomtitle=1mm, toprule=.15mm, leftrule=.75mm, bottomrule=.15mm, colbacktitle=quarto-callout-tip-color!10!white, left=2mm, opacitybacktitle=0.6]

\begin{align*}
\text{(a) Pythagoras: } (x + 5)^2 + (x - 2)^2 = (2x - 1)^2 \\ x^2 -5x -14 = 0 \\ x^2 - 5x - 14 = 0 \\ \text{(b) } x = \dfrac{-(-5) \pm \sqrt{(-5)^2 - 4(1)(-14)}}{2} \\ x = \dfrac{5 \pm 9}{2} \\ x = 7 \text{ or } x = -2 \\ \text{Taking positive root: } x = 7 \\ \text{(c) Legs: } 12 \text{ cm and } 5 \text{ cm} \\ \text{Area } = \tfrac{1}{2} \times 12 \times 5 = 30 \text{ cm}^2
\end{align*}

\textbf{Answer:} \(x = 7,\; \text{Area} = 30 \text{ cm}^2\)

\end{tcolorbox}

\end{exercise}

\begin{exercise}[]\protect\hypertarget{exr-u6-30-1}{}\label{exr-u6-30-1}

~

x+33x−224

The area of the L-shaped region is 234 cm². The big rectangle is
\((3x-2)\) cm by \((x+3)\) cm with a \(2\) cm \(\times\) \(4\) cm
rectangle removed.

\textbf{(a)} Show that \(3x^2 + 7x - 248 = 0\).

\textbf{(b)} Solve to find \(x\).

\begin{tcolorbox}[enhanced jigsaw, arc=.35mm, breakable, toptitle=1mm, title=\textcolor{quarto-callout-tip-color}{\faLightbulb}\hspace{0.5em}{🔍 View Solution}, colframe=quarto-callout-tip-color-frame, rightrule=.15mm, opacityback=0, titlerule=0mm, colback=white, coltitle=black, bottomtitle=1mm, toprule=.15mm, leftrule=.75mm, bottomrule=.15mm, colbacktitle=quarto-callout-tip-color!10!white, left=2mm, opacitybacktitle=0.6]

\begin{align*}
\text{(a) Area: } (3x-2)(x+3) - 8 = 234 \\ 3x^2 +7x - 6 - 8 = 234 \\ 3x^2 +7x - 248 = 0 \\ 3x^2 + 7x - 248 = 0 \\ \text{(b) } x = \dfrac{-(7) \pm \sqrt{(7)^2 - 4(3)(-248)}}{6} \\ x = \dfrac{-7 \pm 55}{6} \\ x = 8 \text{ or } x = -\dfrac{31}{3} \\ \text{Taking positive root: } x = 8
\end{align*}

\textbf{Answer:} \(x = 8\)

\end{tcolorbox}

\end{exercise}

\begin{exercise}[]\protect\hypertarget{exr-u6-30-2}{}\label{exr-u6-30-2}

~

x+73x−221

The area of the L-shaped region is 206 cm². The big rectangle is
\((3x-2)\) cm by \((x+7)\) cm with a \(2\) cm \(\times\) \(1\) cm
rectangle removed.

\textbf{(a)} Show that \(3x^2 + 19x - 222 = 0\).

\textbf{(b)} Solve to find \(x\).

\begin{tcolorbox}[enhanced jigsaw, arc=.35mm, breakable, toptitle=1mm, title=\textcolor{quarto-callout-tip-color}{\faLightbulb}\hspace{0.5em}{🔍 View Solution}, colframe=quarto-callout-tip-color-frame, rightrule=.15mm, opacityback=0, titlerule=0mm, colback=white, coltitle=black, bottomtitle=1mm, toprule=.15mm, leftrule=.75mm, bottomrule=.15mm, colbacktitle=quarto-callout-tip-color!10!white, left=2mm, opacitybacktitle=0.6]

\begin{align*}
\text{(a) Area: } (3x-2)(x+7) - 2 = 206 \\ 3x^2 +19x - 14 - 2 = 206 \\ 3x^2 +19x - 222 = 0 \\ 3x^2 + 19x - 222 = 0 \\ \text{(b) } x = \dfrac{-(19) \pm \sqrt{(19)^2 - 4(3)(-222)}}{6} \\ x = \dfrac{-19 \pm 55}{6} \\ x = 6 \text{ or } x = -\dfrac{37}{3} \\ \text{Taking positive root: } x = 6
\end{align*}

\textbf{Answer:} \(x = 6\)

\end{tcolorbox}

\end{exercise}

\begin{exercise}[]\protect\hypertarget{exr-u6-31-1}{}\label{exr-u6-31-1}

~

A cuboid has dimensions \(x\) cm, \(x\) cm and \((x+5)\) cm. Its surface
area is 544 cm².

\textbf{(a)} Show that \(3x^2 + 10x - 272 = 0\).

\textbf{(b)} Find \(x\).

\textbf{(c)} Find the volume of the cuboid.

\begin{tcolorbox}[enhanced jigsaw, arc=.35mm, breakable, toptitle=1mm, title=\textcolor{quarto-callout-tip-color}{\faLightbulb}\hspace{0.5em}{🔍 View Solution}, colframe=quarto-callout-tip-color-frame, rightrule=.15mm, opacityback=0, titlerule=0mm, colback=white, coltitle=black, bottomtitle=1mm, toprule=.15mm, leftrule=.75mm, bottomrule=.15mm, colbacktitle=quarto-callout-tip-color!10!white, left=2mm, opacitybacktitle=0.6]

\begin{align*}
\text{(a) Surface area: } 2x^2 + 4x(x+5) = 544 \\ 6x^2 + 20x - 544 = 0 \\ 3x^2 + 10x - 272 = 0 \\ \text{(b) } x = \dfrac{-(10) \pm \sqrt{(10)^2 - 4(3)(-272)}}{6} \\ x = \dfrac{-10 \pm 58}{6} \\ x = 8 \text{ or } x = -\dfrac{34}{3} \\ x = 8 \text{ cm} \\ \text{(c) Volume } = 8^2 \cdot (13) = 832 \text{ cm}^3
\end{align*}

\textbf{Answer:} \(x = 8,\; V = 832 \text{ cm}^3\)

\end{tcolorbox}

\end{exercise}

\begin{exercise}[]\protect\hypertarget{exr-u6-31-2}{}\label{exr-u6-31-2}

~

A cuboid has dimensions \(x\) cm, \(x\) cm and \((x+6)\) cm. Its surface
area is 840 cm².

\textbf{(a)} Show that \(x^2 + 4x - 140 = 0\).

\textbf{(b)} Find \(x\).

\textbf{(c)} Find the volume of the cuboid.

\begin{tcolorbox}[enhanced jigsaw, arc=.35mm, breakable, toptitle=1mm, title=\textcolor{quarto-callout-tip-color}{\faLightbulb}\hspace{0.5em}{🔍 View Solution}, colframe=quarto-callout-tip-color-frame, rightrule=.15mm, opacityback=0, titlerule=0mm, colback=white, coltitle=black, bottomtitle=1mm, toprule=.15mm, leftrule=.75mm, bottomrule=.15mm, colbacktitle=quarto-callout-tip-color!10!white, left=2mm, opacitybacktitle=0.6]

\begin{align*}
\text{(a) Surface area: } 2x^2 + 4x(x+6) = 840 \\ 6x^2 + 24x - 840 = 0 \\ x^2 + 4x - 140 = 0 \\ \text{(b) } x = \dfrac{-(4) \pm \sqrt{(4)^2 - 4(1)(-140)}}{2} \\ x = \dfrac{-4 \pm 24}{2} \\ x = 10 \text{ or } x = -14 \\ x = 10 \text{ cm} \\ \text{(c) Volume } = 10^2 \cdot (16) = 1600 \text{ cm}^3
\end{align*}

\textbf{Answer:} \(x = 10,\; V = 1600 \text{ cm}^3\)

\end{tcolorbox}

\end{exercise}

\begin{exercise}[]\protect\hypertarget{exr-u6-32-1}{}\label{exr-u6-32-1}

~

x2x−228

A cuboid has dimensions \(x\) cm, \((2x - 22)\) cm and \(8\) cm. Its
volume is 416 cm³.

\textbf{(a)} Show that \(2x^2 - 22x - 52 = 0\).

\textbf{(b)} Find \(x\).

\begin{tcolorbox}[enhanced jigsaw, arc=.35mm, breakable, toptitle=1mm, title=\textcolor{quarto-callout-tip-color}{\faLightbulb}\hspace{0.5em}{🔍 View Solution}, colframe=quarto-callout-tip-color-frame, rightrule=.15mm, opacityback=0, titlerule=0mm, colback=white, coltitle=black, bottomtitle=1mm, toprule=.15mm, leftrule=.75mm, bottomrule=.15mm, colbacktitle=quarto-callout-tip-color!10!white, left=2mm, opacitybacktitle=0.6]

\begin{align*}
\text{(a) Volume: } 8 \times x \times (2x-22) = 416 \\ 8x(2x-22) = 416 \\ 2x^2 - 22x - 52 = 0 \\ \text{(b) } x = \dfrac{-(-22) \pm \sqrt{(-22)^2 - 4(2)(-52)}}{4} \\ x = \dfrac{22 \pm 30}{4} \\ x = 13 \text{ or } x = -2 \\ \text{Taking positive root: } x = 13 \text{ cm}
\end{align*}

\textbf{Answer:} \(x = 13\)

\end{tcolorbox}

\end{exercise}

\begin{exercise}[]\protect\hypertarget{exr-u6-32-2}{}\label{exr-u6-32-2}

~

x2x−2215

A cuboid has dimensions \(x\) cm, \((2x - 22)\) cm and \(15\) cm. Its
volume is 3780 cm³.

\textbf{(a)} Show that \(2x^2 - 22x - 252 = 0\).

\textbf{(b)} Find \(x\).

\begin{tcolorbox}[enhanced jigsaw, arc=.35mm, breakable, toptitle=1mm, title=\textcolor{quarto-callout-tip-color}{\faLightbulb}\hspace{0.5em}{🔍 View Solution}, colframe=quarto-callout-tip-color-frame, rightrule=.15mm, opacityback=0, titlerule=0mm, colback=white, coltitle=black, bottomtitle=1mm, toprule=.15mm, leftrule=.75mm, bottomrule=.15mm, colbacktitle=quarto-callout-tip-color!10!white, left=2mm, opacitybacktitle=0.6]

\begin{align*}
\text{(a) Volume: } 15 \times x \times (2x-22) = 3780 \\ 15x(2x-22) = 3780 \\ 2x^2 - 22x - 252 = 0 \\ \text{(b) } x = \dfrac{-(-22) \pm \sqrt{(-22)^2 - 4(2)(-252)}}{4} \\ x = \dfrac{22 \pm 50}{4} \\ x = 18 \text{ or } x = -7 \\ \text{Taking positive root: } x = 18 \text{ cm}
\end{align*}

\textbf{Answer:} \(x = 18\)

\end{tcolorbox}

\end{exercise}

\begin{exercise}[]\protect\hypertarget{exr-u6-33}{}\label{exr-u6-33}

\emph{Solve. State any values of \(x\) that must be excluded.}

\textbf{a)} \(\displaystyle \dfrac{4}{x} + \dfrac{4}{x - 3} = 2\)

\textbf{b)} \(\displaystyle \dfrac{5}{x} + \dfrac{10}{x - 6} = 3\)

\textbf{c)} \(\displaystyle \dfrac{6}{x} + \dfrac{12}{x - 5} = 3\)

\textbf{d)} \(\displaystyle \dfrac{6}{x} + \dfrac{12}{x - 5} = 3\)

\begin{tcolorbox}[enhanced jigsaw, arc=.35mm, breakable, toptitle=1mm, title=\textcolor{quarto-callout-tip-color}{\faLightbulb}\hspace{0.5em}{🔍 View Solutions}, colframe=quarto-callout-tip-color-frame, rightrule=.15mm, opacityback=0, titlerule=0mm, colback=white, coltitle=black, bottomtitle=1mm, toprule=.15mm, leftrule=.75mm, bottomrule=.15mm, colbacktitle=quarto-callout-tip-color!10!white, left=2mm, opacitybacktitle=0.6]

\textbf{a)} \(\displaystyle \dfrac{4}{x} + \dfrac{4}{x - 3} = 2\)

\begin{align*}
\text{Multiply through by } x(x - 3): \\ 4(x - 3) + 4x = 2x(x - 3) \\ 2x^2 - 14x + 12 = 0 \\ (x - 6)(x - 1) = 0 \\ x = 6 \text{ or } x = 1
\end{align*}

\textbf{b)} \(\displaystyle \dfrac{5}{x} + \dfrac{10}{x - 6} = 3\)

\begin{align*}
\text{Multiply through by } x(x - 6): \\ 5(x - 6) + 10x = 3x(x - 6) \\ 3x^2 - 33x + 30 = 0 \\ (x - 10)(x - 1) = 0 \\ x = 10 \text{ or } x = 1
\end{align*}

\textbf{c)} \(\displaystyle \dfrac{6}{x} + \dfrac{12}{x - 5} = 3\)

\begin{align*}
\text{Multiply through by } x(x - 5): \\ 6(x - 5) + 12x = 3x(x - 5) \\ 3x^2 - 33x + 30 = 0 \\ (x - 10)(x - 1) = 0 \\ x = 10 \text{ or } x = 1
\end{align*}

\textbf{d)} \(\displaystyle \dfrac{6}{x} + \dfrac{12}{x - 5} = 3\)

\begin{align*}
\text{Multiply through by } x(x - 5): \\ 6(x - 5) + 12x = 3x(x - 5) \\ 3x^2 - 33x + 30 = 0 \\ (x - 10)(x - 1) = 0 \\ x = 10 \text{ or } x = 1
\end{align*}

\end{tcolorbox}

\end{exercise}

\begin{exercise}[]\protect\hypertarget{exr-u6-34}{}\label{exr-u6-34}

\emph{Solve. State any values of \(x\) that must be excluded.}

\textbf{a)} \(\displaystyle \dfrac{-4}{x + 5} + \dfrac{2}{x - 1} = -2\)

\textbf{b)} \(\displaystyle \dfrac{-3}{x + 1} - \dfrac{7}{x - 3} = 2\)

\textbf{c)} \(\displaystyle \dfrac{6}{x + 1} - \dfrac{1}{x - 1} = 1\)

\textbf{d)} \(\displaystyle \dfrac{8}{x + 4} + \dfrac{8}{x - 2} = 2\)

\begin{tcolorbox}[enhanced jigsaw, arc=.35mm, breakable, toptitle=1mm, title=\textcolor{quarto-callout-tip-color}{\faLightbulb}\hspace{0.5em}{🔍 View Solutions}, colframe=quarto-callout-tip-color-frame, rightrule=.15mm, opacityback=0, titlerule=0mm, colback=white, coltitle=black, bottomtitle=1mm, toprule=.15mm, leftrule=.75mm, bottomrule=.15mm, colbacktitle=quarto-callout-tip-color!10!white, left=2mm, opacitybacktitle=0.6]

\textbf{a)} \(\displaystyle \dfrac{-4}{x + 5} + \dfrac{2}{x - 1} = -2\)

\begin{align*}
\text{Multiply by } (x + 5)(x - 1): \\ -4(x - 1) + 2(x + 5) = -2(x + 5)(x - 1) \\ -2x^2 - 6x - 4 = 0 \\ (x + 2)(x + 1) = 0 \\ x = -2 \text{ or } x = -1
\end{align*}

\textbf{b)} \(\displaystyle \dfrac{-3}{x + 1} - \dfrac{7}{x - 3} = 2\)

\begin{align*}
\text{Multiply by } (x + 1)(x - 3): \\ -3(x - 3) - 7(x + 1) = 2(x + 1)(x - 3) \\ 2x^2 + 6x - 8 = 0 \\ (x + 4)(x -1) = 0 \\ x = -4 \text{ or } x = 1
\end{align*}

\textbf{c)} \(\displaystyle \dfrac{6}{x + 1} - \dfrac{1}{x - 1} = 1\)

\begin{align*}
\text{Multiply by } (x + 1)(x - 1): \\ 6(x - 1) - 1(x + 1) = 1(x + 1)(x - 1) \\ x^2 - 5x + 6 = 0 \\ (x -2)(x -3) = 0 \\ x = 2 \text{ or } x = 3
\end{align*}

\textbf{d)} \(\displaystyle \dfrac{8}{x + 4} + \dfrac{8}{x - 2} = 2\)

\begin{align*}
\text{Multiply by } (x + 4)(x - 2): \\ 8(x - 2) + 8(x + 4) = 2(x + 4)(x - 2) \\ 2x^2 - 12x - 32 = 0 \\ (x -8)(x + 2) = 0 \\ x = 8 \text{ or } x = -2
\end{align*}

\end{tcolorbox}

\end{exercise}

\begin{exercise}[]\protect\hypertarget{exr-u6-35}{}\label{exr-u6-35}

\emph{Solve using the substitution \(u = x^2\).}

\textbf{a)} \(\displaystyle x^4 -10x^2 + 16 = 0\)

\textbf{b)} \(\displaystyle x^4 -6x^2 + 5 = 0\)

\textbf{c)} \(\displaystyle x^4 -5x^2 + 4 = 0\)

\textbf{d)} \(\displaystyle x^4 -5x^2 + 4 = 0\)

\begin{tcolorbox}[enhanced jigsaw, arc=.35mm, breakable, toptitle=1mm, title=\textcolor{quarto-callout-tip-color}{\faLightbulb}\hspace{0.5em}{🔍 View Solutions}, colframe=quarto-callout-tip-color-frame, rightrule=.15mm, opacityback=0, titlerule=0mm, colback=white, coltitle=black, bottomtitle=1mm, toprule=.15mm, leftrule=.75mm, bottomrule=.15mm, colbacktitle=quarto-callout-tip-color!10!white, left=2mm, opacitybacktitle=0.6]

\textbf{a)} \(\displaystyle x^4 -10x^2 + 16 = 0\)

\begin{align*}
\text{Let } u = x^2: \quad u^2 -10u + 16 = 0 \\ (u -2)(u -8) = 0 \\ u = 2 \text{ or } u = 8 \\ x^2 = 2 \Rightarrow x = \pm\sqrt{2} \\ x^2 = 8 \Rightarrow x = \pm\sqrt{8}
\end{align*}

\textbf{b)} \(\displaystyle x^4 -6x^2 + 5 = 0\)

\begin{align*}
\text{Let } u = x^2: \quad u^2 -6u + 5 = 0 \\ (u -1)(u -5) = 0 \\ u = 1 \text{ or } u = 5 \\ x^2 = 1 \Rightarrow x = \pm1 \\ x^2 = 5 \Rightarrow x = \pm\sqrt{5}
\end{align*}

\textbf{c)} \(\displaystyle x^4 -5x^2 + 4 = 0\)

\begin{align*}
\text{Let } u = x^2: \quad u^2 -5u + 4 = 0 \\ (u -4)(u -1) = 0 \\ u = 4 \text{ or } u = 1 \\ x^2 = 4 \Rightarrow x = \pm 2 \\ x^2 = 1 \Rightarrow x = \pm 1
\end{align*}

\textbf{d)} \(\displaystyle x^4 -5x^2 + 4 = 0\)

\begin{align*}
\text{Let } u = x^2: \quad u^2 -5u + 4 = 0 \\ (u -4)(u -1) = 0 \\ u = 4 \text{ or } u = 1 \\ x^2 = 4 \Rightarrow x = \pm 2 \\ x^2 = 1 \Rightarrow x = \pm 1
\end{align*}

\end{tcolorbox}

\end{exercise}

\begin{exercise}[]\protect\hypertarget{exr-u6-36}{}\label{exr-u6-36}

\emph{Solve using an appropriate substitution.}

\textbf{a)} \(\displaystyle (x + 5)^2 + 8(x + 5) + 12 = 0\)

\textbf{b)} \(\displaystyle (x + 2)^2 - 5(x + 2) - 6 = 0\)

\textbf{c)} \(\displaystyle (x + 1)^2 + 7(x + 1) + 10 = 0\)

\textbf{d)} \(\displaystyle (x + 4)^2 + 5(x + 4) + 4 = 0\)

\begin{tcolorbox}[enhanced jigsaw, arc=.35mm, breakable, toptitle=1mm, title=\textcolor{quarto-callout-tip-color}{\faLightbulb}\hspace{0.5em}{🔍 View Solutions}, colframe=quarto-callout-tip-color-frame, rightrule=.15mm, opacityback=0, titlerule=0mm, colback=white, coltitle=black, bottomtitle=1mm, toprule=.15mm, leftrule=.75mm, bottomrule=.15mm, colbacktitle=quarto-callout-tip-color!10!white, left=2mm, opacitybacktitle=0.6]

\textbf{a)} \(\displaystyle (x + 5)^2 + 8(x + 5) + 12 = 0\)

\begin{align*}
\text{Let } u = (x + 5): \quad u^2 + 8u + 12 = 0 \\ (u + 2)(u + 6) = 0 \\ u = -2 \Rightarrow x + 5 = -2 \Rightarrow x = -7 \\ u = -6 \Rightarrow x + 5 = -6 \Rightarrow x = -11
\end{align*}

\textbf{b)} \(\displaystyle (x + 2)^2 - 5(x + 2) - 6 = 0\)

\begin{align*}
\text{Let } u = (x + 2): \quad u^2 - 5u - 6 = 0 \\ (u + 1)(u -6) = 0 \\ u = -1 \Rightarrow x + 2 = -1 \Rightarrow x = -3 \\ u = 6 \Rightarrow x + 2 = 6 \Rightarrow x = 4
\end{align*}

\textbf{c)} \(\displaystyle (x + 1)^2 + 7(x + 1) + 10 = 0\)

\begin{align*}
\text{Let } u = (x + 1): \quad u^2 + 7u + 10 = 0 \\ (u + 2)(u + 5) = 0 \\ u = -2 \Rightarrow x + 1 = -2 \Rightarrow x = -3 \\ u = -5 \Rightarrow x + 1 = -5 \Rightarrow x = -6
\end{align*}

\textbf{d)} \(\displaystyle (x + 4)^2 + 5(x + 4) + 4 = 0\)

\begin{align*}
\text{Let } u = (x + 4): \quad u^2 + 5u + 4 = 0 \\ (u + 1)(u + 4) = 0 \\ u = -1 \Rightarrow x + 4 = -1 \Rightarrow x = -5 \\ u = -4 \Rightarrow x + 4 = -4 \Rightarrow x = -8
\end{align*}

\end{tcolorbox}

\end{exercise}

\begin{exercise}[]\protect\hypertarget{exr-u6-37}{}\label{exr-u6-37}

\emph{Solve. Check for extraneous roots.}

\textbf{a)} \(\displaystyle \sqrt{x + 74} = x + 2\)

\textbf{b)} \(\displaystyle \sqrt{x + 29} = x - 1\)

\textbf{c)} \(\displaystyle \sqrt{x + 41} = x - 1\)

\textbf{d)} \(\displaystyle \sqrt{x + 8} = x + 2\)

\begin{tcolorbox}[enhanced jigsaw, arc=.35mm, breakable, toptitle=1mm, title=\textcolor{quarto-callout-tip-color}{\faLightbulb}\hspace{0.5em}{🔍 View Solutions}, colframe=quarto-callout-tip-color-frame, rightrule=.15mm, opacityback=0, titlerule=0mm, colback=white, coltitle=black, bottomtitle=1mm, toprule=.15mm, leftrule=.75mm, bottomrule=.15mm, colbacktitle=quarto-callout-tip-color!10!white, left=2mm, opacitybacktitle=0.6]

\textbf{a)} \(\displaystyle \sqrt{x + 74} = x + 2\)

\begin{align*}
\text{Square both sides: } x + 74 = (x + 2)^2 \\ x^2 + 3x - 70 = 0 \\ (x -7)(x + 10) = 0 \\ x = 7 \text{ or } x = -10 \\ \text{Check both in original equation:} \\ \text{Check } x=-10: 8 = 8 \neq -8 \text{ (extraneous)} \\ \text{Valid: } x = 7
\end{align*}

\textbf{b)} \(\displaystyle \sqrt{x + 29} = x - 1\)

\begin{align*}
\text{Square both sides: } x + 29 = (x - 1)^2 \\ x^2 - 3x - 28 = 0 \\ (x -7)(x + 4) = 0 \\ x = 7 \text{ or } x = -4 \\ \text{Check both in original equation:} \\ \text{Check } x=-4: 5 = 5 \neq -5 \text{ (extraneous)} \\ \text{Valid: } x = 7
\end{align*}

\textbf{c)} \(\displaystyle \sqrt{x + 41} = x - 1\)

\begin{align*}
\text{Square both sides: } x + 41 = (x - 1)^2 \\ x^2 - 3x - 40 = 0 \\ (x -8)(x + 5) = 0 \\ x = 8 \text{ or } x = -5 \\ \text{Check both in original equation:} \\ \text{Check } x=-5: 6 = 6 \neq -6 \text{ (extraneous)} \\ \text{Valid: } x = 8
\end{align*}

\textbf{d)} \(\displaystyle \sqrt{x + 8} = x + 2\)

\begin{align*}
\text{Square both sides: } x + 8 = (x + 2)^2 \\ x^2 + 3x - 4 = 0 \\ (x -1)(x + 4) = 0 \\ x = 1 \text{ or } x = -4 \\ \text{Check both in original equation:} \\ \text{Check } x=-4: 2 = 2 \neq -2 \text{ (extraneous)} \\ \text{Valid: } x = 1
\end{align*}

\end{tcolorbox}

\end{exercise}

\begin{exercise}[]\protect\hypertarget{exr-u6-38}{}\label{exr-u6-38}

\emph{Solve. Check for extraneous roots.}

\textbf{a)} \(\displaystyle \sqrt{x + 3} - \sqrt{7 - x} = 2\)

\textbf{b)} \(\displaystyle \sqrt{x + 5} - \sqrt{8 - x} = 1\)

\textbf{c)} \(\displaystyle \sqrt{x + 1} - \sqrt{4 - x} = 1\)

\begin{tcolorbox}[enhanced jigsaw, arc=.35mm, breakable, toptitle=1mm, title=\textcolor{quarto-callout-tip-color}{\faLightbulb}\hspace{0.5em}{🔍 View Solutions}, colframe=quarto-callout-tip-color-frame, rightrule=.15mm, opacityback=0, titlerule=0mm, colback=white, coltitle=black, bottomtitle=1mm, toprule=.15mm, leftrule=.75mm, bottomrule=.15mm, colbacktitle=quarto-callout-tip-color!10!white, left=2mm, opacitybacktitle=0.6]

\textbf{a)} \(\displaystyle \sqrt{x + 3} - \sqrt{7 - x} = 2\)

\begin{align*}
\text{Isolate: } \sqrt{x + 3} = 2 + \sqrt{7-x} \\ \text{Square: } x + 3 = 2^2 + 2(2)\sqrt{7-x} + (7-x) \\ 2(2)\sqrt{7-x} = 2x - 8 \\ \text{Square again: } 4(4)(7-x) = (2x -8)^2 \\ 4x^2 -16x -48 = 0 \\ x = 6 \\ \text{(check both roots in original)}
\end{align*}

\textbf{b)} \(\displaystyle \sqrt{x + 5} - \sqrt{8 - x} = 1\)

\begin{align*}
\text{Isolate: } \sqrt{x + 5} = 1 + \sqrt{8-x} \\ \text{Square: } x + 5 = 1^2 + 2(1)\sqrt{8-x} + (8-x) \\ 2(1)\sqrt{8-x} = 2x - 4 \\ \text{Square again: } 4(1)(8-x) = (2x -4)^2 \\ 4x^2 -12x -16 = 0 \\ x = 4 \\ \text{(check both roots in original)}
\end{align*}

\textbf{c)} \(\displaystyle \sqrt{x + 1} - \sqrt{4 - x} = 1\)

\begin{align*}
\text{Isolate: } \sqrt{x + 1} = 1 + \sqrt{4-x} \\ \text{Square: } x + 1 = 1^2 + 2(1)\sqrt{4-x} + (4-x) \\ 2(1)\sqrt{4-x} = 2x - 4 \\ \text{Square again: } 4(1)(4-x) = (2x -4)^2 \\ 4x^2 -12x +0 = 0 \\ x = 3 \\ \text{(check both roots in original)}
\end{align*}

\end{tcolorbox}

\end{exercise}

\begin{exercise}[]\protect\hypertarget{exr-u6-39}{}\label{exr-u6-39}

\emph{Solve. (Hint: use the index law \((a^m)^n = a^{mn}\), equate
exponents, then solve the resulting quadratic.)}

\textbf{a)} \(\displaystyle \left(2^{x - 2}\right)^{x - 2} = 16\)

\textbf{b)} \(\displaystyle \left(2^{x}\right)^{x - 4} = 32\)

\textbf{c)} \(\displaystyle \left(2^{x - 3}\right)^{x - 3} = 16\)

\textbf{d)} \(\displaystyle \left(2^{x - 4}\right)^{x} = 32\)

\begin{tcolorbox}[enhanced jigsaw, arc=.35mm, breakable, toptitle=1mm, title=\textcolor{quarto-callout-tip-color}{\faLightbulb}\hspace{0.5em}{🔍 View Solutions}, colframe=quarto-callout-tip-color-frame, rightrule=.15mm, opacityback=0, titlerule=0mm, colback=white, coltitle=black, bottomtitle=1mm, toprule=.15mm, leftrule=.75mm, bottomrule=.15mm, colbacktitle=quarto-callout-tip-color!10!white, left=2mm, opacitybacktitle=0.6]

\textbf{a)} \(\displaystyle \left(2^{x - 2}\right)^{x - 2} = 16\)

\begin{align*}
\text{Use index law: } (a^m)^n = a^{mn} \\ 2^{(x - 2)(x - 2)} = 2^{4} \\ \text{Equate exponents: } (x - 2)(x - 2) = 4 \\ x^2 -4x + 4 = 4 \\ x^2 -4x  = 0 \\ x = \dfrac{-(-4) \pm \sqrt{(-4)^2 - 4(1)(0)}}{2} \\ x = \dfrac{4 \pm 4}{2} \\ x = 4 \text{ or } x = 0
\end{align*}

\textbf{b)} \(\displaystyle \left(2^{x}\right)^{x - 4} = 32\)

\begin{align*}
\text{Use index law: } (a^m)^n = a^{mn} \\ 2^{(x)(x - 4)} = 2^{5} \\ \text{Equate exponents: } (x)(x - 4) = 5 \\ x^2 -4x  = 5 \\ x^2 -4x -5 = 0 \\ x = \dfrac{-(-4) \pm \sqrt{(-4)^2 - 4(1)(-5)}}{2} \\ x = \dfrac{4 \pm 6}{2} \\ x = 5 \text{ or } x = -1
\end{align*}

\textbf{c)} \(\displaystyle \left(2^{x - 3}\right)^{x - 3} = 16\)

\begin{align*}
\text{Use index law: } (a^m)^n = a^{mn} \\ 2^{(x - 3)(x - 3)} = 2^{4} \\ \text{Equate exponents: } (x - 3)(x - 3) = 4 \\ x^2 -6x + 9 = 4 \\ x^2 -6x + 5 = 0 \\ x = \dfrac{-(-6) \pm \sqrt{(-6)^2 - 4(1)(5)}}{2} \\ x = \dfrac{6 \pm 4}{2} \\ x = 5 \text{ or } x = 1
\end{align*}

\textbf{d)} \(\displaystyle \left(2^{x - 4}\right)^{x} = 32\)

\begin{align*}
\text{Use index law: } (a^m)^n = a^{mn} \\ 2^{(x - 4)(x)} = 2^{5} \\ \text{Equate exponents: } (x - 4)(x) = 5 \\ x^2 -4x  = 5 \\ x^2 -4x -5 = 0 \\ x = \dfrac{-(-4) \pm \sqrt{(-4)^2 - 4(1)(-5)}}{2} \\ x = \dfrac{4 \pm 6}{2} \\ x = 5 \text{ or } x = -1
\end{align*}

\end{tcolorbox}

\end{exercise}

∑ × π ÷ √ ◆ ∞ ± ∫ ≠ Δ

\section{⚡ Test Your Knowledge}\label{test-your-knowledge-1}

\chapter{Unit 7: Linear \& Quadratic
Functions}\label{unit-7-linear-quadratic-functions}

\section{🎯 Learning Targets}\label{learning-targets-6}

\begin{longtable}[]{@{}
  >{\raggedright\arraybackslash}p{(\linewidth - 4\tabcolsep) * \real{0.3333}}
  >{\raggedright\arraybackslash}p{(\linewidth - 4\tabcolsep) * \real{0.3333}}
  >{\raggedright\arraybackslash}p{(\linewidth - 4\tabcolsep) * \real{0.3333}}@{}}
\toprule\noalign{}
\begin{minipage}[b]{\linewidth}\raggedright
\textbf{\#}
\end{minipage} & \begin{minipage}[b]{\linewidth}\raggedright
\textbf{Learning Target}
\end{minipage} & \begin{minipage}[b]{\linewidth}\raggedright
\textbf{I can\ldots{}}
\end{minipage} \\
\midrule\noalign{}
\endhead
\bottomrule\noalign{}
\endlastfoot
\textbf{LT1} & Relations \& Functions & Represent relations as sets,
tables, and arrow diagrams, and determine whether a relation is a
function \\
\textbf{LT2} & Representations & Move fluently between tables, graphs,
and equations for linear functions \\
\textbf{LT3} & Domain \& Range & Identify the domain and range of a
relation from a set, table, or graph \\
\textbf{LT4} & Linear Functions & Identify, evaluate, and interpret
linear functions in context \\
\textbf{LT5} & Slope \& Rate of Change & Calculate slope from two points
or a table and interpret it as a rate of change \\
\textbf{LT6} & Linear Equations \& Graphs & Write and graph linear
equations; find equations of parallel and perpendicular lines \\
\textbf{LT7} & Linear Modeling & Form, solve, and interpret linear
models for real-world situations \\
\textbf{LT8} & Quadratic Functions & Identify quadratic functions;
evaluate them; convert between vertex, standard, and factored forms \\
\textbf{LT9} & Key Features of Graphs & Determine the vertex, roots, and
axis of symmetry of a parabola from its equation or graph \\
\textbf{LT10} & Algebraic Manipulation & \emph{(Covered in Unit 6 ---
Quadratic Equations)} \\
\textbf{LT11} & Solving Quadratics & \emph{(Covered in Unit 6 ---
Quadratic Equations)} \\
\textbf{LT12} & Domain \& Range (Quadratic) & Determine and interpret
the domain and range of a quadratic function from its equation or
graph \\
\textbf{LT13} & Quadratic Modeling & Model and interpret real-world
situations using quadratic functions (projectile motion, revenue) \\
\end{longtable}

\section{🗂️ Formula Card}\label{formula-card-3}

\subsection*{Linear \& Quadratic Functions --- Key
Facts}\label{linear-quadratic-functions-key-facts}
\addcontentsline{toc}{subsection}{Linear \& Quadratic Functions --- Key
Facts}

\subsubsection*{Relations \& Functions}\label{relations-functions}
\addcontentsline{toc}{subsubsection}{Relations \& Functions}

\textbf{💡 Key Test} A relation is a \textbf{function} if every input
(\(x\)-value) maps to \textbf{exactly one} output (\(y\)-value). Use the
vertical line test on a graph.

\textbf{Function notation}: \(f(x) = mx + c\) means ``the output when
the input is \(x\)''

\textbf{Domain}: the set of all valid \textbf{inputs} (\(x\)-values)

\textbf{Range}: the set of all valid \textbf{outputs} (\(y\)-values)

\textbf{Examples}:
\[f(x) = 2x + 1 \quad \Rightarrow \quad f(3) = 2(3)+1 = 7\]

\subsubsection*{Linear Functions}\label{linear-functions}
\addcontentsline{toc}{subsubsection}{Linear Functions}

\textbf{📝 Slope Formula}
\[m = \frac{y_2 - y_1}{x_2 - x_1} = \frac{\Delta y}{\Delta x}\]

\textbf{Slope-intercept form}: \(y = mx + c\), where \(m\) = slope,
\(c\) = \(y\)-intercept

\textbf{Parallel lines} share the \textbf{same slope}: \(m_1 = m_2\)

\textbf{Perpendicular lines}: slopes are \textbf{negative reciprocals}:
\(m_1 \times m_2 = -1\)

\textbf{Examples}:
\[\text{Through }(2,5)\text{ with slope }3: \quad y - 5 = 3(x-2) \quad \Rightarrow \quad y = 3x - 1\]
\[\text{Perpendicular to }y=2x+1: \quad m_\perp = -\tfrac{1}{2}\]

\subsubsection*{Quadratic Functions --- Three
Forms}\label{quadratic-functions-three-forms}
\addcontentsline{toc}{subsubsection}{Quadratic Functions --- Three
Forms}

\textbf{⚠️ Conversion Tip} Each form reveals different information.
Convert using expansion or completing the square.

\begin{longtable}[]{@{}lll@{}}
\toprule\noalign{}
Form & Equation & Reveals \\
\midrule\noalign{}
\endhead
\bottomrule\noalign{}
\endlastfoot
\textbf{Standard} & \(f(x) = ax^2 + bx + c\) & \(y\)-intercept
\((0, c)\) \\
\textbf{Vertex} & \(f(x) = a(x-h)^2 + k\) & Vertex \((h, k)\) \\
\textbf{Factored} & \(f(x) = a(x-r_1)(x-r_2)\) & Roots
\(x = r_1, r_2\) \\
\end{longtable}

\textbf{Vertex formula} (from standard form):
\[x_{\text{vertex}} = \frac{-b}{2a}\]

\textbf{Axis of symmetry}: \(x = \dfrac{-b}{2a} = h\)

\subsubsection*{Converting Between
Forms}\label{converting-between-forms}
\addcontentsline{toc}{subsubsection}{Converting Between Forms}

\textbf{Vertex → Standard} (expand):
\[2(x-3)^2 + 1 = 2(x^2-6x+9)+1 = 2x^2 - 12x + 19\]

\textbf{Vertex → Factored} (find roots first):
\[\textstyle(x-3)^2 - 4 = 0 \Rightarrow x = 3\pm2 \Rightarrow f(x)=(x-1)(x-5)\]

\textbf{Standard → Vertex} (complete the square):
\[x^2 - 6x + 5 = (x-3)^2 - 9 + 5 = (x-3)^2 - 4\]

\textbf{Factored → Standard} (expand): \[(x-1)(x-5) = x^2 - 6x + 5\]

\subsubsection*{Domain \& Range}\label{domain-range}
\addcontentsline{toc}{subsubsection}{Domain \& Range}

\textbf{🎯 Quick Rule} Every quadratic has domain \(x \in \mathbb{R}\).
The range depends on whether the parabola opens up or down.

\begin{longtable}[]{@{}lll@{}}
\toprule\noalign{}
Function type & Domain & Range \\
\midrule\noalign{}
\endhead
\bottomrule\noalign{}
\endlastfoot
Linear \(f(x) = mx+c\) & \(x \in \mathbb{R}\) & \(y \in \mathbb{R}\) \\
Quadratic, \(a > 0\) (opens up) & \(x \in \mathbb{R}\) & \(y \geq k\) \\
Quadratic, \(a < 0\) (opens down) & \(x \in \mathbb{R}\) &
\(y \leq k\) \\
\end{longtable}

\section{📝 Practice Problems}\label{practice-problems-4}

\subsection{🟢 Easy}

\begin{exercise}[]\protect\hypertarget{exr-u7-1}{}\label{exr-u7-1}

\emph{Analyse the arrow diagram.}

The arrow diagram shows a relation.

Domain Range 1 6 3 4 2 6

\textbf{(a)} State the domain and range of this relation.

\textbf{(b)} Is this relation a function? Justify your answer.

\begin{tcolorbox}[enhanced jigsaw, arc=.35mm, breakable, toptitle=1mm, title=\textcolor{quarto-callout-tip-color}{\faLightbulb}\hspace{0.5em}{🔍 View Solution}, colframe=quarto-callout-tip-color-frame, rightrule=.15mm, opacityback=0, titlerule=0mm, colback=white, coltitle=black, bottomtitle=1mm, toprule=.15mm, leftrule=.75mm, bottomrule=.15mm, colbacktitle=quarto-callout-tip-color!10!white, left=2mm, opacitybacktitle=0.6]

\[
\begin{aligned}
\text{Check: does each domain element have exactly one arrow?} \\[4pt] \text{No — one input maps to more than one output.}
\end{aligned}
\]

\end{tcolorbox}

{Function definition} {Domain \& range} {Arrow diagrams}

\end{exercise}

\begin{exercise}[]\protect\hypertarget{exr-u7-2}{}\label{exr-u7-2}

\emph{Analyse the relation.}

Consider the relation:
\(\{\, (2,\ -1), (3,\ 4), (-4,\ -1), (5,\ -3), (5,\ -2) \,\}\)

\textbf{(a)} State the domain and range.

\textbf{(b)} Is this a function? Explain.

\begin{tcolorbox}[enhanced jigsaw, arc=.35mm, breakable, toptitle=1mm, title=\textcolor{quarto-callout-tip-color}{\faLightbulb}\hspace{0.5em}{🔍 View Solution}, colframe=quarto-callout-tip-color-frame, rightrule=.15mm, opacityback=0, titlerule=0mm, colback=white, coltitle=black, bottomtitle=1mm, toprule=.15mm, leftrule=.75mm, bottomrule=.15mm, colbacktitle=quarto-callout-tip-color!10!white, left=2mm, opacitybacktitle=0.6]

\[
\begin{aligned}
\text{Domain: } \{-4, 2, 3, 5\} \\[4pt] \text{Range: } \{-3, -2, -1, 4\} \\[4pt] \text{Function? No — one $x$-value maps to two outputs.}
\end{aligned}
\]

\end{tcolorbox}

{Function definition} {Domain \& range} {Arrow diagrams}

\end{exercise}

\begin{exercise}[]\protect\hypertarget{exr-u7-3}{}\label{exr-u7-3}

\emph{Find the equation of the linear function shown in the table.}

\textbf{a)}
\(\displaystyle \begin{array}{c|cccc} x & -2 & -1 & 3 & 4 \\ \hline y & 3 & 0 & -12 & -15 \end{array}\)

\textbf{b)}
\(\displaystyle \begin{array}{c|cccc} x & -3 & -2 & 0 & 3 \\ \hline y & -3 & -2 & 0 & 3 \end{array}\)

\textbf{c)}
\(\displaystyle \begin{array}{c|cccc} x & -3 & 0 & 1 & 3 \\ \hline y & -9 & 0 & 3 & 9 \end{array}\)

\textbf{d)}
\(\displaystyle \begin{array}{c|cccc} x & -3 & 0 & 2 & 3 \\ \hline y & -5 & 4 & 10 & 13 \end{array}\)

\begin{tcolorbox}[enhanced jigsaw, arc=.35mm, breakable, toptitle=1mm, title=\textcolor{quarto-callout-tip-color}{\faLightbulb}\hspace{0.5em}{🔍 View Solutions}, colframe=quarto-callout-tip-color-frame, rightrule=.15mm, opacityback=0, titlerule=0mm, colback=white, coltitle=black, bottomtitle=1mm, toprule=.15mm, leftrule=.75mm, bottomrule=.15mm, colbacktitle=quarto-callout-tip-color!10!white, left=2mm, opacitybacktitle=0.6]

\textbf{a)}

\[
\begin{aligned}
\text{Slope: } m = \dfrac{\Delta y}{\Delta x} = \dfrac{-3}{1} = -3 \\[4pt] \text{Substitute } (-2, 3): 3 = -3(-2) + c \\[4pt] c = -3 \\[4pt] y = -3x - 3
\end{aligned}
\] \textbf{b)}

\[
\begin{aligned}
\text{Slope: } m = \dfrac{\Delta y}{\Delta x} = \dfrac{1}{1} = 1 \\[4pt] \text{Substitute } (-3, -3): -3 = 1(-3) + c \\[4pt] c = 0 \\[4pt] y = x
\end{aligned}
\] \textbf{c)}

\[
\begin{aligned}
\text{Slope: } m = \dfrac{\Delta y}{\Delta x} = \dfrac{9}{3} = 3 \\[4pt] \text{Substitute } (-3, -9): -9 = 3(-3) + c \\[4pt] c = 0 \\[4pt] y = 3x
\end{aligned}
\] \textbf{d)}

\[
\begin{aligned}
\text{Slope: } m = \dfrac{\Delta y}{\Delta x} = \dfrac{9}{3} = 3 \\[4pt] \text{Substitute } (-3, -5): -5 = 3(-3) + c \\[4pt] c = 4 \\[4pt] y = 3x + 4
\end{aligned}
\]

\end{tcolorbox}

{Table of values} {Linear equations} {Representations}

\end{exercise}

\begin{exercise}[]\protect\hypertarget{exr-u7-4}{}\label{exr-u7-4}

\emph{Complete the table of values.}

\textbf{a)}
\(\displaystyle y = -2x + 2 \qquad \begin{array}{c|cccc} x & -1 & 0 & 1 & 4 \\ \hline y & \square & 2 & 0 & \square \end{array}\)

\textbf{b)}
\(\displaystyle y = -x \qquad \begin{array}{c|cccc} x & -2 & -1 & 3 & 4 \\ \hline y & 2 & 1 & \square & \square \end{array}\)

\textbf{c)}
\(\displaystyle y = 2x + 2 \qquad \begin{array}{c|cccc} x & -1 & 0 & 2 & 3 \\ \hline y & 0 & \square & 6 & \square \end{array}\)

\textbf{d)}
\(\displaystyle y = -x - 4 \qquad \begin{array}{c|cccc} x & 0 & 1 & 3 & 4 \\ \hline y & -4 & \square & \square & -8 \end{array}\)

\begin{tcolorbox}[enhanced jigsaw, arc=.35mm, breakable, toptitle=1mm, title=\textcolor{quarto-callout-tip-color}{\faLightbulb}\hspace{0.5em}{🔍 View Solutions}, colframe=quarto-callout-tip-color-frame, rightrule=.15mm, opacityback=0, titlerule=0mm, colback=white, coltitle=black, bottomtitle=1mm, toprule=.15mm, leftrule=.75mm, bottomrule=.15mm, colbacktitle=quarto-callout-tip-color!10!white, left=2mm, opacitybacktitle=0.6]

\textbf{a)}

\[
\begin{aligned}
x = -1: \quad y = -2(-1) + 2 = 4 \\[4pt] x = 4: \quad y = -2(4) + 2 = -6
\end{aligned}
\] \textbf{b)}

\[
\begin{aligned}
x = 3: \quad y = -1(3) + 0 = -3 \\[4pt] x = 4: \quad y = -1(4) + 0 = -4
\end{aligned}
\] \textbf{c)}

\[
\begin{aligned}
x = 0: \quad y = 2(0) + 2 = 2 \\[4pt] x = 3: \quad y = 2(3) + 2 = 8
\end{aligned}
\] \textbf{d)}

\[
\begin{aligned}
x = 1: \quad y = -1(1) - 4 = -5 \\[4pt] x = 3: \quad y = -1(3) - 4 = -7
\end{aligned}
\]

\end{tcolorbox}

{Table of values} {Linear equations} {Representations}

\end{exercise}

\begin{exercise}[]\protect\hypertarget{exr-u7-5}{}\label{exr-u7-5}

\emph{State the domain and range of each relation.}

\textbf{a)}
\(\displaystyle \{\,(1,\ -1), (3,\ 2), (0,\ -2), (-1,\ 4), (4,\ 0)\,\}\)

\textbf{b)}
\(\displaystyle \{\,(2,\ 2), (1,\ -1), (-2,\ 4), (3,\ -3), (-4,\ -3)\,\}\)

\textbf{c)}
\(\displaystyle \{\,(2,\ -3), (3,\ 4), (0,\ 4), (-4,\ -3), (4,\ 0)\,\}\)

\textbf{d)}
\(\displaystyle \{\,(1,\ -3), (0,\ 2), (-2,\ 3), (-4,\ 0), (4,\ -2)\,\}\)

\begin{tcolorbox}[enhanced jigsaw, arc=.35mm, breakable, toptitle=1mm, title=\textcolor{quarto-callout-tip-color}{\faLightbulb}\hspace{0.5em}{🔍 View Solutions}, colframe=quarto-callout-tip-color-frame, rightrule=.15mm, opacityback=0, titlerule=0mm, colback=white, coltitle=black, bottomtitle=1mm, toprule=.15mm, leftrule=.75mm, bottomrule=.15mm, colbacktitle=quarto-callout-tip-color!10!white, left=2mm, opacitybacktitle=0.6]

\textbf{a)}

\[
\begin{aligned}
\text{Domain (all }x\text{-values): } \{-1, 0, 1, 3, 4\} \\[4pt] \text{Range (all }y\text{-values): } \{-2, -1, 0, 2, 4\}
\end{aligned}
\] \textbf{b)}

\[
\begin{aligned}
\text{Domain (all }x\text{-values): } \{-4, -2, 1, 2, 3\} \\[4pt] \text{Range (all }y\text{-values): } \{-3, -1, 2, 4\}
\end{aligned}
\] \textbf{c)}

\[
\begin{aligned}
\text{Domain (all }x\text{-values): } \{-4, 0, 2, 3, 4\} \\[4pt] \text{Range (all }y\text{-values): } \{-3, 0, 4\}
\end{aligned}
\] \textbf{d)}

\[
\begin{aligned}
\text{Domain (all }x\text{-values): } \{-4, -2, 0, 1, 4\} \\[4pt] \text{Range (all }y\text{-values): } \{-3, -2, 0, 2, 3\}
\end{aligned}
\]

\end{tcolorbox}

{Domain} {Range} {Set notation}

\end{exercise}

\begin{exercise}[]\protect\hypertarget{exr-u7-6}{}\label{exr-u7-6}

\emph{State the domain and range of each function.}

\textbf{a)} \(\displaystyle f(x) = 3x - 1\)

\textbf{b)} \(\displaystyle f(x) = 3x + 3\)

\textbf{c)} \(\displaystyle f(x) = 2x - 2\)

\textbf{d)} \(\displaystyle f(x) = -2x\)

\begin{tcolorbox}[enhanced jigsaw, arc=.35mm, breakable, toptitle=1mm, title=\textcolor{quarto-callout-tip-color}{\faLightbulb}\hspace{0.5em}{🔍 View Solutions}, colframe=quarto-callout-tip-color-frame, rightrule=.15mm, opacityback=0, titlerule=0mm, colback=white, coltitle=black, bottomtitle=1mm, toprule=.15mm, leftrule=.75mm, bottomrule=.15mm, colbacktitle=quarto-callout-tip-color!10!white, left=2mm, opacitybacktitle=0.6]

\textbf{a)}

\[
\begin{aligned}
\text{Linear functions are defined for all real numbers.} \\[4pt] \text{Domain: } x \in \mathbb{R} \\[4pt] \text{Range: } y \in \mathbb{R}
\end{aligned}
\] \textbf{b)}

\[
\begin{aligned}
\text{Linear functions are defined for all real numbers.} \\[4pt] \text{Domain: } x \in \mathbb{R} \\[4pt] \text{Range: } y \in \mathbb{R}
\end{aligned}
\] \textbf{c)}

\[
\begin{aligned}
\text{Linear functions are defined for all real numbers.} \\[4pt] \text{Domain: } x \in \mathbb{R} \\[4pt] \text{Range: } y \in \mathbb{R}
\end{aligned}
\] \textbf{d)}

\[
\begin{aligned}
\text{Linear functions are defined for all real numbers.} \\[4pt] \text{Domain: } x \in \mathbb{R} \\[4pt] \text{Range: } y \in \mathbb{R}
\end{aligned}
\]

\end{tcolorbox}

{Domain} {Range} {Set notation}

\end{exercise}

\begin{exercise}[]\protect\hypertarget{exr-u7-7}{}\label{exr-u7-7}

\emph{Identify which expression represents a linear function.}

\textbf{a)}
\(\displaystyle \text{A:}\; |t| + 1 \qquad \text{B:}\; \dfrac{1}{t} + 4 \qquad \text{C:}\; 2t^2 + t + 1 \qquad \text{D:}\; -2t + 3\)

\textbf{b)}
\(\displaystyle \text{A:}\; \dfrac{1}{y} + 1 \qquad \text{B:}\; |y| + 5 \qquad \text{C:}\; -2y - 5 \qquad \text{D:}\; 2y^2 - 3y + 3\)

\textbf{c)}
\(\displaystyle \text{A:}\; 2n - 3 \qquad \text{B:}\; 3n^2 - 2n + 4 \qquad \text{C:}\; \dfrac{1}{n} + 2 \qquad \text{D:}\; |n| + 3\)

\textbf{d)}
\(\displaystyle \text{A:}\; 2m^2 - 2m - 1 \qquad \text{B:}\; |m| + 3 \qquad \text{C:}\; \dfrac{1}{m} + 1 \qquad \text{D:}\; m + 3\)

\begin{tcolorbox}[enhanced jigsaw, arc=.35mm, breakable, toptitle=1mm, title=\textcolor{quarto-callout-tip-color}{\faLightbulb}\hspace{0.5em}{🔍 View Solutions}, colframe=quarto-callout-tip-color-frame, rightrule=.15mm, opacityback=0, titlerule=0mm, colback=white, coltitle=black, bottomtitle=1mm, toprule=.15mm, leftrule=.75mm, bottomrule=.15mm, colbacktitle=quarto-callout-tip-color!10!white, left=2mm, opacitybacktitle=0.6]

\textbf{a)}

\[
\begin{aligned}
\text{A linear function has degree 1 (form } f(x) = mx + c\text{).} \\[4pt] \text{Answer: } D \text{ — } f(x) = -2t + 3 \text{ has degree 1. }\checkmark
\end{aligned}
\] \textbf{b)}

\[
\begin{aligned}
\text{A linear function has degree 1 (form } f(x) = mx + c\text{).} \\[4pt] \text{Answer: } C \text{ — } f(x) = -2y - 5 \text{ has degree 1. }\checkmark
\end{aligned}
\] \textbf{c)}

\[
\begin{aligned}
\text{A linear function has degree 1 (form } f(x) = mx + c\text{).} \\[4pt] \text{Answer: } A \text{ — } f(x) = 2n - 3 \text{ has degree 1. }\checkmark
\end{aligned}
\] \textbf{d)}

\[
\begin{aligned}
\text{A linear function has degree 1 (form } f(x) = mx + c\text{).} \\[4pt] \text{Answer: } D \text{ — } f(x) = m + 3 \text{ has degree 1. }\checkmark
\end{aligned}
\]

\end{tcolorbox}

{Linear functions} {Slope} {y-intercept}

\end{exercise}

\begin{exercise}[]\protect\hypertarget{exr-u7-8}{}\label{exr-u7-8}

\emph{Evaluate the function at each given value.}

\textbf{a)}
\(\displaystyle f(x) = 2x - 3 \quad \text{Find: } f(-3), f(4), f(2)\)

\textbf{b)}
\(\displaystyle f(x) = 3x - 6 \quad \text{Find: } f(0), f(-3), f(3)\)

\textbf{c)}
\(\displaystyle f(x) = -x - 5 \quad \text{Find: } f(0), f(-1), f(2)\)

\textbf{d)}
\(\displaystyle f(x) = x - 3 \quad \text{Find: } f(-1), f(2), f(1)\)

\begin{tcolorbox}[enhanced jigsaw, arc=.35mm, breakable, toptitle=1mm, title=\textcolor{quarto-callout-tip-color}{\faLightbulb}\hspace{0.5em}{🔍 View Solutions}, colframe=quarto-callout-tip-color-frame, rightrule=.15mm, opacityback=0, titlerule=0mm, colback=white, coltitle=black, bottomtitle=1mm, toprule=.15mm, leftrule=.75mm, bottomrule=.15mm, colbacktitle=quarto-callout-tip-color!10!white, left=2mm, opacitybacktitle=0.6]

\textbf{a)}

\[
\begin{aligned}
f(-3) = 2 \cdot (-3) - 3 = -9 \\[4pt] f(4) = 2 \cdot (4) - 3 = 5 \\[4pt] f(2) = 2 \cdot (2) - 3 = 1
\end{aligned}
\] \textbf{b)}

\[
\begin{aligned}
f(0) = 3 \cdot (0) - 6 = -6 \\[4pt] f(-3) = 3 \cdot (-3) - 6 = -15 \\[4pt] f(3) = 3 \cdot (3) - 6 = 3
\end{aligned}
\] \textbf{c)}

\[
\begin{aligned}
f(0) = -1 \cdot (0) - 5 = -5 \\[4pt] f(-1) = -1 \cdot (-1) - 5 = -4 \\[4pt] f(2) = -1 \cdot (2) - 5 = -7
\end{aligned}
\] \textbf{d)}

\[
\begin{aligned}
f(-1) = 1 \cdot (-1) - 3 = -4 \\[4pt] f(2) = 1 \cdot (2) - 3 = -1 \\[4pt] f(1) = 1 \cdot (1) - 3 = -2
\end{aligned}
\]

\end{tcolorbox}

{Linear functions} {Slope} {y-intercept}

\end{exercise}

\begin{exercise}[]\protect\hypertarget{exr-u7-10}{}\label{exr-u7-10}

\emph{Find the slope of the line through each pair of points.}

\textbf{a)} \(\displaystyle (3,\ 2) \text{ and } (1,\ -1)\)

\textbf{b)} \(\displaystyle (-3,\ 2) \text{ and } (-4,\ 5)\)

\textbf{c)} \(\displaystyle (2,\ 3) \text{ and } (6,\ -1)\)

\textbf{d)} \(\displaystyle (-3,\ -4) \text{ and } (-2,\ -2)\)

\textbf{e)} \(\displaystyle (-3,\ -1) \text{ and } (-5,\ -2)\)

\textbf{f)} \(\displaystyle (3,\ -2) \text{ and } (4,\ -4)\)

\begin{tcolorbox}[enhanced jigsaw, arc=.35mm, breakable, toptitle=1mm, title=\textcolor{quarto-callout-tip-color}{\faLightbulb}\hspace{0.5em}{🔍 View Solutions}, colframe=quarto-callout-tip-color-frame, rightrule=.15mm, opacityback=0, titlerule=0mm, colback=white, coltitle=black, bottomtitle=1mm, toprule=.15mm, leftrule=.75mm, bottomrule=.15mm, colbacktitle=quarto-callout-tip-color!10!white, left=2mm, opacitybacktitle=0.6]

\textbf{a)}

\[
\begin{aligned}
m = \dfrac{y_2 - y_1}{x_2 - x_1} = \dfrac{-1 - (2)}{1 - (3)} = \dfrac{-3}{-2} = \dfrac{3}{2}
\end{aligned}
\] \textbf{b)}

\[
\begin{aligned}
m = \dfrac{y_2 - y_1}{x_2 - x_1} = \dfrac{5 - (2)}{-4 - (-3)} = \dfrac{3}{-1} = -3
\end{aligned}
\] \textbf{c)}

\[
\begin{aligned}
m = \dfrac{y_2 - y_1}{x_2 - x_1} = \dfrac{-1 - (3)}{6 - (2)} = \dfrac{-4}{4} = -1
\end{aligned}
\] \textbf{d)}

\[
\begin{aligned}
m = \dfrac{y_2 - y_1}{x_2 - x_1} = \dfrac{-2 - (-4)}{-2 - (-3)} = \dfrac{2}{1} = 2
\end{aligned}
\] \textbf{e)}

\[
\begin{aligned}
m = \dfrac{y_2 - y_1}{x_2 - x_1} = \dfrac{-2 - (-1)}{-5 - (-3)} = \dfrac{-1}{-2} = \dfrac{1}{2}
\end{aligned}
\] \textbf{f)}

\[
\begin{aligned}
m = \dfrac{y_2 - y_1}{x_2 - x_1} = \dfrac{-4 - (-2)}{4 - (3)} = \dfrac{-2}{1} = -2
\end{aligned}
\]

\end{tcolorbox}

{Rate of change} {Slope formula} {Interpretation}

\end{exercise}

\begin{exercise}[]\protect\hypertarget{exr-u7-11}{}\label{exr-u7-11}

\emph{Find the rate of change from the table.}

\textbf{a)}
\(\displaystyle \begin{array}{c|cccc} x & -2 & 1 & 4 & 5 \\ \hline y & -1 & 5 & 11 & 13 \end{array}\)

\textbf{b)}
\(\displaystyle \begin{array}{c|cccc} x & -2 & -1 & 3 & 5 \\ \hline y & 0 & -2 & -10 & -14 \end{array}\)

\textbf{c)}
\(\displaystyle \begin{array}{c|cccc} x & -2 & 0 & 1 & 4 \\ \hline y & 4 & 2 & 1 & -2 \end{array}\)

\textbf{d)}
\(\displaystyle \begin{array}{c|cccc} x & -1 & 0 & 4 & 5 \\ \hline y & 0 & 2 & 10 & 12 \end{array}\)

\begin{tcolorbox}[enhanced jigsaw, arc=.35mm, breakable, toptitle=1mm, title=\textcolor{quarto-callout-tip-color}{\faLightbulb}\hspace{0.5em}{🔍 View Solutions}, colframe=quarto-callout-tip-color-frame, rightrule=.15mm, opacityback=0, titlerule=0mm, colback=white, coltitle=black, bottomtitle=1mm, toprule=.15mm, leftrule=.75mm, bottomrule=.15mm, colbacktitle=quarto-callout-tip-color!10!white, left=2mm, opacitybacktitle=0.6]

\textbf{a)}

\[
\begin{aligned}
\text{Rate of change} = \dfrac{\Delta y}{\Delta x} = \dfrac{6}{3} = 2 \\[4pt] \text{Constant rate → linear, slope } m = 2
\end{aligned}
\] \textbf{b)}

\[
\begin{aligned}
\text{Rate of change} = \dfrac{\Delta y}{\Delta x} = \dfrac{-2}{1} = -2 \\[4pt] \text{Constant rate → linear, slope } m = -2
\end{aligned}
\] \textbf{c)}

\[
\begin{aligned}
\text{Rate of change} = \dfrac{\Delta y}{\Delta x} = \dfrac{-2}{2} = -1 \\[4pt] \text{Constant rate → linear, slope } m = -1
\end{aligned}
\] \textbf{d)}

\[
\begin{aligned}
\text{Rate of change} = \dfrac{\Delta y}{\Delta x} = \dfrac{2}{1} = 2 \\[4pt] \text{Constant rate → linear, slope } m = 2
\end{aligned}
\]

\end{tcolorbox}

{Rate of change} {Slope formula} {Interpretation}

\end{exercise}

\begin{exercise}[]\protect\hypertarget{exr-u7-18}{}\label{exr-u7-18}

\emph{Identify which expression represents a quadratic function.}

\textbf{a)}
\(\displaystyle \text{A:}\; 3t^2 - 3t - 3 \qquad \text{B:}\; 3t^3 - 1t^2 + 3 \qquad \text{C:}\; 2 \qquad \text{D:}\; -2t + 4\)

\textbf{b)}
\(\displaystyle \text{A:}\; 2p^3 + 2p^2 + 3 \qquad \text{B:}\; -p^2 + 3p + 4 \qquad \text{C:}\; 2 \qquad \text{D:}\; -p + 4\)

\textbf{c)}
\(\displaystyle \text{A:}\; 3n^3 - 1n^2 + 3 \qquad \text{B:}\; -2n^2 + 4n - 2 \qquad \text{C:}\; -n - 1 \qquad \text{D:}\; 9\)

\textbf{d)}
\(\displaystyle \text{A:}\; -3x + 1 \qquad \text{B:}\; -2x^2 + 2x \qquad \text{C:}\; 2x^3 x^2 + 2 \qquad \text{D:}\; 8\)

\begin{tcolorbox}[enhanced jigsaw, arc=.35mm, breakable, toptitle=1mm, title=\textcolor{quarto-callout-tip-color}{\faLightbulb}\hspace{0.5em}{🔍 View Solutions}, colframe=quarto-callout-tip-color-frame, rightrule=.15mm, opacityback=0, titlerule=0mm, colback=white, coltitle=black, bottomtitle=1mm, toprule=.15mm, leftrule=.75mm, bottomrule=.15mm, colbacktitle=quarto-callout-tip-color!10!white, left=2mm, opacitybacktitle=0.6]

\textbf{a)}

\[
\begin{aligned}
\text{A quadratic has highest degree 2.} \\[4pt] \text{Answer: } A \text{ — } f(x) = 3t^2 - 3t - 3 \text{ has degree 2. }\checkmark
\end{aligned}
\] \textbf{b)}

\[
\begin{aligned}
\text{A quadratic has highest degree 2.} \\[4pt] \text{Answer: } B \text{ — } f(x) = -p^2 + 3p + 4 \text{ has degree 2. }\checkmark
\end{aligned}
\] \textbf{c)}

\[
\begin{aligned}
\text{A quadratic has highest degree 2.} \\[4pt] \text{Answer: } B \text{ — } f(x) = -2n^2 + 4n - 2 \text{ has degree 2. }\checkmark
\end{aligned}
\] \textbf{d)}

\[
\begin{aligned}
\text{A quadratic has highest degree 2.} \\[4pt] \text{Answer: } B \text{ — } f(x) = -2x^2 + 2x \text{ has degree 2. }\checkmark
\end{aligned}
\]

\end{tcolorbox}

{Quadratic functions} {Form conversions} {Vertex form}

\end{exercise}

\begin{exercise}[]\protect\hypertarget{exr-u7-19}{}\label{exr-u7-19}

\emph{Evaluate each quadratic function at the given values.}

\textbf{a)}
\(\displaystyle f(t) = -2t^2 - 3t - 2 \quad \text{Find: } f(-2), f(0), f(-1)\)

\textbf{b)}
\(\displaystyle f(n) = 2n^2 - 2n - 4 \quad \text{Find: } f(0), f(-2), f(3)\)

\textbf{c)}
\(\displaystyle f(n) = -n^2 - n - 3 \quad \text{Find: } f(2), f(1), f(-1)\)

\textbf{d)}
\(\displaystyle f(m) = -m^2 - 2m - 2 \quad \text{Find: } f(-2), f(-1), f(0)\)

\begin{tcolorbox}[enhanced jigsaw, arc=.35mm, breakable, toptitle=1mm, title=\textcolor{quarto-callout-tip-color}{\faLightbulb}\hspace{0.5em}{🔍 View Solutions}, colframe=quarto-callout-tip-color-frame, rightrule=.15mm, opacityback=0, titlerule=0mm, colback=white, coltitle=black, bottomtitle=1mm, toprule=.15mm, leftrule=.75mm, bottomrule=.15mm, colbacktitle=quarto-callout-tip-color!10!white, left=2mm, opacitybacktitle=0.6]

\textbf{a)}

\[
\begin{aligned}
f(-2) = -2(-2)^2 + 6 - 2 = -4 \\[4pt] f(0) = -2(0)^2 + 0 - 2 = -2 \\[4pt] f(-1) = -2(-1)^2 + 3 - 2 = -1
\end{aligned}
\] \textbf{b)}

\[
\begin{aligned}
f(0) = 2(0)^2 + 0 - 4 = -4 \\[4pt] f(-2) = 2(-2)^2 + 4 - 4 = 8 \\[4pt] f(3) = 2(3)^2 - 6 - 4 = 8
\end{aligned}
\] \textbf{c)}

\[
\begin{aligned}
f(2) = -1(2)^2 - 2 - 3 = -9 \\[4pt] f(1) = -1(1)^2 - 1 - 3 = -5 \\[4pt] f(-1) = -1(-1)^2 + 1 - 3 = -3
\end{aligned}
\] \textbf{d)}

\[
\begin{aligned}
f(-2) = -1(-2)^2 + 4 - 2 = -2 \\[4pt] f(-1) = -1(-1)^2 + 2 - 2 = -1 \\[4pt] f(0) = -1(0)^2 + 0 - 2 = -2
\end{aligned}
\]

\end{tcolorbox}

{Quadratic functions} {Form conversions} {Vertex form}

\end{exercise}

\begin{exercise}[]\protect\hypertarget{exr-u7-26}{}\label{exr-u7-26}

\emph{State the equation of the axis of symmetry.}

\textbf{a)} \(\displaystyle f(x) = -x^2 - 4x - 4\)

\textbf{b)} \(\displaystyle f(x) = 2x^2 - 20x + 55\)

\textbf{c)} \(\displaystyle f(x) = -x^2 + 8x - 13\)

\textbf{d)} \(\displaystyle f(x) = x^2 - 2x + 3\)

\textbf{e)} \(\displaystyle f(x) = -2x^2 - 4x - 4\)

\textbf{f)} \(\displaystyle f(x) = -x^2 + 2x + 5\)

\begin{tcolorbox}[enhanced jigsaw, arc=.35mm, breakable, toptitle=1mm, title=\textcolor{quarto-callout-tip-color}{\faLightbulb}\hspace{0.5em}{🔍 View Solutions}, colframe=quarto-callout-tip-color-frame, rightrule=.15mm, opacityback=0, titlerule=0mm, colback=white, coltitle=black, bottomtitle=1mm, toprule=.15mm, leftrule=.75mm, bottomrule=.15mm, colbacktitle=quarto-callout-tip-color!10!white, left=2mm, opacitybacktitle=0.6]

\textbf{a)}

\[
\begin{aligned}
x = \dfrac{-b}{2a} = \dfrac{-(-4)}{2(-1)} = -2
\end{aligned}
\] \textbf{b)}

\[
\begin{aligned}
x = \dfrac{-b}{2a} = \dfrac{-(-20)}{2(2)} = 5
\end{aligned}
\] \textbf{c)}

\[
\begin{aligned}
x = \dfrac{-b}{2a} = \dfrac{-(8)}{2(-1)} = 4
\end{aligned}
\] \textbf{d)}

\[
\begin{aligned}
x = \dfrac{-b}{2a} = \dfrac{-(-2)}{2(1)} = 1
\end{aligned}
\] \textbf{e)}

\[
\begin{aligned}
x = \dfrac{-b}{2a} = \dfrac{-(-4)}{2(-2)} = -1
\end{aligned}
\] \textbf{f)}

\[
\begin{aligned}
x = \dfrac{-b}{2a} = \dfrac{-(2)}{2(-1)} = 1
\end{aligned}
\]

\end{tcolorbox}

{Vertex} {Axis of symmetry} {Roots}

\end{exercise}

\begin{exercise}[]\protect\hypertarget{exr-u7-34}{}\label{exr-u7-34}

\emph{State whether each graph represents a function. Write Yes or No.}

\textbf{a)}

x y

\textbf{b)}

x y

\textbf{c)}

x y

\textbf{d)}

x y

\begin{tcolorbox}[enhanced jigsaw, arc=.35mm, breakable, toptitle=1mm, title=\textcolor{quarto-callout-tip-color}{\faLightbulb}\hspace{0.5em}{🔍 View Solutions}, colframe=quarto-callout-tip-color-frame, rightrule=.15mm, opacityback=0, titlerule=0mm, colback=white, coltitle=black, bottomtitle=1mm, toprule=.15mm, leftrule=.75mm, bottomrule=.15mm, colbacktitle=quarto-callout-tip-color!10!white, left=2mm, opacitybacktitle=0.6]

✅ \textbf{(a)} \textbf{Yes} --- passes the Vertical Line Test\\
❌ \textbf{(b)} \textbf{No} --- fails the Vertical Line Test\\
❌ \textbf{(c)} \textbf{No} --- fails the Vertical Line Test\\
✅ \textbf{(d)} \textbf{Yes} --- passes the Vertical Line Test

\end{tcolorbox}

{Function definition} {Domain \& range} {Arrow diagrams}

\end{exercise}

∑ × π ÷ √ ◆ ∞ ± ∫ ≠ Δ

\subsection{🟡 Medium}

\begin{exercise}[]\protect\hypertarget{exr-u7-9}{}\label{exr-u7-9}

\emph{Read and interpret the linear function.}

A phone plan costs \$30 per month plus \$3 per GB of data.

\textbf{(a)} Write an equation for \(T\) in terms of \(g\).

\textbf{(b)} Interpret the slope and intercept in context.

\textbf{(c)} Find the value when \(g = 5\).

\begin{tcolorbox}[enhanced jigsaw, arc=.35mm, breakable, toptitle=1mm, title=\textcolor{quarto-callout-tip-color}{\faLightbulb}\hspace{0.5em}{🔍 View Solution}, colframe=quarto-callout-tip-color-frame, rightrule=.15mm, opacityback=0, titlerule=0mm, colback=white, coltitle=black, bottomtitle=1mm, toprule=.15mm, leftrule=.75mm, bottomrule=.15mm, colbacktitle=quarto-callout-tip-color!10!white, left=2mm, opacitybacktitle=0.6]

\[
\begin{aligned}
\text{(a) } T = 3g + 30 \\[4pt] \text{(b) Slope } = 3: \text{ cost per GB used} \\[4pt] \text{Intercept } = 30: \text{ monthly base cost} \\[4pt] \text{(c) When } g = 5: T = 3(5) + 30 = 45 \text{ \$}
\end{aligned}
\]

\end{tcolorbox}

{Linear functions} {Slope} {y-intercept}

\end{exercise}

\begin{exercise}[]\protect\hypertarget{exr-u7-12}{}\label{exr-u7-12}

\emph{Interpret the slope in context.}

A plumber charges a call-out fee of \$39 plus \$2 per hour. The equation
is \(C = 2h + 39\).

\textbf{(a)} Identify the slope and interpret it in context.

\textbf{(b)} What does the intercept represent?

\begin{tcolorbox}[enhanced jigsaw, arc=.35mm, breakable, toptitle=1mm, title=\textcolor{quarto-callout-tip-color}{\faLightbulb}\hspace{0.5em}{🔍 View Solution}, colframe=quarto-callout-tip-color-frame, rightrule=.15mm, opacityback=0, titlerule=0mm, colback=white, coltitle=black, bottomtitle=1mm, toprule=.15mm, leftrule=.75mm, bottomrule=.15mm, colbacktitle=quarto-callout-tip-color!10!white, left=2mm, opacitybacktitle=0.6]

\[
\begin{aligned}
\text{(a) Slope } = 2 \text{ (C per hour)} \\[4pt] \text{cost per hour of work} \\[4pt] \text{(b) Intercept } = 39: \text{ fixed call-out fee}
\end{aligned}
\]

\end{tcolorbox}

{Rate of change} {Slope formula} {Interpretation}

\end{exercise}

\begin{exercise}[]\protect\hypertarget{exr-u7-13}{}\label{exr-u7-13}

\emph{Find intercepts and analyse the graph.}

Consider the linear function \(y = -2x + 1\).

x y -4 -3 -2 -1 1 2 3 4 (-4,9) (0,1)

\textbf{(a)} Find the \(x\)-intercept and \(y\)-intercept.

\textbf{(b)} Use the graph to verify your intercepts.

\begin{tcolorbox}[enhanced jigsaw, arc=.35mm, breakable, toptitle=1mm, title=\textcolor{quarto-callout-tip-color}{\faLightbulb}\hspace{0.5em}{🔍 View Solution}, colframe=quarto-callout-tip-color-frame, rightrule=.15mm, opacityback=0, titlerule=0mm, colback=white, coltitle=black, bottomtitle=1mm, toprule=.15mm, leftrule=.75mm, bottomrule=.15mm, colbacktitle=quarto-callout-tip-color!10!white, left=2mm, opacitybacktitle=0.6]

\[
\begin{aligned}
y\text{-int: set }x=0: y = 1 \\[4pt] x\text{-int: set }y=0: 0 = -2x + 1 \Rightarrow x = \dfrac{1}{2} \\[4pt] \text{Plot the points: x-int at }(\dfrac{1}{2},\, 0)\text{ and y-int at }(0,\ 1).
\end{aligned}
\]

\end{tcolorbox}

{Graphing} {Intercepts} {Parallel \& perpendicular}

\end{exercise}

\begin{exercise}[]\protect\hypertarget{exr-u7-14}{}\label{exr-u7-14}

\emph{Write the equation of the line.}

A line passes through \((1,\ -7)\) and \((0,\ -5)\).

x y -4 -3 -2 -1 1 2 3 4 (-4,3) (0,-5)

\textbf{(a)} Calculate the slope of the line.

\textbf{(b)} Write the equation in the form \(y = mx + c\).

\begin{tcolorbox}[enhanced jigsaw, arc=.35mm, breakable, toptitle=1mm, title=\textcolor{quarto-callout-tip-color}{\faLightbulb}\hspace{0.5em}{🔍 View Solution}, colframe=quarto-callout-tip-color-frame, rightrule=.15mm, opacityback=0, titlerule=0mm, colback=white, coltitle=black, bottomtitle=1mm, toprule=.15mm, leftrule=.75mm, bottomrule=.15mm, colbacktitle=quarto-callout-tip-color!10!white, left=2mm, opacitybacktitle=0.6]

\[
\begin{aligned}
m = \dfrac{-5 - (-7)}{0 - (1)} = -2 \\[4pt] y - (-7) = -2(x - (1)) \\[4pt] y = -2x - 5
\end{aligned}
\]

\end{tcolorbox}

{Graphing} {Intercepts} {Parallel \& perpendicular}

\end{exercise}

\begin{exercise}[]\protect\hypertarget{exr-u7-20}{}\label{exr-u7-20}

\emph{Analyse the parabola.}

The parabola is given by \(f(x) = 2(x - 3)^2 - 1\).

x y x=3 (3, -1) 2.3 3.7

\textbf{(a)} State the vertex and axis of symmetry.

\textbf{(b)} Does the parabola open upward or downward? Why?

\textbf{(c)} Write the equation in standard form
\(f(x) = ax^2 + bx + c\).

\begin{tcolorbox}[enhanced jigsaw, arc=.35mm, breakable, toptitle=1mm, title=\textcolor{quarto-callout-tip-color}{\faLightbulb}\hspace{0.5em}{🔍 View Solution}, colframe=quarto-callout-tip-color-frame, rightrule=.15mm, opacityback=0, titlerule=0mm, colback=white, coltitle=black, bottomtitle=1mm, toprule=.15mm, leftrule=.75mm, bottomrule=.15mm, colbacktitle=quarto-callout-tip-color!10!white, left=2mm, opacitybacktitle=0.6]

\[
\begin{aligned}
\text{(a) Vertex: } (3,\ -1), \text{ axis: } x = 3 \\[4pt] \text{(b) Opens upward since } a = 2 \text{{ > 0}} \\[4pt] \text{(c) Standard form: } f(x) = 2x^2 - 12x + 17
\end{aligned}
\]

\end{tcolorbox}

{Quadratic functions} {Form conversions} {Vertex form}

\end{exercise}

\begin{exercise}[]\protect\hypertarget{exr-u7-21}{}\label{exr-u7-21}

\emph{Convert from vertex form to standard form.}

\textbf{a)} \(\displaystyle f(x) = (x + 2)^2 + 6\)

\textbf{b)} \(\displaystyle f(x) = -2(x - 2)^2 - 6\)

\textbf{c)} \(\displaystyle f(x) = 2(x + 1)^2 - 3\)

\textbf{d)} \(\displaystyle f(x) = 2(x - 1)^2 - 2\)

\begin{tcolorbox}[enhanced jigsaw, arc=.35mm, breakable, toptitle=1mm, title=\textcolor{quarto-callout-tip-color}{\faLightbulb}\hspace{0.5em}{🔍 View Solutions}, colframe=quarto-callout-tip-color-frame, rightrule=.15mm, opacityback=0, titlerule=0mm, colback=white, coltitle=black, bottomtitle=1mm, toprule=.15mm, leftrule=.75mm, bottomrule=.15mm, colbacktitle=quarto-callout-tip-color!10!white, left=2mm, opacitybacktitle=0.6]

\textbf{a)}

\[
\begin{aligned}
\text{Expand the bracket:} \\[4pt] = 1(x^2 + 4x + 4) + 6 \\[4pt] = x^2 + 4x + 10
\end{aligned}
\] \textbf{b)}

\[
\begin{aligned}
\text{Expand the bracket:} \\[4pt] = -2(x^2 - 4x + 4) - 6 \\[4pt] = -2x^2 + 8x - 14
\end{aligned}
\] \textbf{c)}

\[
\begin{aligned}
\text{Expand the bracket:} \\[4pt] = 2(x^2 + 2x + 1) - 3 \\[4pt] = 2x^2 + 4x - 1
\end{aligned}
\] \textbf{d)}

\[
\begin{aligned}
\text{Expand the bracket:} \\[4pt] = 2(x^2 - 2x + 1) - 2 \\[4pt] = 2x^2 - 4x
\end{aligned}
\]

\end{tcolorbox}

{Quadratic functions} {Form conversions} {Vertex form}

\end{exercise}

\begin{exercise}[]\protect\hypertarget{exr-u7-24}{}\label{exr-u7-24}

\emph{Expand to standard form.}

\textbf{a)} \(\displaystyle f(x) = (x - 1)(x - 5)\)

\textbf{b)} \(\displaystyle f(x) = -(x - 1)(x - 5)\)

\textbf{c)} \(\displaystyle f(x) = (x + 1)(x + 4)\)

\textbf{d)} \(\displaystyle f(x) = 2(x + 5)(x - 4)\)

\begin{tcolorbox}[enhanced jigsaw, arc=.35mm, breakable, toptitle=1mm, title=\textcolor{quarto-callout-tip-color}{\faLightbulb}\hspace{0.5em}{🔍 View Solutions}, colframe=quarto-callout-tip-color-frame, rightrule=.15mm, opacityback=0, titlerule=0mm, colback=white, coltitle=black, bottomtitle=1mm, toprule=.15mm, leftrule=.75mm, bottomrule=.15mm, colbacktitle=quarto-callout-tip-color!10!white, left=2mm, opacitybacktitle=0.6]

\textbf{a)}

\[
\begin{aligned}
\text{Expand} \\[4pt] = 1(x^2 - 6x + 5) \\[4pt] = x^2 - 6x + 5
\end{aligned}
\] \textbf{b)}

\[
\begin{aligned}
\text{Expand} \\[4pt] = -1(x^2 - 6x + 5) \\[4pt] = -x^2 + 6x - 5
\end{aligned}
\] \textbf{c)}

\[
\begin{aligned}
\text{Expand} \\[4pt] = 1(x^2 + 5x + 4) \\[4pt] = x^2 + 5x + 4
\end{aligned}
\] \textbf{d)}

\[
\begin{aligned}
\text{Expand} \\[4pt] = 2(x^2 + x - 20) \\[4pt] = 2x^2 + 2x - 40
\end{aligned}
\]

\end{tcolorbox}

{Quadratic functions} {Form conversions} {Vertex form}

\end{exercise}

\begin{exercise}[]\protect\hypertarget{exr-u7-25}{}\label{exr-u7-25}

\emph{Find the vertex of each quadratic function.}

\textbf{a)} \(\displaystyle f(x) = 3x^2 + 6x + 2\)

\textbf{b)} \(\displaystyle f(x) = 2x^2 + 12x + 22\)

\textbf{c)} \(\displaystyle f(x) = x^2 - 2x + 5\)

\textbf{d)} \(\displaystyle f(x) = -2x^2 + 2\)

\textbf{e)} \(\displaystyle f(x) = x^2 - 4x + 3\)

\textbf{f)} \(\displaystyle f(x) = 2x^2 + 2\)

\begin{tcolorbox}[enhanced jigsaw, arc=.35mm, breakable, toptitle=1mm, title=\textcolor{quarto-callout-tip-color}{\faLightbulb}\hspace{0.5em}{🔍 View Solutions}, colframe=quarto-callout-tip-color-frame, rightrule=.15mm, opacityback=0, titlerule=0mm, colback=white, coltitle=black, bottomtitle=1mm, toprule=.15mm, leftrule=.75mm, bottomrule=.15mm, colbacktitle=quarto-callout-tip-color!10!white, left=2mm, opacitybacktitle=0.6]

\textbf{a)}

\[
\begin{aligned}
x_{\text{vertex}} = \dfrac{-b}{2a} = \dfrac{-(6)}{2(3)} = -1 \\[4pt] y = 3(-1)^2 - 6 + 2 = -1 \\[4pt] \text{Vertex: }\ x=-1,\ y=-1
\end{aligned}
\] \textbf{b)}

\[
\begin{aligned}
x_{\text{vertex}} = \dfrac{-b}{2a} = \dfrac{-(12)}{2(2)} = -3 \\[4pt] y = 2(-3)^2 - 36 + 22 = 4 \\[4pt] \text{Vertex: }\ x=-3,\ y=4
\end{aligned}
\] \textbf{c)}

\[
\begin{aligned}
x_{\text{vertex}} = \dfrac{-b}{2a} = \dfrac{-(-2)}{2(1)} = 1 \\[4pt] y = 1(1)^2 - 2 + 5 = 4 \\[4pt] \text{Vertex: }\ x=1,\ y=4
\end{aligned}
\] \textbf{d)}

\[
\begin{aligned}
x_{\text{vertex}} = \dfrac{-b}{2a} = \dfrac{-(0)}{2(-2)} = 0 \\[4pt] y = -2(0)^2  + 2 = 2 \\[4pt] \text{Vertex: }\ x=0,\ y=2
\end{aligned}
\] \textbf{e)}

\[
\begin{aligned}
x_{\text{vertex}} = \dfrac{-b}{2a} = \dfrac{-(-4)}{2(1)} = 2 \\[4pt] y = 1(2)^2 - 8 + 3 = -1 \\[4pt] \text{Vertex: }\ x=2,\ y=-1
\end{aligned}
\] \textbf{f)}

\[
\begin{aligned}
x_{\text{vertex}} = \dfrac{-b}{2a} = \dfrac{-(0)}{2(2)} = 0 \\[4pt] y = 2(0)^2  + 2 = 2 \\[4pt] \text{Vertex: }\ x=0,\ y=2
\end{aligned}
\]

\end{tcolorbox}

{Vertex} {Axis of symmetry} {Roots}

\end{exercise}

\begin{exercise}[]\protect\hypertarget{exr-u7-27}{}\label{exr-u7-27}

\emph{Read key features from the graph.}

The graph of a quadratic function is shown.

x y x=2 (2, -4) 0 4

\textbf{(a)} State the roots of the function.

\textbf{(b)} State the coordinates of the vertex.

\textbf{(c)} Write the equation of the axis of symmetry.

\begin{tcolorbox}[enhanced jigsaw, arc=.35mm, breakable, toptitle=1mm, title=\textcolor{quarto-callout-tip-color}{\faLightbulb}\hspace{0.5em}{🔍 View Solution}, colframe=quarto-callout-tip-color-frame, rightrule=.15mm, opacityback=0, titlerule=0mm, colback=white, coltitle=black, bottomtitle=1mm, toprule=.15mm, leftrule=.75mm, bottomrule=.15mm, colbacktitle=quarto-callout-tip-color!10!white, left=2mm, opacitybacktitle=0.6]

\[
\begin{aligned}
\text{Roots: } x = 0 \text{ and } x = 4 \\[4pt] \text{Vertex: } (2,\ -4) \\[4pt] \text{Axis of symmetry: } x = 2
\end{aligned}
\]

\end{tcolorbox}

{Vertex} {Axis of symmetry} {Roots}

\end{exercise}

\begin{exercise}[]\protect\hypertarget{exr-u7-28}{}\label{exr-u7-28}

\emph{State the domain and range of each quadratic function.}

\textbf{a)} \(\displaystyle f(x) = 2(x - 3)^2 + 2\)

\textbf{b)} \(\displaystyle f(x) = -(x - 4)^2 + 2\)

\textbf{c)} \(\displaystyle f(x) = -(x + 3)^2 - 1\)

\textbf{d)} \(\displaystyle f(x) = (x + 3)^2 - 2\)

\begin{tcolorbox}[enhanced jigsaw, arc=.35mm, breakable, toptitle=1mm, title=\textcolor{quarto-callout-tip-color}{\faLightbulb}\hspace{0.5em}{🔍 View Solutions}, colframe=quarto-callout-tip-color-frame, rightrule=.15mm, opacityback=0, titlerule=0mm, colback=white, coltitle=black, bottomtitle=1mm, toprule=.15mm, leftrule=.75mm, bottomrule=.15mm, colbacktitle=quarto-callout-tip-color!10!white, left=2mm, opacitybacktitle=0.6]

\textbf{a)}

\[
\begin{aligned}
\text{Domain: all real numbers} \\[4pt] x \in \mathbb{R} \\[4pt] \text{Parabola opens upward; vertex at }(3,\ 2) \\[4pt] \text{Range: } y \geq 2
\end{aligned}
\] \textbf{b)}

\[
\begin{aligned}
\text{Domain: all real numbers} \\[4pt] x \in \mathbb{R} \\[4pt] \text{Parabola opens downward; vertex at }(4,\ 2) \\[4pt] \text{Range: } y \leq 2
\end{aligned}
\] \textbf{c)}

\[
\begin{aligned}
\text{Domain: all real numbers} \\[4pt] x \in \mathbb{R} \\[4pt] \text{Parabola opens downward; vertex at }(-3,\ -1) \\[4pt] \text{Range: } y \leq -1
\end{aligned}
\] \textbf{d)}

\[
\begin{aligned}
\text{Domain: all real numbers} \\[4pt] x \in \mathbb{R} \\[4pt] \text{Parabola opens upward; vertex at }(-3,\ -2) \\[4pt] \text{Range: } y \geq -2
\end{aligned}
\]

\end{tcolorbox}

{Domain \& range} {Vertex} {Quadratic functions}

\end{exercise}

\begin{exercise}[]\protect\hypertarget{exr-u7-29}{}\label{exr-u7-29}

\emph{State the domain and range from the graph.}

The graph of a quadratic function is shown.

x y (-2, -1)

\textbf{(a)} State the domain of the function.

\textbf{(b)} State the range of the function.

\begin{tcolorbox}[enhanced jigsaw, arc=.35mm, breakable, toptitle=1mm, title=\textcolor{quarto-callout-tip-color}{\faLightbulb}\hspace{0.5em}{🔍 View Solution}, colframe=quarto-callout-tip-color-frame, rightrule=.15mm, opacityback=0, titlerule=0mm, colback=white, coltitle=black, bottomtitle=1mm, toprule=.15mm, leftrule=.75mm, bottomrule=.15mm, colbacktitle=quarto-callout-tip-color!10!white, left=2mm, opacitybacktitle=0.6]

\[
\begin{aligned}
\text{(a) Domain: } x \in \mathbb{R} \\[4pt] \text{(b) Vertex at }(-2,\ -1)\text{; opens upward} \\[4pt] \text{Range: } y \geq -1
\end{aligned}
\]

\end{tcolorbox}

{Domain \& range} {Vertex} {Quadratic functions}

\end{exercise}

\begin{exercise}[]\protect\hypertarget{exr-u7-32}{}\label{exr-u7-32}

\emph{Read the graph and find the equation.}

The graph of a quadratic function is shown.

x y

\textbf{(a)} State the roots (x-intercepts) of the function.

\textbf{(b)} State the coordinates of the vertex.

\textbf{(c)} Write the equation of the axis of symmetry.

\textbf{(d)} Write the equation of the quadratic function.

\begin{tcolorbox}[enhanced jigsaw, arc=.35mm, breakable, toptitle=1mm, title=\textcolor{quarto-callout-tip-color}{\faLightbulb}\hspace{0.5em}{🔍 View Solution}, colframe=quarto-callout-tip-color-frame, rightrule=.15mm, opacityback=0, titlerule=0mm, colback=white, coltitle=black, bottomtitle=1mm, toprule=.15mm, leftrule=.75mm, bottomrule=.15mm, colbacktitle=quarto-callout-tip-color!10!white, left=2mm, opacitybacktitle=0.6]

\[
\begin{aligned}
\text{(a) Roots: } x = -1 \text{ and } x = 1 \\[4pt] \text{(b) Vertex: } (0,\ -2) \\[4pt] \text{(c) Axis of symmetry: } x = 0 \\[4pt] \text{(d) Equation: } 2(x + 1)(x - 1) \text{ or } 2x^2 - 2
\end{aligned}
\]

\end{tcolorbox}

{Vertex} {Axis of symmetry} {Roots}

\end{exercise}

∑ × π ÷ √ ◆ ∞ ± ∫ ≠ Δ

\subsection{🔴 Challenging}

\begin{exercise}[]\protect\hypertarget{exr-u7-15}{}\label{exr-u7-15}

\emph{Find the equation of each line described.}

\textbf{a)}
\(\displaystyle \text{Parallel to } y = 3x + 3\text{, through } (-2,\ -3)\)

\textbf{b)}
\(\displaystyle \text{Parallel to } y = 3x - 4\text{, through } (3,\ -2)\)

\textbf{c)}
\(\displaystyle \text{Parallel to } y = -2x + 4\text{, through } (-3,\ 2)\)

\textbf{d)}
\(\displaystyle \text{Parallel to } y = -x + 1\text{, through } (-3,\ 1)\)

\begin{tcolorbox}[enhanced jigsaw, arc=.35mm, breakable, toptitle=1mm, title=\textcolor{quarto-callout-tip-color}{\faLightbulb}\hspace{0.5em}{🔍 View Solutions}, colframe=quarto-callout-tip-color-frame, rightrule=.15mm, opacityback=0, titlerule=0mm, colback=white, coltitle=black, bottomtitle=1mm, toprule=.15mm, leftrule=.75mm, bottomrule=.15mm, colbacktitle=quarto-callout-tip-color!10!white, left=2mm, opacitybacktitle=0.6]

\textbf{a)}

\[
\begin{aligned}
\text{Parallel lines have equal slopes: } m = 3 \\[4pt] \text{Through }(-2,\ -3): -3 = 3(-2) + c \\[4pt] c = 3 \\[4pt] y = 3x + 3
\end{aligned}
\] \textbf{b)}

\[
\begin{aligned}
\text{Parallel lines have equal slopes: } m = 3 \\[4pt] \text{Through }(3,\ -2): -2 = 3(3) + c \\[4pt] c = -11 \\[4pt] y = 3x - 11
\end{aligned}
\] \textbf{c)}

\[
\begin{aligned}
\text{Parallel lines have equal slopes: } m = -2 \\[4pt] \text{Through }(-3,\ 2): 2 = -2(-3) + c \\[4pt] c = -4 \\[4pt] y = -2x - 4
\end{aligned}
\] \textbf{d)}

\[
\begin{aligned}
\text{Parallel lines have equal slopes: } m = -1 \\[4pt] \text{Through }(-3,\ 1): 1 = -1(-3) + c \\[4pt] c = -2 \\[4pt] y = -x - 2
\end{aligned}
\]

\end{tcolorbox}

{Graphing} {Intercepts} {Parallel \& perpendicular}

\end{exercise}

\begin{exercise}[]\protect\hypertarget{exr-u7-16}{}\label{exr-u7-16}

\emph{Find the equation of each line described.}

\textbf{a)}
\(\displaystyle \text{Perpendicular to } y = -x - 4\text{, through } (2,\ -2)\)

\textbf{b)}
\(\displaystyle \text{Perpendicular to } y = 2x + 4\text{, through } (0,\ 0)\)

\textbf{c)}
\(\displaystyle \text{Perpendicular to } y = -x + 1\text{, through } (-1,\ 0)\)

\textbf{d)}
\(\displaystyle \text{Perpendicular to } y = 2x - 2\text{, through } (0,\ 1)\)

\begin{tcolorbox}[enhanced jigsaw, arc=.35mm, breakable, toptitle=1mm, title=\textcolor{quarto-callout-tip-color}{\faLightbulb}\hspace{0.5em}{🔍 View Solutions}, colframe=quarto-callout-tip-color-frame, rightrule=.15mm, opacityback=0, titlerule=0mm, colback=white, coltitle=black, bottomtitle=1mm, toprule=.15mm, leftrule=.75mm, bottomrule=.15mm, colbacktitle=quarto-callout-tip-color!10!white, left=2mm, opacitybacktitle=0.6]

\textbf{a)}

\[
\begin{aligned}
\text{Perpendicular slope: } m_\perp = \dfrac{-1}{-1} = 1 \\[4pt] \text{Through }(2,\ -2): -2 = 1(2) + c \\[4pt] c = -4 \\[4pt] y = x - 4
\end{aligned}
\] \textbf{b)}

\[
\begin{aligned}
\text{Perpendicular slope: } m_\perp = \dfrac{-1}{2} = -\dfrac{1}{2} \\[4pt] \text{Through }(0,\ 0): 0 = -\dfrac{1}{2}(0) + c \\[4pt] c = 0 \\[4pt] y = -\dfrac{1}{2}x
\end{aligned}
\] \textbf{c)}

\[
\begin{aligned}
\text{Perpendicular slope: } m_\perp = \dfrac{-1}{-1} = 1 \\[4pt] \text{Through }(-1,\ 0): 0 = 1(-1) + c \\[4pt] c = 1 \\[4pt] y = x + 1
\end{aligned}
\] \textbf{d)}

\[
\begin{aligned}
\text{Perpendicular slope: } m_\perp = \dfrac{-1}{2} = -\dfrac{1}{2} \\[4pt] \text{Through }(0,\ 1): 1 = -\dfrac{1}{2}(0) + c \\[4pt] c = 1 \\[4pt] y = -\dfrac{1}{2}x + 1
\end{aligned}
\]

\end{tcolorbox}

{Graphing} {Intercepts} {Parallel \& perpendicular}

\end{exercise}

\begin{exercise}[]\protect\hypertarget{exr-u7-17}{}\label{exr-u7-17}

\emph{Form and solve a linear model.}

A phone plan charges \$10 per month plus \$2 per GB of data.

\textbf{(a)} Write an equation for total monthly bill \(T\) after using
\(g\) GB.

\textbf{(b)} Find the bill for 12 GB usage.

\textbf{(c)} How many GB can be used for \$90?

\begin{tcolorbox}[enhanced jigsaw, arc=.35mm, breakable, toptitle=1mm, title=\textcolor{quarto-callout-tip-color}{\faLightbulb}\hspace{0.5em}{🔍 View Solution}, colframe=quarto-callout-tip-color-frame, rightrule=.15mm, opacityback=0, titlerule=0mm, colback=white, coltitle=black, bottomtitle=1mm, toprule=.15mm, leftrule=.75mm, bottomrule=.15mm, colbacktitle=quarto-callout-tip-color!10!white, left=2mm, opacitybacktitle=0.6]

\[
\begin{aligned}
\text{(a) } T = 2g + 10 \\[4pt] \text{(b) } T = 2(12) + 10 = 34 \text{ \$} \\[4pt] \text{(c) } 90 = 2g + 10 \Rightarrow g = 40 \text{ GB}
\end{aligned}
\]

\end{tcolorbox}

{Linear modelling} {Real-world context} {Interpretation}

\end{exercise}

\begin{exercise}[]\protect\hypertarget{exr-u7-22}{}\label{exr-u7-22}

\emph{Convert from vertex form to factored form.}

\textbf{a)} \(\displaystyle f(x) = (x + 1)^2 - 1\)

\textbf{b)} \(\displaystyle f(x) = (x - 2)^2 - 4\)

\textbf{c)} \(\displaystyle f(x) = (x + 3)^2 - 4\)

\textbf{d)} \(\displaystyle f(x) = (x - 4)^2 - 4\)

\begin{tcolorbox}[enhanced jigsaw, arc=.35mm, breakable, toptitle=1mm, title=\textcolor{quarto-callout-tip-color}{\faLightbulb}\hspace{0.5em}{🔍 View Solutions}, colframe=quarto-callout-tip-color-frame, rightrule=.15mm, opacityback=0, titlerule=0mm, colback=white, coltitle=black, bottomtitle=1mm, toprule=.15mm, leftrule=.75mm, bottomrule=.15mm, colbacktitle=quarto-callout-tip-color!10!white, left=2mm, opacitybacktitle=0.6]

\textbf{a)}

\[
\begin{aligned}
\text{Set } f(x) = 0: \\[4pt] (x - (-1))^2 = 1 \\[4pt] x + 1 = \pm 1 \\[4pt] x = -2 \text{ or } x = 0 \\[4pt] \therefore f(x) = (x + 2)(x)
\end{aligned}
\] \textbf{b)}

\[
\begin{aligned}
\text{Set } f(x) = 0: \\[4pt] (x - (2))^2 = 4 \\[4pt] x - 2 = \pm 2 \\[4pt] x = 0 \text{ or } x = 4 \\[4pt] \therefore f(x) = (x)(x - 4)
\end{aligned}
\] \textbf{c)}

\[
\begin{aligned}
\text{Set } f(x) = 0: \\[4pt] (x - (-3))^2 = 4 \\[4pt] x + 3 = \pm 2 \\[4pt] x = -5 \text{ or } x = -1 \\[4pt] \therefore f(x) = (x + 5)(x + 1)
\end{aligned}
\] \textbf{d)}

\[
\begin{aligned}
\text{Set } f(x) = 0: \\[4pt] (x - (4))^2 = 4 \\[4pt] x - 4 = \pm 2 \\[4pt] x = 2 \text{ or } x = 6 \\[4pt] \therefore f(x) = (x - 2)(x - 6)
\end{aligned}
\]

\end{tcolorbox}

{Quadratic functions} {Form conversions} {Vertex form}

\end{exercise}

\begin{exercise}[]\protect\hypertarget{exr-u7-23}{}\label{exr-u7-23}

\emph{Convert to vertex form by completing the square.}

\textbf{a)} \(\displaystyle f(x) = -x^2 + 2x - 7\)

\textbf{b)} \(\displaystyle f(x) = x^2 - 4\)

\textbf{c)} \(\displaystyle f(x) = -x^2 - 6\)

\textbf{d)} \(\displaystyle f(x) = x^2 - 4x + 3\)

\begin{tcolorbox}[enhanced jigsaw, arc=.35mm, breakable, toptitle=1mm, title=\textcolor{quarto-callout-tip-color}{\faLightbulb}\hspace{0.5em}{🔍 View Solutions}, colframe=quarto-callout-tip-color-frame, rightrule=.15mm, opacityback=0, titlerule=0mm, colback=white, coltitle=black, bottomtitle=1mm, toprule=.15mm, leftrule=.75mm, bottomrule=.15mm, colbacktitle=quarto-callout-tip-color!10!white, left=2mm, opacitybacktitle=0.6]

\textbf{a)}

\[
\begin{aligned}
\text{Factor out } -1\text{:} \\[4pt] = -1\left(x^2 - 2x\right) - 7 \\[4pt] = -1\left(x - 1\right)^2 + 1 - 7 \\[4pt] = -(x - 1)^2 - 6
\end{aligned}
\] \textbf{b)}

\[
\begin{aligned}
\text{Complete the square:} \\[4pt] = (x + 0)^2 - 0 - 4 \\[4pt] = (x + 0)^2 - 4 \\[4pt] = (x)^2 - 4
\end{aligned}
\] \textbf{c)}

\[
\begin{aligned}
\text{Factor out } -1\text{:} \\[4pt] = -1\left(x^2 + 0x\right) - 6 \\[4pt] = -1\left(x + 0\right)^2 + 0 - 6 \\[4pt] = -(x)^2 - 6
\end{aligned}
\] \textbf{d)}

\[
\begin{aligned}
\text{Complete the square:} \\[4pt] = (x - 2)^2 - 4 + 3 \\[4pt] = (x - 2)^2 - 1 \\[4pt] = (x - 2)^2 - 1
\end{aligned}
\]

\end{tcolorbox}

{Quadratic functions} {Form conversions} {Vertex form}

\end{exercise}

\begin{exercise}[]\protect\hypertarget{exr-u7-30}{}\label{exr-u7-30}

\emph{Analyse the projectile model.}

A ball is thrown upward. Its height (m) after \(t\) seconds is
\(h(t) = -5t^2 + 15t\).

t y x=3/2 (3/2, 11) 0 3

\textbf{(a)} What is the initial height of the ball?

\textbf{(b)} Find the maximum height and when it occurs.

\textbf{(c)} Set up the equation to find when the ball hits the ground.

\begin{tcolorbox}[enhanced jigsaw, arc=.35mm, breakable, toptitle=1mm, title=\textcolor{quarto-callout-tip-color}{\faLightbulb}\hspace{0.5em}{🔍 View Solution}, colframe=quarto-callout-tip-color-frame, rightrule=.15mm, opacityback=0, titlerule=0mm, colback=white, coltitle=black, bottomtitle=1mm, toprule=.15mm, leftrule=.75mm, bottomrule=.15mm, colbacktitle=quarto-callout-tip-color!10!white, left=2mm, opacitybacktitle=0.6]

\[
\begin{aligned}
\text{(a) } h(0) = 0 \text{ m — initial height} \\[4pt] \text{(b) Vertex: } t = \dfrac{-15}{2(-5)} = \dfrac{3}{2} \text{ s} \\[4pt] \text{Max height } = 11 \text{ m} \\[4pt] \text{(c) Set } h = 0: -5t^2 + 15t = 0
\end{aligned}
\]

\end{tcolorbox}

{Quadratic modelling} {Projectile motion} {Revenue}

\end{exercise}

\begin{exercise}[]\protect\hypertarget{exr-u7-31}{}\label{exr-u7-31}

\emph{Analyse the revenue model.}

A company sells items at price \(p\) dollars. Demand is \(100 - 2p\)
units.

\textbf{(a)} Write an equation for revenue \(R\) in terms of \(p\).

\textbf{(b)} Find the price that maximises revenue and state the maximum
revenue.

\begin{tcolorbox}[enhanced jigsaw, arc=.35mm, breakable, toptitle=1mm, title=\textcolor{quarto-callout-tip-color}{\faLightbulb}\hspace{0.5em}{🔍 View Solution}, colframe=quarto-callout-tip-color-frame, rightrule=.15mm, opacityback=0, titlerule=0mm, colback=white, coltitle=black, bottomtitle=1mm, toprule=.15mm, leftrule=.75mm, bottomrule=.15mm, colbacktitle=quarto-callout-tip-color!10!white, left=2mm, opacitybacktitle=0.6]

\[
\begin{aligned}
\text{(a) Demand } = 100 - 2p \\[4pt] R = p(100 - 2p) = 100p - 2p^2 \\[4pt] \text{(b) Vertex: } p = \dfrac{100}{2 \times 2} = 25 \\[4pt] R_{\text{max}} = 1250
\end{aligned}
\]

\end{tcolorbox}

{Quadratic modelling} {Projectile motion} {Revenue}

\end{exercise}

\begin{exercise}[]\protect\hypertarget{exr-u7-33}{}\label{exr-u7-33}

\emph{Find the intersection of the two functions.}

\(f(x) = x^2 - 9x + 6\) and \(g(x) = -2x - 4\)

\textbf{(a)} Set \(f(x) = g(x)\) and form a quadratic equation.

\textbf{(b)} Solve to find the \(x\)-coordinates of the intersection
points.

\textbf{(c)} State the coordinates of the intersection points.

\begin{tcolorbox}[enhanced jigsaw, arc=.35mm, breakable, toptitle=1mm, title=\textcolor{quarto-callout-tip-color}{\faLightbulb}\hspace{0.5em}{🔍 View Solution}, colframe=quarto-callout-tip-color-frame, rightrule=.15mm, opacityback=0, titlerule=0mm, colback=white, coltitle=black, bottomtitle=1mm, toprule=.15mm, leftrule=.75mm, bottomrule=.15mm, colbacktitle=quarto-callout-tip-color!10!white, left=2mm, opacitybacktitle=0.6]

\[
\begin{aligned}
\textbf{(a)} \quad x^2 - 9x + 6 = -2x - 4 \implies x^2 - 7x + 10 = 0 \\[6pt] \textbf{(b)} \quad (x - 2)(x - 5) = 0 \implies x = 2 \text{ or } x = 5 \\[6pt] \textbf{(c)} \quad (2,\ -8) \text{ and } (5,\ -14)
\end{aligned}
\]

\end{tcolorbox}

{Vertex} {Axis of symmetry} {Roots}

\end{exercise}

∑ × π ÷ √ ◆ ∞ ± ∫ ≠ Δ

\section{⚡ Test Your Knowledge}\label{test-your-knowledge-2}




\end{document}
